\chapter{Bar-cobar adjunction and curved CDG comparison}
\label{chap:bar-cobar-adjunction}

% Local macros (provide-command idempotent; required for cross-chapter
% references to thm:koszul-reflection in theorem_A_infinity_2.tex).
\providecommand{\Kosz}{\mathsf{Kosz}}
\providecommand{\co}{\mathrm{co}}

% Bicoloured primitive package macros, imported for cross-volume
% coherence with Vol II (\texttt{chiral-bar-cobar-vol2:main.tex}).
% The chiral algebra $\cA$ of this chapter is recovered as
% $A_b = \mathrm{End}_{\openFactCat}(b)$ for a boundary chart $b$
% (Vol II, Definition~\ref{def:bicoloured-primitive-package}).
\providecommand{\openFactCat}{\cC^{\mathrm{open}}}
\providecommand{\logCurve}{(X, D, \tau)}
\providecommand{\kappaChHodge}{\kappa_{\mathrm{ch}}^{\mathrm{Hodge}}}
\providecommand{\primPkgBicolour}{\bigl(\openFactCat|_{\logCurve},\, \cD^{\mathrm{cl}}|_{X};\, b,\, A_b,\, F^{\mathrm{cl}};\, \HalfBraid,\, \mathrm{tr}^{\mathrm{cl}}_{X};\, \TraceC\bigr)}
\providecommand{\bulkChirHoch}[1]{\mathrm{HH}^\bullet_{\mathrm{ch}}\bigl(#1\bigr)}
\providecommand{\HalfBraid}{\beta^{\mathrm{half}}}
\providecommand{\TraceC}{\mathrm{tr}^{\mathrm{op}}_{\openFactCat}}

\paragraph{Chart projection.}
The chiral algebra $\cA$ of this chapter is throughout the endomorphism
algebra
$\cA \;=\; A_b \;=\; \mathrm{End}_{\openFactCat}(b)$
of a boundary chart $b$ in the open factorisation dg-category
$\openFactCat$ on the tangentially-decorated curve
$\logCurve = (X, D, \tau)$, in the bicoloured primitive package of
Vol~II
\textup{(}\texttt{chiral-bar-cobar-vol2:chapters/theory/foundations.tex:def:bicoloured-primitive-package}\textup{)}.
The
bar-cobar inversion
$\Omegach(\barBch(\cA))\xrightarrow{\sim}\cA$
proved here is therefore an inversion \emph{at the level of the chart
$b$}, not at the level of the operadic primitive
$\primPkgBicolour$. On the universal stage chain
$\mathsf{P}\xrightarrow{\alpha}\mathsf{C}\xrightarrow{\beta}
 \mathsf{S}\xrightarrow{\gamma}\mathsf{Z}\xrightarrow{\delta}\mathsf{A}$
the present chapter operates at level~$\mathsf{C}$ (chart):
$\cA=A_b$ is the endomorphism object of one boundary vacuum,
$\barBch(\cA)$ is its level-$\mathsf{S}$ shadow,
and the level-$\mathsf{P}$ primitive $\primPkgBicolour$ together
with the level-$\mathsf{Z}$ derived chiral centre
$\bulkChirHoch{A_b}\simeq\Zderch{A_b}$ are the subject of Vol~II. The
choice of $b$ is a chart (licensing type~$\alpha$); changing $b$
produces Morita-equivalent endomorphism algebras whose bar complexes
are quasi-isomorphic but whose ordered structure can shift. The
inversion proved here is invariant under this choice, but the
identification $\cA = A_b$ should be read as a projection from the
bicoloured nine-tuple to its open-colour endomorphism object, in the
sense of Beilinson--Drinfeld~\cite{BD04} and
Lurie~\cite[\S5.5]{HA}; the homotopical machinery of the
adjunction itself is that of Loday--Vallette~\cite{LV12}.

The bar-cobar adjunction $\Omegach \dashv \barBch$ represents the
space of twisting morphisms
$\operatorname{Tw}(\cC, \cA) = \MC(\operatorname{Conv}(\cC, \cA))$
by a pair of functors between chiral algebras and conilpotent
chiral coalgebras on $\Ran(X)$
(Theorem~\ref{thm:bar-cobar-adjunction}). The universal twisting
morphism
$\tau_\cA \in \MC\bigl(\operatorname{Conv}(\barBch(\cA), \cA)\bigr)$
is the genus-$0$ symmetric projection of the ordered MC class
$\Theta_\cA^{\Eone}\in\MC({\gAmod}^{\Eone})$. In the ordered
convolution dg Lie algebra it satisfies
\[
D\Theta_\cA^{\Eone}
+\tfrac{1}{2}[\Theta_\cA^{\Eone},\Theta_\cA^{\Eone}]=0,
\]
and applying $\operatorname{av}\colon{\gAmod}^{\Eone}\to\gAmod$
gives the modular equation
\[
D\Theta_\cA+\tfrac{1}{2}[\Theta_\cA,\Theta_\cA]=0.
\]
The degree-$2$ symmetric projection of $\Theta_\cA$ is the modular
characteristic $\kappa(\cA)$.
Theorem~B (Theorem~\ref{thm:bar-cobar-inversion-qi}) says that
$\tau_\cA$ is \emph{acyclic} on the Koszul locus: the twisted tensor
product $\cA \otimes_{\tau_\cA} \barBch(\cA)$ has trivial cohomology,
so the counit
$\psi \colon \Omegach(\barBch(\cA)) \xrightarrow{\;\sim\;} \cA$
is a quasi-isomorphism. The MC equation supplies the identities; the
acyclicity of the induced twisting morphism supplies reconstruction.
Neither assertion identifies bar-cobar inversion with Verdier duality
or with the formation of the Koszul dual algebra~$\cA^!$.

Not every chiral
algebra is quadratic: Virasoro, $\mathcal{W}_N$, and
$\mathcal{W}_\infty$ all require curved or filtered presentations in
which $\mu_0 \neq 0$ and $\mu_1^2 = [\mu_0, -]_{\mu_2}$. On this
extended locus the adjunction survives in the language of curved
$\Ainf$~algebras, $I$-adic completions, and coderived categories.
The strict convolution dg Lie algebra
$\Convstr(\barBch(\cA), \cA)$ is replaced by its homotopy-invariant
envelope $\Convinf(\barBch(\cA), \cA)$, and the twisting morphism
mediating the adjunction becomes a Maurer--Cartan element in this
larger algebra. On the curved locus a bar-cobar comparison is an
ordinary chain quasi-isomorphism only after a curvature-killing
MC twist or after passing to the square-zero total differential.
Without such a square-zero model the comparison is read in the
second-kind coderived category: the cone is required to be coacyclic
and the equivalence is detected on comodules over the completed
CDG-coalgebra.

The precise content of Theorem~A
(Theorem~\ref{thm:bar-cobar-adjunction}) is this: there is an
adjunction on the homotopy/derived chiral categories
$\Omegach \dashv \barBch$
between chiral algebras and conilpotent chiral coalgebras on
$\Ran(X)$, with the unit
$\eta \colon \cC \to \barBch(\Omegach(\cC))$
and counit
$\psi \colon \Omegach(\barBch(\cA)) \to \cA$.
Verdier duality enters through a separate finite-type and completion
comparison:
\[
\mathbb{D}_{\Ran}\,\barBch(\cA) \simeq \cA^!_\infty ,
\]
where the left-hand side is a factorization \emph{algebra}, not a
coalgebra. This comparison constructs the homotopy Koszul-dual
algebra; it is not one of the unit or counit maps and it does not
recover~$\cA$.
The unified statement is the \emph{Koszul reflection}
\textup{(}Theorem~\ref{thm:koszul-reflection}\textup{)}:
on the Koszul locus the algebra-side counit
$\psi_\cA$ and the coalgebra-side unit $\eta_\cC$ are
quasi-isomorphisms \textup{(}clause~\ref{KR-ii}\textup{)}; off that
locus the corresponding cones are required to vanish only after the
completed coderived localization
$D^{\co}(\barBch\text{-}\mathrm{CoFact})$
\textup{(}clause~\ref{KR-iii}\textup{)}. The notation $K$ in the
reflection theorem denotes the induced Verdier/Koszul reflection on
the localized category; it is not a reversal of the canonical
adjunction direction $\Omegach\dashv\barBch$, and it is not a raw
chain-level assertion that $K^2=\mathrm{id}$ on curved CDG objects.
The adjunction is mediated by the canonical twisting morphism
$\tau \colon \barBch(\cA) \to \cA$, which is a Maurer--Cartan
element in the strict convolution dg~Lie algebra
$\Convstr(\barBch(\cA), \cA)$
(Definition~\ref{def:convolution-dg-lie}); at the homotopy level,
$\tau$ is an MC element in $\Convinf$ and therefore invariant under
quasi-isomorphic replacement
(Theorem~\ref{thm:operadic-homotopy-convolution}). On the strict
quadratic locus the adjunction is classical; off it, the
curvature $\mu_0 \in A^2$ forces $\mu_1^2(a) =
\mu_2(\mu_0, a) - \mu_2(a, \mu_0) = [\mu_0, a]_{\mu_2}$
(Theorem~\ref{thm:curvature-central}), so the raw fiberwise
differential is a CDG differential rather than an ordinary
differential. A square-zero differential is present only in the
uncurved/gauge-trivial case or in the corrected total model
$\Dg{g}^{\,2}=0$; otherwise the bar object is used through its
coderived comodule category. The four-regime
hierarchy organizes the successively stronger completions needed:
\begin{enumerate}[label=\textup{(\roman*)}]
\item \emph{Quadratic} ($d_0^{\,2} = 0$): strict $d_{\mathrm{bar}}^2 = 0$,
 no completion needed (Heisenberg, free fields, lattice VOAs);
\item \emph{Curved-central} ($d_0^{\,2} \neq 0$,
 $\mu_0 \in Z(A)$): completed bar-cobar with chain-level recovery
 only after curvature cancellation or in the square-zero total
 model; the raw fiberwise operator has geometric curvature
 $d_{\mathrm{fib}}^2 = \kappa \cdot \omega_g$ (affine algebras, Virasoro);
\item \emph{Filtered-complete} (non-quadratic OPE):
 $I$-adic completion required, work in the weight-completed
 coderived category
 ($\Walg_N$ at finite~$N$);
\item \emph{Programmatic} (infinite generators):
 MC4 completion with the same square-zero/coderived distinction
 (proved; Theorem~\ref{thm:completed-bar-cobar-strong}).
\end{enumerate}
This hierarchy is codified in Convention~\ref{conv:regime-tags} and
governs every theorem statement in this chapter and the next.

The fixed-curve assertions in this chapter have a separate
relative-family gate. Theorem~$A^{\infty,2}$
(Theorem~\ref{thm:A-infinity-2}) proves the conilpotent-complete
fixed-curve adjunction. Transport over a smooth family requires the
base-change obstruction
$[\alpha_{\mathrm{BC}}]$ of
Remark~\ref{rem:A-infinity-2-modular-family-scope} to vanish; extension
over the Deligne--Mumford boundary requires the nodal sewing
obstruction $[\beta_{\mathrm{nodal}}]$ of the same remark to vanish.
The logarithmic collision-extension hypothesis
Lemma~\ref{lem:R-twisted-descent}\textup{\ref{R3}} is only the
fixed-curve Fulton--MacPherson condition needed for ordered-to-symmetric
descent. It does not replace either relative base change or nodal
sewing.

The first part
(\S\ref{sec:curved-koszul-i-adic}--\S\ref{sec:reconstruction-duality-casimir})
develops curved $\Ainf$~structures, $I$-adic completion, and the
regime hierarchy (quadratic, curved-central, filtered-complete,
programmatic) that determines which completions are needed. The
second part
(\S\ref{sec:filtered-vs-curved-comprehensive}--\S\ref{sec:cech-hca})
proves inversion, establishes the categorical logarithm reading of
bar-cobar ($\barBch = \log$, $\Omegach = \exp$), and lifts the
whole story to the $\Ran$-space via the \v{C}ech resolution.
Once the adjunction and its inversion are established, the
genus-$0$ bar complex is a functor, and the passage to higher genus
(Chapter~\ref{chap:higher-genus}) amounts to deforming the
differential by curvature. The curved extension provides the
framework in which $\kappa(\cA) \cdot \omega_g$ and the shadow
obstruction tower $\Theta_\cA^{\leq r}$ operate.

\begin{remark}[The ordered bar as primitive]
\label{rem:bar-ordered-primacy}
\index{ordered bar complex!as primitive object}
\index{averaging map!av@$\operatorname{av}$}
The bar complex $\barBch(\cA)$ presented in this chapter is the
\emph{symmetric} bar: it lives on
$(s^{-1}\bar\cA)^{\otimes n}_{\Sigma_n}$ with the coshuffle
coproduct, a cocommutative coalgebra on unordered configurations.
The natural primitive is the \emph{ordered} bar complex
$\barB^{\mathrm{ord}}(\cA) = T^c(s^{-1}\bar\cA)$ of
Chapter~\ref{ch:ordered-associative-chiral-kd}, the cofree
conilpotent coalgebra with deconcatenation coproduct. The two
are related by a surjective map
\begin{equation}\label{eq:bar-ordered-to-symmetric}
\barB^{\mathrm{ord}}(\cA)
\;\xrightarrow{\;\operatorname{av}\;}
\barBch(\cA),
\qquad
\operatorname{av}\bigl(a_1 \otimes \cdots \otimes a_n\bigr)
=
\frac{1}{n!}\sum_{\sigma \in \Sigma_n}
a_{\sigma(1)} \otimes \cdots \otimes a_{\sigma(n)},
\end{equation}
which is the $\Sigma_n$-coinvariant projection at each degree.
This map is a quasi-isomorphism of chain complexes in
characteristic~$0$, but it is \emph{not} an isomorphism of
coalgebras: the deconcatenation coproduct on
$\barB^{\mathrm{ord}}$ is coassociative but not cocommutative,
while the coshuffle coproduct on $\barBch$ is cocommutative.
Passing to coinvariants destroys the ordered coproduct and with it
three layers of structure:
\begin{enumerate}[label=\textup{(\roman*)}]
\item \emph{The $R$-matrix.} The genus-$0$, degree-$2$ component
 of the ordered MC element
 $\Theta_\cA^{\Eone}$ is the classical $r$-matrix
 $r(z) \in \End(V^{\otimes 2}) \otimes \cO(*\Delta)$, a
 matrix-valued meromorphic function satisfying the classical
 Yang--Baxter equation. Its $\Sigma_2$-coinvariant is the scalar
 $\kappa(\cA) = \operatorname{av}(r(z))$ for abelian algebras
 (for non-abelian Kac--Moody,
 $\kappa = \operatorname{av}(r(z)) + \dim(\fg)/2$;
 see equation~\eqref{eq:intro-kappa-from-r}): the
 full spectral data collapses to a single number.
\item \emph{The braided monoidal category.} The deconcatenation
 coproduct on $\barB^{\mathrm{ord}}(\cA)$ controls the tensor
 product on the category of ordered factorization modules
 $\Factord(\cA)$. The $R$-matrix provides the braiding.
 After coinvariants, only the symmetric monoidal shadow survives.
\item \emph{The Swiss-cheese directionality.} The
 $\mathrm{SC}^{\mathrm{ch},\mathrm{top}}$~picture couples
 holomorphic collision data on~$\FM_k(\C)$ with topological
 interval splitting on~$\Conf_k(\R)$. On the ordered bar these
 appear as the bar differential and the deconcatenation
 coproduct, but they do \emph{not} make
 $\barB^{\mathrm{ord}}(\cA)$ itself an
 $\mathrm{SC}^{\mathrm{ch},\mathrm{top}}$ object: the ordered
 bar remains a single-coloured $\Eone$-chiral coalgebra, and the
 actual Swiss-cheese datum emerges only on the derived center
 pair $(C^\bullet_{\mathrm{ch}}(\cA,\cA),\,\cA)$.
\end{enumerate}
The five main theorems~A--H are invariants of the symmetric bar
and therefore survive averaging. The Yangian $R$-matrix, the
Drinfeld--Kohno comparison, and the full braided monoidal category
of line operators live in the ordered bar and are invisible to the
symmetric theory. This chapter develops the $\operatorname{av}$-image
of the ordered theory: every construction here has a richer ordered
ancestor in Chapter~\ref{ch:ordered-associative-chiral-kd}, and the
passage from ordered to symmetric is the unique source of information
loss \textup{(}Theorem~\ref{thm:e1-primacy}\textup{)}.
\end{remark}

\section{Curved Koszul duality and I-adic completion}
\label{sec:curved-koszul-i-adic}

The Virasoro algebra cannot be quadratic: the Virasoro OPE
$T(z)T(w) \sim (c/2)/(z-w)^4 + 2T(w)/(z-w)^2 + \partial T(w)/(z-w)$
involves a fourth-order pole, and the stress tensor $T$ appears on
both sides of the relation, so no finite-dimensional generating space
and quadratic relation space suffice. The bar complex
$\barB(\mathrm{Vir}_c)$ is curved: its
differential squares not to zero but to the commutator
$\mu_1^2(a) = [\mu_0, a]$, where $\mu_0 = \kappa \cdot \omega$ is
the curvature element encoding the central extension. The failure
of $d^2=0$ is the algebraic trace of the
genus-$1$ boundary $\delta_0 \in \partial\overline{\mathcal{M}}_{1,1}$:
nodal degeneration introduces curvature into every genus
tower.

The same phenomenon intensifies for $\mathcal{W}_N$ (composite
fields, higher-order poles) and becomes qualitatively new for
$\mathcal{W}_\infty$ (infinitely many generators, no finite
presentation). In each case, the categorical logarithm $\barB$
still produces a well-defined object, but recovering the original
algebra by exponentiation $\Omega \circ \barB \xrightarrow{\sim} \cA$
requires increasingly refined completions. Three regimes emerge:
\begin{itemize}
\item \emph{Curved structures}: $\mu_0 \neq 0$, the differential
 satisfies $\mu_1^2 = [\mu_0, -]$ (commutator, with minus sign).
\item \emph{$I$-adic completion}: the bar complex must be completed
 against the augmentation ideal to ensure convergence of infinite sums.
\item \emph{Filtered structures}: more general than curved; the
 filtration controls the infinite-dimensional hierarchy of generators
 (Gui--Li--Zeng~\cite{GLZ22}).
\end{itemize}

\subsection{Curved $A_\infty$ algebras: definitions}
\label{sec:curved-ainfty-definition}

\begin{definition}[Curved \texorpdfstring{$A_\infty$}{A-infinity} algebra]
\label{def:curved-ainfty}
\label{ex:curved-ainfty-complete}
A \emph{curved $A_\infty$ algebra} is a graded vector space $A$ with operations:
\begin{equation}
\{\mu_n: A^{\otimes n} \to A\}_{n \geq 0}
\end{equation}
of degree $2-n$, satisfying the \emph{curved $A_\infty$ relations}:
\begin{equation}\label{eq:curved-ainfty-relations}
\sum_{\substack{r+s+t=n \\ r,t \geq 0,\; s \geq 0}} (-1)^{rs+t}\, \mu_{r+1+t}(\mathrm{id}^{\otimes r} \otimes \mu_s \otimes \mathrm{id}^{\otimes t}) = 0
\end{equation}

The summation begins at $n = 0$: the operation
$\mu_0 \colon \mathbb{C} \to A$ is the \emph{curvature}.
\end{definition}

\begin{remark}[Consequences of the curved relations at low $n$]
\label{rem:curved-ainfty-low-n}
Setting $n = 1$ in~\eqref{eq:curved-ainfty-relations} gives
$\mu_1^2 = [\mu_0, -]_{\mu_2}$, so $\mu_1$ is a differential
only when $\mu_0$ is central with respect to~$\mu_2$.
For $n \geq 2$ the curved relations do \emph{not} coincide with
the ordinary $A_\infty$ relations unless $\mu_0=0$ or the
curvature insertions vanish in the chosen model. The next relation
already contains the terms
$\mu_3(\mu_0,a,b)$, $\mu_3(a,\mu_0,b)$, and
$\mu_3(a,b,\mu_0)$, with the Koszul signs of
\eqref{eq:curved-ainfty-relations}. Thus centrality of~$\mu_0$
kills only the binary commutator in $\mu_1^2$; it does not by
itself uncurve the higher $A_\infty$ structure.
\end{remark}

\begin{remark}[Curvature, backreaction, and \texorpdfstring{$d^2$}{d\textasciicircum 2}]\label{rem:curvature-backreaction-terminology}
\begin{enumerate}[label=(\roman*)]
\item The \emph{curvature} is $m_0 := \mu_0(1) \in A^2$.
\item $\mu_1^2(a) = \mu_2(m_0, a) - \mu_2(a, m_0)$: the
 square-zero assertion belongs to the corrected total bar
 differential or to a curvature-killed model. The raw curved
 CDG differential has square controlled by its curvature.
\item In physics, \emph{backreaction} is the same phenomenon: the BRST operator $Q = \mu_1$ satisfies $Q^2 \neq 0$, with the obstruction measured by $m_0$.
\item Centrality of $m_0$ gives $\mu_1^2=0$ for the internal operator
 (Theorem~\ref{thm:curvature-central}); it does not uncurve the
 algebra. A square-zero uncurved model is obtained by MC twisting only
 when the completed curvature class is gauge-trivial.
\end{enumerate}
\end{remark}

%/C31: family-specific scalar couplings checked against landscape_census.tex and bershadsky_polyakov.tex:
% kappa(\widehat{\mathfrak{g}}_k)=(k+h^\vee)\dim\mathfrak{g}/(2h^\vee), kappa(\mathrm{Vir}_c)=c/2, kappa(\mathcal{B}^k)=c/6.
\begin{remark}[Gravity dictionary for curved \texorpdfstring{$A_\infty$}{A-infinity}]
\label{rem:curved-ainfty-gravity-dictionary}
For genus $g \geq 1$, the curved fiberwise bar model carries the scalar curvature element
\[
 m_0^{(g)} = \kappa(\mathcal{A}) \cdot \omega_g \cdot \mathbf{1},
 \qquad
 \omega_g = c_1(\lambda) \in H^2(\overline{\mathcal{M}}_g),
\]
where $\lambda = \det \mathbb{E}$ is the Hodge line bundle on the moduli space. The fiberwise differential satisfies
\[
 d_{\mathrm{fib}}^2 = \kappa(\mathcal{A}) \cdot \omega_g,
\]
while the period-corrected total differential satisfies
$\Dg{g}^{\,2}=0$ in the square-zero total model. This gives the
gravity dictionary:
\begin{enumerate}[label=(\roman*)]
\item $m_0^{(g)} \neq 0$ if and only if $\kappa(\mathcal{A}) \neq 0$, equivalently if and only if $d_{\mathrm{fib}}^2 \neq 0$. This is the precise sense in which gravity is turned on. The scalar curvature term $m_0^{(g)}$ is the cosmological-constant term of the genus-$g$ curved $A_\infty$ model. This statement is fiberwise. For the standard families $m_0^{(g)}$ is central, so the internal commutator formula $m_1^2 = [m_0,-]$ may vanish even when $d_{\mathrm{fib}}^2 \neq 0$.
\item $\kappa(\mathcal{A})$ is the coefficient of the Hodge class in $d_{\mathrm{fib}}^2$, hence the strength of the genus-$g$ gravitational backreaction. In that dictionary, $\kappa(\mathcal{A})$ is the Newton coupling. It is family-specific: $\kappa(\widehat{\mathfrak{g}}_k) = (k+h^\vee)\dim\mathfrak{g}/(2h^\vee)$ for affine Kac--Moody, $\kappa(\mathrm{Vir}_c) = c/2$ for Virasoro, and $\kappa(\mathcal{B}^k) = c/6$ for Bershadsky--Polyakov.
\item The critical gravitational coupling is the fixed point of the duality involution on the scalar lane. When the scalar complementarity constant is written as $\kappa(\mathcal{A}) + \kappa(\mathcal{A}^!) = K_\kappa(\mathcal{A})$, the fixed point is $\kappa_{\mathrm{crit}} = K_\kappa(\mathcal{A})/2$. This constant is family-specific. For affine Kac--Moody, Heisenberg, and free-field families one has $K_\kappa = 0$, so $\kappa_{\mathrm{crit}} = 0$. For Virasoro one has $K_\kappa = 13$, so $\kappa_{\mathrm{crit}} = 13/2$, equivalently $c = 13$. For the Bershadsky--Polyakov family the manuscript uses $\kappa(\mathcal{B}^k) = c/6$ together with the central-charge conductor $c(\mathcal{B}^k) + c((\mathcal{B}^k)^!) = 196$; hence the self-dual central charge is $c = 98$, and the induced scalar critical coupling is $\kappa_{\mathrm{crit}} = 49/3$.
\end{enumerate}
\end{remark}

\begin{theorem}[Curvature as \texorpdfstring{$\mu_1$}{mu1}-cycle; \ClaimStatusProvedHere]
\label{thm:curvature-central}
For a curved $A_\infty$ algebra $(A, \mu_0, \mu_1, \mu_2, \ldots)$, the curvature element satisfies:
\begin{enumerate}
\item $\mu_1(\mu_0) = 0$ \quad (the curvature is a $\mu_1$-cycle);
\item $\mu_1^2 = [\mu_0, -]_{\mu_2}$, where $[\mu_0, a]_{\mu_2} := \mu_2(\mu_0, a) - \mu_2(a, \mu_0)$ is the inner derivation.
\end{enumerate}
In particular, $\mu_1$ fails to be a differential precisely when $\mu_0$
is \emph{not} central with respect to $\mu_2$. For the standard
positive-energy families considered here, the genus-$g$ curvature is
$\mu_0 = \kappa(\cA)\cdot\omega_g\cdot\mathbf{1}$: a scalar multiple of
the vacuum vector, which is central in the binary bar operation
$\mu_2$, so $[\mu_0, -]_{\mu_2} = 0$ automatically. The curvature
therefore enters not through $\mu_2$ but through higher operations
($\mu_3, \mu_4, \ldots$) and through the fiberwise Hodge bundle
structure on $\overline{\cM}_{g,n}$.
\end{theorem}

\begin{proof}
The curved $A_\infty$ relations $\sum_{r+s+t=n}(-1)^{rs+t}\mu_{r+1+t}(\mathrm{id}^{\otimes r}\otimes \mu_s \otimes \mathrm{id}^{\otimes t})=0$ give, at each level:
\begin{itemize}
\item $n=0$: $\mu_1(\mu_0) = 0$ \quad (the curvature is a $\mu_1$-cycle).
\item $n=1$: $\mu_1^2(a) - \mu_2(\mu_0, a) + \mu_2(a, \mu_0) = 0$ for all $a$ \quad (failure of $\mu_1^2 = 0$).
\end{itemize}

Rearranging the $n=1$ relation gives
\[
\mu_1^2(a)
=\mu_2(\mu_0,a)-\mu_2(a,\mu_0)
=[\mu_0,a]_{\mu_2}.
\]
No higher relation is needed for the two assertions of the theorem.
The $n=2$ relation is a curved Leibniz identity with
$\mu_0$-insertions in~$\mu_3$; after a curvature-killing MC twist, or
on an associated model in which those insertions vanish, it reduces to
the ordinary Leibniz rule.
\end{proof}

\subsection{I-adic completion: topology and convergence}
\label{sec:i-adic-completion}

\begin{definition}[I-adic topology]
\label{def:i-adic-topology}
\index{nilpotent completion}
Let $A$ be a curved $A_\infty$ algebra with augmentation ideal $I = \ker(\varepsilon: A \to \mathbb{C})$.

The \emph{I-adic completion} of $A$ is:
\begin{align}
\hat{A} &:= \varprojlim_n A/I^n \notag \\
 &= \bigl\{a \in {\textstyle\prod_{n \geq 0}} A/I^n :
 \text{compatible}\bigr\}.
\end{align}

An element $a \in \hat{A}$ can be written as a formal series:
\begin{equation}
a = a_0 + a_1 + a_2 + \cdots \quad \text{where } a_n \in I^n/I^{n+1}
\end{equation}
\end{definition}

\begin{theorem}[When completion is necessary {\cite{Positselski11,GLZ22}}; \ClaimStatusConditional]
\label{thm:completion-necessity}
Completion $A \to \hat{A}$ is necessary when:
\begin{enumerate}
\item \emph{Infinite sums}: Operations $\mu_n$ produce infinite sums not convergent in $A$
\item \emph{Failure of finite-type completion hypotheses}: the raw
 direct-sum bar object does not carry the finite-dimensional
 pro-conilpotent topology needed for continuous bar-cobar operations
\item \emph{Non-quadratic}: Relations involve infinitely many generators
\end{enumerate}

For example, the Virasoro algebra and $W_\infty$ need completion,
while Heisenberg and the finite-type Kac--Moody lanes do not.

The key datum is the filtration structure that guarantees the completed
bar-cobar round-trip remains a quasi-isomorphism. The answer is the strong
filtration axiom
$\mu_r(F^{i_1},\dots,F^{i_r}) \subset F^{i_1+\cdots+i_r}$
(Definition~\ref{def:strong-completion-tower}), which ensures that the
bar differential on each finite quotient $\cA_{\le N}$ is a
\emph{finite} sum (Lemma~\ref{lem:degree-cutoff}). Virasoro admits a
completed same-family shadow used throughout Parts~\ref{part:characteristic-datum}--\ref{part:standard-landscape};
$W_\infty$ is handled by the strong completion-tower theorem
(Theorem~\ref{thm:completed-bar-cobar-strong}), which supplies sharp
completion hypotheses.
\end{theorem}

\begin{proof}[Illustration: Virasoro finite windows]
The Virasoro algebra has generators $\{L_n\}_{n \in \mathbb{Z}}$ with:
\begin{equation}
[L_m, L_n] = (m-n)L_{m+n} + \frac{c}{12}m(m^2-1)\delta_{m+n,0}
\end{equation}

For a fixed bar word, the bracket is finite. For example,
$\omega=L_{-2}\otimes L_2$ gives
\begin{equation}
d(\omega) = [L_{-2}, L_2] = -4L_0 - \tfrac{c}{2}.
\end{equation}
The completion issue is therefore not deconcatenation
nonconilpotency of a finite word. It is the universal propagator and
continuous duality surface. The compatible family
\[
\xi_N=\sum_{1\leq |k|\leq N}s^{-1}L_{-k}\otimes s^{-1}L_k
\]
has finite image in every conformal-weight quotient but no value in
the discrete direct-sum bar complex. Its differential records the
central coefficients
$\tfrac{c}{12}k(k^2-1)$ and is again a compatible finite-window
family. Thus the operations are finite on each weight window, while
the inverse system itself lives in
\[
\widehat{\bar B}^{\mathrm{wt}}(\mathrm{Vir}_c)
=\varprojlim_N\bar B(\mathrm{Vir}_{c,\leq N}),
\]
not in the raw direct sum. The conformal-weight completion is exactly
the topology in which these compatible finite-window expressions,
their curvature corrections, and the continuous Verdier dual are
defined.
\end{proof}

\begin{remark}[Scope of the proof]\label{rem:completion-necessity-scope}
The Virasoro example above establishes necessity for retaining the
finite-window inverse system behind condition~\textup{(1)}; it does
not claim that the bar differential of a single finite Virasoro word is
an infinite sum.
Condition~(2) is demonstrated by the $\beta\gamma$-system: the
raw direct-sum presentation has infinite-rank graded pieces, so the
finite-type pro-conilpotent Positselski surface is not available
until the weight topology is retained
(see~\S\ref{subsec:betagamma-all-genera} in
Chapter~\ref{ch:genus-expansions}).
Condition~(3) is illustrated by affine $\widehat{\mathfrak{sl}}_2$ at non-integral level, where the Serre-type relations involve infinitely many generators (see~\S\ref{sec:three-theorems-sl2} in Chapter~\ref{ch:genus-expansions}).
\end{remark}

\begin{proposition}[Coacyclic-cone criterion for curved bar comparison in the
second-kind ambient; \ClaimStatusConditional]
\label{prop:curved-bar-acyclicity}
\index{coacyclicity!curved bar complex}
\index{Positselski!acyclicity}
Let $\mathcal{A}$ be a Koszul chiral algebra with nonzero obstruction
coefficient $\kappa(\mathcal{A}) \neq 0$. Suppose the genus-$g$ bar
data have been placed in an explicitly defined complete curved
CDG-coalgebra/factorization model whose curvature is
$h=m_0^{(g)}$ and whose coderived category is formed by Positselski's
second-kind totalizations. Assume, in addition, that the finite-window
bar-cobar comparison maps
\[
\epsilon_N\colon \Omega(C_N)\longrightarrow \mathcal{A}_{\leq N}
\]
have cones in the finite-stage coacyclic subcategories and that the
cone tower is strict Mittag--Leffler in the sense of
Definition~\ref{def:strict-ml-completed-chiral-coalgebra}. Then the
completed cone of
\[
\widehat\epsilon\colon
\widehat\Omega(\widehat C)\longrightarrow \widehat{\mathcal A}
\]
is coacyclic in the completed second-kind coderived category.
The assertion is about the comparison cone. Nonzero curvature alone
does not imply that the underlying curved bar object
$\bar{B}^{(g)}(\mathcal{A})$ is zero in the coderived category, and it
does not imply ordinary contractibility after forgetting curvature and
filtration.
\end{proposition}

\begin{proof}
At genus $g \geq 1$, the bar differential satisfies
$\dfib^{\,2} = \mcurv{g} \cdot \mathrm{id}$ where $\mcurv{g} = \kappa(\mathcal{A}) \cdot \lambda_g \neq 0$
(Theorem~\ref{thm:genus-universality}; Convention~\ref{conv:higher-genus-differentials}). The fiberwise differential $\dfib$ is curved, but the total corrected differential $\Dg{g}$ satisfies
$\Dg{g}^{\,2} = 0$. For the standard families the scalar curvature
$m_0^{(g)}=\kappa(\cA)\omega_g\mathbf 1$ is central for the binary
operation, so one must not infer coacyclicity from a nonzero
commutator $[m_0,-]_{m_2}$; that commutator may vanish. The curvature
acts through the second-kind CDG totalization and the completed
factorization structure.

Positselski's second-kind formalism~\cite[\S\S3--4]{Positselski11}
defines weak equivalences by coacyclic cones inside the exact
category of curved comodules. By hypothesis each finite window cone
\[
K_N:=\operatorname{Cone}(\epsilon_N)
\]
belongs to the finite-stage coacyclic subcategory. The transition maps
are strict and the cone tower is Mittag--Leffler, so products,
finite-stage totalizations, and the admitted direct sums used to
generate coacyclic objects commute with the inverse limit. Therefore
\[
\widehat K:=\varprojlim_N K_N
=\operatorname{Cone}(\widehat\epsilon)
\]
lies in the completed coacyclic subcategory. This proves the curved
bar-cobar comparison in the coderived ambient.

The proof deliberately uses the cone tower. A scalar class
$\mcurv{g}=\kappa(\mathcal A)\lambda_g$ may be nonzero without being
invertible and without giving a contracting homotopy for the bar
object itself. Thus the theorem supplies the strongest true
comparison available in this ambient: coacyclicity of the comparison
cone under explicit finite-window and completion hypotheses.
\end{proof}

\begin{remark}[Necessity of exotic derived categories at higher genus]
\label{rem:positselski-acyclicity}
At genus $g \geq 1$ with $\kappa \neq 0$, coacyclicity of
$\operatorname{Cone}(\widehat\epsilon)$ is a statement about the
bar-cobar comparison map, not about the disappearance of the bar object
itself. The raw curved bar coalgebra is still the object over which the
associated comodule homological algebra is formed. The correct
framework for that algebra is Positselski's coderived category
$D^{\mathrm{co}}(\bar{B}^{(g)}\text{-}\mathrm{CoMod})$~\cite[\S\S3--4]{Positselski11},
which is why Theorem~\ref{thm:positselski-chiral} requires the
coderived/contraderived setting.
\end{remark}

\begin{remark}[Derived--coderived reduction for chiral CDG-coalgebras]
\label{rem:derived-coderived-chiral-CDG}
\index{coderived category!derived reduction}
\index{Positselski!derived--coderived reduction}
The chiral bar complex $C = \bar{B}^{\mathrm{ch}}(\cA)$ is a CDG-coalgebra
with curvature $h = m_0^{(g)}$. Three hypotheses govern the relationship
between its coderived and ordinary derived categories:
\begin{enumerate}[label=\textup{(H\arabic*)}]
\item \emph{Conilpotency}: $C$ is conilpotent, i.e., the coaugmentation
 filtration is exhaustive. This holds by
 Theorem~\ref{thm:coalgebra-via-NAP}(4): the conformal weight grading
 on $\bar{B}^{\mathrm{ch}}(\cA)$ provides an exhaustive
 $\mathbb{N}$-filtration with $C = \bigcup_n F^n C$.
\item \emph{Finite-type weight spaces}: each graded piece of $C$ is
 finite-dimensional. This holds because $\cA$ is finitely generated
 with finite-dimensional conformal weight spaces (the positive-energy
 axiom), so $\bar{B}^n(\cA)$ in each weight is a finite-dimensional
 vector space.
\item \emph{Compact generation}: (H1) and (H2) together imply that the
 coderived category $D^{\mathrm{co}}(C\text{-}\mathrm{comod}^{\mathrm{ch}})$
 is compactly generated by finite-dimensional comodules.
\end{enumerate}

\smallskip\noindent
\emph{Genus $0$ \textup{(}$h = 0$\textup{)}.}
By the conilpotent reduction theorem of Positselski~\cite[Corollary~4.4]{Positselski11}
(cf.\ Remark~\ref{rem:exotic-vs-ordinary}),
(H1) and $h = 0$ imply that every acyclic CDG-comodule is coacyclic, so
\[
D^{\mathrm{co}}(C\text{-}\mathrm{comod}^{\mathrm{ch}})
\;\simeq\;
D(C\text{-}\mathrm{comod}^{\mathrm{ch}}).
\]
The coderived category reduces to the ordinary derived category.
This is the finite strict-ML regime of Theorem~B; the displayed
bar-cobar inversion is the counit face of Theorem~A on the Koszul locus.

\smallskip\noindent
\emph{Genus $g \geq 1$ with $\kappa(\cA) \neq 0$
\textup{(}$h \neq 0$\textup{)}.}
The raw higher-genus fiber operator satisfies
$\dfib^{\,2} = \mcurv{g} \cdot \mathrm{id}$ and therefore lives on the
curved CDG surface. Proposition~\ref{prop:curved-bar-acyclicity}
concerns the second-kind curved totalization: the relevant cone is
coacyclic after the curved coalgebra and its completed exact ambient
have been fixed. It does not assert that the period-corrected
square-zero complex $(\barB^{(g)}(\cA),\Dg{g})$ has zero ordinary
cohomology; that ordinary cohomology is the corrected genus-$g$
invariant. The
coderived--contraderived equivalence
(Theorem~\ref{thm:positselski-chiral-proved}) is nontrivial.
The Positselski equivalence
$D^{\mathrm{co}}(\barB^{\mathrm{ch}}(\cA)\text{-}\mathrm{comod}) \simeq
D^{\mathrm{ctr}}(\barB^{\mathrm{ch}}(\cA)\text{-}\mathrm{contra})$
of Theorem~\ref{thm:positselski-chiral-proved} provides the
correct framework. Its proof uses (H1)--(H3) at three
points: conilpotency for cofree resolutions, finite-type for the
duality between contramodules and complete modules over the graded
dual bar-dual algebra, and compact generation for the triangulated
comparison functor.
See Remark~\ref{rem:curvature-coderived} for the connection to
bar-cobar duality and
\S\ref{subsec:chiral-coderived-contraderived} for the chiral
coderived and contraderived categories.

\smallskip\noindent
\emph{Summary.} Each finite-window bar complex of a positive-energy
Koszul chiral algebra satisfies (H1)--(H3); class-$\mathsf M$ raw
direct sums acquire this surface only after weight completion. At
genus $0$, the uncurved conilpotent finite-type lane yields
$D^{\mathrm{co}} \simeq D$ (ordinary derived category). At genus
$g \geq 1$, the completed curved lane yields
$D^{\mathrm{co}} \simeq D^{\mathrm{ctr}}$ (comodule--contramodule
correspondence). The three models of the genus-$g$ bar complex
(flat associated graded in $D$; corrected holomorphic in $D$; curved
geometric in $D^{\mathrm{co}}$) are linked by these identifications.
\end{remark}

\subsection{Filtered versus curved structures}
\label{sec:filtered-vs-curved}

\begin{theorem}[Filtered cooperads (Gui--Li--Zeng~\cite{GLZ22}); \ClaimStatusProvedElsewhere]
\label{thm:filtered-cooperads}
A \emph{filtered cooperad} $\mathcal{C}$ is more general than a curved cooperad:
\begin{equation}
\mathcal{C} = \bigcup_{n=0}^\infty F^n \mathcal{C}
\end{equation}
where $F^n \mathcal{C} \subset F^{n+1} \mathcal{C}$ is an increasing filtration, with:
\begin{enumerate}
\item Comultiplication: $\Delta(F^n) \subset \bigoplus_{i+j=n} F^i \otimes F^j$
\item Counit: $\varepsilon(F^{>0}) = 0$
\end{enumerate}

The filtered structure does not reduce to a single curvature element.
\end{theorem}

\begin{remark}[Provenance and citation]
This theorem is imported and treated as \ClaimStatusProvedElsewhere. The
filtered-cooperad framework used here is sourced from \cite{GLZ22}; subsequent
internal reductions in this manuscript are organized by
Theorem~\ref{thm:filtered-to-curved} and Theorem~\ref{thm:bar-convergence}.
\end{remark}

\begin{example}[W-algebras: filtered but not simply curved]
\label{ex:w-algebra-filtered-comprehensive}
For the $\mathcal{W}_3$ algebra, the filtration is by conformal weight:
\begin{equation}
F^n \mathcal{W}_3 = \text{span}\{T, W, \partial T, TT, \partial W, \ldots \text{ up to
weight } n\}
\end{equation}

This does not come from a single curvature element $\mu_0$. Instead, there are curvature
contributions at each weight:
\begin{align}
\mu_0^{(2)} &= T \quad \text{(weight 2)}\\
\mu_0^{(3)} &= W \quad \text{(weight 3)}\\
\mu_0^{(4)} &= TT + \text{regular} \quad \text{(weight 4)}\\
&\vdots
\end{align}

The full structure requires the complete filtered cooperad.
\end{example}

\begin{theorem}[When filtered reduces to curved; \ClaimStatusProvedHere]
\label{thm:filtered-to-curved}
Let $\mathcal{C}$ be a filtered chiral cooperad with exhaustive, separated, complete
filtration $F_\bullet\mathcal{C}$. Assume:
\begin{enumerate}
\item $\dim(F_n\mathcal{C}/F_{n-1}\mathcal{C})<\infty$ for every $n$;
\item $\gr(\mathcal{C})$ is generated in degree $1$ with quadratic relations;
\item higher relation classes differ from quadratic consequences by central
filtration-$\ge2$ corrections.
\end{enumerate}
Then there exists a curved cooperad $\mathcal{C}_{\mathrm{curv}}$ with curvature
$\mu_0\in F_2\mathcal{C}_{\mathrm{curv}}$ and a filtered quasi-isomorphism
\[
\mathcal{C}_{\mathrm{curv}}\xrightarrow{\;\simeq_f\;}\mathcal{C}.
\]
If $\gr(\mathcal{C})$ has infinitely many nonzero graded pieces, this curved
description is only defined after completion.
\end{theorem}

\begin{proof}
Set $\mathcal{A}:=\mathcal{C}^\vee$ (continuous dual with respect to
$F_\bullet$). By assumption (1), filtration pieces are finite-dimensional, so
duality exchanges filtered cooperads and filtered chiral algebras without loss
of filtration data.

Assumptions (2)--(3) are exactly the hypotheses of
Proposition~\ref{prop:filtered-to-curved-fc} for $\mathcal{A}$. Hence there is a
curved model $\mathcal{A}_{\mathrm{curv}}$ with curvature
$\mu_0\in F_2\mathcal{A}_{\mathrm{curv}}$ and a filtered quasi-isomorphism
$\mathcal{A}_{\mathrm{curv}}\to\mathcal{A}$.

Dualizing back gives a curved cooperad
$\mathcal{C}_{\mathrm{curv}}:=\mathcal{A}_{\mathrm{curv}}^\vee$ and a filtered
quasi-isomorphism $\mathcal{C}_{\mathrm{curv}}\to\mathcal{C}$. The curved identity
$d^2=[\mu_0,-]$ is preserved under this dualization because all filtration
quotients are finite-dimensional.

If $\gr(\mathcal{C})$ has infinitely many nonzero graded pieces, the curvature
decomposes into infinitely many filtration components. Then the bar/cobar
operations are no longer finite strict and convergence requires completion,
as in Theorem~\ref{thm:conilpotency-convergence} and
Theorem~\ref{thm:bar-convergence}.
\end{proof}

\begin{remark}[Dependencies]
This proof uses:
\begin{enumerate}[label=(D\arabic*)]
\item the filtered algebra-side reduction in
Proposition~\ref{prop:filtered-to-curved-fc};
\item finite-type filtered duality between filtered algebras and cooperads;
\item convergence control from Theorem~\ref{thm:conilpotency-convergence} and
Theorem~\ref{thm:bar-convergence}.
\end{enumerate}
\end{remark}

\subsection{Conilpotency and convergence without completion}
\label{sec:conilpotency-convergence}

\begin{definition}[Conilpotent coalgebra]
\label{def:conilpotent-complete}
\index{conilpotent!coalgebra}
\index{conilpotent!filtration}
A coalgebra $C$ is \emph{conilpotent} if for each $c \in C$, there exists $N$ such that:
\begin{equation}
\Delta^{(N)}(c) = 0
\end{equation}
where $\Delta^{(N)}$ is the $N$-fold iterated comultiplication.
\end{definition}

\begin{theorem}[Conilpotency controls algebraic convergence; \ClaimStatusProvedHere]
\label{thm:conilpotency-convergence}
\label{thm:conilpotency-bar}
\label{thm:koszul-conilpotent}
\index{bar construction!convergence}
Let $C=\bar{B}(A)$. Suppose either that $C$ is an uncurved
conilpotent dg coalgebra in the discrete finite-type lane, or that
the curved data have first been replaced by a square-zero total model
whose finite output windows receive only finitely many input
contributions. Then:
\begin{enumerate}
\item all algebraic sums in the cobar differential terminate on each
 element;
\item the bar-cobar composition $\Omega \circ \bar{B}(A) \to A$ is
 defined without completion in the discrete finite-type lane;
\item The bar, cobar, and Verdier-dual objects entering Koszul
 duality are defined without taking $\hat{A}$ only on that
 finite-type square-zero surface.
\end{enumerate}
For class-$\mathsf M$ raw direct-sum presentations, finite bar-length
conilpotence gives the first conclusion on each finite word but does
not supply the finite-type Positselski surface, the infinite OPE mode
sums, or the completed bar-cobar comparison. Those require the
weight-completed tower of
Theorem~\ref{thm:chiral-positselski-weight-completed}.
\end{theorem}

\begin{proof}
Write $C = \bar{B}(A)$ and $\bar{C} = \ker(\varepsilon)$ for the augmentation
ideal. The cobar construction is the tensor algebra
$\Omega(C) = T(s^{-1}\bar{C}) = \bigoplus_{k \ge 0} (s^{-1}\bar{C})^{\otimes k}$
equipped with the differential $d_\Omega$ determined on generators by the reduced
comultiplication $\bar{\Delta}\colon \bar{C} \to \bar{C} \otimes \bar{C}$:
\begin{equation}\label{eq:cobar-diff-generators}
d_\Omega(s^{-1}c) \;=\; -\sum (-1)^{|c'|}\, s^{-1}c' \otimes s^{-1}c''
\end{equation}
where $\bar{\Delta}(c) = \sum c' \otimes c''$ (Sweedler notation), extended to all
of $T(s^{-1}\bar{C})$ as a derivation:
\begin{equation}\label{eq:cobar-diff-derivation}
d_\Omega(s^{-1}c_1 \otimes \cdots \otimes s^{-1}c_k)
\;=\; \sum_{i=1}^{k} (-1)^{\epsilon_i}\,
s^{-1}c_1 \otimes \cdots \otimes d_\Omega(s^{-1}c_i) \otimes \cdots \otimes s^{-1}c_k
\end{equation}
with signs $\epsilon_i = |s^{-1}c_1| + \cdots + |s^{-1}c_{i-1}|$.

\emph{Step 1 (conilpotent case).}
Suppose $C$ is conilpotent (Definition~\ref{def:conilpotent-complete}): for each
$c \in \bar{C}$, there exists $N = N(c)$ such that the iterated reduced
comultiplication $\bar{\Delta}^{(N)}(c) = 0$. Then $\bar{\Delta}(c)$ is a
\emph{finite} sum in $\bar{C} \otimes \bar{C}$, so~\eqref{eq:cobar-diff-generators}
is a finite sum. On a general element
$\xi = s^{-1}c_1 \otimes \cdots \otimes s^{-1}c_k$,
the derivation rule~\eqref{eq:cobar-diff-derivation} applies the comultiplication to
each factor $c_i$ individually; conilpotency of each~$c_i$ ensures each summand is
finite. The outer sum has exactly $k$ terms. Therefore $d_\Omega(\xi)$ is a finite
sum in $T(s^{-1}\bar{C})$, and no completion or topology is needed.

\emph{Step 2 ($d_\Omega^2 = 0$ in the uncurved or square-zero total
case).}
Since $d_\Omega$ is a derivation of the free algebra $T(s^{-1}\bar{C})$, the
Leibniz rule gives $d_\Omega^2(x \otimes y) = d_\Omega^2(x) \otimes y + x \otimes
d_\Omega^2(y)$, so it suffices to verify $d_\Omega^2 = 0$ on generators
$s^{-1}c$. Write $\bar{\Delta}(c) = \sum c' \otimes c''$ in Sweedler notation.
Then by~\eqref{eq:cobar-diff-generators}:
\begin{align*}
d_\Omega^2(s^{-1}c)
&= d_\Omega\Bigl(-\sum (-1)^{|c'|}\,s^{-1}c' \otimes s^{-1}c''\Bigr) \\
&= -\sum (-1)^{|c'|}\bigl(d_\Omega(s^{-1}c') \otimes s^{-1}c''
 + (-1)^{|s^{-1}c'|}\,s^{-1}c' \otimes d_\Omega(s^{-1}c'')\bigr).
\end{align*}
The first group of terms applies $\bar{\Delta}$ to the left tensorand:
writing $\bar{\Delta}(c') = \sum c'_{(1)} \otimes c'_{(2)}$, it produces
$+\sum \pm\, s^{-1}c'_{(1)} \otimes s^{-1}c'_{(2)} \otimes s^{-1}c''$.
The second group applies $\bar{\Delta}$ to the right tensorand:
writing $\bar{\Delta}(c'') = \sum c''_{(1)} \otimes c''_{(2)}$, it produces
$+\sum \pm\, s^{-1}c' \otimes s^{-1}c''_{(1)} \otimes s^{-1}c''_{(2)}$.
Coassociativity of~$\bar{\Delta}$ states
$(\bar{\Delta} \otimes \mathrm{id}) \circ \bar{\Delta}
= (\mathrm{id} \otimes \bar{\Delta}) \circ \bar{\Delta}$,
so these two sums are identical up to sign. Tracking the Koszul signs from the
desuspension (each $s^{-1}$ shifts degree by $-1$), the two contributions cancel
exactly: the sign discrepancy between applying $\bar{\Delta}$ on the left versus
the right factor is $(-1)^{|s^{-1}c'|} = (-1)^{|c'|-1}$, and this combines with
the outer sign $(-1)^{|c'|}$ to produce opposite signs on the two coassociativity
terms. Therefore $d_\Omega^2(s^{-1}c) = 0$. All sums involved are finite by
Step~1 (cf.~also \cite{LV12}, \S2.2).

\emph{Step 3 (bar-cobar composition).}
The bar-cobar composition $\Omega(\bar{B}(A)) = T(s^{-1}\overline{\bar{B}(A)})$
inherits conilpotency from $\bar{B}(A)$: every element of $\bar{B}^n(A)$ lies in the
tensor algebra on $A$, and the comultiplication is the deconcatenation coproduct,
which satisfies $\bar{\Delta}^{(n)}(\omega) = 0$ for $\omega \in \bar{B}^n(A)$.
All cobar operations therefore terminate after finitely many steps, and the
counit $\varepsilon\colon \Omega(\bar{B}(A)) \to A$ is well-defined without
completion.

\emph{Step 4 (completed case).}
If $C$ is not conilpotent but is \emph{cocomplete} (pro-conilpotent), i.e., the
natural map $C \xrightarrow{\;\sim\;} \varprojlim C/F_N C$ is an isomorphism,
where $F_N C = \ker(\bar{\Delta}^{(N)})$ is the coradical filtration, then elements
of~$C$ are coherent sequences $(c_N)_{N \ge 0}$ with $c_N \in C/F_N C$ and the
canonical projections $\pi_{N+1,N}\colon C/F_{N+1}C \to C/F_NC$ satisfying
$\pi_{N+1,N}(c_{N+1}) = c_N$. The completed cobar construction is
$\hat{\Omega}(C) = \varprojlim_N \Omega(C/F_N C)$. Each quotient $C/F_N C$ is
conilpotent by construction ($\bar{\Delta}^{(N)} = 0$ on $C/F_NC$), so the cobar
differential $d_\Omega$ is well-defined on each $\Omega(C/F_N C)$ by Steps~1--2.
Since $d_\Omega$ respects the coradical filtration: if
$c \in F_N C$, then $\bar{\Delta}(c) \in \sum_{i+j=N} F_i C \otimes F_j C$
(since the coradical filtration is a coalgebra filtration), so the differentials on
$\Omega(C/F_N C)$ form a compatible system; the projection
$\pi_{N+1,N}$ intertwines $d_\Omega$ at level $N+1$ with $d_\Omega$ at level~$N$.
The completed differential $d_{\hat{\Omega}}$ is therefore well-defined on the
inverse limit: for an element $(c_N)_{N \ge 0} \in \hat{\Omega}(C)$, the image
$d_{\hat{\Omega}}((c_N)) = (d_\Omega(c_N))_{N \ge 0}$ is a coherent sequence
because each $d_\Omega(c_N)$ is well-defined (Step~1) and the compatibility
$\pi_{N+1,N}(d_\Omega(c_{N+1})) = d_\Omega(\pi_{N+1,N}(c_{N+1})) =
d_\Omega(c_N)$ holds. The identity $d_{\hat{\Omega}}^2 = 0$ follows from
$d_\Omega^2 = 0$ at each finite level (Step~2), since a compatible system of
zeros has zero inverse limit.

This proves all three conclusions in the stated finite-type
square-zero ambient: the cobar differential terminates on each element,
the bar-cobar composition is defined without completion, and the
objects used in Koszul duality are defined via the direct
(uncompleted) bar-cobar adjunction. It proves only algebraic
bar-length conilpotence for a raw class-$\mathsf M$ direct sum. The
topological comparison for class~$\mathsf M$ is a separate completed
statement: the finite-window tower must satisfy the continuous
Mittag--Leffler hypotheses before the coderived/comodule-detected
comparison is available.
\end{proof}

\begin{example}[Heisenberg: conilpotent]
\label{ex:heisenberg-conilpotent-complete}
The Heisenberg algebra $\mathcal{H}_\kappa$ has bar complex:
\begin{equation}
\bar{B}^n(\mathcal{H}_\kappa) = \mathcal{H}_\kappa^{\otimes n} \otimes \Omega^n
\end{equation}

For $\omega = a_{n_1} \otimes \cdots \otimes a_{n_k} \otimes \eta_{ij}$, the
comultiplication is:
\begin{equation}
\Delta(\omega) = \sum_{\text{splittings}} \omega_L \otimes \omega_R
\end{equation}

After $k$ iterations, $\Delta^{(k)}(\omega) = 0$ because we run out of tensor factors.
Thus $\bar{B}(\mathcal{H}_\kappa)$ is conilpotent, and no completion is needed.
\end{example}

\begin{example}[Virasoro: raw direct-sum topology is not the finite-type Positselski surface]
\label{ex:virasoro-not-conilpotent}
The cofree tensor bar coalgebra on the Virasoro state space is
conilpotent by bar length on every finite word. The obstruction is
different: the raw direct-sum topology does not contain the infinite
mode sums produced by the OPE and by the harmonic curvature
corrections. For instance the compatible finite-weight expressions
\[
\xi_N=\sum_{\lvert k\rvert\leq N}
 s^{-1}L_{-k}\otimes s^{-1}L_k
\]
define a Cauchy family in the conformal-weight completion, but not an
element of the discrete direct-sum bar complex. The same phenomenon
appears in the central terms of
$[L_m,L_{-m}]$ and in the higher harmonic corrections of the
genus-$g$ curved differential. Hence the raw direct-sum presentation
is not the finite-type pro-conilpotent Positselski surface; the
weight-completed/pro object is required. The obstruction is
finite-type completion, not deconcatenation nonconilpotency of~$L_0$.
\end{example}

\subsection{Examples: computing Koszul duals with completion}
\label{sec:koszul-duals-completion-examples}

\begin{example}[Virasoro Koszul dual]
\label{ex:virasoro-koszul-dual}

Since the raw direct-sum $\bar{B}(\mathrm{Vir})$ is not the
finite-type pro-conilpotent Positselski surface
(Example~\ref{ex:virasoro-not-conilpotent}), one forms the
completed bar complex and takes continuous linear dual.
Writing $\mathrm{Vir}_c = \mathcal{W}^k(\mathfrak{sl}_2)$
with $h^\vee = 2$, level-shifting duality
(Corollary~\ref{cor:level-shifting-part1}) gives
$k' = -k - 2h^\vee = -k-4$, so
\begin{equation}
\mathrm{Vir}_c^! \cong \mathrm{Vir}_{26-c},
\end{equation}
since $c(k) + c(k') = 2\operatorname{rank}(\mathfrak{sl}_2) + 4h^\vee\dim\mathfrak{sl}_2 = 2 + 24 = 26$ by Theorem~\ref{thm:central-charge-complementarity}(b) (the $\mathcal{W}$-algebra central-charge sum; this is the Virasoro-family formula, not the affine KM formula $c + c' = 2\dim\mathfrak{g}$). Physically, matter--ghost duality pairs a matter system at central charge~$c$ with ghosts at $c_{\mathrm{ghost}} = 26 - c$. In particular, $\mathrm{Vir}_{26}^! \cong \mathrm{Vir}_0$, consistent with bosonic string BRST cohomology at ghost number zero.

At the M/S level, the completed dual has infinitely many primitive cumulants ($\Delta_{\mathrm{Vir}}(t) = t^3 + 2t^5 + \cdots$; Definition~\ref{def:primitive-defect-series}). The MC4 structural framework is proved (Theorem~\ref{thm:completed-bar-cobar-strong}); the residual H-level target identification is the subject of Example~\ref{ex:winfty-completion-frontier}.
\end{example}

The infinite primitive-cumulant tail in
Example~\ref{ex:virasoro-koszul-dual} is the genus-$0$ shadow of
a more concrete analytic object: the spectral R-matrix of
$\mathrm{Vir}_c$ acting on a primary state. On a rank-$1$
primary line the path-ordered exponential of the collision
residue closes into an elementary function.

\begin{computation}[Virasoro spectral R-matrix on primary states;
\ClaimStatusProvedHere]
\label{comp:virasoro-spectral-r-matrix}
\index{Virasoro spectral R-matrix!closed form|textbf}
\index{R-matrix!Virasoro closed form}
For the Virasoro algebra at central charge $c$ acting on a
primary state of conformal weight $h$, the spectral R-matrix
admits the closed form
\begin{equation}
\label{eq:virasoro-r-matrix-closed}
R(z) \;=\; z^{2h}\,\exp\!\left(-\frac{c}{4\,z^{2}}\right)
 \;=\; z^{2h}\,\sum_{k=0}^{\infty}
 \frac{(-c/4)^{k}}{k!\,z^{2k}}.
\end{equation}
The first four series coefficients are
\[
 R_0 \;=\; z^{2h},
 \qquad
 R_2 \;=\; -\frac{c}{4},
 \qquad
 R_4 \;=\; \frac{c^{2}}{32},
 \qquad
 R_6 \;=\; -\frac{c^{3}}{384}.
\]
Only even powers of $1/z$ appear, reflecting bosonic parity.
The non-terminating character of this series is the genus-$0$
signature of class~$M$ (infinite shadow depth): unlike
Heisenberg ($R(z) = z^{k}$, terminating after the leading
power) or affine Kac--Moody
($R(z) = z^{\Omega/(k+h^{\vee})}$, terminating), the Virasoro
R-matrix has infinitely many corrections in $1/z^{2}$.
\end{computation}

\begin{proof}
The Virasoro OPE is
\[
 T(z)T(w) \;\sim\; \frac{c/2}{(z-w)^{4}}
 \;+\; \frac{2\,T(w)}{(z-w)^{2}}
 \;+\; \frac{\partial T(w)}{z-w}.
\]
By the $d\log$ absorption principle, the collision residue has
poles one order lower than the OPE:
\[
 r^{\mathrm{coll}}(z) \;=\; \frac{c/2}{z^{3}} \;+\; \frac{2\,T}{z}.
\]
On a primary state $\lvert h\rangle$ the modes act as
$T_{(1)}\lvert h\rangle = h\,\lvert h\rangle$ (by
$L_{0}\lvert h\rangle = h\lvert h\rangle$), while $T_{(3)}$
acts as the central scalar $c/2$. The path-ordered exponential
\[
 R(z) \;=\; \mathcal{P}\exp\!\left(\oint\;\sum_{n\geq 0}
 r_{n}\,z^{-n-1}\,dz\right)
\]
on the primary sector is a scalar path-ordered exponential of
commuting generators ($\mathbbm{1}$ and $L_{0}$), hence reduces
to the ordinary exponential
\[
 R(z)\,\lvert h\rangle
 \;=\; \exp\!\left(2h\,\log z \;-\; \frac{c/2}{2\,z^{2}}\right)
 \lvert h\rangle
 \;=\; z^{2h}\,\exp\!\left(-\frac{c}{4\,z^{2}}\right)
 \lvert h\rangle.
\]
The $z^{2h}$ factor comes from integrating $2h/z$ along the
contour; the exponential of the central term produces
$\exp(-c/(4z^{2}))$. Expanding the exponential yields
\eqref{eq:virasoro-r-matrix-closed}.
\end{proof}

\begin{remark}[Independent checks]
\label{rem:virasoro-r-matrix-paths}
The closed form
$R(z) = z^{2h}\,\exp(-c/(4z^{2}))$ is checked by $91$ tests in
\texttt{compute/lib/theorem\_virasoro\_spectral\_r\_matrix\_engine.py},
including:
\begin{itemize}
\item direct OPE mode computation;
\item path-ordered exponential numerical integration;
\item series coefficient checks at $c = 1, 13, 25, 26$;
\item comparison with the BPZ connection on conformal blocks;
\item bosonic parity (only even $1/z$ powers).
\end{itemize}
The $R_{2} = -c/4$ first correction is the genus-$0$
non-formality witness: it is the coefficient from which the
infinite primitive-cumulant series
$\Delta_{\mathrm{Vir}}(t) = t^{3} + 2t^{5} + \cdots$ of
Example~\ref{ex:virasoro-koszul-dual} descends via the shadow
obstruction tower.
\end{remark}

The Virasoro example exposes the shape of the general completion
problem: finite-type truncations $W_N$ lie in the proved Koszul
regime, and the challenge is to assemble these finite-stage
quasi-isomorphisms into a completed quasi-isomorphism for the full
infinite-type algebra in the square-zero total model, or into a
coacyclic-cone comparison in the curved CDG ambient. The obstacle is
analytic (convergence of the inverse system) together with the choice
of square-zero versus coderived ambient, and it admits a clean
algebraic resolution once the correct filtration axiom is identified.

\begin{example}[\texorpdfstring{$W_\infty$}{W-infinity}: MC4 completion criterion]
\label{ex:winfty-completion-frontier}

The tower $W_\infty = \varprojlim_N W_N$ is the archetype of the
programmatic regime: each finite stage~$W_N$ is principal finite
type with proved bar-cobar quasi-isomorphism, but infinitely many
generators prevent a direct application of the finite-stage
theorem. The strong completion-tower theorem
(Theorem~\ref{thm:completed-bar-cobar-strong}) resolves this for
models satisfying the strong filtration axiom:
the strong filtration axiom
$\mu_r(F^{i_1},\dots,F^{i_r}) \subset F^{i_1+\cdots+i_r}$
makes the inverse-limit differential continuous and the
Mittag--Leffler condition automatic, so the completed bar-cobar
round-trip is a quasi-isomorphism in the square-zero total model and
a coacyclic-cone equivalence in the curved CDG ambient.
For the standard spin-truncation tower, which need not satisfy this
strong axiom on the algebra itself, the replacement input is the
bar-level weight-cutoff criterion of
Proposition~\ref{prop:mc4-weight-cutoff}.

The H-level comparison requires:
\begin{enumerate}
\item identify the separated complete H-level/factorization target whose finite quotients recover $W_N$;
\item verify the exact residue identities of Proposition~\ref{prop:winfty-ds-residue-identity-criterion};
\item prove finite detection on $\mathcal{I}_N$, after which Corollary~\ref{cor:winfty-hlevel-comparison-criterion} makes the comparison formal.
\end{enumerate}
\end{example}

At each finite stage~$N$, the bar-cobar adjunction
(Theorem~\ref{thm:bar-cobar-isomorphism-main}) is mediated by the
universal twisting morphism
$\tau_N \colon \bar B^{\mathrm{ch}}(\cA_{\le N}) \to \cA_{\le N}$,
which is a Maurer--Cartan element in the convolution algebra
$\operatorname{Conv}(\bar B^{\mathrm{ch}}(\cA_{\le N}),
\cA_{\le N})$
(Definition~\ref{def:twisting-morphism}).
The MC4 completion problem asks: \emph{when do these finite-stage
MC elements assemble to an MC element
$\widehat\tau \in
\MC\bigl(\widehat{\operatorname{Conv}}(\widehat{\bar B}^{\mathrm{ch}}(\cA),
\cA)\bigr)$
in the completed convolution algebra?}
For a strong completion tower, a single axiom on the filtration (that
the $A_\infty$~operations respect the additive weight) makes
the MC equation $\partial(\widehat\tau) + \widehat\tau \star
\widehat\tau = 0$ converge in the inverse limit. Each component
of this equation involves only finitely many terms
(Lemma~\ref{lem:degree-cutoff}); the passage to the completion is then
quotientwise finite.

\begin{definition}[Strong completion tower]
\label{def:strong-completion-tower}
\index{strong completion tower|textbf}
\index{completion closure!strong completion tower}
An augmented curved chiral $\Ainf$-algebra $\cA$ on a curve $X$ is a
\emph{strong completion tower} if it carries a descending filtration
\[
\cA = F^0\cA \supset F^1\cA \supset F^2\cA \supset \cdots,
\qquad
\textstyle\bigcap_{N \ge 0} F^{N+1}\cA = 0,
\]
such that:
\begin{enumerate}
\item $\cA$ is separated and complete:
 $\cA \cong \varprojlim_N \cA_{\le N}$, where
 $\cA_{\le N} := \cA/F^{N+1}\cA$;
\item every quotient $\cA_{\le N}$ is finite type and lies in the proved
 bar-cobar regime;
\item $\bar\cA = F^1\cA$ (the augmentation ideal is the first filtration
 piece); and
\item all chiral $\Ainf$-operations are filtration-nondecreasing:
 \begin{equation}\label{eq:strong-filtration-condition}
 \mu_r(F^{i_1}\cA, \dots, F^{i_r}\cA)
 \subset F^{i_1+\cdots+i_r}\cA.
 \end{equation}
\end{enumerate}
In particular, each projection $p_N\colon \cA_{\le N+1} \twoheadrightarrow
\cA_{\le N}$ is a strict morphism of curved chiral $\Ainf$-algebras,
and the finite-stage bar constructions are compatible with the tower.
\end{definition}

\begin{lemma}[Degree cutoff: finite MC equation at each stage;
\ClaimStatusProvedHere]
\label{lem:degree-cutoff}
\index{degree cutoff}
For a strong completion tower, the MC equation
$\partial(\tau_N) + \tau_N \star \tau_N = 0$
in $\operatorname{Conv}(\bar B^{\mathrm{ch}}(\cA_{\le N}),
\cA_{\le N})$ involves only degrees $r \le N$:
the bar differential on $\cA_{\le N}$ is a finite sum.
\end{lemma}

\begin{proof}
If all inputs lie in $\bar\cA = F^1\cA$, then
$\mu_r(\bar\cA^{\otimes r}) \subset F^r\cA$ by
\eqref{eq:strong-filtration-condition}. Modulo $F^{N+1}\cA$, any term
with $r \ge N+1$ vanishes. So on $\cA_{\le N}$ the bar differential is
the finite sum $b_1 + \cdots + b_N$ of coderivations.
\end{proof}

\begin{theorem}[MC element lifts to the completed convolution algebra]
\label{thm:completed-bar-cobar-strong}
\index{completed bar-cobar!strong completion tower}
\index{MC4!completion closure theorem}
\index{twisting morphism!completed}
Let $\cA$ be a strong completion tower
\textup{(}Definition~\textup{\ref{def:strong-completion-tower})}.
Write $\tau_N \in \MC\bigl(\operatorname{Conv}(\bar B^{\mathrm{ch}}
(\cA_{\le N}), \cA_{\le N})\bigr)$ for the universal twisting
morphism at stage~$N$. The quasi-isomorphism assertions below are
ordinary chain assertions only for the square-zero total differential
model. In the genuinely curved CDG model, the same finite-stage
construction gives a coacyclic-cone statement after the second-kind
coacyclicity hypotheses are satisfied. Then:
\begin{enumerate}
\item \emph{Completed coalgebra.}
 The completed bar construction
 $\widehat{\bar B}^{\mathrm{ch}}(\cA)
 := \varprojlim_N \bar B^{\mathrm{ch}}(\cA_{\le N})$
 exists as a separated complete pronilpotent curved dg chiral
 coalgebra with continuous differential.
 It is the domain of the completed twisting morphism
 $\widehat\tau\colon \widehat{\bar B}^{\mathrm{ch}}(\cA) \to \cA$.
 At each stage,
 $\widehat{\bar B}^{\mathrm{ch}}(\cA)/F^{N+1}
 \cong \bar B^{\mathrm{ch}}(\cA_{\le N})$.
\item \emph{Completed algebra.}
 The completed cobar object
 $\widehat\Omega^{\mathrm{ch}}\bigl(
 \widehat{\bar B}^{\mathrm{ch}}(\cA)\bigr)
 := \varprojlim_N \Omega^{\mathrm{ch}}(
 \bar B^{\mathrm{ch}}(\cA_{\le N}))$
 is the twisted tensor product
 $\cA \otimes_{\widehat\tau}
 \widehat{\bar B}^{\mathrm{ch}}(\cA)$:
 a separated complete curved dg chiral algebra.
\item \emph{MC element is acyclic in the square-zero total model.}
 The completed twisting morphism $\widehat\tau$ is acyclic:
 the counit
 \[
 \widehat\epsilon \colon
 \widehat\Omega^{\mathrm{ch}}\bigl(
 \widehat{\bar B}^{\mathrm{ch}}(\cA)\bigr)
 \longrightarrow \cA
 \]
 is a quasi-isomorphism when both sides carry the square-zero total
 differential. At each quotient this is the
 finite-stage acyclicity of~$\tau_N$; the passage to the
 limit uses the Milnor exact sequence with the
 Mittag--Leffler condition guaranteed by the strong
 filtration axiom.
\item \emph{Dual acyclicity in the square-zero total model.}
 If $C = \varprojlim_N C_{\le N}$ is a separated complete
 pronilpotent curved dg chiral coalgebra with every finite
 quotient in the proved finite-stage regime, then the completed unit
 $\widehat\eta \colon C \to
 \widehat{\bar B}^{\mathrm{ch}}\bigl(
 \widehat\Omega^{\mathrm{ch}}(C)\bigr)$
 is a quasi-isomorphism when the completed total differential squares
 to zero; in a curved CDG ambient the corresponding assertion is that
 its cone is coacyclic in the specified second-kind category.
\item \emph{Completed bar-cobar inversion.}
 In the square-zero total model, completed bar and completed cobar
 are quasi-inverse equivalences on the full subcategory of strong
 completion towers whose finite quotients lie in the proved finite-stage
 regime. In the genuinely curved CDG ambient, the same sentence means
 that the completed algebra-side counit and coalgebra-side unit have
 coacyclic cones in the specified second-kind coderived category; it
 is not a raw chain-level inverse statement. The completed twisting
 morphism
 $\widehat\tau = \varprojlim_N \tau_N$ is a Maurer--Cartan
 element in the completed convolution algebra
 $\widehat{\operatorname{Conv}} :=
 \varprojlim_N \operatorname{Conv}(\bar B^{\mathrm{ch}}
 (\cA_{\le N}), \cA_{\le N})$,
 and this MC element represents the completed bar-cobar
 equivalence by twisting representability
 \textup{(}Theorem~\textup{%
 \ref{thm:completed-twisting-representability})}.
 The Verdier/Koszul-dual algebra $\cA^!$ is obtained separately
 from $\mathbb{D}_{\Ran}\widehat{\bar B}^{\mathrm{ch}}(\cA)$ under
 the finite-type/perfectness hypotheses of
 Theorem~\ref{thm:reconstruction-vs-duality}; it is not the object
 recovered by $\widehat\Omega\widehat{\bar B}$.
\end{enumerate}
\emph{This resolves the MC4 completion problem for strong completion
towers. When the strong-filtration axiom holds, the degree cutoff
\textup{(}Lemma~\textup{\ref{lem:degree-cutoff})} makes each
component of the MC equation a finite sum.
For towers that do not satisfy the strong-filtration axiom
\textup{(}e.g.\ $W_\infty$ with its spin-truncation tower\textup{)},
the same conclusion holds by a different route: the weight-cutoff
criterion \textup{(}Proposition~\textup{\ref{prop:mc4-weight-cutoff})}
together with inverse-limit continuity
\textup{(}Proposition~\textup{\ref{prop:inverse-limit-differential-continuity})}
supplies the Mittag--Leffler and convergence inputs needed by the
reduction principle
\textup{(}Proposition~\textup{\ref{prop:mc4-reduction-principle})}.}
\ClaimStatusConditional{}
\end{theorem}

\begin{proof}
\emph{Step~1: the completed coalgebra
 \textup{(}domain of $\widehat\tau$\textup{)}.}
Each quotient map $p_N\colon \cA_{\le N+1} \twoheadrightarrow
\cA_{\le N}$ induces a morphism of finite-stage bar coalgebras
$q_N\colon \bar B^{\mathrm{ch}}(\cA_{\le N+1}) \to
\bar B^{\mathrm{ch}}(\cA_{\le N})$ that commutes with the bar
differential, coproduct, coaugmentation, and curvature (by functoriality
of the finite-stage bar construction and strictness of~$p_N$). So the
inverse limit $\widehat{\bar B}^{\mathrm{ch}}(\cA)$ exists in the dg
category of chiral coalgebras with structure maps defined componentwise:
$b = \varprojlim_N b_N$, $\Delta = \varprojlim_N \Delta_N$,
$h = \varprojlim_N h_N$. The curved coalgebra identities hold
quotientwise and hence on the limit. The completed twisting morphism
$\widehat\tau := \varprojlim_N \tau_N$ is well-defined as a
degree-$+1$ map $\widehat{\bar B}^{\mathrm{ch}}(\cA) \to \cA$.

\emph{Continuity: the MC equation converges.}
By Lemma~\ref{lem:degree-cutoff}, modulo $F^{N+1}$ only degrees
$\le N$ contribute to the convolution product $\widehat\tau \star
\widehat\tau$. The MC equation
$\partial(\widehat\tau) + \widehat\tau \star \widehat\tau = 0$
therefore holds modulo $F^{N+1}$ (where it reduces to the
finite-stage equation for~$\tau_N$), hence on the completed
coalgebra. The bar differential $\widehat{b}$ is continuous
because the inverse-limit topology is initial for the projections.

\emph{Pronilpotence.}
The reduced coproduct on bar-length-$\le m$ elements is nilpotent on each finite
quotient (since bar length is bounded). Passing to the inverse limit
gives topological pronilpotence.

\emph{Step~2: the completed cobar object.}
Each $\Omega^{\mathrm{ch}}(\bar B^{\mathrm{ch}}(\cA_{\le N}))$ is a
curved dg algebra, and the transition maps are compatible, so the
inverse limit is again a complete curved dg algebra. The
inverse-limit description ensures continuity automatically: everything
is defined quotientwise.

\emph{Step~3: counit in the square-zero total model.}
The quotient of $\widehat\epsilon$ modulo $F^{N+1}$ is exactly the
finite-stage counit $\epsilon_N\colon
\Omega^{\mathrm{ch}}(\bar B^{\mathrm{ch}}(\cA_{\le N}))
\xrightarrow{\sim} \cA_{\le N}$, which is a quasi-isomorphism by the
finite-stage theorem. This is an ordinary cohomology statement only
for the square-zero total differential. In that case apply the
complete filtered comparison lemma
(Lemma~\ref{lem:complete-filtered-comparison}; Milnor exact sequence +
surjective quotient tower):
\[
0 \to \varprojlim\nolimits^1_N H^{m-1}(\operatorname{Cone}(\epsilon_N))
\to H^m(\operatorname{Cone}(\widehat\epsilon))
\to \varprojlim_N H^m(\operatorname{Cone}(\epsilon_N)) \to 0.
\]
The right term vanishes because each $\epsilon_N$ is a quasi-isomorphism.
The left term vanishes because the surjective quotient tower is
Mittag--Leffler. Hence $\widehat\epsilon$ is a quasi-isomorphism in
the square-zero total model. For a curved CDG realization replace this
last sentence by: the completed cone is coacyclic after the
second-kind coacyclicity hypotheses of
Proposition~\ref{prop:curved-bar-acyclicity} hold.

\emph{Step~4: unit on the coalgebra side.}
By the same quotientwise argument: the quotient of $\widehat\eta$ modulo
$F^{N+1}$ is the finite-stage unit, which is a quasi-isomorphism.
Lemma~\ref{lem:complete-filtered-comparison} applies identically in
the square-zero total model. In a curved CDG realization this proves
the stated unit only after the completed cone is shown coacyclic in
the relevant second-kind category.
\end{proof}

\begin{proposition}[Standard weight truncations and the induced bar filtration]
\label{prop:standard-strong-filtration}
\index{bar complex!weight truncations for standard families|textbf}
Let $\cA$ be one of the following chiral algebras:
\begin{enumerate}[label=\textup{(\alph*)}]
\item $V_k(\fg)$ for any simple~$\fg$ and $k \neq -h^\vee$;
\item $\mathrm{Vir}_c$ for any~$c$;
\item $\cW^k(\fg, f_{\mathrm{prin}})$ for any simple~$\fg$,
 principal nilpotent~$f$, and non-critical~$k$;
\item $V_\Lambda$ for any positive-definite even lattice~$\Lambda$.
\end{enumerate}
Write
\[
\cA = \bigoplus_{h \geq 0} \cA_h,
\qquad
\cA_{\le N} := \bigoplus_{0 \le h \le N} \cA_h,
\qquad
\widehat{\cA}^{\mathrm{wt}} := \prod_{h \geq 0} \cA_h .
\]
Then:
\begin{enumerate}[label=\textup{(\roman*)}]
\item the inverse limit of the finite weight truncations is the weight
 completion,
 \[
 \varprojlim_N \cA_{\le N} \cong \widehat{\cA}^{\mathrm{wt}},
 \]
 and the direct sum~$\cA$ embeds into $\widehat{\cA}^{\mathrm{wt}}$ as
 a dense subspace; in particular, the direct-sum algebra~$\cA$ is not
 identified with this inverse limit unless it is already weight-complete;
\item the finite truncations $\{\cA_{\le N}\}_{N \ge 0}$ form an inverse
 system of finite-type curved chiral $\Ainf$-algebras;
\item for the reduced bar complexes
 \[
 C_N := \bar B(\cA_{\le N}),
 \]
 the total conformal-weight filtration
 \[
 F_{\le w}C_N
 := \operatorname{span}\left\{
 s^{-1}a_1|\cdots|s^{-1}a_r \,\middle|\,
 \sum_{j=1}^r \operatorname{wt}(a_j) \le w
 \right\}
 \]
 is an exhaustive increasing filtration preserved by the bar
 differential;
\item for each weight bound~$w$ and each $N \ge w$, the transition map
 $C_{N+1} \to C_N$ restricts to an isomorphism
 \[
 F_{\le w}C_{N+1} \xrightarrow{\sim} F_{\le w}C_N.
 \]
\end{enumerate}
Consequently these standard families supply the bar-level hypotheses of
Proposition~\ref{prop:mc4-weight-cutoff}. What they do not supply is the
strong-filtration axiom~\eqref{eq:strong-filtration-condition} on
$\cA$ itself for the descending conformal-weight filtration.
\ClaimStatusProvedHere{}
\end{proposition}

\begin{proof}
Each family carries the standard nonnegative conformal-weight
decomposition
\[
\cA = \bigoplus_{h \ge 0} \cA_h
\]
with finite-dimensional weight spaces. For affine Kac--Moody,
Virasoro, and principal~$\mathcal{W}$-algebras this is the usual
$L_0$-grading by conformal weight; for a positive-definite even lattice
vertex algebra, finite-dimensionality of each~$\cA_h$ follows because
only finitely many lattice vectors have a given norm.

The finite truncation $\cA_{\le N}$ is therefore finite-dimensional.
The transition map $\cA_{\le N+1} \twoheadrightarrow \cA_{\le N}$ is
the quotient by the weight-$(N+1)$ summand, so the inverse system
$\{\cA_{\le N}\}_N$ is well defined. Its inverse limit is the product
\[
\varprojlim_N \cA_{\le N}
\cong \prod_{h \ge 0} \cA_h
= \widehat{\cA}^{\mathrm{wt}},
\]
whereas the original chiral algebra is the direct sum
$\bigoplus_{h \ge 0} \cA_h$. This is the completion issue in the
finding: bounded-below grading gives an inverse limit only after
passing to the weight completion.

Now let $C_N = \bar B(\cA_{\le N})$. The weight decomposition on
$\cA_{\le N}$ induces the total conformal weight on bar words by
\[
\operatorname{wt}(s^{-1}a_1|\cdots|s^{-1}a_r)
:= \sum_{j=1}^r \operatorname{wt}(a_j).
\]
Desuspension changes cohomological degree, not conformal weight. A
summand of the bar differential applies one of the chiral operations
$\mu_r$ to a consecutive block of homogeneous inputs
$a_1,\dots,a_r$ and extracts singular modes. For a single OPE step,
\[
\operatorname{wt}(a_{(n)}b)
= \operatorname{wt}(a) + \operatorname{wt}(b) - n - 1
\le \operatorname{wt}(a) + \operatorname{wt}(b)
\qquad (n \ge 0),
\]
so residue extraction never increases conformal weight. Iterating this
estimate through the $A_\infty$ terms shows that every summand of the
bar differential weakly decreases total conformal weight. Equivalently,
\begin{equation}\label{eq:bar-weight-preservation}
d_{\mathrm{bar}}(F_{\le w}C_N) \subset F_{\le w}C_N .
\end{equation}
This is the correct bar-level statement corresponding to the OPE weight
formula. It does \emph{not} imply the strong-filtration axiom
\eqref{eq:strong-filtration-condition} on~$\cA$ itself; for instance,
for affine currents one has
$\operatorname{wt}(J^a_{(0)}J^b)=1<2$.

Finally, if a homogeneous bar word has total conformal weight at most
$w$, then each letter occurring in it has weight at most~$w$. Hence for
$N \ge w$ every such word already lies in $\bar B(\cA_{\le N})$, and
the transition map $C_{N+1} \to C_N$ is an isomorphism on
$F_{\le w}$. This proves~\textup{(}iv\textup{)}, and the conclusion
about Proposition~\ref{prop:mc4-weight-cutoff} is immediate.
\end{proof}

\begin{proposition}[Reduction of MC4 to finite-stage compatibility;
\ClaimStatusProvedHere]
\label{prop:mc4-reduction-principle}
Let $\{C_N\}_{N \ge 0}$ be the inverse system of finite-stage bar
complexes
\[
C_N := \bar B(\cA_{\le N}),
\qquad
\widehat{C} := \varprojlim_N C_N,
\]
for a tower whose finite stages $\cA_{\le N}$ lie in the proved
bar-cobar regime. Assume:
\begin{enumerate}
\item each transition map $C_{N+1} \to C_N$ is a morphism of complexes;
\item for every cohomological degree~$m$, the inverse system
 $\{H^m(C_N)\}_N$ satisfies the Mittag--Leffler condition; and
\item the finite-stage bar-cobar quasi-isomorphisms in the
 square-zero total model, respectively coacyclic-cone equivalences in
 the curved CDG model,
 $\epsilon_N \colon \Omega(C_N) \xrightarrow{\sim} \cA_{\le N}$ are
 compatible with the tower maps.
\end{enumerate}
Then:
\begin{enumerate}
\item the completed cohomology is computed by the limit of the
 finite-stage cohomologies,
 \[
 H^m(\widehat{C}) \cong \varprojlim_N H^m(C_N);
 \]
\item the induced map
 \[
 \widehat{\epsilon}\colon \Omega(\widehat{C}) \longrightarrow
 \varprojlim_N \cA_{\le N}
 \]
is a quasi-isomorphism whenever the completed cobar differential is
continuous and square-zero; in the genuinely curved CDG ambient the
same map has coacyclic cone in the specified second-kind category.
\end{enumerate}
Consequently, the remaining content of MC4 for a concrete tower is to
verify continuity, pro-nilpotence, and the Mittag--Leffler/convergence
conditions; no new formal obstruction remains once these are checked.
\end{proposition}

\begin{proof}
For any inverse system of cochain complexes there is a Milnor exact
sequence
\[
0 \to \varprojlim\nolimits^1_N H^{m-1}(C_N)
\to H^m(\widehat{C})
\to \varprojlim_N H^m(C_N) \to 0.
\]
Assumption~\textup{(2)} kills the derived limit term
$\varprojlim^1$, so the first claim follows.

For the second claim, the maps $\epsilon_N$ assemble by
Assumption~\textup{(3)} into a morphism of inverse systems. Passing to
limits therefore gives a map
$\widehat{\epsilon}\colon \Omega(\widehat{C}) \to \varprojlim_N \cA_{\le N}$.
The finite-stage maps are quasi-isomorphisms in the square-zero total
model, so after applying the first part to the bar complexes and the
same Milnor argument to the compatible cobar tower,
$\widehat{\epsilon}$ induces an isomorphism on cohomology. The
continuity and square-zero hypotheses ensure that the completed cobar
differential is well defined as an ordinary differential, so
$\widehat{\epsilon}$ is a quasi-isomorphism. In the curved CDG
ambient this argument is applied to the completed cone and proves
coacyclicity rather than ordinary acyclicity.
\end{proof}

% \label{rem:mc4-concrete-checklist} removed (echo of Proposition~\ref{prop:mc4-reduction-principle})

\begin{corollary}[Degreewise stabilization criterion for MC4;
\ClaimStatusProvedHere]
\label{cor:mc4-degreewise-stabilization}
In the setting of Proposition~\ref{prop:mc4-reduction-principle},
assume in addition that for every cohomological degree~$m$ there exists
$N(m)$ such that the transition maps
\[
H^m(C_{N+1}) \longrightarrow H^m(C_N)
\]
are isomorphisms for all $N \ge N(m)$. Then:
\begin{enumerate}
\item the inverse system $\{H^m(C_N)\}_N$ is Mittag--Leffler;
\item the completed cohomology stabilizes at the finite stage,
 \[
 H^m(\widehat{C}) \cong H^m(C_{N(m)});
 \]
\item every compatible system of finite-stage bar-cobar
 quasi-isomorphisms induces a completed quasi-isomorphism
 \[
 \Omega(\widehat{C}) \xrightarrow{\sim} \varprojlim_N \cA_{\le N},
 \]
 provided the completed cobar differential is continuous and
 square-zero; in the curved CDG ambient this is replaced by a
 coacyclic-cone statement.
\end{enumerate}
Thus eventual stabilization of the finite-stage bar cohomology is a
sufficient MC4 input.
\end{corollary}

\begin{proof}
Eventual constancy implies the Mittag--Leffler condition, so
Proposition~\ref{prop:mc4-reduction-principle} applies. Since the
transition maps are isomorphisms for $N \ge N(m)$, the inverse limit
$\varprojlim_N H^m(C_N)$ identifies with the stable value
$H^m(C_{N(m)})$. The second and third claims are then exactly the
conclusions of
Proposition~\ref{prop:mc4-reduction-principle}.
\end{proof}

\begin{corollary}[Finite-dimensional surjectivity criterion for MC4;
\ClaimStatusProvedHere]
\label{cor:mc4-surjective-criterion}
In the setting of Proposition~\ref{prop:mc4-reduction-principle},
assume in addition that for every cohomological degree~$m$:
\begin{enumerate}
\item each $H^m(C_N)$ is finite-dimensional; and
\item there exists $N_{\mathrm{surj}}(m)$ such that the transition maps
 \[
 H^m(C_{N+1}) \longrightarrow H^m(C_N)
 \]
 are surjective for all $N \ge N_{\mathrm{surj}}(m)$.
\end{enumerate}
Then the inverse system $\{H^m(C_N)\}_N$ is Mittag--Leffler, and every
compatible system of finite-stage bar-cobar quasi-isomorphisms induces a
completed quasi-isomorphism
\[
\Omega(\widehat{C}) \xrightarrow{\sim} \varprojlim_N \cA_{\le N},
\]
provided the completed cobar differential is continuous and
square-zero; in the curved CDG ambient this is replaced by a
coacyclic-cone statement.
Thus MC4 follows from eventual surjectivity on finite-dimensional
weight slices; one does not need a priori eventual constancy of the full
cohomology tower.
\end{corollary}

\begin{proof}
Fix $m$. For $N \ge N_{\mathrm{surj}}(m)$ and every $r \ge 0$, the image
of the composite map
\[
H^m(C_{N+r}) \longrightarrow H^m(C_N)
\]
is all of $H^m(C_N)$. Hence the descending image filtration in
$H^m(C_N)$ stabilizes immediately, so the inverse system is
Mittag--Leffler. Proposition~\ref{prop:mc4-reduction-principle} then
gives the completed cohomology comparison and the completed bar-cobar
quasi-isomorphism.
\end{proof}

\begin{proposition}[Weight-cutoff criterion for MC4;
\ClaimStatusProvedHere]
\label{prop:mc4-weight-cutoff}
In the setting of Proposition~\ref{prop:mc4-reduction-principle},
assume in addition that each finite-stage bar complex $C_N$ carries an
exhaustive increasing auxiliary-weight filtration
\[
F_{\le 0}C_N \subset F_{\le 1}C_N \subset \cdots \subset C_N
\]
such that:
\begin{enumerate}
\item the differential and the transition maps preserve the filtration;
\item for each weight bound $w$, the filtered piece $F_{\le w}C_N$ is
 finite-dimensional in every cohomological degree; and
\item for each $w$ and every $N \ge w$, the transition map induces an
 isomorphism of filtered subcomplexes
 \[
 F_{\le w}C_{N+1} \xrightarrow{\sim} F_{\le w}C_N.
 \]
\end{enumerate}
Then for every cohomological degree $m$, weight bound $w$, and
$N \ge w$, the induced map
\[
H^m(F_{\le w}C_{N+1}) \xrightarrow{\sim} H^m(F_{\le w}C_N)
\]
is an isomorphism. Consequently the associated graded weight slices of
cohomology stabilize, and if the cohomology splits as a direct sum of
finite-dimensional weight slices, then the maps
\[
H^m(C_{N+1})_w \xrightarrow{\sim} H^m(C_N)_w
\]
are isomorphisms for all $N \ge w$. Hence the eventual-surjectivity
input of Corollary~\ref{cor:mc4-surjective-criterion} is automatic once
the tower is defined by genuine weight truncations.
\end{proposition}

\begin{proof}
Assumption~\textup{(3)} identifies the filtered subcomplexes
$F_{\le w}C_{N+1}$ and $F_{\le w}C_N$ for every $N \ge w$, so their
cohomology groups are canonically isomorphic. Passing to successive
quotients gives stabilization of the associated graded weight pieces.
If the cohomology splits by weight, the direct-summand description of
the weight slices identifies $H^m(C_N)_w$ with the corresponding
graded piece, so the same stabilization holds on the nose. In
particular the transition maps are eventually surjective on each
finite-dimensional weight slice, and
Corollary~\ref{cor:mc4-surjective-criterion} applies.
\end{proof}

\begin{proposition}[\texorpdfstring{$W_\infty$}{W_infty} criterion from principal finite-type stages;
\ClaimStatusProvedHere]
\label{prop:winfty-mc4-criterion}
Let
\[
W_\infty = \varprojlim_N W_N
\]
be a separated complete infinite-generator $\mathcal{W}$-tower whose
finite stages $W_N$ lie in the proved principal finite-type
$\mathcal{W}$ regime. Assume:
\begin{enumerate}
\item the transition maps $W_{N+1} \to W_N$ are compatible with the
 finite-stage bar complexes and bar-cobar quasi-isomorphisms;
\item for every cohomological degree~$m$, the inverse system
 $\{H^m(\bar B(W_N))\}_N$ is eventually constant; and
\item the transition maps define an inverse system of curved dg
 coalgebras on the finite-stage bar complexes.
\end{enumerate}
Then the completed bar object
\[
\widehat{\bar B}(W_\infty) := \varprojlim_N \bar B(W_N)
\]
is a separated complete pro-nilpotent curved dg coalgebra, and the
canonical comparison map
\[
\Omega\bigl(\widehat{\bar B}(W_\infty)\bigr) \xrightarrow{\sim} W_\infty
\]
is a quasi-isomorphism in the square-zero total model. In the
genuinely curved CDG ambient, the same comparison has coacyclic cone
in the completed second-kind category. Thus the M-level completed
bar-cobar package for $W_\infty$ reduces to compatible inverse-system
dg structure and stabilization of the principal finite-type tower,
together with the square-zero or coacyclic-cone hypothesis appropriate
to the chosen ambient. The stronger H-level realization problem
belongs to MC4; any physics-facing identification built on top of that
realization is downstream at MC5. By
Corollary~\ref{cor:mc4-surjective-criterion}, the eventual-constancy
hypothesis may be replaced by eventual surjectivity on finite-dimensional
weight slices, and Proposition~\ref{prop:mc4-weight-cutoff} shows that
this surjectivity is formal for the standard principal-stage tower.
\end{proposition}

\begin{proof}
The principal finite-type stages supply the finite-stage bar-cobar
quasi-isomorphisms required by
Proposition~\ref{prop:mc4-reduction-principle}. Assumption~\textup{(2)}
is exactly the degreewise stabilization hypothesis of
Corollary~\ref{cor:mc4-degreewise-stabilization}, and
Assumption~\textup{(3)} lets one apply
Proposition~\ref{prop:inverse-limit-differential-continuity} to the
inverse-limit bar and cobar differentials. Applying
Corollary~\ref{cor:mc4-degreewise-stabilization} to the tower
$\{\bar B(W_N)\}_N$ yields the stated completed quasi-isomorphism.
In the curved CDG ambient, the same inverse-limit argument applied to
the completed cone yields the coacyclic-cone comparison.
\end{proof}

\begin{corollary}[Standard principal-stage cutoff for
\texorpdfstring{$W_\infty$}{W_infty}; \ClaimStatusProvedHere]
\label{cor:winfty-weight-cutoff}
Let $W_N$ be the standard principal finite-type stage of $W_\infty$,
so that $W_N$ retains exactly the generators of conformal weights
$2,3,\dots,N$ and the transition map $W_{N+1}\to W_N$ forgets the
generator of weight $N+1$. Then the tower of bar complexes
$\{\bar B(W_N)\}_N$ satisfies the hypotheses of
Proposition~\ref{prop:mc4-weight-cutoff} for the total conformal-weight
filtration. In particular, for every cohomological degree $m$ and
conformal weight $w$, the map
\[
H^m(\bar B(W_{N+1}))_w \xrightarrow{\sim} H^m(\bar B(W_N))_w
\]
is an isomorphism for all $N \ge w$.

Therefore the standard principal-stage tower already supplies the
surjectivity/stabilization part of MC4. Together with
Proposition~\ref{prop:inverse-limit-differential-continuity}, this
reduces the standard $W_\infty$ tower to the standard-tower completed
M-level package of Corollary~\ref{cor:winfty-standard-mc4-package}.
\end{corollary}

\begin{proof}
A bar chain of total conformal weight $w$ can involve only generators
whose conformal weights are at most $w$. Once $N \ge w$, all such
generators already occur in $W_N$, so the filtered bar subcomplex of
weights $\le w$ is identical in every later stage. The bar
differential preserves the total conformal-weight filtration used
throughout the \(W\)-algebra computations, so
Proposition~\ref{prop:mc4-weight-cutoff} applies.
\end{proof}

\begin{remark}[MC4 splitting: positive vs.\ resonant]
\label{rem:mc4-positive-vs-resonant}
Proposition~\ref{prop:mc4-weight-cutoff} and
Corollary~\ref{cor:winfty-weight-cutoff} show that positive towers
(those carrying a strict positive weight grading preserved by the bar
differential) are settled under those hypotheses:
the weight-$w$ summand stabilizes at finite stage~$w$.
This resolves the ``MC4$^+$'' sub-problem, in the positive-tower
ambient, for
$\mathcal{W}_{1+\infty}$, affine Yangians, and positive RTT towers.

The residual ``MC4$^0$'' sub-problem concerns algebras with
weight-zero resonance: a finite-dimensional piece $R_\cA$ on which
higher operations can preserve filtration degree~$0$
(Definition~\ref{def:resonance-rank}). The resonance rank
$\rho(\cA) := \dim H^*(R_\cA, m_1^{R_\cA})$ classifies the
completion difficulty: $\rho = 0$ (positive weight cutoff),
$0 < \rho < \infty$ (finite resonance, tractable via
Theorem~\ref{thm:resonance-filtered-bar-cobar}), or
$\rho = \infty$ (wild).

For Virasoro and non-quadratic $\mathcal{W}_N$-algebras, the
proved regime is $0 < \rho < \infty$
(Theorem~\ref{thm:platonic-completion}).
The same-family shadow $\mathrm{Vir}_{26-c}$ is the depth-zero
resonance shadow; the genuine $\mathcal{W}_\infty$-type dual
requires the full resonance-filtered completion.
See Appendix~\ref{app:nilpotent-completion},
\S\ref{sec:mc4-splitting} for the complete development.
\end{remark}

\begin{proposition}[Continuity of inverse-limit bar and cobar differentials;
\ClaimStatusProvedHere]
\label{prop:inverse-limit-differential-continuity}
Let $\{C_N,d_N,\Delta_N\}_{N \ge 0}$ be an inverse system of finite-stage
curved dg coalgebras with transition maps
$\pi_{N+1,N}\colon C_{N+1}\to C_N$ commuting with the differentials and
coproducts. Endow
\[
\widehat C := \varprojlim_N C_N
\]
with the inverse-limit topology, and let
\[
\Omega(\widehat C)
:=
\prod_{m \ge 0} (s^{-1}\widehat C)^{\widehat\otimes m}
\]
carry the completed tensor-product topology. Then:
\begin{enumerate}
\item the componentwise differential
 $\widehat d((x_N)_N):=(d_Nx_N)_N$ on $\widehat C$ is well defined and
 continuous;
\item the inverse-limit coproduct $\widehat\Delta$ is continuous; and
\item the completed cobar differential on $\Omega(\widehat C)$ induced
 by $\widehat d$ and $\widehat\Delta$ is continuous.
\end{enumerate}
Consequently, once an MC4 tower is realized as a genuine inverse system
of finite-stage curved dg coalgebras, continuity is formal for the
inverse-limit topology.
\end{proposition}

\begin{proof}
Let $p_N\colon \widehat C \to C_N$ be the projection. Because the
transition maps commute with the differentials,
\[
p_N \circ \widehat d = d_N \circ p_N
\]
for every $N$, so the componentwise formula defines a map
$\widehat d\colon \widehat C \to \widehat C$. The inverse-limit
topology on $\widehat C$ is initial with respect to the projections
$p_N$, hence $\widehat d$ is continuous as soon as each composite
$p_N \circ \widehat d$ is continuous. This holds because each $C_N$ is
given the finite-stage topology and $d_N$ is a morphism of complexes.

The same argument applies to the coproduct. Writing
$p_N^{\otimes 2}\colon \widehat C \widehat\otimes \widehat C \to
C_N\otimes C_N$ for the induced projection, one has
\[
p_N^{\otimes 2} \circ \widehat\Delta = \Delta_N \circ p_N,
\]
so $\widehat\Delta$ is well defined and continuous.

The completed cobar differential on $\Omega(\widehat C)$ is the sum of
the internal part coming from $\widehat d$ and the bar part coming from
$\widehat\Delta$. On each tensor-length summand
$(s^{-1}\widehat C)^{\widehat\otimes m}$, each contribution is a finite
sum of maps obtained by inserting either $\widehat d$ or
$\widehat\Delta$ in one slot. Completed tensor products preserve
continuity of continuous multilinear maps, and finite sums of
continuous maps remain continuous. The product topology on
$\Omega(\widehat C)$ is defined componentwise by the tensor-length
summands, so the total cobar differential is continuous.
\end{proof}

\begin{corollary}[Standard principal-stage
\texorpdfstring{$W_\infty$}{W_infty} tower satisfies the M-level MC4
package; \ClaimStatusProvedHere]
\label{cor:winfty-standard-mc4-package}
For the standard principal finite-type stages $W_N$ of $W_\infty$, the
completed bar object
\[
\widehat{\bar B}(W_\infty) := \varprojlim_N \bar B(W_N)
\]
has continuous bar and cobar differentials, and the canonical
comparison map
\[
\Omega\bigl(\widehat{\bar B}(W_\infty)\bigr) \xrightarrow{\sim} W_\infty
\]
is a quasi-isomorphism in the square-zero total model, and a
coacyclic-cone equivalence in the completed curved CDG ambient.
Thus the standard principal-stage tower realizes the full
standard-tower M-level completed bar-cobar package for $W_\infty$; the
remaining problem is
the H-level comparison with stronger factorization-theoretic or
physical realizations.
\end{corollary}

\begin{proof}
The finite-stage bar complexes form an inverse system of curved dg
coalgebras by Proposition~\ref{prop:winfty-mc4-criterion},
Assumption~\textup{(3)}, so
Proposition~\ref{prop:inverse-limit-differential-continuity} gives the
continuity of the completed bar and cobar differentials. The
finite-stage bar-cobar quasi-isomorphisms are provided by the proved
principal finite-type regime, and
Corollary~\ref{cor:winfty-weight-cutoff} supplies the
Mittag--Leffler/stabilization input. Since
$W_\infty = \varprojlim_N W_N$ by definition, applying
Proposition~\ref{prop:mc4-reduction-principle} to the standard tower
gives the stated quasi-isomorphism.
In the curved CDG ambient, the same application gives the
coacyclic-cone version.
\end{proof}

\begin{proposition}[Comparison with a completed target by compatible
finite quotients; \ClaimStatusProvedHere]
\label{prop:completed-target-comparison}
Let $\{\cA_{\le N}\}_{N \ge 0}$ be an inverse system of dg algebras with
surjective transition maps and completed limit
\[
\widehat{\cA}:=\varprojlim_N \cA_{\le N}.
\]
Let $\cB$ be a separated complete dg algebra with descending filtration
by closed dg ideals
\[
F^1\cB \supset F^2\cB \supset \cdots.
\]
Assume:
\begin{enumerate}
\item the canonical map
 \[
 \cB \xrightarrow{\sim} \varprojlim_N \cB/F^{N+1}\cB
 \]
 is an isomorphism of dg algebras;
\item for each $N$ there is a dg algebra quasi-isomorphism
 \[
 f_N\colon \cB/F^{N+1}\cB \xrightarrow{\sim} \cA_{\le N};
 \]
\item the maps $f_N$ commute with the quotient and transition maps; and
\item for every cohomological degree $m$, the inverse systems
 $\{H^m(\cB/F^{N+1}\cB)\}_N$ and $\{H^m(\cA_{\le N})\}_N$ satisfy the
 Mittag--Leffler condition.
\end{enumerate}
Then the induced map
\[
\widehat f\colon \cB \longrightarrow \widehat{\cA}
\]
is a quasi-isomorphism.

If, in addition, the quotient tower $\{\cB/F^{N+1}\cB\}_N$ carries
bar-cobar comparison maps compatible with the tower
$\{\cA_{\le N}\}_N$ and the completed bar and cobar differentials on
$\cB$ are continuous, then $\cB$ and $\widehat{\cA}$ define
quasi-isomorphic completed M-level bar-cobar objects.
\end{proposition}

\begin{proof}
Assumption~\textup{(1)} identifies $\cB$ with the inverse limit of its
finite quotients. By Assumptions~\textup{(2)} and~\textup{(3)}, the
maps $f_N$ assemble into a morphism of inverse systems and hence into a
map
\[
\widehat f\colon
\cB \cong \varprojlim_N \cB/F^{N+1}\cB \longrightarrow
\varprojlim_N \cA_{\le N}=\widehat{\cA}.
\]

Apply the Milnor exact sequence to the quotient tower of $\cB$ and to
the tower $\{\cA_{\le N}\}_N$. Assumption~\textup{(4)} kills the
derived-limit terms, so for every cohomological degree $m$ one has
\[
H^m(\cB)\cong \varprojlim_N H^m(\cB/F^{N+1}\cB),
\qquad
H^m(\widehat{\cA})\cong \varprojlim_N H^m(\cA_{\le N}).
\]
Since each $f_N$ is a quasi-isomorphism, the maps
$H^m(f_N)\colon H^m(\cB/F^{N+1}\cB)\xrightarrow{\sim}
H^m(\cA_{\le N})$ are isomorphisms compatible with the inverse systems.
Passing to inverse limits shows that $H^m(\widehat f)$ is an isomorphism
for every $m$, so $\widehat f$ is a quasi-isomorphism.

For the final claim, apply the same inverse-limit comparison to the
compatible bar towers. Continuity of the completed bar and cobar
differentials makes the completed bar objects well defined, and
Proposition~\ref{prop:mc4-reduction-principle} identifies both
completed bar-cobar packages with the same inverse-limit finite-stage
data. Hence the resulting completed M-level bar-cobar objects are
quasi-isomorphic.
\end{proof}

\begin{corollary}[H-level comparison criterion for
\texorpdfstring{$W_\infty$}{W_infty}; \ClaimStatusConditional]
\label{cor:winfty-hlevel-comparison-criterion}
Let $\mathcal{W}^{\mathrm{ht}}$ be a separated complete dg model for a
factorization-theoretic or physical completion of $W_\infty$, equipped
with a descending conformal-weight filtration by closed dg ideals.
Assume:
\begin{enumerate}
\item $\mathcal{W}^{\mathrm{ht}}$ is the inverse limit of its finite quotients;
\item each finite quotient
 $\mathcal{W}^{\mathrm{ht}}/F^{N+1}\mathcal{W}^{\mathrm{ht}}$ is dg
 quasi-isomorphic to the principal finite-type stage $W_N$,
 compatibly with the truncation tower; and
\item the quotient tower inherits the same finite-stage bar-cobar data
 and continuity hypotheses as in
 Corollary~\ref{cor:winfty-standard-mc4-package}.
\end{enumerate}
Then the induced comparison map
\[
\mathcal{W}^{\mathrm{ht}} \xrightarrow{\sim} W_\infty
\]
is a quasi-isomorphism, and $\mathcal{W}^{\mathrm{ht}}$ carries the same
completed M-level bar-cobar object as the standard principal-stage
completion. Thus the remaining $\mathcal{W}_\infty$ problem is to
construct such a filtered H-level target and identify its finite
quotients by the exact residue identities
$\mathsf{C}^{\mathrm{res}}_{s,t;u;m,n}(N)
=\mathsf{C}^{\mathrm{DS}}_{s,t;u;m,n}(N)$,
discharged on the finite-detection packet $\mathcal{I}_N$, whose first
live higher-spin stages are reduced in this chapter to
the stage-$3$ fifteen-coefficient packet and the exact stage-$4$
six-entry identity packet on~$\mathcal{I}_4$, organized as three local
OPE blocks and splitting into four higher-spin channels together with
two proved principal Virasoro-target values
$\mathsf{C}^{\mathrm{DS}}_{4,4;2;0,6}=2$ and
$\mathsf{C}^{\mathrm{DS}}_{3,4;2;0,5}=0$, with the corresponding
residue-side contraction supplied only under the normalized two-point
package of Corollary~\ref{cor:winfty-stage4-residue-four-channel}, not
another stabilization theorem.
\end{corollary}

\begin{proof}
Apply Proposition~\ref{prop:completed-target-comparison} with
$\cA_{\le N}=W_N$. The Mittag--Leffler input is supplied by the same
weight-cutoff argument as in Corollary~\ref{cor:winfty-weight-cutoff},
and the final statement follows from the bar-cobar comparison part of
Proposition~\ref{prop:completed-target-comparison}.
\end{proof}

\begin{definition}[Principal-stage compatible
\texorpdfstring{$W_\infty$}{W_infty} target]
\label{def:winfty-principal-stage-compatible}
A separated complete dg algebra $\mathcal{W}^{\mathrm{ht}}$ is called a
\emph{principal-stage compatible $W_\infty$ target} if it is equipped
with a descending conformal-weight filtration by closed dg ideals such
that:
\begin{enumerate}
\item $\mathcal{W}^{\mathrm{ht}} \cong
 \varprojlim_N \mathcal{W}^{\mathrm{ht}}/F^{N+1}\mathcal{W}^{\mathrm{ht}}$;
\item each finite quotient
 $\mathcal{W}^{\mathrm{ht}}/F^{N+1}\mathcal{W}^{\mathrm{ht}}$ is dg
 quasi-isomorphic to the principal finite-type stage $W_N$;
\item the quotient maps intertwine the principal-stage truncation tower;
 and
\item the quotient tower carries bar-cobar comparison maps compatible
 with the finite-type $W_N$ bar complexes.
\end{enumerate}
This is the H-level input singled out by
Corollary~\ref{cor:winfty-hlevel-comparison-criterion}.
\end{definition}

\begin{definition}[Principal-stage quotient system for a factorization
\texorpdfstring{$W_\infty$}{W_infty} completion]
\label{def:winfty-quotient-system}
Let $\mathcal{W}^{\mathrm{fact}}_\infty$ be a separated complete
factorization-theoretic completion of the principal $W$-tower. A
\emph{principal-stage quotient system} on
$\mathcal{W}^{\mathrm{fact}}_\infty$ consists of dg quotient maps
\[
q_N\colon \mathcal{W}^{\mathrm{fact}}_\infty \longrightarrow
\mathcal{W}^{\mathrm{fact}}_{\le N}
\]
such that:
\begin{enumerate}
\item the kernels of the $q_N$ define a descending separated complete
 conformal-weight filtration;
\item each quotient object
 $\mathcal{W}^{\mathrm{fact}}_{\le N}$ is dg quasi-isomorphic to the
 principal finite-type stage $W_N$;
\item the transition maps between the quotients intertwine the
 principal-stage truncation tower; and
\item the induced bar complexes agree with the finite-type bar-cobar
 package already proved for $W_N$.
\end{enumerate}
\end{definition}

\begin{proposition}[Formal descent criterion for the
\texorpdfstring{$W_\infty$}{W_infty} factorization target;
\ClaimStatusConditional]
\label{prop:winfty-quotient-system-criterion}
A separated complete factorization-theoretic completion equipped with a
principal-stage quotient system
(Definition~\ref{def:winfty-quotient-system}) is a principal-stage
compatible $W_\infty$ target
(Definition~\ref{def:winfty-principal-stage-compatible}).
Bar-cobar compatibility further implies
Corollary~\ref{cor:winfty-hlevel-comparison-criterion}.
\end{proposition}

\begin{proof}
The quotient-system data package precisely the inputs of Definition~\ref{def:winfty-principal-stage-compatible}.
\end{proof}

\begin{proposition}[Factorization-envelope criterion for principal
stages; \ClaimStatusConditional]
\label{prop:winfty-factorization-envelope-criterion}
Let $\mathcal{F}_\infty$ be a separated complete factorization algebra
on a curve with quotient maps
$\mathcal{F}_\infty \twoheadrightarrow \mathcal{F}_{\le N}$
whose kernels define a descending separated complete conformal-weight
filtration. If each $\mathcal{F}_{\le N}$ admits a BD chiral envelope
quasi-isomorphic to $W_N$ and the quotient maps intertwine the
principal-stage tower, then the chiral envelope of
$\mathcal{F}_\infty$ is a principal-stage compatible $W_\infty$ target.
Bar-data compatibility further implies
Proposition~\ref{prop:winfty-quotient-system-criterion} and
Corollary~\ref{cor:winfty-hlevel-comparison-criterion}.
\end{proposition}

\begin{proof}
The chiral-envelope functor produces compatible quotients recovering the principal stages; Proposition~\ref{prop:winfty-quotient-system-criterion} applies.
\end{proof}

%% Completion closure theory (raeeznotes85 absorption)

\begin{definition}[Completion closure of the finite-type regime]
\label{def:completion-closure}
\index{completion closure|textbf}
Let $\Fft$ denote the class of finite-type augmented curved chiral
$\Ainf$-algebras for which the finite-stage bar-cobar theorem is proved.
The \emph{completion closure} $\CompCl(\Fft)$ is the class of strong
completion towers $\cA = \varprojlim_N \cA_{\le N}$ with every quotient
$\cA_{\le N} \in \Fft$.

Dually, define $\Cstr$ to be the class of separated complete pronilpotent
curved dg chiral coalgebras $C = \varprojlim_N C_{\le N}$ such that each
$C_{\le N}$ lies in the finite-stage essential image.
\end{definition}

\begin{proposition}[Essential image of the completed bar functor;
\ClaimStatusConditional]
\label{prop:completion-closure-essential-image}
\index{bar functor!completed essential image}
A separated complete pronilpotent curved dg chiral coalgebra
$C=\varprojlim_N C_{\le N}$ lies in the essential image of
$\widehat{\bar B}^{\mathrm{ch}}\colon\CompCl(\Fft)\to\Cstr$ in the
localized square-zero/coderived sense if and only if the following
finite-window conditions hold.
\begin{enumerate}[label=\textup{(\roman*)}]
\item Each quotient $C_{\le N}$ lies in the finite-stage essential
 image: there is $\cA_{\le N}\in\Fft$ and a finite-stage weak
 equivalence
 $C_{\le N}\simeq \bar B^{\mathrm{ch}}(\cA_{\le N})$.
\item The algebras $\{\cA_{\le N}\}_N$ and the coalgebras
 $\{C_{\le N}\}_N$ form compatible strict inverse systems with
 surjective transition maps, continuous curvature, coproduct,
 differential, and twisting morphisms.
\item The cone towers of the finite-stage units and counits are
 Mittag--Leffler and are respectively acyclic in the square-zero total
 model or coacyclic in the curved second-kind model.
\end{enumerate}
Under these hypotheses
$\cA:=\varprojlim_N\cA_{\le N}$ belongs to $\CompCl(\Fft)$ and
$C\simeq\widehat{\bar B}^{\mathrm{ch}}(\cA)$ after localization.
No object enters the essential image merely because its underlying
graded coalgebra is complete or conilpotent; the finite-stage
twisting morphisms and the ML cone control are part of the datum.
\end{proposition}

\begin{proof}
If $C\simeq\widehat{\bar B}^{\mathrm{ch}}(\cA)$ for
$\cA\in\CompCl(\Fft)$, quotienting by $F^{N+1}$ gives the finite-stage
bar coalgebras
$C_{\le N}\simeq\bar B^{\mathrm{ch}}(\cA_{\le N})$, and
Theorem~\ref{thm:completed-bar-cobar-strong} gives the compatible
unit, counit, and cone conditions.

Conversely, assume \textup{(i)--(iii)}. The strict inverse system
$\{\cA_{\le N}\}_N$ defines an object
$\cA=\varprojlim_N\cA_{\le N}$ of $\CompCl(\Fft)$, and the compatible
finite-stage equivalences assemble to a map
\[
C\longrightarrow \widehat{\bar B}^{\mathrm{ch}}(\cA).
\]
The cone of this map is the inverse limit of the finite-stage cones.
The Mittag--Leffler hypothesis kills the Milnor
$\varprojlim^1$ term in the square-zero total model; in the curved
CDG model the same strict cone tower lies in the completed coacyclic
subcategory by Proposition~\ref{prop:curved-bar-acyclicity}. Hence the
assembled map is an equivalence in the chosen localization.
\end{proof}

\begin{corollary}[Homotopy-categorical equivalence on the completion
closure; \ClaimStatusConditional]
\label{cor:completion-closure-equivalence}
\index{completion closure!homotopy equivalence}
If one localizes strong completion towers and objects in $\Cstr$ by
quotientwise quasi-isomorphisms in the square-zero total model, and
by proved coacyclic-cone equivalences in the genuinely curved CDG
model, then
\[
\widehat\Omega^{\mathrm{ch}} \;\dashv\; \widehat{\bar B}^{\mathrm{ch}}
\]
induces an equivalence of homotopy categories. That is,
Theorem~\textup{\ref{thm:completed-bar-cobar-strong}} gives quasi-inverse
functors on the localized categories.
\end{corollary}

\begin{proof}
The algebra-side counit and coalgebra-side unit are quasi-isomorphisms,
respectively coacyclic-cone equivalences in the curved CDG ambient, by
Theorem~\ref{thm:completed-bar-cobar-strong}(3)--(4). Both pass to
the localized categories.
\end{proof}

\begin{remark}[Adjunction direction convention]\label{rem:adjunction-direction-convention}
The canonical adjunction direction throughout this monograph is
$\Omegach \dashv \barBch$ (cobar left adjoint to bar), as stated in
the chapter opening: the Hom-set identity is
$\operatorname{Hom}_{\mathrm{Alg}}(\Omegach(\cC), \cA)
\cong \operatorname{Tw}(\cC, \cA)
\cong \operatorname{Hom}_{\mathrm{Coalg}}(\cC, \barBch(\cA))$,
following the Loday--Vallette convention~\textup{\cite{LV12}}.
The displayed formula above and occasional later passages must not be
read with the bar functor as the left adjoint. When a display lists
the bar functor first, it is only listing the right adjoint before the
left adjoint, not changing the adjunction.
The mathematical content (unit
$\eta \colon \cC \to \barBch(\Omegach(\cC))$, counit
$\psi \colon \Omegach(\barBch(\cA)) \to \cA$) is
unambiguous and consistent throughout.
\end{remark}

\begin{definition}[Window-stability]
\label{def:window-stability}
\index{window-stability|textbf}
For a tower $\{\cA_{\le N}\}_{N \ge 0}$ and each pair of bounds
$(m,q)$ (cohomological degree, total conformal weight), let
\[
W_{m,q}(\cA_{\le N})
\subset \bar B^{\mathrm{ch}}(\cA_{\le N})
\]
be the finite-dimensional subcomplex spanned by bar chains of
cohomological degree~$m$ and total conformal weight~$\le q$. The tower
is \emph{window-stable} if for every $(m,q)$ there exists $N_0(m,q)$
such that for all $N \ge N_0(m,q)$ the transition maps identify
$W_{m,q}(\cA_{\le N+1}) \cong W_{m,q}(\cA_{\le N})$ as subcomplexes
(differentials agree under the identification).
\end{definition}

\begin{theorem}[Coefficient-stability criterion; \ClaimStatusConditional]
\label{thm:coefficient-stability-criterion}
\index{coefficient-stability criterion|textbf}
Let $\{\cA_{\le N}\}_{N \ge 0}$ be a compatible tower of finite-type
stages presented by generators and OPE/$\Ainf$-structure constants.
Suppose that for every finite generator set~$S$, bar degree bound~$m$,
and total-weight bound~$q$, the finitely many OPE/$\Ainf$-coefficients
entering the bar differential on the window $W_{m,q}(\cA_{\le N})$ are
eventually independent of~$N$. Then:
\begin{enumerate}
\item the tower is window-stable;
\item the inverse limit $\cA = \varprojlim_N \cA_{\le N}$ carries a
 unique strong completion-tower structure; and
\item Theorem~\textup{\ref{thm:completed-bar-cobar-strong}} applies.
\end{enumerate}
\end{theorem}

\begin{proof}
Fix $(m,q)$. The window $W_{m,q}(\cA_{\le N})$ involves finitely many
basis vectors. The bar differential on that finite-dimensional space is
given by a finite matrix whose entries are universal polynomial
expressions in finitely many OPE/$\Ainf$-coefficients. By hypothesis,
those coefficients stabilize for $N \gg 0$, so every matrix entry
stabilizes. This gives window-stability.

Eventual constancy of the quotient structure tensors on every
$(m,q)$-window provides a compatible inverse system of multilinear
operations on all finite quotients. Passing to the inverse limit
defines continuous $\Ainf$-operations on~$\cA$. The $\Ainf$ identities
are polynomial on each quotient and pass to the limit entrywise.
Continuity of the completed bar differential follows because on each
weight quotient only finitely many degrees survive and all matrices are
stable. Theorem~\ref{thm:completed-bar-cobar-strong} then applies.
\end{proof}

\begin{theorem}[Completed twisting representability; \ClaimStatusProvedHere]
\label{thm:completed-twisting-representability}
\index{twisting morphism!completed|textbf}
For $\cA \in \CompCl(\Fft)$ and $C \in \Cstr$, define $\Twcts(C,\cA)$
to be the set of degree-$(-1)$ continuous twisting morphisms from~$C$
to~$\cA$. Then
\[
\operatorname{Hom}_{\mathrm{alg}}(\widehat\Omega^{\mathrm{ch}}C,\,\cA)
\;\cong\; \Twcts(C,\cA)
\;\cong\; \operatorname{Hom}_{\mathrm{coalg}}(C,\,
\widehat{\bar B}^{\mathrm{ch}}\cA).
\]
\end{theorem}

\begin{proof}
A continuous map into or out of a complete object is equivalent to a
compatible family of maps on finite quotients/windows. On each finite
quotient, ordinary bar/cobar representability holds by the finite-stage
regime. Passing to compatible inverse limits preserves the representing
property because the windows are finite and stabilized.
\end{proof}

\begin{theorem}[MC-twisting closure of the completion closure;
\ClaimStatusConditional]
\label{thm:mc-twisting-closure}
\index{Maurer--Cartan!twisting closure}
Let $\cA$ be a strong completion tower and let
$\alpha \in F^1\cA^1$ be a continuous Maurer--Cartan element. Then the
twisted algebra $\cA_\alpha$ is again a strong completion tower, and
\[
\widehat\Omega^{\mathrm{ch}}\bigl(\widehat{\bar B}^{\mathrm{ch}}
(\cA_\alpha)\bigr) \simeq \cA_\alpha.
\]
\end{theorem}

\begin{proof}
The twisted operations are
\[
\mu_n^\alpha(a_1,\dots,a_n)
=
\sum_{i_0,\ldots,i_n \ge 0}
\mu_{n+i_0+\cdots+i_n}
(\alpha^{\otimes i_0},a_1,\alpha^{\otimes i_1},
\ldots,a_n,\alpha^{\otimes i_n}).
\]
Since $\alpha \in F^1\cA$, each inserted copy of~$\alpha$ raises
filtration degree by at least one. Modulo $F^{N+1}$, only multiindices
with $i_0+\cdots+i_n \le N$ can contribute. Thus the twist is
quotientwise finite and continuous, preserves the strong
completion-tower axioms, and Theorem~\ref{thm:completed-bar-cobar-strong}
applies.
\end{proof}

\begin{theorem}[Uniform PBW bridge from MC1 to MC4; \ClaimStatusConditional]
\label{thm:uniform-pbw-bridge}
\index{PBW!bridge to MC4}
Let $\cA = \varprojlim_N \cA_{\le N}$ be a strong completion tower, and
assume each finite stage carries a PBW filtration~$P$. Suppose:
\begin{enumerate}
\item for each quotient level~$q$, the associated-graded quotients
 $\operatorname{gr}_P(\cA_{\le N}/F^{q+1})$ stabilize for $N \gg 0$;
 and
\item for each~$q$, the bar spectral sequence of $\cA_{\le N}/F^{q+1}$
 degenerates at a page $E^{r(q)}$ independent of sufficiently
 large~$N$.
\end{enumerate}
Then the quotient bar cohomology stabilizes for each~$q$, and completed
bar-cobar duality holds for~$\cA$.
\end{theorem}

\begin{proof}
For fixed~$q$, the PBW filtration on the quotient bar complex is finite.
Stabilization of the associated graded gives stabilization of the $E^0$
and $E^1$ pages for large~$N$. Uniform degeneration at page $E^{r(q)}$
implies stabilization of~$E^\infty$. Since the filtration is finite,
stabilized $E^\infty$ implies stabilized total cohomology of the quotient
bar complex. Theorem~\ref{thm:completed-bar-cobar-strong} then applies
quotientwise.
\end{proof}

\begin{remark}[Completion closure: the correct MC4 envelope]
\label{rem:completion-closure-envelope}
Theorem~\ref{thm:completed-bar-cobar-strong} together with the
coefficient-stability criterion
(Theorem~\ref{thm:coefficient-stability-criterion}) and the uniform PBW
bridge (Theorem~\ref{thm:uniform-pbw-bridge}) give the finite-window
envelope for MC4. Continuity, pro-nilpotence, Mittag--Leffler control,
and inverse-limit identification are consequences of the strong
completion-tower axiom.

The genuine example-specific residue for $\mathcal{W}_\infty$ and
Yangian towers is now reduced to: \emph{verify coefficient stabilization
on finite windows and identify the inverse-limit target}. For positive
towers (weight-graded by a strict positive grading), the
weight-cutoff criterion of Proposition~\ref{prop:mc4-weight-cutoff} makes
this automatic. For towers with weight-zero resonance, the
finite-resonance criterion of
Theorem~\ref{thm:resonance-filtered-bar-cobar} applies once the
resonance rank is shown to be finite.
\end{remark}

%% End completion closure theory

\begin{proposition}[Factorization realization criterion for
\texorpdfstring{$W_\infty$}{W_infty}; \ClaimStatusConditional]
\label{prop:winfty-factorization-package}
Let $\mathcal{W}^{\mathrm{fact}}_\infty$ be a
factorization-theoretic completion of the principal $W$-tower. Suppose:
\begin{enumerate}
\item $\mathcal{W}^{\mathrm{fact}}_\infty$ carries a separated complete
 conformal-weight filtration;
\item its finite quotients recover the proved principal finite-type
 stages $W_N$;
\item the induced bar-cobar data on those quotients agree with the
 standard finite-stage bar complexes; and
\item the resulting comparison map
 \[
 \mathcal{W}^{\mathrm{fact}}_\infty \xrightarrow{\sim} W_\infty
 \]
 identifies the factorization-theoretic completion with the standard
 principal-stage completed M-level package.
\end{enumerate}
then $\mathcal{W}^{\mathrm{fact}}_\infty$ is a principal-stage
compatible $W_\infty$ target in the sense of
Definition~\ref{def:winfty-principal-stage-compatible}. The existence
input is Theorem~\ref{thm:winfty-factorization-kd}: the factorization
target is constructed there via $\Einf$-sectorwise convergence at each
finite stage, with the DS--KD intertwining
(Theorem~\ref{thm:ds-koszul-intertwine}) providing the
finite-stage identifications and the Mittag--Leffler property
(Corollary~\ref{cor:winfty-standard-mc4-package}) controlling
the inverse limit.
\end{proposition}

\begin{remark}[Reduction route for the
\texorpdfstring{$W_\infty$}{W_infty} factorization target]
\label{rem:winfty-factorization-route}
The reduction chain from
Proposition~\ref{prop:winfty-factorization-package} proceeds through
higher-spin ideals
(Proposition~\ref{prop:winfty-higher-spin-ideal-criterion}),
spin-triangular OPE closure
(Proposition~\ref{prop:winfty-spin-triangular-ideals}),
coefficient-level DS matching
(Proposition~\ref{prop:winfty-ds-local-coefficient-criterion}),
residue identity
(Proposition~\ref{prop:winfty-ds-residue-identity-criterion}),
generator seeds
(Proposition~\ref{prop:winfty-ds-generator-seed},
Corollary~\ref{cor:winfty-ds-finite-seed-set}),
and the explicit stage-$3$ $W_3$ packet
(Proposition~\ref{prop:winfty-ds-stage3-explicit-packet}).
The first genuinely new higher-spin packet is stage~$4$.
\end{remark}

\begin{proposition}[Higher-spin ideal criterion for principal-stage
\texorpdfstring{$W_\infty$}{W_infty} quotients; \ClaimStatusConditional]
\label{prop:winfty-higher-spin-ideal-criterion}
Let $\mathcal{F}_\infty$ be a separated complete factorization algebra
on a curve, topologically generated by fields
$W^{(s)}$ of conformal weights $s \ge 2$. For each $N \ge 2$, let
$I_{>N}\mathcal{F}_\infty$ denote the closed factorization ideal
generated by the fields $W^{(s)}$ with $s > N$. Assume:
\begin{enumerate}
\item the differential, translation operator, and factorization
 products preserve each ideal $I_{>N}\mathcal{F}_\infty$;
\item modulo $I_{>N}\mathcal{F}_\infty$, the OPE of the generators
 $W^{(2)},\dots,W^{(N)}$ closes among those same generators;
\item the quotient factorization algebra
 $\mathcal{F}_{\le N}:=\mathcal{F}_\infty/I_{>N}\mathcal{F}_\infty$
 admits a chiral envelope or equivalent dg model quasi-isomorphic to
 the principal stage $W_N$; and
\item the quotient maps
 $\mathcal{F}_{\le N+1}\twoheadrightarrow \mathcal{F}_{\le N}$
 intertwine the principal-stage truncation tower.
\end{enumerate}
Then the quotients $\{\mathcal{F}_{\le N}\}_{N \ge 2}$ define a
principal-stage quotient system in the sense of
Definition~\ref{def:winfty-quotient-system}. Consequently,
Proposition~\ref{prop:winfty-factorization-envelope-criterion},
Proposition~\ref{prop:winfty-quotient-system-criterion}, and
Corollary~\ref{cor:winfty-hlevel-comparison-criterion} apply.

In particular, the remaining $\mathcal{W}_\infty$ construction problem
is to prove that the higher-spin generators above stage $N$ define
stable factorization ideals whose quotients recover the proved
principal stages.
\end{proposition}

\begin{proof}
Assumptions~\textup{(1)--(4)} supply closed factorization ideals, closed OPE, finite-stage identification with
$W_N$, and compatibility with the
principal-stage tower. This is exactly the data required by
Definition~\ref{def:winfty-quotient-system}. The stated comparison
results then follow from
Proposition~\ref{prop:winfty-factorization-envelope-criterion},
Proposition~\ref{prop:winfty-quotient-system-criterion}, and
Corollary~\ref{cor:winfty-hlevel-comparison-criterion}.
\end{proof}

\begin{proposition}[Spin-triangular OPE criterion for the
\texorpdfstring{$W_\infty$}{W_infty} factorization ideals;
\ClaimStatusConditional]
\label{prop:winfty-spin-triangular-ideals}
Let $\mathcal{F}_\infty$ be a separated complete factorization algebra
on a curve, topologically generated by fields $W^{(s)}$ of conformal
weight $s \ge 2$. Assume:
\begin{enumerate}
\item for every pair of generators, the coefficient of the pole
 $(z-w)^{-n}$ in the local OPE
 $W^{(s)}(z)W^{(t)}(w)$ is a linear combination of derivatives of
 fields $W^{(u)}$ with $u \le s+t-n$;
\item the translation operator and the residue/bar operations are
 obtained from these OPE coefficients by the standard factorization
 formulas, so no operation increases the maximal conformal weight
 appearing in an expression; and
\item for each $N$, after quotienting by the fields $W^{(s)}$ with
 $s>N$, the induced OPE formulas on
 $W^{(2)},\dots,W^{(N)}$ agree with the principal finite-type stage
 $W_N$.
\end{enumerate}
Then the closed factorization ideals
\[
I_{>N}\mathcal{F}_\infty
= \langle W^{(s)} \mid s>N\rangle_{\mathrm{fact}}
\]
form a compatible principal-stage quotient system. In particular,
Proposition~\ref{prop:winfty-higher-spin-ideal-criterion} applies, and
the resulting quotients recover the principal finite-type stages.
\end{proposition}

\begin{proof}
Spin-triangularity (1) closes the quotient OPE on generators $\le N$; (2) preserves the ideal under all operations; (3) identifies the quotient with $W_N$. Nesting of ideals gives compatibility for varying $N$.
\end{proof}

\begin{proposition}[Coefficient-level DS criterion for principal-stage
\texorpdfstring{$W_\infty$}{W_infty} quotients;
\ClaimStatusConditional]
\label{prop:winfty-ds-coefficient-criterion}
Fix $N \ge 2$ and let
$\mathcal{F}_{\le N}:=\mathcal{F}_\infty/I_{>N}\mathcal{F}_\infty$
be a quotient from
Proposition~\ref{prop:winfty-spin-triangular-ideals}.
If the quotient is strongly generated by $W^{(2)},\dots,W^{(N)}$ with
OPE coefficients matching the principal DS formulas for $W_N$, no
additional local relations, and induced bar operations coinciding with
the finite-stage package, then $\mathcal{F}_{\le N}$ identifies with
$W_N$ and
Proposition~\ref{prop:winfty-higher-spin-ideal-criterion} applies.
\end{proposition}

\begin{proof}
The hypotheses supply the same generators, OPE, and bar data as $W_N$; the reconstruction principle gives the identification.
\end{proof}

\begin{proposition}[Local-coefficient criterion for principal-stage
\texorpdfstring{$W_\infty$}{W_infty} quotients;
\ClaimStatusConditional]
\label{prop:winfty-ds-local-coefficient-criterion}
Fix $N \ge 2$ and let
$\mathcal{F}_{\le N}:=\mathcal{F}_\infty/I_{>N}\mathcal{F}_\infty$
be as in
Proposition~\ref{prop:winfty-spin-triangular-ideals}. If the local
OPE coefficients $C_{s,t}^{u;m,n}(N)$ agree with the principal DS data,
translation descendants generate all local relations, and bar operations
match, then Proposition~\ref{prop:winfty-ds-coefficient-criterion}
applies and the quotient identifies with $W_N$.
\end{proposition}

\begin{proof}
Immediate from Proposition~\ref{prop:winfty-ds-coefficient-criterion}.
\end{proof}

\begin{proposition}[Residue-coefficient identity criterion for
principal-stage \texorpdfstring{$W_\infty$}{W_infty} quotients;
\ClaimStatusConditional]
\label{prop:winfty-ds-residue-identity-criterion}
Fix $N \ge 2$ and let
$\mathcal{F}_{\le N}:=\mathcal{F}_\infty/I_{>N}\mathcal{F}_\infty$
be as in Proposition~\ref{prop:winfty-spin-triangular-ideals}.
If $\mathsf{C}^{\mathrm{res}}_{s,t;u;m,n}(N)
= \mathsf{C}^{\mathrm{DS}}_{s,t;u;m,n}(N)$
for every admissible tuple and the quotient OPE and bar operations are
entirely determined by these residue coefficients, then
Proposition~\ref{prop:winfty-ds-local-coefficient-criterion} applies
and the quotient identifies with $W_N$.
\end{proposition}

\begin{proof}
Immediate from Proposition~\ref{prop:winfty-ds-local-coefficient-criterion}.
\end{proof}

\begin{proposition}[Generator-seed criterion for principal-stage
\texorpdfstring{$W_\infty$}{W_infty} residue identities;
\ClaimStatusConditional]
\label{prop:winfty-ds-generator-seed}
Fix $N \ge 2$ and let
\[
\mathcal{F}_{\le N}:=\mathcal{F}_\infty/I_{>N}\mathcal{F}_\infty
\]
be as in
Proposition~\ref{prop:winfty-ds-residue-identity-criterion}. Assume:
\begin{enumerate}
\item $\mathcal{F}_{\le N}$ is strongly generated by
 $W^{(2)},\dots,W^{(N)}$;
\item for each pair of generators $W^{(s)},W^{(t)}$ with
 $2 \le s,t \le N$, the primary residue coefficients
 \[
 \mathsf{C}^{\mathrm{res}}_{s,t;u;0,n}(N)
 \]
 agree with the principal Drinfeld--Sokolov coefficients
 $\mathsf{C}^{\mathrm{DS}}_{s,t;u;0,n}(N)$ for every visible pole order
 and target spin~$u$;
\item translation commutes with the residue extraction formulas, so all
 descendant coefficients are obtained from the primary ones by
 repeated application of $\partial$; and
\item the bar operations are generated from the same residue calculus
 and therefore are determined by those translated coefficient
 families.
\end{enumerate}
Then the full mode-by-mode identities
\[
\mathsf{C}^{\mathrm{res}}_{s,t;u;m,n}(N)
=\mathsf{C}^{\mathrm{DS}}_{s,t;u;m,n}(N)
\]
hold for all admissible indices. Consequently the hypotheses of
Proposition~\ref{prop:winfty-ds-residue-identity-criterion} are
satisfied.
\end{proposition}

\begin{proof}
Primary generator agreement propagates to all descendants by translation (assumption~(3)), and strong generation (assumption~(1)) plus bar recovery (assumption~(4)) identify the full coefficient family, verifying Proposition~\ref{prop:winfty-ds-residue-identity-criterion}.
\end{proof}

\begin{corollary}[Finite primary seed set for principal-stage
\texorpdfstring{$W_\infty$}{W_infty} comparison;
\ClaimStatusConditional]
\label{cor:winfty-ds-finite-seed-set}
Fix $N \ge 2$ and define the finite index set
\[
\mathcal{I}_N:=
\left\{
(s,t,u,n)\ \middle|\
\begin{array}{l}
2 \le s \le t \le N,\\
2 \le u \le \min(N,s+t-1),\\
1 \le n \le s+t-u
\end{array}
\right\}.
\]
Assume the hypotheses of
Proposition~\ref{prop:winfty-ds-generator-seed}. Then it is enough to
check the primary identities
\[
\mathsf{C}^{\mathrm{res}}_{s,t;u;0,n}(N)
=\mathsf{C}^{\mathrm{DS}}_{s,t;u;0,n}(N)
\qquad\text{for all } (s,t,u,n)\in\mathcal{I}_N.
\]
These finitely many checks imply the full residue/OPE identity package
required by
Proposition~\ref{prop:winfty-ds-residue-identity-criterion}.
\end{corollary}

\begin{proof}
Spin-triangularity implies that for fixed $s,t \le N$ only target spins
$u \le s+t-1$ and pole orders $n \le s+t-u$ can appear. Since the
stage quotient discards spins $>N$, the admissible tuples are exactly
those in $\mathcal{I}_N$, hence finite.

Proposition~\ref{prop:winfty-ds-generator-seed} shows that agreement on
these primary generator coefficients propagates by translation to every
descendant coefficient and therefore to the full residue/bar package.
\end{proof}

\begin{corollary}[First principal-stage seed packets for
\texorpdfstring{$W_\infty$}{W_infty} comparison;
\ClaimStatusConditional]
\label{cor:winfty-ds-lowstage-seeds}
The finite seed sets of
Corollary~\ref{cor:winfty-ds-finite-seed-set} begin as follows:
\[
\mathcal{I}_2
=\{(2,2,2,1),(2,2,2,2)\},
\]
and
\begin{align*}
\mathcal{I}_3
= {} & \{(2,2,2,1),(2,2,2,2),(2,2,3,1)\} \\
& \sqcup\{(2,3,2,1),(2,3,2,2),(2,3,2,3),(2,3,3,1),(2,3,3,2)\} \\
& \sqcup\{(3,3,2,1),(3,3,2,2),(3,3,2,3),(3,3,2,4), \\
& \qquad (3,3,3,1),(3,3,3,2),(3,3,3,3)\}.
\end{align*}
In particular, the first genuinely new infinite-generator stage beyond
the proved Virasoro block is the stage-$3$ packet of fifteen
primary coefficients: the two Virasoro coefficients, the single
vanishing $W^{(2)}W^{(2)}\to W^{(3)}$ coefficient, five mixed
$W^{(2)}$/$W^{(3)}$ coefficients, and seven $W^{(3)}$ self-coupling
coefficients.
\end{corollary}

\begin{proof}
For $N=2$, the defining inequalities for $\mathcal{I}_N$ force
$s=t=u=2$ and $1\le n\le 2$, giving exactly the two displayed tuples.
For $N=3$, the admissible ordered pairs of generator spins are
$(2,2)$, $(2,3)$, and $(3,3)$. The inequalities
\[
2\le u\le \min(3,s+t-1),
\qquad
1\le n\le s+t-u
\]
then give the displayed decomposition into the three corresponding
blocks. The final sentence is simply a repackaging of those tuples by
source pair and target spin.
\end{proof}

\begin{proposition}[Incremental interacting packet from stage
\texorpdfstring{$N$}{N} to stage \texorpdfstring{$N{+}1$}{N+1}]
\label{prop:winfty-ds-stage-growth-packet}
Fix $N \ge 3$ and let $\mathcal{I}_N$ and $\mathcal{I}_{N+1}$ be the
seed packets of Corollary~\ref{cor:winfty-ds-finite-seed-set}. Write
\[
\mathcal{I}_{N+1}^{\mathrm{Vir,new}}
:=
\mathcal{I}_{N+1}^{\mathrm{Vir}}\setminus \mathcal{I}_N
=
\{(2,N{+}1,u,n)\in \mathcal{I}_{N+1}\},
\]
and define the incremental interacting packet
\[
\mathcal{J}_{N+1}
:=
\mathcal{I}_{N+1}\setminus
\bigl(\mathcal{I}_N\sqcup \mathcal{I}_{N+1}^{\mathrm{Vir,new}}\bigr).
\]
Assume the hypotheses of
Proposition~\ref{prop:winfty-ds-generator-seed}. Then:
\begin{enumerate}
\item the coefficients indexed by $\mathcal{I}_N$ are exactly the
 stage-$N$ packet identities;
\item the coefficients indexed by
 $\mathcal{I}_{N+1}^{\mathrm{Vir,new}}$ are determined by the exact
 stress-tensor action on the new generator $W^{(N+1)}$;
\item the genuinely new interacting stage-$(N+1)$ packet is therefore
 exactly $\mathcal{J}_{N+1}$, and it decomposes as the disjoint union
 \[
 \mathcal{J}_{N+1}
 =
 \mathcal{J}_{N+1}^{\mathrm{tail}}
 \sqcup
 \mathcal{J}_{N+1}^{\mathrm{mix}}
 \sqcup
 \mathcal{J}_{N+1}^{\mathrm{self}},
 \]
 where
 \[
 \mathcal{J}_{N+1}^{\mathrm{tail}}
 :=
 \{(s,t,N{+}1,n)\in\mathcal{I}_{N+1}\mid 3\le s\le t\le N\},
 \]
 \[
 \mathcal{J}_{N+1}^{\mathrm{mix}}
 :=
 \{(s,N{+}1,u,n)\in\mathcal{I}_{N+1}\mid 3\le s\le N,\ 2\le u\le N{+}1\},
 \]
 \[
 \mathcal{J}_{N+1}^{\mathrm{self}}
 :=
 \{(N{+}1,N{+}1,u,n)\in\mathcal{I}_{N+1}\mid 2\le u\le N{+}1\}.
 \]
\end{enumerate}
Consequently, once the stage-$N$ packet identities on $\mathcal{I}_N$
are known, the stage-$(N+1)$ comparison is equivalent to matching the
coefficients on the incremental interacting packet $\mathcal{J}_{N+1}$.
For $N=3$, this is exactly the stage-$4$ residual packet of
Proposition~\ref{prop:winfty-ds-stage4-residual-packet}.
\ClaimStatusConditional{}
\end{proposition}

\begin{proof}
Since $\mathcal{I}_N\subset \mathcal{I}_{N+1}$, new tuples have at least one index equal to $N{+}1$. Those with $s=2$ form $\mathcal{I}_{N+1}^{\mathrm{Vir,new}}$; the complement $\mathcal{J}_{N+1}$ (with $s\geq 3$) splits exhaustively into the three displayed blocks by whether the new index appears in $u$, $t$, or both.
\end{proof}

\begin{corollary}[Top-pole/parity reduction of the incremental
\texorpdfstring{$W_\infty$}{W_infty} stage-growth packet;
\ClaimStatusConditional]
\label{cor:winfty-ds-stage-growth-top-parity}
Fix $N \ge 3$ and assume the hypotheses of
Proposition~\ref{prop:winfty-ds-stage-growth-packet},
Proposition~\ref{prop:winfty-ds-primary-top-pole}, and
Proposition~\ref{prop:winfty-ds-self-ope-parity}. Let
\[
\mathcal{J}_{N+1}^{\mathrm{top}}
:=
\{(s,t,u,s+t-u)\in \mathcal{J}_{N+1}\},
\]
and remove the odd self-OPE top-pole entries by setting
\[
\mathcal{J}_{N+1}^{\mathrm{red}}
:=
\mathcal{J}_{N+1}^{\mathrm{top}}
\setminus
\bigl\{(s,s,u,2s-u)\in\mathcal{J}_{N+1}^{\mathrm{top}}
\mid 2s-u \text{ odd}\bigr\}.
\]
Then, once the stage-$N$ packet identities on $\mathcal{I}_N$ are
known, the stage-$(N+1)$ comparison is equivalent to the identities on
the reduced incremental packet
$\mathcal{J}_{N+1}^{\mathrm{red}}$.

For $N=3$, one has
\[
\mathcal{J}_4^{\mathrm{red}}
=
\mathcal{J}_4^{\mathrm{par}},
\]
so this is exactly the six-entry stage-$4$ block of
Corollary~\ref{cor:winfty-ds-stage4-parity-packet}. The additional
contraction of that six-entry block to the four higher-spin channels is
controlled by the stage-$4$ Ward-inheritance input.
\end{corollary}

\begin{proof}
Top-pole primaryity (Proposition~\ref{prop:winfty-ds-primary-top-pole}) and odd self-OPE parity (Proposition~\ref{prop:winfty-ds-self-ope-parity}) reduce $\mathcal{J}_{N+1}$ to $\mathcal{J}_{N+1}^{\mathrm{red}}$. For $N=3$ this recovers the stage-$4$ parity packet.
\end{proof}

\begin{corollary}[First reduced incremental packet beyond
\texorpdfstring{$\mathcal{I}_4$}{I4}; \ClaimStatusConditional]
\label{cor:winfty-ds-stage5-reduced-packet}
Assume the hypotheses of
Corollary~\ref{cor:winfty-ds-stage-growth-top-parity} with $N=4$.
Then the first reduced incremental packet beyond the stage-$4$ block is
\[
\mathcal{J}_5^{\mathrm{red}}
=
\Bigl\{
(3,4;5;0,2)
\Bigr\}
\sqcup
\Bigl\{
(3,5;2;0,6),\,
(3,5;3;0,5),\,
(3,5;4;0,4),\,
(3,5;5;0,3)
\Bigr\}
\]
\[
\sqcup
\Bigl\{
(4,5;2;0,7),\,
(4,5;3;0,6),\,
(4,5;4;0,5),\,
(4,5;5;0,4)
\Bigr\}
\sqcup
\Bigl\{
(5,5;2;0,8),\,
(5,5;4;0,6)
\Bigr\}.
\]
In particular,
\[
|\mathcal{J}_5^{\mathrm{red}}|=11,
\]
decomposed as one tail channel, two four-entry mixed blocks, and two
self-coupling channels.
\end{corollary}

\begin{proof}
Enumerate the tail, mixed, and self blocks of $\mathcal{J}_5$ from Proposition~\ref{prop:winfty-ds-stage-growth-packet}; apply top-pole primaryity and odd self-OPE parity to each. The surviving entries are exactly the displayed eleven tuples.
\end{proof}

\begin{proposition}[Primary top-pole criterion for generator seed
packets; \ClaimStatusConditional]
\label{prop:winfty-ds-primary-top-pole}
Fix $N \ge 2$ and assume the hypotheses of
Proposition~\ref{prop:winfty-ds-generator-seed}. Assume that
for each $2 \le u \le N$ the strong generator $W^{(u)}$ is primary of
conformal weight $u$ with respect to the Virasoro field
$T:=W^{(2)}$, and that the local OPE is decomposed into conformal
families of these primary generators and their translation
descendants.

Then for every visible triple $(s,t,u)$ with
\[
2 \le s \le t \le N,
\qquad
2 \le u \le \min(N,s+t-1),
\]
the primary generator $W^{(u)}$ can appear only at the top pole order
\[
n_{\max}(s,t;u):=s+t-u.
\]
Equivalently,
\[
\mathsf{C}^{\mathrm{res}}_{s,t;u;0,n}(N)
\;=\;
\mathsf{C}^{\mathrm{DS}}_{s,t;u;0,n}(N)
\;=\;0
\qquad\text{for every } 1\le n < s+t-u.
\]
Consequently the primary seed set of
Corollary~\ref{cor:winfty-ds-finite-seed-set} reduces further to the
top-pole subset
\[
\mathcal{I}_N^{\mathrm{top}}
:=\{(s,t,u,s+t-u)\in \mathcal{I}_N\}.
\]
\end{proposition}

\begin{proof}
Because $W^{(s)}$, $W^{(t)}$, and $W^{(u)}$ are primary fields of
weights $s$, $t$, and $u$, conformal covariance organizes the
$W^{(u)}$-family inside the OPE
$W^{(s)}(z)W^{(t)}(w)$ as one primary coefficient multiplying
$W^{(u)}(w)$ at pole order $s+t-u$, together with the translation
descendants $\partial^m W^{(u)}(w)$ at the lower poles
$s+t-u-m$. The undifferentiated primary field therefore occurs only at
the top pole order $n=s+t-u$.

In the notation of
Proposition~\ref{prop:winfty-ds-generator-seed}, the coefficients
\[
\mathsf{C}^{\mathrm{res}}_{s,t;u;0,n}(N),
\qquad
\mathsf{C}^{\mathrm{DS}}_{s,t;u;0,n}(N)
\]
record the coefficients of the undifferentiated primary generator
$W^{(u)}$ itself. Hence those coefficients vanish for all lower pole
orders $n<s+t-u$, both on the residue side and on the principal
Drinfeld--Sokolov side. This proves the displayed vanishing and
identifies the reduced top-pole seed set $\mathcal{I}_N^{\mathrm{top}}$.
\end{proof}

\begin{proposition}[Odd top-pole vanishing for identical even
generators; \ClaimStatusConditional]
\label{prop:winfty-ds-self-ope-parity}
Fix $N \ge 2$ and assume the hypotheses of
Proposition~\ref{prop:winfty-ds-primary-top-pole}. Assume further
that the strong generators $W^{(2)},\dots,W^{(N)}$ are even fields.
Then for every visible self-coupling triple $(s,s,u)$ one has
\[
\mathsf{C}^{\mathrm{res}}_{s,s;u;0,\,2s-u}(N)
=
\mathsf{C}^{\mathrm{DS}}_{s,s;u;0,\,2s-u}(N)
=0
\qquad\text{whenever } 2s-u \text{ is odd}.
\]
Equivalently, for identical even generators the top-pole primary
coefficient can be nonzero only when the top pole order is even.
\end{proposition}

\begin{proof}
Let $A:=W^{(s)}$ and set $m:=2s-u-1$, so the top-pole primary
coefficient of the target generator $W^{(u)}$ in the self-OPE
$A(z)A(w)$ is the $A_{(m)}A$ term. Because $A$ is even, the
vertex-algebra skew-symmetry formula gives
\[
A_{(m)}A
=
\sum_{j\ge 0}\frac{(-1)^{m+j+1}}{j!}\,\partial^j(A_{(m+j)}A).
\]
By Proposition~\ref{prop:winfty-ds-primary-top-pole}, the
undifferentiated primary generator channel occurs only at the top pole,
so the terms with $j>0$ do not contribute to the same primary
coefficient. Therefore the primary coefficient satisfies
\[
A_{(m)}A = (-1)^{m+1}A_{(m)}A = (-1)^{2s-u}A_{(m)}A.
\]
If $2s-u$ is odd, this becomes $A_{(m)}A=-A_{(m)}A$, hence the
coefficient vanishes. The same skew-symmetry identity holds on the
principal Drinfeld--Sokolov side, so both displayed coefficients
vanish.
\end{proof}

\begin{proposition}[Stage-\texorpdfstring{$3$}{3} principal packet from the explicit
\texorpdfstring{$W_3$}{W3} OPE; \ClaimStatusConditional]
\label{prop:winfty-ds-stage3-explicit-packet}
Fix $N=3$ and write
\[
T:=W^{(2)},
\qquad
W:=W^{(3)}.
\]
For the principal stage $W_3$, the primary residue coefficients on the
seed set $\mathcal{I}_3$ are completely determined by the explicit
$W_3$ OPE package (Definition~\ref{def:w3-algebra};
Theorem~\ref{thm:w-w-ope-complete} for the complete $W$--$W$
expansion).
More precisely, among the fifteen tuples in $\mathcal{I}_3$, only the
three primary coefficients
\[
\mathsf{C}^{\mathrm{DS}}_{2,2;2;0,2}(3)=2,
\qquad
\mathsf{C}^{\mathrm{DS}}_{2,3;3;0,2}(3)=3,
\qquad
\mathsf{C}^{\mathrm{DS}}_{3,3;2;0,4}(3)=2
\]
are nonzero; the remaining twelve coefficients in $\mathcal{I}_3$
vanish. Consequently the stage-$3$ comparison problem is equivalent to
matching these three numbers and the twelve forced vanishings, after
which translation recovers all descendant coefficients.
\end{proposition}

\begin{proof}
The seed set $\mathcal{I}_3$ was written explicitly in
Corollary~\ref{cor:winfty-ds-lowstage-seeds}. We read its coefficients
off from the explicit $W_3$ OPE package.

For the pair $(T,T)$, the only primary generator appearing in the
Virasoro OPE is $T$ itself at pole order $2$, with coefficient $2$.
The pole-order-$1$ term is $\partial T$, hence a descendant, and there
is no primary target of spin $3$. Therefore
\[
\mathsf{C}^{\mathrm{DS}}_{2,2;2;0,2}(3)=2,
\qquad
\mathsf{C}^{\mathrm{DS}}_{2,2;2;0,1}(3)=
\mathsf{C}^{\mathrm{DS}}_{2,2;3;0,1}(3)=0.
\]

For the mixed pair $(T,W)$, the primary-field OPE shows that the only
primary target is $W$ at pole order $2$, with coefficient $3$. The
pole-order-$1$ term is $\partial W$, hence a descendant, and no primary
target of spin $2$ appears. Thus
\[
\mathsf{C}^{\mathrm{DS}}_{2,3;3;0,2}(3)=3,
\]
and all four remaining mixed primary coefficients in $\mathcal{I}_3$
vanish.

For the pair $(W,W)$, the complete $W$-$W$ OPE has primary generator
target $T$ only at pole order $4$, again with coefficient $2$. The
terms at pole orders $3,2,1$ are the descendants
$\partial T$, $\partial^2 T$, and $\partial^3 T$, while the
$\Lambda$-terms lie in conformal weight $4$ and therefore outside the
generator seed set $\mathcal{I}_3$. No primary target of spin $3$
appears. Hence
\[
\mathsf{C}^{\mathrm{DS}}_{3,3;2;0,4}(3)=2,
\]
and the other seven $(W,W)$-entries in $\mathcal{I}_3$ vanish.

This accounts for all fifteen tuples in $\mathcal{I}_3$. The final
sentence is then exactly the translation-propagation statement of
Proposition~\ref{prop:winfty-ds-generator-seed}.
\end{proof}

\begin{proposition}[Stage-\texorpdfstring{$4$}{4} residual packet after the proved
\texorpdfstring{$W_3$}{W3} sector]
\label{prop:winfty-ds-stage4-residual-packet}
Fix $N=4$ and let $\mathcal{I}_4$ be the primary seed set of
Corollary~\ref{cor:winfty-ds-finite-seed-set}. Write
\[
\mathcal{I}_4^{\mathrm{Vir}}
:=\{(s,t,u,n)\in\mathcal{I}_4 \mid s=2\},
\qquad
\mathcal{I}_4^{W_3}
:=\{(3,3,u,n)\in\mathcal{I}_4 \mid u\le 3\}.
\]
Then every primary coefficient indexed by
\[
\mathcal{I}_4^{\mathrm{Vir}}\sqcup \mathcal{I}_4^{W_3}
\]
is already determined by the exact stress-tensor action on
$W^{(2)},W^{(3)},W^{(4)}$ together with the explicit stage-$3$
\texorpdfstring{$W_3$}{W3} packet of
Proposition~\ref{prop:winfty-ds-stage3-explicit-packet}. The only
genuinely new stage-$4$ primary coefficients are therefore the residual
packet
\[
\mathcal{J}_4
:=\mathcal{I}_4\setminus
\bigl(\mathcal{I}_4^{\mathrm{Vir}}\sqcup \mathcal{I}_4^{W_3}\bigr),
\]
which has cardinality $29$ and decomposes as
\[
\mathcal{J}_4
=\{(3,3,4,1),(3,3,4,2)\}
\sqcup \mathcal{J}_4^{3,4}
\sqcup \mathcal{J}_4^{4,4},
\]
where
\begin{align*}
\mathcal{J}_4^{3,4}
={} & \{(3,4,2,1),(3,4,2,2),(3,4,2,3),(3,4,2,4),(3,4,2,5)\} \\
& \sqcup \{(3,4,3,1),(3,4,3,2),(3,4,3,3),(3,4,3,4)\} \\
& \sqcup \{(3,4,4,1),(3,4,4,2),(3,4,4,3)\},
\end{align*}
and
\begin{align*}
\mathcal{J}_4^{4,4}
={} & \{(4,4,2,1),(4,4,2,2),(4,4,2,3),(4,4,2,4),(4,4,2,5),(4,4,2,6)\} \\
& \sqcup \{(4,4,3,1),(4,4,3,2),(4,4,3,3),(4,4,3,4),(4,4,3,5)\} \\
& \sqcup \{(4,4,4,1),(4,4,4,2),(4,4,4,3),(4,4,4,4)\}.
\end{align*}
Consequently, the stage-$4$ comparison problem is equivalent to
matching these $29$ coefficients. Once that packet agrees with the
principal Drinfeld--Sokolov coefficients, translation again recovers
all descendant coefficients.
\ClaimStatusConditional{}
\end{proposition}

\begin{proof}
All source-$2$ tuples lie in $\mathcal{I}_4^{\mathrm{Vir}}$ (stress-tensor action); the $(3,3)$ tuples with $u\le 3$ lie in $\mathcal{I}_4^{W_3}$ (stage-$3$ packet). Enumerating the remaining admissible tuples for $(3,3,4)$, $(3,4)$, and $(4,4)$ gives the displayed decomposition with cardinality $2+12+15=29$.
\end{proof}

\begin{computation}[Stage-\texorpdfstring{$4$}{4} seed set and residual packet check]
\label{comp:winfty-stage4-seed-check}
\index{MC4!W-infinity stage-4 check}
The seed set cardinalities and residual decomposition of
Proposition~\ref{prop:winfty-ds-stage4-residual-packet} have been
checked computationally
(\texttt{compute/\allowbreak lib/\allowbreak w4\_stage4\_coefficients.py},
$74$~tests):
\begin{itemize}
\item $|\mathcal{I}_4| = 54$,
 $|\mathcal{I}_4^{\mathrm{Vir}}| = 18$,
 $|\mathcal{I}_4^{W_3}| = 7$,
 $|\mathcal{J}_4| = 29$;
\item Residual decomposition: $(3,3)$-tail of size~$2$,
 $\mathcal{J}_4^{3,4}$ of size~$12$ (sub-decomposition $5+4+3$),
 $\mathcal{J}_4^{4,4}$ of size~$15$ (sub-decomposition $6+5+4$);
\item Stage-$3$ packet reproduced:
 $\mathsf{C}^{\mathrm{DS}}_{2,2;2;0,2} = 2$,
 $\mathsf{C}^{\mathrm{DS}}_{2,3;3;0,2} = 3$,
 $\mathsf{C}^{\mathrm{DS}}_{3,3;2;0,4} = 2$ (three nonzero,
 twelve vanishing);
\item Complementarity sum:
 $c(k) + c(-k-8) = 246$ (constant in~$k$);
\item Universal $T$-coupling pattern:
 $\mathsf{C}_{s,s;2;0,2s-2} = 2$ confirmed for $s = 2, 3$;
 predicted for $s = 4$.
\end{itemize}
\end{computation}

\begin{corollary}[Stage-\texorpdfstring{$4$}{4} top-pole packet after primaryity;
\ClaimStatusConditional]
\label{cor:winfty-ds-stage4-top-pole-packet}
Assume the hypotheses of
Proposition~\ref{prop:winfty-ds-stage4-residual-packet} and
Proposition~\ref{prop:winfty-ds-primary-top-pole}. Then the residual
stage-$4$ packet $\mathcal{J}_4$ reduces further to the top-pole subset
\[
\mathcal{J}_4^{\mathrm{top}}
=\{(3,3,4,2),(3,4,2,5),(3,4,3,4),(3,4,4,3),
(4,4,2,6),(4,4,3,5),(4,4,4,4)\},
\]
which has cardinality $7$. Equivalently, among the $29$ residual
stage-$4$ primary coefficients, only these seven top-pole coefficients
can be nonzero; the remaining twenty-two coefficients in
$\mathcal{J}_4$ vanish by primaryity.

Consequently, the genuinely new principal stage-$4$ comparison problem
is equivalent to matching these seven coefficients with the
Drinfeld--Sokolov data.
\end{corollary}

\begin{proof}
Proposition~\ref{prop:winfty-ds-primary-top-pole} retains only the top pole $n=s+t-u$ for each triple, giving the seven displayed tuples; all lower-pole entries vanish.
\end{proof}

\begin{corollary}[Stage-\texorpdfstring{$4$}{4} parity-compressed packet;
\ClaimStatusConditional]
\label{cor:winfty-ds-stage4-parity-packet}
Assume the hypotheses of
Corollary~\ref{cor:winfty-ds-stage4-top-pole-packet} and
Proposition~\ref{prop:winfty-ds-self-ope-parity}. Then the stage-$4$
top-pole packet reduces further to
\[
\mathcal{J}_4^{\mathrm{par}}
=\{(3,3,4,2),(3,4,2,5),(3,4,3,4),(3,4,4,3),
(4,4,2,6),(4,4,4,4)\},
\]
which has cardinality $6$. Equivalently, among the $7$ top-pole
coefficients of
Corollary~\ref{cor:winfty-ds-stage4-top-pole-packet}, the odd self-OPE
entry $(4,4,3,5)$ vanishes by skew-symmetry. Hence the residual
stage-$4$ comparison problem is equivalent to matching these six
coefficients, together with the forced vanishing of the remaining $23$
entries in the original residual packet $\mathcal{J}_4$.
\end{corollary}

\begin{proof}
Among the seven top-pole entries, the only odd self-OPE top pole is $(4,4,3,5)$; Proposition~\ref{prop:winfty-ds-self-ope-parity} kills it, leaving six.
\end{proof}

\begin{corollary}[Stage-\texorpdfstring{$4$}{4} packet as three local OPE blocks;
\ClaimStatusConditional]
\label{cor:winfty-ds-stage4-ope-blocks}
Assume the hypotheses of
Corollary~\ref{cor:winfty-ds-stage4-parity-packet}. Then the residual
principal stage-$4$ comparison problem is equivalent to matching the
six coefficients in the three top-pole OPE blocks
\begin{align*}
W^{(3)}(z)W^{(3)}(w)
&\sim
\frac{c_{334}\,W^{(4)}(w)}{(z-w)^2}
\quad\text{mod the proved $W_3$ sector and descendants},\\
W^{(3)}(z)W^{(4)}(w)
&\sim
\frac{c_{342}\,W^{(2)}(w)}{(z-w)^5}
+\frac{c_{343}\,W^{(3)}(w)}{(z-w)^4}
+\frac{c_{344}\,W^{(4)}(w)}{(z-w)^3}\\
&\qquad
\quad\text{mod descendants},\\
W^{(4)}(z)W^{(4)}(w)
&\sim
\frac{c_{442}\,W^{(2)}(w)}{(z-w)^6}
+\frac{c_{444}\,W^{(4)}(w)}{(z-w)^4}\\
&\qquad
\quad\text{mod the proved Virasoro sector and descendants}.
\end{align*}
More precisely,
\[
c_{334}=\mathsf{C}_{3,3;4;0,2},
\qquad
(c_{342},c_{343},c_{344})
\!=\!
\bigl(
\mathsf{C}_{3,4;2;0,5},
\mathsf{C}_{3,4;3;0,4},
\mathsf{C}_{3,4;4;0,3}
\bigr),
\]
and
\[
(c_{442},c_{444})
=(\mathsf{C}_{4,4;2;0,6},\mathsf{C}_{4,4;4;0,4}),
\]
where either the residue coefficients
$\mathsf{C}^{\mathrm{res}}$ or the principal Drinfeld--Sokolov
coefficients $\mathsf{C}^{\mathrm{DS}}$ may be used. Equivalently, the
remaining stage-$4$ problem is one new coefficient in the $(3,3)$ block,
three in the mixed $(3,4)$ block, and two in the even $(4,4)$ block.
\end{corollary}

\begin{proof}
Group the six parity-compressed tuples by source pair: $(3,3)$ contributes one, $(3,4)$ three, $(4,4)$ two.
\end{proof}

\begin{corollary}[Stage-\texorpdfstring{$4$}{4} reduction to one mixed block and three
self-coupling scalars; \ClaimStatusConditional]
\label{cor:winfty-ds-stage4-mixed-self-split}
Assume the hypotheses of
Corollary~\ref{cor:winfty-ds-stage4-ope-blocks}. Then the residual
principal stage-$4$ comparison decomposes as:
\begin{enumerate}[label=\textup{(\roman*)}]
\item a self-coupling sector consisting of the three scalar top-pole
 coefficients
 \[
 c_{334},\qquad c_{442},\qquad c_{444};
 \]
\item a mixed higher-spin sector consisting of the single
 $W^{(3)}$-$W^{(4)}$ block
 \[
 (c_{342},c_{343},c_{344}).
 \]
\end{enumerate}
Equivalently, the only genuinely mixed stage-$4$ higher-spin data are
the three coefficients in the $W^{(3)}(z)W^{(4)}(w)$ block; the
remaining three live coefficients come from the self-couplings
$W^{(3)}$-$W^{(3)}$ and $W^{(4)}$-$W^{(4)}$.
\end{corollary}

\begin{proof}
Regroup: $(3,3)$ and $(4,4)$ give the self-coupling sector; $(3,4)$ is the mixed sector.
\end{proof}

\begin{proposition}[Mixed top-pole swap parity for even generators;
\ClaimStatusConditional]
\label{prop:winfty-ds-mixed-top-pole-swap}
Fix $N \ge 2$ and assume the hypotheses of
Proposition~\ref{prop:winfty-ds-primary-top-pole}. Assume further
that the strong generators $W^{(2)},\dots,W^{(N)}$ are even fields.
For every visible mixed triple $(s,t,u)$ with $2 \le s < t \le N$ and
$u \le \min(N,s+t-1)$, write
\[
\mathsf{C}^{\mathrm{res}}_{s,t;u}(N)
:=
\mathsf{C}^{\mathrm{res}}_{s,t;u;0,\,s+t-u}(N),
\qquad
\mathsf{C}^{\mathrm{DS}}_{s,t;u}(N)
:=
\mathsf{C}^{\mathrm{DS}}_{s,t;u;0,\,s+t-u}(N)
\]
for the top-pole primary coefficients. Then the reversed-order
coefficients satisfy
\[
\mathsf{C}^{\mathrm{res}}_{t,s;u}(N)
=
(-1)^{s+t-u}\mathsf{C}^{\mathrm{res}}_{s,t;u}(N),
\qquad
\mathsf{C}^{\mathrm{DS}}_{t,s;u}(N)
=
(-1)^{s+t-u}\mathsf{C}^{\mathrm{DS}}_{s,t;u}(N).
\]
Equivalently, for even primary generators the top-pole mixed
$W^{(u)}$-channel is swap-even when $s+t-u$ is even and swap-odd when
$s+t-u$ is odd.
\end{proposition}

\begin{proof}
The vertex-algebra skew-symmetry formula for even fields, combined with top-pole primaryity (Proposition~\ref{prop:winfty-ds-primary-top-pole}) to eliminate $j>0$ terms, gives $\mathsf{C}_{s,t;u}^{\mathrm{res}}=(-1)^{s+t-u}\mathsf{C}_{t,s;u}^{\mathrm{res}}$. The same holds on the DS side.
\end{proof}

\begin{corollary}[Stage-\texorpdfstring{$4$}{4} mixed block split by swap parity;
\ClaimStatusConditional]
\label{cor:winfty-ds-stage4-mixed-swap-parity}
Assume the hypotheses of
Corollary~\ref{cor:winfty-ds-stage4-mixed-self-split} and
Proposition~\ref{prop:winfty-ds-mixed-top-pole-swap}. Write
\[
W^{(4)}(z)W^{(3)}(w)
\sim
\frac{\widetilde c_{432}\,W^{(2)}(w)}{(z-w)^5}
+\frac{\widetilde c_{433}\,W^{(3)}(w)}{(z-w)^4}
+\frac{\widetilde c_{434}\,W^{(4)}(w)}{(z-w)^3}
\quad\text{mod descendants}
\]
for the reversed-order mixed block. Then
\[
\widetilde c_{432}=-c_{342},
\qquad
\widetilde c_{433}=c_{343},
\qquad
\widetilde c_{434}=-c_{344}.
\]
Equivalently, the remaining mixed stage-$4$ problem splits into one
swap-even channel $c_{343}$ and two swap-odd channels
$c_{342},c_{344}$.
\end{corollary}

\begin{proof}
Apply Proposition~\ref{prop:winfty-ds-mixed-top-pole-swap} to $(3,4,u)$ for $u=2,3,4$; the signs $(-1)^{7-u}$ give the displayed identities.
\end{proof}

\begin{proposition}[Formal mixed Virasoro-target vanishing under a
normalized two-point package; \ClaimStatusProvedHere]
\label{prop:winfty-formal-mixed-virasoro-zero}
Fix $N \ge 3$ and consider a stage-$N$ local OPE calculus with visible
primary generators $W^{(2)},\dots,W^{(N)}$, where
$T:=W^{(2)}$ is the Virasoro field. Assume:
\begin{enumerate}[label=\textup{(\roman*)}]
\item each visible generator $W^{(u)}$ has conformal weight $u$, the
 local OPE is decomposed into conformal families of these primary
 generators and their translation descendants, and the Virasoro Ward
 identity holds for insertions of $T$ against each visible generator;
\item for every mixed pair $2 \le s < t \le N$, the vacuum two-point
 function vanishes:
 \[
 \langle W^{(s)}(z)W^{(t)}(0)\rangle=0;
 \]
\item the stress tensor is normalized by
 \[
 \langle T(z)T(0)\rangle=\frac{c/2}{z^4}
 \qquad\text{with } c\neq 0.
 \]
\end{enumerate}
Let
\[
\mathsf{C}_{s,t;u;0,n}(N)
\]
denote the undifferentiated primary coefficient of the target
$W^{(u)}$ at pole order $n$ in this local OPE calculus. Then for every
visible mixed pair $2 \le s < t \le N$, the top-pole stress-tensor
coefficient vanishes:
\[
\mathsf{C}_{s,t;2;0,\,s+t-2}(N)=0.
\]
Equivalently, under this normalized two-point package the
undifferentiated $W^{(2)}=T$ channel is absent from the top pole of
every mixed OPE $W^{(s)}(z)W^{(t)}(w)$ with $s\ne t$.
\end{proposition}

\begin{proof}
Mixed-weight orthogonality (assumption~\textup{(ii)}) makes $\langle T(x)A(z)B(0)\rangle = 0$. The top-pole $T$-channel contributes $\frac{c}{2}\mathsf{C}_{s,t;2;0,s+t-2}$ to the $x^{-4}z^{-(s+t-2)}$ coefficient; since $c\neq 0$, the coefficient vanishes.
\end{proof}

\begin{proposition}[Principal Drinfeld--Sokolov vanishing of the mixed
Virasoro target; \ClaimStatusConditional]
\label{prop:winfty-ds-mixed-virasoro-ds-zero}
Fix $N \ge 3$ and consider the principal Drinfeld--Sokolov stage
$\mathcal{W}_N$ at a generic level with nonzero central charge $c$.
Assume the Zamolodchikov normalization and mixed-weight orthogonality
package of Theorem~\ref{thm:wn-obstruction}. Then for every visible
mixed pair $2 \le s < t \le N$, the top-pole stress-tensor coefficient
vanishes:
\[
\mathsf{C}^{\mathrm{DS}}_{s,t;2;0,\,s+t-2}(N)=0.
\]
Equivalently, on the principal Drinfeld--Sokolov side the
undifferentiated $W^{(2)}=T$ channel is absent from the top pole of
every mixed OPE $W^{(s)}(z)W^{(t)}(w)$ with $s \ne t$.
\end{proposition}

\begin{proof}
Apply Proposition~\ref{prop:winfty-formal-mixed-virasoro-zero} to the
principal Drinfeld--Sokolov stage. Theorem~\ref{thm:wn-obstruction}
supplies the required normalized two-point package: the visible
generators are primary, the mixed two-point functions vanish for
distinct conformal weights, and
$\langle T(z)T(0)\rangle=\frac{c/2}{z^4}$ with $c\neq 0$ at generic
level.
\end{proof}

\begin{corollary}[Stage-\texorpdfstring{$4$}{4} mixed block as one vanishing channel and a
parity pair; \ClaimStatusConditional]
\label{cor:winfty-ds-stage4-mixed-two-channel}
Assume the hypotheses of
Corollary~\ref{cor:winfty-ds-stage4-mixed-swap-parity} and
Proposition~\ref{prop:winfty-ds-mixed-virasoro-ds-zero}. Then on the
principal Drinfeld--Sokolov side one has
\[
\mathsf{C}^{\mathrm{DS}}_{3,4;2;0,5}=0,
\qquad
\mathsf{C}^{\mathrm{DS}}_{4,3;2;0,5}=0.
\]
Consequently, the stage-$4$ mixed comparison problem reduces to:
\begin{enumerate}[label=\textup{(\roman*)}]
\item proving the corresponding residue-side vanishing in the
 $W^{(2)}$ target channel;
\item matching the remaining swap-even coefficient
 $\mathsf{C}_{3,4;3;0,4}$;
\item matching the remaining swap-odd coefficient
 $\mathsf{C}_{3,4;4;0,3}$.
\end{enumerate}
Equivalently, the principal mixed stage-$4$ data consist of one
proved-zero $W^{(2)}$ target channel together with a
swap-even / swap-odd pair targeting $W^{(3)}$ and $W^{(4)}$.
\end{corollary}

\begin{proof}
Proposition~\ref{prop:winfty-ds-mixed-virasoro-ds-zero} gives the first vanishing; swap parity (Corollary~\ref{cor:winfty-ds-stage4-mixed-swap-parity}) gives the second.
\end{proof}

\begin{proposition}[Formal self-coupling stress-tensor coefficient
under a normalized two-point package; \ClaimStatusProvedHere]
\label{prop:winfty-formal-self-t-coefficient}
Fix $N \ge 2$ and consider a stage-$N$ local OPE calculus with visible
primary generators $W^{(2)},\dots,W^{(N)}$, where
$T:=W^{(2)}$ is the Virasoro field. Assume:
\begin{enumerate}[label=\textup{(\roman*)}]
\item each visible generator $W^{(u)}$ has conformal weight $u$, the
 local OPE is decomposed into conformal families of these primary
 generators and their translation descendants, and the Virasoro Ward
 identity holds for insertions of $T$ against each visible generator;
\item for every visible generator $W^{(s)}$ one has the normalized
 two-point function
 \[
 \langle W^{(s)}(z)W^{(s)}(0)\rangle=\frac{c/s}{z^{2s}}
 \qquad\text{with } c\neq 0;
 \]
\item the stress tensor is normalized by
 \[
 \langle T(z)T(0)\rangle=\frac{c/2}{z^4}.
 \]
\end{enumerate}
Let
\[
\mathsf{C}_{s,s;u;0,n}(N)
\]
denote the undifferentiated primary coefficient of the target
$W^{(u)}$ at pole order $n$ in this local OPE calculus. Then for every
visible generator $W^{(s)}$ with $2 \le s \le N$, the top-pole
stress-tensor coefficient in the self-OPE is
\[
\mathsf{C}_{s,s;2;0,\,2s-2}(N)=2.
\]
Equivalently, under this normalized two-point package the
undifferentiated $T=W^{(2)}$ channel appears in every self-OPE
$W^{(s)}(z)W^{(s)}(w)$ with the universal coefficient $2$.
\end{proposition}

\begin{proof}
Write $N_s:=c/s$. The Virasoro Ward identity for $\langle T(x)A(z)A(0)\rangle$ with $A:=W^{(s)}$ has $x^{-4}z^{-2s+2}$ coefficient equal to $sN_s=c$. The only undifferentiated $T$-contribution at the top pole $2s-2$ gives $\mathsf{C}_{s,s;2;0,2s-2}\cdot c/2$. Matching and dividing by $c\neq 0$ yields $\mathsf{C}_{s,s;2;0,2s-2}=2$.
\end{proof}

\begin{proposition}[Formal converse: the universal self-coupling
\texorpdfstring{$T$}{T}-coefficient forces the normalized two-point
function; \ClaimStatusProvedHere]
\label{prop:winfty-formal-self-normalization-from-t}
Fix $N \ge 2$ and consider a stage-$N$ local OPE calculus with visible
primary generators $W^{(2)},\dots,W^{(N)}$, where
$T:=W^{(2)}$ is the Virasoro field. Assume:
\begin{enumerate}[label=\textup{(\roman*)}]
\item each visible generator $W^{(u)}$ has conformal weight $u$, the
 local OPE is decomposed into conformal families of these primary
 generators and their translation descendants, and the Virasoro Ward
 identity holds for insertions of $T$ against each visible generator;
\item the stress tensor is normalized by
 \[
 \langle T(z)T(0)\rangle=\frac{c/2}{z^4}
 \qquad\text{with } c\neq 0;
 \]
\item for some visible generator $W^{(s)}$ with $2 \le s \le N$, the
 top-pole stress-tensor coefficient in the self-OPE is
 \[
 \mathsf{C}_{s,s;2;0,\,2s-2}(N)=2.
 \]
\end{enumerate}
Then the self-pairing of $W^{(s)}$ is normalized by
\[
\langle W^{(s)}(z)W^{(s)}(0)\rangle=\frac{c/s}{z^{2s}}.
\]
Equivalently, under the visible Virasoro package the universal
self-coupling coefficient $2$ and the Zamolodchikov normalization of
$W^{(s)}$ are equivalent.
\end{proposition}

\begin{proof}
Global conformal invariance forces $\langle A(z)A(0)\rangle = N_s z^{-2s}$. The Ward-identity computation of Proposition~\ref{prop:winfty-formal-self-t-coefficient} gives $x^{-4}z^{-2s+2}$ coefficient $sN_s$ on one side and $\mathsf{C}_{s,s;2;0,2s-2}\cdot c/2 = c$ on the other. Matching gives $N_s = c/s$.
\end{proof}

\begin{proposition}[Principal Drinfeld--Sokolov self-coupling
stress-tensor coefficient; \ClaimStatusConditional]
\label{prop:winfty-ds-self-t-coefficient}
Fix $N \ge 2$ and consider the principal Drinfeld--Sokolov stage
$\mathcal{W}_N$ at a generic level with nonzero central charge $c$.
Assume the Zamolodchikov normalization package of
Theorem~\ref{thm:wn-obstruction}. Then for every visible generator
$W^{(s)}$ with $2 \le s \le N$, the top-pole stress-tensor coefficient
in the self-OPE is
\[
\mathsf{C}^{\mathrm{DS}}_{s,s;2;0,\,2s-2}(N)=2.
\]
Equivalently, on the principal Drinfeld--Sokolov side the
undifferentiated $T=W^{(2)}$ channel appears in every self-OPE
$W^{(s)}(z)W^{(s)}(w)$ with the universal coefficient $2$.
\end{proposition}

\begin{proof}
Apply Proposition~\ref{prop:winfty-formal-self-t-coefficient} to the
principal Drinfeld--Sokolov stage. Theorem~\ref{thm:wn-obstruction}
supplies the required normalized two-point package, with
\[
\langle W^{(s)}(z)W^{(s)}(0)\rangle=\frac{c/s}{z^{2s}}
\qquad\text{and}\qquad
\langle T(z)T(0)\rangle=\frac{c/2}{z^4}
\]
at generic level.
\end{proof}

\begin{corollary}[Principal stage-\texorpdfstring{$4$}{4} self-coupling
\texorpdfstring{$W^{(4)}$-$W^{(4)}\to T$}{W4-W4 to T}
normalization; \ClaimStatusConditional]
\label{cor:winfty-ds-stage4-self-t-normalization}
In the setting of Proposition~\ref{prop:winfty-ds-self-t-coefficient},
one has
\[
\mathsf{C}^{\mathrm{DS}}_{4,4;2;0,6}=2.
\]
Equivalently, the principal stage-$4$ self-coupling
$W^{(4)}(z)W^{(4)}(w)\to T(w)$ is not free data: its top-pole
coefficient is fixed to $2$ in the principal Drinfeld--Sokolov
normalization.
\end{corollary}

\begin{proof}
Apply Proposition~\ref{prop:winfty-ds-self-t-coefficient} with $s=4$.
\end{proof}

\begin{corollary}[Stage-\texorpdfstring{$4$}{4} principal target packet after proved
Virasoro-target elimination; \ClaimStatusConditional]
\label{cor:winfty-ds-stage4-five-plus-zero}
Assume the hypotheses of
Corollary~\ref{cor:winfty-ds-stage4-mixed-two-channel} and
Corollary~\ref{cor:winfty-ds-stage4-self-t-normalization}. Then the
residual principal stage-$4$ target packet decomposes as:
\begin{enumerate}[label=\textup{(\roman*)}]
\item two higher-spin self-coupling targets
 \[
 c_{334},\qquad c_{444};
 \]
\item the mixed swap-even higher-spin target
 \[
 \mathsf{C}_{3,4;3;0,4};
 \]
\item the mixed swap-odd higher-spin target
 \[
 \mathsf{C}_{3,4;4;0,3};
 \]
\item the proved principal Virasoro-target value
 \[
 \mathsf{C}^{\mathrm{DS}}_{4,4;2;0,6}=2;
 \]
\item the proved principal Virasoro-target value
 \[
 \mathsf{C}^{\mathrm{DS}}_{3,4;2;0,5}=0.
 \]
\end{enumerate}
Equivalently, the standard stage-$4$ identity packet on~$\mathcal{I}_4$
splits into four higher-spin targets
\[
c_{334},\qquad c_{444},\qquad
\mathsf{C}_{3,4;3;0,4},\qquad \mathsf{C}_{3,4;4;0,3},
\]
together with two Virasoro-target channels whose principal values are
fixed by the preceding propositions:
\[
\mathsf{C}^{\mathrm{DS}}_{4,4;2;0,6}=2,
\qquad
\mathsf{C}^{\mathrm{DS}}_{3,4;2;0,5}=0,
\]
and whose residue-side contraction to four channels is the additional
package isolated in
Corollary~\ref{cor:winfty-stage4-residue-four-channel}.
\end{corollary}

\begin{proof}
Combine the mixed-self split (Corollary~\ref{cor:winfty-ds-stage4-mixed-self-split}), the mixed two-channel reduction (Corollary~\ref{cor:winfty-ds-stage4-mixed-two-channel}), and the universal $T$-coupling normalization (Corollary~\ref{cor:winfty-ds-stage4-self-t-normalization}).
\end{proof}

\begin{proposition}[MC4 completion packet for the standard
\texorpdfstring{$W_\infty$}{W_infty} tower]
\label{prop:winfty-mc4-frontier-package}
Let $\mathcal{W}^{\mathrm{ht}}$ be a separated complete H-level target
for $W_\infty$ satisfying the quotient-system hypotheses of
Corollary~\ref{cor:winfty-hlevel-comparison-criterion}. Assume
that the quotient residue calculus is translation compatible
and that the strong generators are even Virasoro-primary fields, so
that
Proposition~\ref{prop:winfty-ds-generator-seed},
Corollary~\ref{cor:winfty-ds-finite-seed-set},
Proposition~\ref{prop:winfty-ds-stage3-explicit-packet}, and
Corollary~\ref{cor:winfty-ds-stage4-five-plus-zero} apply. Then:
\begin{enumerate}
\item for every finite stage $N$, the remaining H-level comparison with
 the principal Drinfeld--Sokolov stage $W_N$ reduces to the finite
 primary seed packet $\mathcal{I}_N$ of residue identities
 \[
 \mathsf{C}^{\mathrm{res}}_{s,t;u;0,n}(N)
 =\mathsf{C}^{\mathrm{DS}}_{s,t;u;0,n}(N);
 \]
\item at stage~$3$, that packet is already fixed by the
 explicit $W_3$ OPE: three nonzero primary coefficients and twelve
 forced vanishing statements;
\item at stage~$4$, after removing the proved Virasoro and
 $W_3$ sectors, the first genuinely open higher-spin packet is exactly
 the exact six-entry identity packet on~$\mathcal{I}_4$
 \[
 \mathsf{C}^{\mathrm{res}}_{s,t;u;m,n}(4)
 =
 \mathsf{C}^{\mathrm{DS}}_{s,t;u;m,n}(4)
 \]
 in the channels
 \[
 \begin{gathered}
 (3,3;4;0,2),\qquad
 (4,4;4;0,4),\\
 (3,4;3;0,4),\qquad
 (3,4;4;0,3),\qquad
 (4,4;2;0,6),\\
 (3,4;2;0,5).
 \end{gathered}
 \]
 with the first four entries forming the residual higher-spin
 comparison packet, while the last two are Virasoro-target channels
 whose principal values are already fixed by
 Corollary~\ref{cor:winfty-ds-stage4-five-plus-zero}.
\end{enumerate}
In particular, once the H-level target and quotient system are
constructed, the standard $\mathcal{W}_\infty$ MC4 problem is not a
further stabilization theorem. It is the finite stagewise
residue-identification problem on the packets $\mathcal{I}_N$, with the
first genuinely open higher-spin packet occurring at stage~$4$ in the
displayed six-entry identity form on~$\mathcal{I}_4$.
\ClaimStatusConditional{}
\end{proposition}

\begin{proof}
\textup{(1)}: Corollary~\ref{cor:winfty-hlevel-comparison-criterion} plus Proposition~\ref{prop:winfty-ds-generator-seed} and Corollary~\ref{cor:winfty-ds-finite-seed-set}. \textup{(2)}: Proposition~\ref{prop:winfty-ds-stage3-explicit-packet}. \textup{(3)}: Corollary~\ref{cor:winfty-ds-stage4-five-plus-zero}.
\end{proof}

\begin{corollary}[Minimal closure criterion for the standard
\texorpdfstring{$W_\infty$}{W_infty} MC4 completion;
\ClaimStatusConditional]
\label{cor:winfty-stage4-closure-criterion}
Under the hypotheses of
Proposition~\ref{prop:winfty-mc4-frontier-package}, the standard
\texorpdfstring{$W_\infty$}{W_infty} H-level comparison is reduced to
the stagewise packet identities on~$\mathcal{I}_N$. The first
genuinely open higher-spin packet is the exact six-entry stage-$4$
identity packet on~$\mathcal{I}_4$. Within that packet, the residual
higher-spin channels still to be matched are
\[
\begin{gathered}
(3,3;4;0,2),\qquad
(4,4;4;0,4),\qquad
(3,4;3;0,4),\qquad
(3,4;4;0,3),
\end{gathered}
\]
while the corresponding principal Virasoro-target values are already
fixed to
\[
\mathsf{C}^{\mathrm{DS}}_{4,4;2;0,6}=2,
\qquad
\mathsf{C}^{\mathrm{DS}}_{3,4;2;0,5}=0.
\]
\end{corollary}

\begin{proof}
Corollary~\ref{cor:winfty-hlevel-comparison-criterion} upgrades
stagewise packet matching on~$\mathcal{I}_N$ to the H-level
comparison, and
Proposition~\ref{prop:winfty-mc4-frontier-package} identifies the first
open stage as the six-entry block on~$\mathcal{I}_4$.
Corollary~\ref{cor:winfty-ds-stage4-five-plus-zero} fixes the two
principal Virasoro-target values, leaving the four displayed
higher-spin channels as the residual higher-spin packet.
\end{proof}

\begin{corollary}[Constructing the completed chiral Koszul-dual
candidate for \texorpdfstring{$W_\infty$}{W_infty};
\ClaimStatusConditional]
\label{cor:winfty-dual-candidate-construction}
Assume the hypotheses of
Proposition~\ref{prop:winfty-mc4-frontier-package}. Then the completed
bar object of the corresponding H-level target
$\mathcal{W}^{\mathrm{ht}}$ is quasi-isomorphic to the inverse-limit
bar object of the principal stages:
\[
\widehat{\bar B}\bigl(\mathcal{W}^{\mathrm{ht}}\bigr)
\xrightarrow{\sim}
\varprojlim_N \bar B(W_N).
\]
Its completed cobar comparison recovers
$\mathcal{W}^{\mathrm{ht}}$ up to quasi-isomorphism, and any two
principal-stage compatible H-level targets satisfying the same
stagewise packet identities determine quasi-isomorphic completed
bar-cobar objects.

Consequently, the foundational route to the chiral Koszul-dual object
of \texorpdfstring{$W_\infty$}{W_infty} is exactly:
\begin{enumerate}
\item construct a principal-stage compatible H-level or factorization
 target;
\item prove the finite packet identities on~$\mathcal{I}_N$; and
\item invoke
 Corollary~\ref{cor:winfty-hlevel-comparison-criterion} to lift the
 standard principal-stage completed bar object to that target.
\end{enumerate}
No separate infinite-generator quadratic presentation is required at
this stage.
\end{corollary}

\begin{proof}
Under the stated hypotheses,
Proposition~\ref{prop:winfty-mc4-frontier-package} identifies the
finite-stage quotient problem with the packet identities on
$\mathcal{I}_N$, while
Corollary~\ref{cor:winfty-hlevel-comparison-criterion} identifies
$\mathcal{W}^{\mathrm{ht}}$ with the standard principal-stage completed
M-level package. The bar-cobar comparison part of that corollary
therefore gives the displayed quasi-isomorphism and the cobar recovery.
The same comparison shows that any two such H-level targets define the
same completed bar-cobar object up to quasi-isomorphism.
\end{proof}

\begin{corollary}[Stage-\texorpdfstring{$4$}{4} \texorpdfstring{$W_\infty$}{W_infty}
reduction on the Ward-normalized H-level locus;
\ClaimStatusConditional]
\label{cor:winfty-stage4-residue-four-channel}
Assume the hypotheses of
Proposition~\ref{prop:winfty-mc4-frontier-package}. Assume further
that $\mathcal{W}^{\mathrm{ht}}$ is stage-$4$ Ward-normalized in the
sense of Definition~\ref{def:winfty-stage4-ward-normalized}. Then the
two
Virasoro-target residue channels are fixed by the Ward-normalized
package:
\[
\mathsf{C}^{\mathrm{res}}_{3,4;2;0,5}(4)=0,
\qquad
\mathsf{C}^{\mathrm{res}}_{4,4;2;0,6}(4)=2.
\]
Consequently the first genuinely open stage-$4$
\texorpdfstring{$W_\infty$}{W_infty} packet contracts from the six-entry
identity packet of Proposition~\ref{prop:winfty-mc4-frontier-package} to
the four-channel higher-spin identity packet
\[
\mathsf{C}^{\mathrm{res}}_{s,t;u;0,n}(4)
=
\mathsf{C}^{\mathrm{DS}}_{s,t;u;0,n}(4)
\]
in the channels
\[
\begin{gathered}
(3,3;4;0,2),\qquad
(4,4;4;0,4),\\
(3,4;3;0,4),\qquad
(3,4;4;0,3).
\end{gathered}
\]
\end{corollary}

\begin{proof}
The Ward-normalized input supplies the hypotheses of Propositions~\ref{prop:winfty-formal-mixed-virasoro-zero} and~\ref{prop:winfty-formal-self-t-coefficient}, fixing the two Virasoro-target channels; removing them leaves four higher-spin identities.
\end{proof}

\begin{proposition}[Exact missing input for the conditional
\texorpdfstring{$W_\infty$}{W_infty} stage-\texorpdfstring{$4$}{4} contraction;
\ClaimStatusConditional]
\label{prop:winfty-stage4-visible-pairing-gap}
Assume the hypotheses of
Proposition~\ref{prop:winfty-mc4-frontier-package}. Assume further
that the stage-$4$ quotient residue calculus on the visible generators
$W^{(2)},W^{(3)},W^{(4)}$ is decomposed into conformal families of
Virasoro-primary generators and that the Virasoro field $T:=W^{(2)}$
acts on them, including on itself, by the standard stage-$4$
Virasoro Ward identity. Then the only additional input from the
normalized residue two-point/Ward package used to deduce
\[
\mathsf{C}^{\mathrm{res}}_{3,4;2;0,5}(4)=0,
\qquad
\mathsf{C}^{\mathrm{res}}_{4,4;2;0,6}(4)=2
\]
is the visible weight-$4$ self-pairing normalization
\[
\langle W^{(4)}(z)W^{(4)}(0)\rangle=\frac{c/4}{z^{8}},
\]
with the same nonzero central charge $c$ as on the principal side.
In particular, once the visible Virasoro Ward action is fixed, the
remaining stage-$4$ Ward-inheritance input is exactly inheritance of
this weight-$4$ normalization, equivalently the single scalar identity
\[
\mathsf{C}^{\mathrm{res}}_{4,4;2;0,6}(4)=2.
\]
\end{proposition}

\begin{proof}
Mixed-weight orthogonality and $W^{(3)}$ normalization are already forced (Propositions~\ref{prop:winfty-stage4-visible-orthogonality} and~\ref{prop:winfty-stage4-visible-w3-normalization}). The mixed vanishing follows from Proposition~\ref{prop:winfty-formal-mixed-virasoro-zero}; the self-coupling identity reduces via Proposition~\ref{prop:winfty-formal-self-t-coefficient} to the weight-$4$ normalization.
\end{proof}

\begin{definition}[Stage-\texorpdfstring{$4$}{4} Ward-normalized standard
\texorpdfstring{$W_\infty$}{W_infty} H-level target]
\label{def:winfty-stage4-ward-normalized}
Assume the setting of
Proposition~\ref{prop:winfty-mc4-frontier-package}. We say that the
H-level target $\mathcal{W}^{\mathrm{ht}}$ is
\emph{stage-$4$ Ward-normalized} if its induced stage-$4$ quotient
residue calculus satisfies the visible Virasoro Ward action together
with the visible weight-$4$ self-pairing normalization
\[
\langle W^{(4)}(z)W^{(4)}(0)\rangle=\frac{c/4}{z^{8}}
\]
with the same nonzero central charge $c$ as on the principal
Drinfeld--Sokolov stage. Equivalently, by
Proposition~\ref{prop:winfty-stage4-visible-orthogonality}, the
visible generators are mixed-weight orthogonal, and by
Proposition~\ref{prop:winfty-stage4-visible-w3-normalization} the
visible $W^{(3)}$ self-pairing is already normalized as well, so the
full visible package required by
Proposition~\ref{prop:winfty-stage4-visible-pairing-gap} holds.
\end{definition}

\begin{remark}[Stage-\texorpdfstring{$4$}{4} four-channel reduction]
\label{rem:winfty-stage4-four-channel-gap}
The MC4 structural framework is proved
(Theorem~\ref{thm:completed-bar-cobar-strong}). The stage-$4$
four-channel packet reduces to the H-level input of
Definition~\ref{def:winfty-stage4-ward-normalized}. The
quotient-system hypotheses of
Corollary~\ref{cor:winfty-hlevel-comparison-criterion} and
Proposition~\ref{prop:winfty-mc4-frontier-package} supply compatible
finite quotients, translation compatibility, and Virasoro-primary
strong generators, but they do not by themselves identify the
visible weight-$4$ self-pairing inside that quotient residue calculus.
Mixed-weight orthogonality is forced by the Virasoro Ward
identity
(Proposition~\ref{prop:winfty-stage4-visible-orthogonality}). The
full six-entry block of
Proposition~\ref{prop:winfty-mc4-frontier-package} is reduced by the
principal-stage identification input of
Theorem~\ref{thm:winfty-all-stages-rigidity-closure}, inside the
strong completion-tower framework of
Theorem~\ref{thm:completed-bar-cobar-strong}.
The reduction path is:
Proposition~\ref{prop:winfty-stage4-visible-pairing-gap}, once the
visible Virasoro Ward action is fixed, has exact remaining input
Proposition~\ref{prop:winfty-stage4-visible-diagonal-normalization},
equivalently the single scalar identity
\[
\mathsf{C}^{\mathrm{res}}_{4,4;2;0,6}(4)=2.
\]
Its package form is
Proposition~\ref{prop:winfty-stage4-ward-inheritance}.
\end{remark}

\begin{proposition}[Stage-\texorpdfstring{$4$}{4} visible mixed-weight orthogonality from
the Virasoro Ward identity; \ClaimStatusConditional]
\label{prop:winfty-stage4-visible-orthogonality}
Let $\mathcal{W}^{\mathrm{fact}}_\infty$ be a principal-stage
compatible factorization target in the sense of
Proposition~\ref{prop:winfty-factorization-package}, and let
$\mathcal{F}_{\le 4}$ be its stage-$4$ quotient. Assume:
\begin{enumerate}[label=\textup{(\roman*)}]
\item the visible generators $W^{(2)},W^{(3)},W^{(4)}$ in
 $\mathcal{F}_{\le 4}$ are even Virasoro-primary fields, and the local
 OPE is decomposed into their conformal families and translation
 descendants;
\item the Virasoro field $T:=W^{(2)}$ acts on these visible generators
 by the standard stage-$4$ Virasoro Ward identity;
\end{enumerate}
Then the quotient vacuum functional is mixed-weight orthogonal on the
visible generators:
\[
\langle W^{(s)}(z)W^{(t)}(0)\rangle = 0
\qquad (s\neq t,\ 2\le s,t\le 4).
\]
\end{proposition}

\begin{proof}
The global conformal Ward identities force
$\langle W^{(s)}(z)W^{(t)}(w)\rangle = C_{s,t}(z-w)^{-(s+t)}$;
the special conformal equation gives $(s-t)C_{s,t}=0$,
so $C_{s,t}=0$ whenever $s\neq t$.
\end{proof}

\begin{proposition}[Stage-\texorpdfstring{$4$}{4} visible
\texorpdfstring{$W^{(3)}$}{W3} normalization from the proved
\texorpdfstring{$W_3$}{W3} packet; \ClaimStatusConditional]
\label{prop:winfty-stage4-visible-w3-normalization}
Assume the hypotheses of
Proposition~\ref{prop:winfty-mc4-frontier-package}. Assume further
that the stage-$4$ quotient residue calculus on the visible generators
$W^{(2)},W^{(3)},W^{(4)}$ is decomposed into conformal families of
Virasoro-primary generators and that the Virasoro field $T:=W^{(2)}$
acts on them, including on itself, by the standard stage-$4$
Virasoro Ward identity. Then the visible stage-$4$ coefficient
\[
\mathsf{C}^{\mathrm{res}}_{3,3;2;0,4}(4)=2
\]
is already fixed by the stage-$3$ packet, and hence
\[
\langle W^{(3)}(z)W^{(3)}(0)\rangle=\frac{c/3}{z^{6}}.
\]
\end{proposition}

\begin{proof}
The tuple $(3,3,2,4)\in\mathcal{I}_4^{W_3}$ is fixed by the
proved stage-$3$ packet
(Proposition~\ref{prop:winfty-ds-stage3-explicit-packet}):
$\mathsf{C}^{\mathrm{res}}_{3,3;2;0,4}(4)=2$.
Proposition~\ref{prop:winfty-formal-self-normalization-from-t} with
$s=3$ then gives the displayed normalization.
\end{proof}

\begin{proposition}[Stage-\texorpdfstring{$4$}{4} visible
weight-\texorpdfstring{$4$}{4} normalization criterion]
\label{prop:winfty-stage4-visible-diagonal-normalization}
\ClaimStatusConditional{}
Let $\mathcal{W}^{\mathrm{fact}}_\infty$ be a principal-stage
compatible factorization target in the sense of
Proposition~\ref{prop:winfty-factorization-package}, and let
$\mathcal{F}_{\le 4}$ be its stage-$4$ quotient. Assume:
\begin{enumerate}[label=\textup{(\roman*)}]
\item the visible generators $W^{(2)},W^{(3)},W^{(4)}$ in
 $\mathcal{F}_{\le 4}$ are even Virasoro-primary fields, and the local
 OPE is decomposed into their conformal families and translation
 descendants;
\item the Virasoro field $T:=W^{(2)}$ acts on these visible generators
 by the standard stage-$4$ Virasoro Ward identity, including on
 itself.
\item the stage-$4$ quotient is identified, with the chosen
 strong-generator signs and two-point normalization, with the
 principal Drinfeld--Sokolov $W_4$ stage of
 Proposition~\ref{prop:w4-ds-ope-explicit}.
\end{enumerate}
Then the remaining visible self-pairing is normalized by
\[
\langle W^{(4)}(z)W^{(4)}(0)\rangle = \frac{c/4}{z^{8}}
\]
with the same nonzero central charge $c$ as on the principal
Drinfeld--Sokolov stage. Equivalently,
\[
\mathsf{C}^{\mathrm{res}}_{4,4;2;0,6}(4)=2.
\]
\end{proposition}

\begin{proof}
The added principal-stage identification supplies the normalized
Drinfeld--Sokolov two-point function
$\langle W^{(4)}(z)W^{(4)}(0)\rangle=(c/4)z^{-8}$. Then
Proposition~\ref{prop:winfty-formal-self-t-coefficient} with $s=4$
gives
$\mathsf{C}^{\mathrm{res}}_{4,4;2;0,6}(4)=2$. Conversely,
Proposition~\ref{prop:winfty-formal-self-normalization-from-t} with
$s=4$ recovers the displayed two-point normalization from that scalar.
\end{proof}

\begin{proposition}[Stage-\texorpdfstring{$4$}{4} Ward normalization
criterion from the visible Virasoro package]
\label{prop:winfty-stage4-ward-inheritance}
\ClaimStatusConditional{}
Let $\mathcal{W}^{\mathrm{fact}}_\infty$ be a principal-stage
compatible factorization target in the sense of
Proposition~\ref{prop:winfty-factorization-package}, and let
$\mathcal{F}_{\le 4}$ be its stage-$4$ quotient. Assume:
\begin{enumerate}[label=\textup{(\roman*)}]
\item the visible generators $W^{(2)},W^{(3)},W^{(4)}$ in
 $\mathcal{F}_{\le 4}$ are even Virasoro-primary fields, and the local
 OPE is decomposed into their conformal families and translation
 descendants;
\item the Virasoro field $T:=W^{(2)}$ acts on these visible generators
 by the standard stage-$4$ Virasoro Ward identity, including on
 itself;
\item the visible weight-$4$ self-pairing normalization of
 Proposition~\ref{prop:winfty-stage4-visible-diagonal-normalization}
 holds, equivalently
 $\mathsf{C}^{\mathrm{res}}_{4,4;2;0,6}(4)=2$.
\end{enumerate}
Then the induced H-level target is stage-$4$ Ward-normalized in the
sense of Definition~\ref{def:winfty-stage4-ward-normalized}. Consequently
Corollary~\ref{cor:winfty-stage4-residue-four-channel} applies, so the
first live standard-tower stage-$4$ packet contracts from the
unconditional six-entry block of
Proposition~\ref{prop:winfty-mc4-frontier-package} to the four
higher-spin channels
\[
(3,3;4;0,2),\qquad
(4,4;4;0,4),\qquad
(3,4;3;0,4),\qquad
(3,4;4;0,3).
\]
\end{proposition}

\begin{proof}
Definition~\ref{def:winfty-stage4-ward-normalized} is exactly the
visible Virasoro Ward package together with the weight-$4$ self-pairing
normalization. The last hypothesis supplies that normalization, so the
target is stage-$4$ Ward-normalized. The contraction from the six-entry
packet to the four higher-spin channels is then precisely
Corollary~\ref{cor:winfty-stage4-residue-four-channel}.
\end{proof}

\begin{corollary}[Equivalent exact forms of the remaining
\texorpdfstring{$W_\infty$}{W_infty} stage-\texorpdfstring{$4$}{4} input;
\ClaimStatusConditional]
\label{cor:winfty-stage4-single-scalar-equivalent}
Assume the hypotheses of
Proposition~\ref{prop:winfty-mc4-frontier-package}. Assume further
that the stage-$4$ quotient residue calculus on the visible generators
$W^{(2)},W^{(3)},W^{(4)}$ is decomposed into conformal families of
Virasoro-primary generators and that the Virasoro field $T:=W^{(2)}$
acts on these visible generators by the standard stage-$4$ Virasoro
Ward identity, including on itself. Then the following are
equivalent:
\begin{enumerate}[label=\textup{(\roman*)}]
\item the induced H-level target is stage-$4$ Ward-normalized in the
 sense of Definition~\ref{def:winfty-stage4-ward-normalized};
\item the visible weight-$4$ self-pairing is normalized by
 \[
 \langle W^{(4)}(z)W^{(4)}(0)\rangle = \frac{c/4}{z^{8}};
 \]
\item the single scalar residue identity
 \[
 \mathsf{C}^{\mathrm{res}}_{4,4;2;0,6}(4)=2
 \]
 holds.
\end{enumerate}
Under these equivalent conditions,
Corollary~\ref{cor:winfty-stage4-residue-four-channel} applies, so the
first live standard-tower stage-$4$ packet contracts from the
unconditional six-entry block to the four higher-spin channels
\[
(3,3;4;0,2),\qquad
(4,4;4;0,4),\qquad
(3,4;3;0,4),\qquad
(3,4;4;0,3).
\]
\end{corollary}

\begin{proof}
The equivalence \textup{(i)} \(\Longleftrightarrow\) \textup{(ii)} is
exactly Definition~\ref{def:winfty-stage4-ward-normalized}. Under the
same visible Virasoro package,
Proposition~\ref{prop:winfty-formal-self-t-coefficient} with $s=4$
gives \textup{(ii)} \(\Longrightarrow\) \textup{(iii)}, and
Proposition~\ref{prop:winfty-formal-self-normalization-from-t} with
$s=4$ gives \textup{(iii)} \(\Longrightarrow\) \textup{(ii)}.
The final sentence is exactly
Corollary~\ref{cor:winfty-stage4-residue-four-channel}.
\end{proof}

\begin{proposition}[Stage-\texorpdfstring{$4$}{4} swap-even residue channel from a visible
invariant pairing; \ClaimStatusConditional]
\label{prop:winfty-stage4-residue-pairing-reduction}
Assume the hypotheses of
Corollary~\ref{cor:winfty-stage4-single-scalar-equivalent}. Assume
further that the visible stage-$4$ quotient residue calculus carries a
bilinear pairing \((-, -)_{\mathrm{vis}}\) on the span of
\(W^{(2)},W^{(3)},W^{(4)}\) such that:
\begin{enumerate}[label=\textup{(\roman*)}]
\item the visible generators are pairwise orthogonal for distinct
 conformal weights, and the self-pairings are normalized by
 \[
 (W^{(2)},W^{(2)})_{\mathrm{vis}}=\frac{c}{2},\qquad
 (W^{(3)},W^{(3)})_{\mathrm{vis}}=\frac{c}{3},\qquad
 (W^{(4)},W^{(4)})_{\mathrm{vis}}=\frac{c}{4}
 \]
 with the same nonzero central charge~\(c\);
\item for the visible primary field \(W^{(3)}\), the pairing satisfies
 the standard primary invariant-pairing identity
 \[
 \bigl(W^{(3)}_{(1)}a,b\bigr)_{\mathrm{vis}}
 =
 -\bigl(a,W^{(3)}_{(3)}b\bigr)_{\mathrm{vis}}
 \]
 whenever \(a,b\) are visible generators.
\end{enumerate}
Then the swap-even residue channel is forced by the \((3,3)\)-block:
\[
\mathsf{C}^{\mathrm{res}}_{3,4;3;0,4}(4)
=
-\frac{3}{4}\,
\mathsf{C}^{\mathrm{res}}_{3,3;4;0,2}(4).
\]
Equivalently,
\[
\bigl(\mathsf{C}^{\mathrm{res}}_{3,4;3;0,4}(4)\bigr)^2
=
\frac{9}{16}\,
\bigl(\mathsf{C}^{\mathrm{res}}_{3,3;4;0,2}(4)\bigr)^2.
\]
\end{proposition}

\begin{proof}
Apply the $W^{(3)}$ invariant-pairing identity with $a=W^{(3)}$, $b=W^{(4)}$: mode $1$ produces $\mathsf{C}^{\mathrm{res}}_{3,3;4;0,2}\cdot c/4$, the complementary mode $3$ produces $\mathsf{C}^{\mathrm{res}}_{3,4;3;0,4}\cdot c/3$, and the sign is negative ($W^{(3)}$ odd weight). Since $c\neq 0$, rearranging gives the displayed identity.
\end{proof}

\begin{corollary}[Stage-\texorpdfstring{$4$}{4} residue packet as three higher-spin
channels on the visible pairing locus; \ClaimStatusConditional]
\label{cor:winfty-stage4-residue-three-channel}
Assume the hypotheses of
Proposition~\ref{prop:winfty-stage4-residue-pairing-reduction}. Then,
on the stage-$4$ Ward-normalized visible pairing locus, the first
genuinely open stage-$4$ higher-spin residue packet contracts from the
four-channel block of
Corollary~\ref{cor:winfty-stage4-residue-four-channel} to the
three-channel block
\[
\begin{gathered}
(3,3;4;0,2),\qquad
(4,4;4;0,4),\\
(3,4;4;0,3),
\end{gathered}
\]
with the swap-even channel \((3,4;3;0,4)\) recovered formally from
\((3,3;4;0,2)\) by
Proposition~\ref{prop:winfty-stage4-residue-pairing-reduction}.
\end{corollary}

\begin{proof}
Proposition~\ref{prop:winfty-stage4-residue-pairing-reduction} forces $(3,4;3;0,4)$ from $(3,3;4;0,2)$; removing it leaves three channels.
\end{proof}

\begin{corollary}[Stage-\texorpdfstring{$4$}{4} higher-spin comparison as a
primitive-plus-transport square triple on the visible pairing locus;
\ClaimStatusConditional]
\label{cor:winfty-stage4-primitive-transport-square-triple}
Assume the hypotheses of
Proposition~\ref{prop:winfty-stage4-residue-pairing-reduction}. Then,
on the stage-$4$ Ward-normalized visible pairing locus, the four
higher-spin identities of
Corollary~\ref{cor:winfty-stage4-residue-four-channel} are equivalent
at signless square-class level to the three identities
\[
\bigl(\mathsf{C}^{\mathrm{res}}_{3,3;4;0,2}(4)\bigr)^2
=
c_{334}^{\,2},
\qquad
\bigl(\mathsf{C}^{\mathrm{res}}_{4,4;4;0,4}(4)\bigr)^2
=
c_{444}^{\,2},
\]
\[
\bigl(\mathsf{C}^{\mathrm{res}}_{3,4;4;0,3}(4)\bigr)^2
=
\tfrac{5}{7}\,c_{334}^{\,2}.
\]
Equivalently, on that locus the swap-even square identity
\[
\bigl(\mathsf{C}^{\mathrm{res}}_{3,4;3;0,4}(4)\bigr)^2
=
\tfrac{9}{16}\,c_{334}^{\,2}
\]
is automatic once the \((3,3;4;0,2)\) square identity holds. In
particular, the signless stage-$4$ higher-spin problem on the visible
pairing locus consists of the two primitive self-coupling squares
\[
\bigl(\mathsf{C}^{\mathrm{res}}_{3,3;4;0,2}(4)\bigr)^2,
\qquad
\bigl(\mathsf{C}^{\mathrm{res}}_{4,4;4;0,4}(4)\bigr)^2
\]
together with the swap-odd \(W^{(4)}\)-transport square
\[
\bigl(\mathsf{C}^{\mathrm{res}}_{3,4;4;0,3}(4)\bigr)^2.
\]
\end{corollary}

\begin{proof}
On the visible pairing locus, Proposition~\ref{prop:winfty-stage4-residue-pairing-reduction} gives
$(\mathsf{C}^{\mathrm{res}}_{3,4;3;0,4})^2 = \tfrac{9}{16}(\mathsf{C}^{\mathrm{res}}_{3,3;4;0,2})^2$.
On the principal side,
Corollary~\ref{cor:w4-ds-stage4-square-class-reduction} gives
\[
\mathsf{C}_{3,4;3;0,4}^{\,2}
=
\tfrac{9}{16}\,c_{334}^{\,2},
\qquad
\mathsf{C}_{3,4;4;0,3}^{\,2}
=
\tfrac{5}{7}\,c_{334}^{\,2}.
\]
Therefore the swap-even square identity is automatic once
\(
\bigl(\mathsf{C}^{\mathrm{res}}_{3,3;4;0,2}(4)\bigr)^2=c_{334}^{\,2}
\)
holds, while the remaining two displayed identities are exactly the
\((4,4)\) self-coupling square identity and the swap-odd
\(W^{(4)}\)-transport square identity. The converse implication is
immediate.
\end{proof}

\begin{remark}[Exact next square-class gap after the visible pairing
reduction]
\label{rem:winfty-stage4-primitive-transport-gap}
Corollary~\ref{cor:winfty-stage4-primitive-transport-square-triple}
shows that the visible pairing package brings the residue side to the
same square-class profile as the principal side up to one additional
transport relation: the principal DS packet still forces
\[
\mathsf{C}_{3,4;4;0,3}^{\,2}
=
\tfrac{5}{7}\,c_{334}^{\,2},
\]
while on the residue side the corresponding relation
\[
\bigl(\mathsf{C}^{\mathrm{res}}_{3,4;4;0,3}(4)\bigr)^2
=
\tfrac{5}{7}\,
\bigl(\mathsf{C}^{\mathrm{res}}_{3,3;4;0,2}(4)\bigr)^2
\]
is not yet formal from the visible Ward plus pairing package alone.
Establishing this visible top-pole Borcherds transport relation would
reduce the residue higher-spin packet to the same two primitive
self-coupling square classes
\(
\bigl(\mathsf{C}^{\mathrm{res}}_{3,3;4;0,2}(4)\bigr)^2
\)
and
\(
\bigl(\mathsf{C}^{\mathrm{res}}_{4,4;4;0,4}(4)\bigr)^2
\)
as on the principal side.
\end{remark}

\begin{proposition}[Stage-\texorpdfstring{$4$}{4} visible top-pole
Borcherds transport criterion]
\label{prop:winfty-stage4-visible-borcherds-transport}
\ClaimStatusConditional{}
Assume the hypotheses of
Proposition~\ref{prop:winfty-stage4-residue-pairing-reduction}. Assume,
in addition, that the visible stage-$4$ quotient is identified with the
principal Drinfeld--Sokolov $W_4$ stage, with the same central charge
and chosen strong-generator signs. Then the visible stage-$4$ quotient
residue calculus satisfies the transport square identity
\[
\bigl(\mathsf{C}^{\mathrm{res}}_{3,4;4;0,3}(4)\bigr)^2
=
\tfrac{5}{7}\,
\bigl(\mathsf{C}^{\mathrm{res}}_{3,3;4;0,2}(4)\bigr)^2.
\]
Equivalently, on the stage-$4$ Ward-normalized visible pairing locus
the swap-odd \(W^{(4)}\)-transport square is forced by the
\((3,3)\)-square.
\end{proposition}

\begin{proof}
The principal Drinfeld--Sokolov computation
(Proposition~\ref{prop:w4-ds-ope-explicit},
Computation~\ref{comp:w4-ds-ope-extraction}, $121$~tests) gives all
stage-$4$ OPE coefficients as explicit rational functions of~$c$;
the inter-coefficient relation
$\mathsf{C}_{3,4;4;0,3}^{\,2} = \tfrac{5}{7}c_{334}^{\,2}$
as an algebraic consequence of the DS structure
(Corollary~\ref{cor:w4-ds-stage4-square-class-reduction}). The added
identification hypothesis transports this DS identity to the visible
residue calculus. Without that identification, the visible Ward and
pairing package leaves the transport identity as the proof obligation
isolated in Remark~\ref{rem:winfty-stage4-primitive-transport-gap}.
\end{proof}

\begin{corollary}[Equivalent exact forms of the remaining stage-\texorpdfstring{$4$}{4}
higher-spin transport input on the visible pairing locus;
\ClaimStatusConditional]
\label{cor:winfty-stage4-visible-borcherds-two-primitive}
Assume the hypotheses of
Proposition~\ref{prop:winfty-stage4-residue-pairing-reduction}. Then
the following are equivalent:
\begin{enumerate}[label=\textup{(\roman*)}]
\item the visible top-pole Borcherds transport relation
 \[
 \bigl(\mathsf{C}^{\mathrm{res}}_{3,4;4;0,3}(4)\bigr)^2
 =
 \tfrac{5}{7}\,
 \bigl(\mathsf{C}^{\mathrm{res}}_{3,3;4;0,2}(4)\bigr)^2;
 \]
\item on the stage-$4$ Ward-normalized visible pairing locus, the
 higher-spin comparison of
 Corollary~\ref{cor:winfty-stage4-residue-four-channel} is equivalent
 at signless square-class level to the primitive pair
 \[
 \bigl(\mathsf{C}^{\mathrm{res}}_{3,3;4;0,2}(4)\bigr)^2
 =
 c_{334}^{\,2},
 \qquad
 \bigl(\mathsf{C}^{\mathrm{res}}_{4,4;4;0,4}(4)\bigr)^2
 =
 c_{444}^{\,2}.
 \]
\end{enumerate}
Under these equivalent conditions, the visible pairing locus has the
same two-primitive stage-$4$ square-class profile as the principal
Drinfeld--Sokolov side.
\end{corollary}

\begin{proof}
\textup{(i)}$\Rightarrow$\textup{(ii)}: the $\tfrac{5}{7}$-ratio from Corollary~\ref{cor:w4-ds-stage4-square-class-reduction} makes the transport square automatic once the $(3,3)$ primitive identity holds. \textup{(ii)}$\Rightarrow$\textup{(i)}: the same principal formula recovers the transport relation.
\end{proof}

\begin{proposition}[Exact local reduction order for the stage-\texorpdfstring{$4$}{4}
\texorpdfstring{$W_\infty$}{W_infty} packet; \ClaimStatusConditional]
\label{prop:winfty-stage4-local-reduction-order}
Assume the hypotheses of
Proposition~\ref{prop:winfty-mc4-frontier-package}. Assume further
that the stage-$4$ quotient residue calculus on the visible generators
\(W^{(2)},W^{(3)},W^{(4)}\) is decomposed into conformal families of
Virasoro-primary generators and that the Virasoro field \(T:=W^{(2)}\)
acts on them by the standard stage-$4$ Virasoro Ward identity,
including on itself. Then the local stage-$4$ comparison separates
into two exact steps:
\begin{enumerate}[label=\textup{(\roman*)}]
\item \emph{Virasoro-target normalization step.}
 The remaining input for the four-channel higher-spin refinement is
 exactly the visible weight-$4$ normalization, equivalently the single
 scalar identity
 \[
 \mathsf{C}^{\mathrm{res}}_{4,4;2;0,6}(4)=2;
 \]
\item \emph{Higher-spin transport step.}
 On the additional visible pairing locus of
 Proposition~\ref{prop:winfty-stage4-residue-pairing-reduction}, the
 remaining input for the principal two-primitive square profile is
 exactly the visible top-pole Borcherds transport relation
 \[
 \bigl(\mathsf{C}^{\mathrm{res}}_{3,4;4;0,3}(4)\bigr)^2
 =
 \tfrac{5}{7}\,
 \bigl(\mathsf{C}^{\mathrm{res}}_{3,3;4;0,2}(4)\bigr)^2.
 \]
\end{enumerate}
Equivalently, after the unconditional six-entry packet of
Proposition~\ref{prop:winfty-mc4-frontier-package}, the local
stage-$4$ problem is first a one-scalar Ward-normalization step and
then, on the visible pairing locus, a one-relation higher-spin
transport problem.
\end{proposition}

\begin{proof}
\textup{(i)}: Corollary~\ref{cor:winfty-stage4-single-scalar-equivalent}. \textup{(ii)}: Corollary~\ref{cor:winfty-stage4-primitive-transport-square-triple} reduces to the primitive-plus-transport triple; Corollary~\ref{cor:winfty-stage4-visible-borcherds-two-primitive} records the equivalent two-primitive form.
\end{proof}

\begin{remark}[Stage-4 Ward inheritance from \texorpdfstring{$\mathcal W_4$}{W4} rigidity]
The principal values are already explicit
(Proposition~\ref{prop:w4-ds-ope-explicit}). The three stage-$4$
identities (visible weight-$4$ normalization
(Proposition~\ref{prop:winfty-stage4-visible-diagonal-normalization}),
Ward inheritance
(Proposition~\ref{prop:winfty-stage4-ward-inheritance}), and
Borcherds transport
(Proposition~\ref{prop:winfty-stage4-visible-borcherds-transport})) follow
from $\mathcal{W}_4$ rigidity only after the visible stage-$4$ quotient
has been identified with the principal Drinfeld--Sokolov $W_4$ stage,
with central charge and strong-generator signs fixed. The unconditional
bar-side output remains the six-entry packet of
Proposition~\ref{prop:winfty-mc4-frontier-package}; the reduction to
the two primitive square classes of
Corollary~\ref{cor:winfty-stage4-visible-borcherds-two-primitive}
requires that principal-stage identification. For $N\geq5$ the same
point is the explicit proof obligation in
Theorem~\ref{thm:winfty-all-stages-rigidity-closure}.
\end{remark}

\begin{proposition}[Uniform Virasoro-target contraction of reduced
incremental packets under the normalized residue package;
\ClaimStatusConditional]
\label{prop:winfty-stage-growth-virasoro-target-contraction}
Fix $N \ge 3$ and assume the hypotheses of
Corollary~\ref{cor:winfty-ds-stage-growth-top-parity}. Assume
further that the stage-$(N+1)$ quotient residue calculus on the
visible generators
$W^{(2)},\dots,W^{(N+1)}$ satisfies the hypotheses of
Propositions~\ref{prop:winfty-formal-mixed-virasoro-zero}
and~\ref{prop:winfty-formal-self-t-coefficient}. Then the
Virasoro-target channels inside
$\mathcal{J}_{N+1}^{\mathrm{red}}$ are exactly
\[
\bigl\{
(s,N{+}1;2;0,s+N-1)\mid 3\le s\le N
\bigr\}
\sqcup
\bigl\{
(N{+}1,N{+}1;2;0,2N)
\bigr\}.
\]
Their residue values are fixed by
\[
\mathsf{C}^{\mathrm{res}}_{s,N+1;2;0,s+N-1}(N+1)=0
\qquad (3\le s\le N),
\]
\[
\mathsf{C}^{\mathrm{res}}_{N+1,N+1;2;0,2N}(N+1)=2.
\]
Consequently the reduced incremental packet contracts to the
higher-spin packet
\[
\mathcal{J}_{N+1}^{\mathrm{hs}}
:=
\mathcal{J}_{N+1}^{\mathrm{red}}
\setminus
\Bigl(
\bigl\{
(s,N{+}1;2;0,s+N-1)\mid 3\le s\le N
\bigr\}
\sqcup
\bigl\{
(N{+}1,N{+}1;2;0,2N)
\bigr\}
\Bigr).
\]
For $N=3$ this recovers
Corollary~\ref{cor:winfty-stage4-residue-four-channel}.
\end{proposition}

\begin{proof}
Proposition~\ref{prop:winfty-formal-mixed-virasoro-zero} kills each mixed target-$2$ channel $(s,N{+}1;2;0,s+N-1)$; Proposition~\ref{prop:winfty-formal-self-t-coefficient} fixes the self target-$2$ channel to $2$. Removing these leaves $\mathcal{J}_{N+1}^{\mathrm{hs}}$.
\end{proof}

\begin{corollary}[First reduced stage beyond
\texorpdfstring{$\mathcal{I}_4$}{I4} under the normalized residue
package; \ClaimStatusConditional]
\label{cor:winfty-stage5-residue-eight-channel}
Assume the hypotheses of
Corollary~\ref{cor:winfty-ds-stage5-reduced-packet}. Assume further
that the stage-$5$ quotient residue calculus on the visible generators
$W^{(2)},W^{(3)},W^{(4)},W^{(5)}$ satisfies the hypotheses of
Propositions~\ref{prop:winfty-formal-mixed-virasoro-zero}
and~\ref{prop:winfty-formal-self-t-coefficient}. Then the three
Virasoro-target channels in $\mathcal{J}_5^{\mathrm{red}}$ are
fixed:
\[
\mathsf{C}^{\mathrm{res}}_{3,5;2;0,6}(5)=0,
\qquad
\mathsf{C}^{\mathrm{res}}_{4,5;2;0,7}(5)=0,
\qquad
\mathsf{C}^{\mathrm{res}}_{5,5;2;0,8}(5)=2.
\]
Consequently the first reduced stage beyond~$\mathcal{I}_4$ contracts
to the eight higher-spin channels
\[
\begin{gathered}
(3,4;5;0,2),\\
(3,5;3;0,5),\qquad
(3,5;4;0,4),\qquad
(3,5;5;0,3),\\
(4,5;3;0,6),\qquad
(4,5;4;0,5),\qquad
(4,5;5;0,4),\\
(5,5;4;0,6).
\end{gathered}
\]
\end{corollary}

\begin{proof}
Apply
Proposition~\ref{prop:winfty-stage-growth-virasoro-target-contraction}
with $N=4$, and use
Corollary~\ref{cor:winfty-ds-stage5-reduced-packet} to identify the
resulting packet explicitly.
\end{proof}

\begin{corollary}[First higher-spin packet beyond
\texorpdfstring{$\mathcal{I}_4$}{I4}; \ClaimStatusConditional]
\label{cor:winfty-stage5-higher-spin-packet}
Assume the hypotheses of
Corollary~\ref{cor:winfty-stage5-residue-eight-channel}. Then the
first higher-spin packet beyond the stage-$4$ four-channel block is
\[
\mathcal{J}_5^{\mathrm{hs}}
=
\Bigl\{
(3,4;5;0,2)
\Bigr\}
\sqcup
\Bigl\{
(3,5;3;0,5),\,
(3,5;4;0,4),\,
(3,5;5;0,3)
\Bigr\}
\]
\[
\sqcup
\Bigl\{
(4,5;3;0,6),\,
(4,5;4;0,5),\,
(4,5;5;0,4)
\Bigr\}
\sqcup
\Bigl\{
(5,5;4;0,6)
\Bigr\}.
\]
In particular,
\[
|\mathcal{J}_5^{\mathrm{hs}}|=8,
\]
and its block decomposition by source pair is
\[
1+3+3+1
\]
for the blocks
$(3,4)$, $(3,5)$, $(4,5)$, $(5,5)$.
\end{corollary}

\begin{proof}
Corollary~\ref{cor:winfty-stage5-residue-eight-channel} identifies the
contracted packet explicitly, and
Proposition~\ref{prop:winfty-stage-growth-virasoro-target-contraction}
defines that contracted packet to be
$\mathcal{J}_5^{\mathrm{hs}}$.
\end{proof}

\begin{remark}[Stage-$5$ higher-spin packet decomposition]
\label{rem:winfty-stage5-higher-spin-subblocks}%
\label{rem:winfty-stage5-entry-transport}%
\label{rem:winfty-stage5-entry-singletons}%
\label{rem:winfty-stage5-entry-mixed-self}%
\label{rem:winfty-stage5-reduced-tail-singleton}%
\label{rem:winfty-stage5-tail-mechanism}%
\label{rem:winfty-stage5-transport-target-ladders}%
\label{rem:winfty-stage5-higher-spin-target-blocks}%
\label{rem:winfty-stage5-target5-corridor}%
\label{rem:winfty-stage5-target5-residual}%
\label{rem:winfty-stage5-target5-transport-mechanism}%
\label{rem:winfty-stage5-target5-transport-singletons}%
\label{rem:winfty-stage5-transport-pole-profiles}%
\label{rem:winfty-stage5-first-subproblems}%
The $8$-element packet $\mathcal{J}_5^{\mathrm{hs}}$ of
Corollary~\ref{cor:winfty-stage5-higher-spin-packet} decomposes by
source pair into blocks $(3,4)$, $(3,5)$, $(4,5)$, $(5,5)$ of sizes
$1+3+3+1$. The entry packet
$\mathcal{J}_5^{\mathrm{entry}}=\{(3,4;5;0,2),(5,5;4;0,6)\}$ carries
the mixed-entry singleton $\mathcal{J}_5^{\mathrm{ent,mix}}=\{(3,4;5;0,2)\}$ and the self-return
singleton $\mathcal{J}_5^{\mathrm{ent,self}}=\{(5,5;4;0,6)\}$; the transport packet
$\mathcal{J}_5^{\mathrm{tr}}$ has $6$ elements decomposing by target
spin into three $2$-element ladders $\mathcal{J}_5^{\mathrm{tr},u}$ for
$u=3,4,5$. By target spin,
$\mathcal{J}_5^{u=3}$, $\mathcal{J}_5^{u=4}$, $\mathcal{J}_5^{u=5}$
have sizes $2,3,3$. The target-$5$ corridor
$\mathcal{J}_5^{u=5}=\{(3,4;5;0,2),(3,5;5;0,3),(4,5;5;0,4)\}$ sits at
consecutive pole orders $2,3,4$; its residual continuation after the
tail singleton consists of the two transport singletons
$(3,5;5;0,3)$ and $(4,5;5;0,4)$, representing the
$W^{(5)}$-projection coefficients
$\mathsf{C}^{\mathrm{res}}_{3,5;5;0,3}(5)$ and
$\mathsf{C}^{\mathrm{res}}_{4,5;5;0,4}(5)$. The reduced tail block
$\mathcal{J}_5^{\mathrm{tail,red}}=\{(3,4;5;0,2)\}$ identifies the
missing mechanism as matching
$\mathsf{C}^{\mathrm{res}}_{3,4;5;0,2}(5)=\mathsf{C}^{\mathrm{DS}}_{3,4;5;0,2}(5)$.
Transport ladders have pole profiles $\{3,4\}$, $\{4,5\}$, $\{5,6\}$
for targets $5,4,3$ respectively.
\end{remark}

\begin{proposition}[Stage-\texorpdfstring{$5$}{5} visible
\texorpdfstring{$W^{(5)}$}{W5} normalization from the proved
\texorpdfstring{$W^{(5)}$-$W^{(5)}\to T$}{W5-W5 to T} coefficient;
\ClaimStatusConditional]
\label{prop:winfty-stage5-visible-w5-normalization}
Assume the hypotheses of
Corollary~\ref{cor:winfty-stage5-residue-eight-channel}. Assume
further that the stage-$5$ quotient residue calculus on the visible
generators $W^{(2)},W^{(3)},W^{(4)},W^{(5)}$ is decomposed into
conformal families of Virasoro-primary generators and that the
Virasoro field $T:=W^{(2)}$ acts on them, including on itself, by the
standard stage-$5$ Virasoro Ward identity. Then the visible stage-$5$
coefficient
\[
\mathsf{C}^{\mathrm{res}}_{5,5;2;0,8}(5)=2
\]
is already fixed by the stated hypotheses, and hence
\[
\langle W^{(5)}(z)W^{(5)}(0)\rangle=\frac{c/5}{z^{10}}.
\]
\end{proposition}

\begin{proof}
Apply Proposition~\ref{prop:winfty-formal-self-normalization-from-t} with $s=5$.
\end{proof}

\begin{proposition}[Stage-\texorpdfstring{$5$}{5} target-\texorpdfstring{$5$}{5} pole-\texorpdfstring{$3$}{3} transport singleton
vanishes on a visible \texorpdfstring{$W^{(5)}$}{W5}-pairing locus;
\ClaimStatusConditional]
\label{prop:winfty-stage5-target5-pole3-pairing-vanishing}
Assume the hypotheses of
Remark~\ref{rem:winfty-stage5-target5-transport-singletons}.
Assume further that the visible stage-$5$ quotient residue calculus
carries a bilinear pairing \((-, -)_{\mathrm{vis}}\) on the span of
\(W^{(2)},W^{(3)},W^{(4)},W^{(5)}\) such that:
\begin{enumerate}[label=\textup{(\roman*)}]
\item the visible generators are even Virasoro-primary fields,
 pairwise orthogonal for distinct conformal weights, and their
 self-pairings are normalized by
 \[
 (W^{(u)},W^{(u)})_{\mathrm{vis}}=\frac{c}{u}
 \qquad (2\le u\le 5)
 \]
 with the same nonzero central charge \(c\);
\item for the visible primary field \(W^{(5)}\), the pairing satisfies
 the standard primary invariant-pairing identity
 \[
 \bigl(W^{(5)}_{(n)}a,b\bigr)_{\mathrm{vis}}
 =
 -\bigl(a,W^{(5)}_{(8-n)}b\bigr)_{\mathrm{vis}}
 \]
 whenever \(a,b\) are visible generators and the displayed modes are
 defined.
\end{enumerate}
Then the first singleton in the residual target-$5$ continuation
vanishes:
\[
\mathsf{C}^{\mathrm{res}}_{3,5;5;0,3}(5)=0.
\]
\end{proposition}

\begin{proof}
The invariant-pairing identity of assumption~\textup{(ii)} with \(n=2\), \(a=W^{(3)}\), \(b=W^{(5)}\) relates the reversed-order coefficient \(\mathsf{C}^{\mathrm{res}}_{5,3;5;0,3}=-\mathsf{C}^{\mathrm{res}}_{3,5;5;0,3}\) (swap parity, \(3+5-5=3\) odd, by Proposition~\ref{prop:winfty-ds-mixed-top-pole-swap}) to the complementary mode \(8-2=6\) self-OPE coefficient \((5,5;3;0,7)\), which vanishes by Proposition~\ref{prop:winfty-ds-self-ope-parity} (odd self-OPE). Since \(c\neq 0\), the displayed vanishing follows.
\end{proof}

\begin{proposition}[Stage-\texorpdfstring{$5$}{5} target-\texorpdfstring{$5$}{5} pole-\texorpdfstring{$4$}{4} transport singleton
from the self-return singleton on a visible
\texorpdfstring{$W^{(5)}$}{W5}-pairing locus; \ClaimStatusConditional]
\label{prop:winfty-stage5-target5-pole4-from-self-return}
Assume the hypotheses of
Remark~\ref{rem:winfty-stage5-target5-transport-singletons} and
Proposition~\ref{prop:winfty-stage5-target5-pole3-pairing-vanishing}.
Then
\[
\mathsf{C}^{\mathrm{res}}_{4,5;5;0,4}(5)
=
-\frac{5}{4}\,
\mathsf{C}^{\mathrm{res}}_{5,5;4;0,6}(5).
\]
\end{proposition}

\begin{proof}
Apply the invariant-pairing technique of Proposition~\ref{prop:winfty-stage5-target5-pole3-pairing-vanishing} with \(n=3\), \(a=W^{(4)}\), \(b=W^{(5)}\). The swap symmetry of Proposition~\ref{prop:winfty-ds-mixed-top-pole-swap} (\(4+5-5=4\) even) gives \(\mathsf{C}^{\mathrm{res}}_{5,4;5;0,4}=\mathsf{C}^{\mathrm{res}}_{4,5;5;0,4}\); the complementary mode \(8-3=5\) produces \(\mathsf{C}^{\mathrm{res}}_{5,5;4;0,6}\); the sign is negative (\(W^{(5)}\) odd weight). Weight orthogonality and \(c\neq 0\) yield the displayed identity.
\end{proof}

\begin{proposition}[Stage-\texorpdfstring{$5$}{5} target-\texorpdfstring{$5$}{5} pole-\texorpdfstring{$4$}{4} transport singleton
vanishes on a visible \texorpdfstring{$W^{(4)}$}{W4}-pairing locus;
\ClaimStatusConditional]
\label{prop:winfty-stage5-target5-pole4-w4-vanishing}
Assume the hypotheses of
Remark~\ref{rem:winfty-stage5-target5-transport-singletons}.
Assume further that the visible stage-$5$ quotient residue calculus
carries a bilinear pairing \((-, -)_{\mathrm{vis}}\) on the span of
\(W^{(2)},W^{(3)},W^{(4)},W^{(5)}\) such that:
\begin{enumerate}[label=\textup{(\roman*)}]
\item the visible generators are even Virasoro-primary fields,
 pairwise orthogonal for distinct conformal weights, and their
 self-pairings are normalized by
 \[
 (W^{(u)},W^{(u)})_{\mathrm{vis}}=\frac{c}{u}
 \qquad (2\le u\le 5)
 \]
 with the same nonzero central charge \(c\);
\item for the visible primary field \(W^{(4)}\), the pairing satisfies
 the standard primary invariant-pairing identity
 \[
 \bigl(W^{(4)}_{(n)}a,b\bigr)_{\mathrm{vis}}
 =
 \bigl(a,W^{(4)}_{(6-n)}b\bigr)_{\mathrm{vis}}
 \]
 whenever \(a,b\) are visible generators and the displayed modes are
 defined.
\end{enumerate}
Then
\[
\mathsf{C}^{\mathrm{res}}_{4,5;5;0,4}(5)=0.
\]
\end{proposition}

\begin{proof}
The invariant-pairing identity with \(n=3\), \(a=W^{(4)}\), \(b=W^{(5)}\) relates \(\mathsf{C}^{\mathrm{res}}_{4,4;5;0,3}\) (killed by odd self-OPE parity, Proposition~\ref{prop:winfty-ds-self-ope-parity}) to \(\mathsf{C}^{\mathrm{res}}_{4,5;5;0,4}\cdot c/4\). Since \(c\neq 0\), the vanishing follows.
\end{proof}

\begin{corollary}[Stage-\texorpdfstring{$5$}{5} self-return singleton vanishes on the
visible \texorpdfstring{$W^{(4)}$}{W4}/\texorpdfstring{$W^{(5)}$}{W5}
pairing locus; \ClaimStatusConditional]
\label{cor:winfty-stage5-self-return-vanishing-on-pairing}
Assume the hypotheses of
Proposition~\ref{prop:winfty-stage5-target5-pole4-from-self-return}
and Proposition~\ref{prop:winfty-stage5-target5-pole4-w4-vanishing}.
Then
\[
\mathsf{C}^{\mathrm{res}}_{5,5;4;0,6}(5)=0.
\]
\end{corollary}

\begin{proof}
Combine Proposition~\ref{prop:winfty-stage5-target5-pole4-w4-vanishing} with
Proposition~\ref{prop:winfty-stage5-target5-pole4-from-self-return}.
\end{proof}

\begin{proposition}[Stage-\texorpdfstring{$5$}{5} reduced tail singleton from a visible
\texorpdfstring{$W^{(3)}$}{W3}-pairing locus; \ClaimStatusConditional]
\label{prop:winfty-stage5-tail-from-w3-pairing}
Under the hypotheses of
Remark~\ref{rem:winfty-stage5-tail-mechanism}, assume further
the visible stage-$5$ pairing package of
Proposition~\ref{prop:winfty-stage5-target5-pole3-pairing-vanishing}
with the $W^{(3)}$ invariant-pairing identity
$\bigl(W^{(3)}_{(n)}a,b\bigr)_{\mathrm{vis}}
=-\bigl(a,W^{(3)}_{(4-n)}b\bigr)_{\mathrm{vis}}$.
Then:
\[
\mathsf{C}^{\mathrm{res}}_{3,4;5;0,2}(5)
=
-\frac{5}{4}\,
\mathsf{C}^{\mathrm{res}}_{3,5;4;0,4}(5).
\]
\end{proposition}

\begin{proof}
Apply the $W^{(3)}$ invariant-pairing identity with \(n=1\), \(a=W^{(4)}\), \(b=W^{(5)}\); the complementary mode \(4-1=3\) produces \(\mathsf{C}^{\mathrm{res}}_{3,5;4;0,4}\cdot c/4\). The sign is negative (\(W^{(3)}\) odd weight). Since \(c\neq 0\), rearranging gives the displayed identity.
\end{proof}

\begin{proposition}[Stage-\texorpdfstring{$5$}{5} reduced tail singleton from a visible
\texorpdfstring{$W^{(4)}$}{W4}-pairing locus; \ClaimStatusConditional]
\label{prop:winfty-stage5-tail-from-w4-pairing}
Under the hypotheses of
Remark~\ref{rem:winfty-stage5-tail-mechanism}, assume the visible stage-$5$ pairing package of
Proposition~\ref{prop:winfty-stage5-target5-pole3-pairing-vanishing}
with the $W^{(4)}$ invariant-pairing identity
$\bigl(W^{(4)}_{(n)}a,b\bigr)_{\mathrm{vis}}
=\bigl(a,W^{(4)}_{(6-n)}b\bigr)_{\mathrm{vis}}$.
Then:
\[
\mathsf{C}^{\mathrm{res}}_{3,4;5;0,2}(5)
=
\frac{5}{3}\,
\mathsf{C}^{\mathrm{res}}_{4,5;3;0,6}(5).
\]
\end{proposition}

\begin{proof}
Apply the $W^{(4)}$ invariant-pairing identity with \(n=1\), \(a=W^{(3)}\), \(b=W^{(5)}\). The swap symmetry of Proposition~\ref{prop:winfty-ds-mixed-top-pole-swap} (\(3+4-5=2\) even) gives \(\mathsf{C}^{\mathrm{res}}_{4,3;5;0,2}=\mathsf{C}^{\mathrm{res}}_{3,4;5;0,2}\); the complementary mode \(6-1=5\) produces \(\mathsf{C}^{\mathrm{res}}_{4,5;3;0,6}\cdot c/3\). The sign is positive (\(W^{(4)}\) even weight). Since \(c\neq 0\), rearranging gives the displayed identity.
\end{proof}

\begin{corollary}[Stage-\texorpdfstring{$5$}{5} tail singleton equates neighboring transport channels;
\ClaimStatusConditional]
\label{cor:winfty-stage5-tail-cross-target-reduction}
Under the hypotheses of
Propositions~\ref{prop:winfty-stage5-tail-from-w3-pairing}
and~\ref{prop:winfty-stage5-tail-from-w4-pairing}:
$\mathsf{C}^{\mathrm{res}}_{3,5;4;0,4}(5)
=
-\frac{4}{3}\,
\mathsf{C}^{\mathrm{res}}_{4,5;3;0,6}(5)$.
\end{corollary}

\begin{proof}
Eliminate \(\mathsf{C}^{\mathrm{res}}_{3,4;5;0,2}(5)\) between the two propositions.
\end{proof}

\begin{corollary}[Stage-\texorpdfstring{$5$}{5} target-\texorpdfstring{$5$}{5} corridor contracts to the tail
singleton; \ClaimStatusConditional]
\label{cor:winfty-stage5-target5-corridor-to-tail}
Under the hypotheses of
Propositions~\ref{prop:winfty-stage5-target5-pole3-pairing-vanishing}
and~\ref{prop:winfty-stage5-target5-pole4-w4-vanishing}:
$\mathsf{C}^{\mathrm{res}}_{3,5;5;0,3}(5)=0$ and
$\mathsf{C}^{\mathrm{res}}_{4,5;5;0,4}(5)=0$.
\end{corollary}

\begin{proof}
The first vanishing is exactly
Proposition~\ref{prop:winfty-stage5-target5-pole3-pairing-vanishing}.
The second vanishing is exactly
Proposition~\ref{prop:winfty-stage5-target5-pole4-w4-vanishing}.
Together they show that the target-$5$ corridor contracts to the
tail singleton and carries no surviving target-$5$ coefficient.
\end{proof}

\begin{corollary}[Stage-\texorpdfstring{$5$}{5} target-\texorpdfstring{$5$}{5} corridor carries no new
independent coefficient; \ClaimStatusConditional]
\label{cor:winfty-stage5-target5-no-new-independent-data}
Under the hypotheses of
Corollaries~\ref{cor:winfty-stage5-target5-corridor-to-tail}
and~\ref{cor:winfty-stage5-tail-cross-target-reduction}:
\[
\mathsf{C}^{\mathrm{res}}_{3,4;5;0,2}(5)
=
-\frac{5}{4}\,\mathsf{C}^{\mathrm{res}}_{3,5;4;0,4}(5)
=
\frac{5}{3}\,\mathsf{C}^{\mathrm{res}}_{4,5;3;0,6}(5),
\]
while the two target-$5$ transport coefficients vanish.
\end{corollary}

\begin{proof}
Corollary~\ref{cor:winfty-stage5-target5-corridor-to-tail}
supplies the vanishing of the two target-$5$ transport coefficients.
Proposition~\ref{prop:winfty-stage5-tail-from-w4-pairing}
expresses
$\mathsf{C}^{\mathrm{res}}_{3,4;5;0,2}(5)$ as
$\frac{5}{3}\mathsf{C}^{\mathrm{res}}_{4,5;3;0,6}(5)$, while
Corollary~\ref{cor:winfty-stage5-tail-cross-target-reduction}
rewrites the same coefficient as
$-\frac{5}{4}\mathsf{C}^{\mathrm{res}}_{3,5;4;0,4}(5)$.
Hence every surviving corridor coefficient is determined by the tail
singleton, so the corridor carries no new independent data.
\end{proof}

\begin{proposition}[Stage-\texorpdfstring{$5$}{5} target-\texorpdfstring{$4$}{4} pole-\texorpdfstring{$5$}{5} transport singleton
vanishes; \ClaimStatusConditional]
\label{prop:winfty-stage5-target4-pole5-w4-vanishing}
Under the hypotheses of
Proposition~\ref{prop:winfty-stage5-target5-pole4-w4-vanishing}:
$\mathsf{C}^{\mathrm{res}}_{4,5;4;0,5}(5)=0$.
\end{proposition}

\begin{proof}
The $W^{(4)}$ invariant-pairing identity with \(n=4\), \(a=W^{(5)}\), \(b=W^{(4)}\) relates \(\mathsf{C}^{\mathrm{res}}_{4,5;4;0,5}\cdot c/4\) to the complementary mode \(6-4=2\) self-OPE coefficient \((4,4;5;0,3)\), which vanishes by odd self-OPE parity (Proposition~\ref{prop:winfty-ds-self-ope-parity}).
\end{proof}

\begin{proposition}[Stage-\texorpdfstring{$5$}{5} target-\texorpdfstring{$3$}{3} pole-\texorpdfstring{$5$}{5} transport singleton
vanishes; \ClaimStatusConditional]
\label{prop:winfty-stage5-target3-pole5-w3-vanishing}
Under the hypotheses of
Proposition~\ref{prop:winfty-stage5-tail-from-w3-pairing}:
$\mathsf{C}^{\mathrm{res}}_{3,5;3;0,5}(5)=0$.
\end{proposition}

\begin{proof}
The $W^{(3)}$ invariant-pairing identity with \(n=4\), \(a=W^{(5)}\), \(b=W^{(3)}\) relates \(\mathsf{C}^{\mathrm{res}}_{3,5;3;0,5}\cdot c/3\) to the complementary mode \(4-4=0\) self-OPE coefficient \((3,3;5;0,1)\), which vanishes by odd self-OPE parity (Proposition~\ref{prop:winfty-ds-self-ope-parity}).
\end{proof}

\begin{proposition}[Stage-\texorpdfstring{$5$}{5} target-\texorpdfstring{$4$}{4}/target-\texorpdfstring{$3$}{3} transport channels
are paired on a visible \texorpdfstring{$W^{(5)}$}{W5}-pairing locus;
\ClaimStatusConditional]
\label{prop:winfty-stage5-transport-cross-target-reduction}
Assume the hypotheses of
Proposition~\ref{prop:winfty-stage5-target5-pole3-pairing-vanishing}.
Then
\[
\mathsf{C}^{\mathrm{res}}_{3,5;4;0,4}(5)
=
-\frac{4}{3}\,
\mathsf{C}^{\mathrm{res}}_{4,5;3;0,6}(5).
\]
\end{proposition}

\begin{proof}
Apply the $W^{(5)}$ invariant-pairing identity with \(n=3\), \(a=W^{(3)}\), \(b=W^{(4)}\). The swap symmetry of Proposition~\ref{prop:winfty-ds-mixed-top-pole-swap} gives \(\mathsf{C}^{\mathrm{res}}_{5,3;4;0,4}=\mathsf{C}^{\mathrm{res}}_{3,5;4;0,4}\) (\(3+5-4=4\) even) and \(\mathsf{C}^{\mathrm{res}}_{5,4;3;0,6}=\mathsf{C}^{\mathrm{res}}_{4,5;3;0,6}\) (\(4+5-3=6\) even). The sign is negative (\(W^{(5)}\) odd weight). Weight orthogonality and \(c\neq 0\) yield the displayed identity.
\end{proof}

\begin{corollary}[Stage-\texorpdfstring{$5$}{5} mixed transport reduction to one
effective independent coefficient; \ClaimStatusConditional]
\label{cor:winfty-stage5-transport-effective-independent-frontier}
Under the hypotheses of
Propositions~\ref{prop:winfty-stage5-target4-pole5-w4-vanishing},
\ref{prop:winfty-stage5-target3-pole5-w3-vanishing}, and
\ref{prop:winfty-stage5-transport-cross-target-reduction}:
$\mathsf{C}^{\mathrm{res}}_{4,5;4;0,5}(5)=0$,
$\mathsf{C}^{\mathrm{res}}_{3,5;3;0,5}(5)=0$, and
$\mathsf{C}^{\mathrm{res}}_{4,5;3;0,6}(5)
=
-\frac{3}{4}\,
\mathsf{C}^{\mathrm{res}}_{3,5;4;0,4}(5)$.
\end{corollary}

\begin{proof}
The first two vanishings are the content of
Propositions~\ref{prop:winfty-stage5-target4-pole5-w4-vanishing}
and~\ref{prop:winfty-stage5-target3-pole5-w3-vanishing}.
The reduction identity follows from
Proposition~\ref{prop:winfty-stage5-transport-cross-target-reduction}.
\end{proof}

\begin{corollary}[Stage-\texorpdfstring{$5$}{5} higher-spin packet reduces to one effective
independent coefficient; \ClaimStatusConditional]
\label{cor:winfty-stage5-effective-independent-frontier}
Under the hypotheses of
Corollaries~\ref{cor:winfty-stage5-self-return-vanishing-on-pairing},
\ref{cor:winfty-stage5-target5-no-new-independent-data}, and
\ref{cor:winfty-stage5-transport-effective-independent-frontier},
set $A_5:=\mathsf{C}^{\mathrm{res}}_{3,5;4;0,4}(5)$.
Then every coefficient in $\mathcal{J}_5^{\mathrm{hs}}$ is either zero or determined by $A_5$:
\[
\mathsf{C}^{\mathrm{res}}_{3,4;5;0,2}(5)=-\tfrac{5}{4}\,A_5,
\quad
\mathsf{C}^{\mathrm{res}}_{4,5;3;0,6}(5)=-\tfrac{3}{4}\,A_5,
\]
with all other channels vanishing.
The effective independent coefficient is the target-$4$
singleton \((3,5;4;0,4)\).
\end{corollary}

\begin{proof}
Combine the self-return vanishing
(Corollary~\ref{cor:winfty-stage5-self-return-vanishing-on-pairing}),
the target-$5$ corridor reduction
(Corollary~\ref{cor:winfty-stage5-target5-no-new-independent-data}),
and the mixed transport reduction
(Corollary~\ref{cor:winfty-stage5-transport-effective-independent-frontier}).
The surviving channels are exactly $\mathsf{C}^{\mathrm{res}}_{3,4;5;0,2}$
and $\mathsf{C}^{\mathrm{res}}_{4,5;3;0,6}$, both determined by
$A_5 = \mathsf{C}^{\mathrm{res}}_{3,5;4;0,4}(5)$.
\end{proof}

\begin{proposition}[Exact local reduction order for the first stage-\texorpdfstring{$5$}{5}
higher-spin packet]
\label{prop:winfty-stage5-local-reduction-order}
Assume the hypotheses of
Remark~\ref{rem:winfty-stage5-entry-transport}. Then the first
stage-$5$ higher-spin problem separates into the following exact local
order:
\begin{enumerate}[label=\textup{(\roman*)}]
\item \emph{Entry step.} The reduction begins with the two-channel entry
 packet \(\mathcal{J}_5^{\mathrm{entry}}\), first through the
 mixed-entry singleton \((3,4;5;0,2)\) and then through the
 self-return singleton \((5,5;4;0,6)\).
\item \emph{Target-$5$ corridor step.} The first local mixed-transport
 strip is the target-$5$ corridor
 \[
 (3,4;5;0,2),\qquad
 (3,5;5;0,3),\qquad
 (4,5;5;0,4),
 \]
 whose residual continuation after the tail singleton
 \((3,4;5;0,2)\) is exactly the two singleton transport ladder
 \[
 (3,5;5;0,3),\qquad
 (4,5;5;0,4),
 \]
 taken in pole order \(3\) then \(4\).
\item \emph{Remaining mixed transport step.} After the target-$5$
 ladder, the remaining transport problem is reduced by fixed target
 in the order
 \[
 \mathcal{J}_5^{\mathrm{tr},5},\qquad
 \mathcal{J}_5^{\mathrm{tr},4},\qquad
 \mathcal{J}_5^{\mathrm{tr},3}.
 \]
\end{enumerate}
On the visible pairing loci of
Propositions~\ref{prop:winfty-stage5-target5-pole3-pairing-vanishing}
--\ref{cor:winfty-stage5-tail-cross-target-reduction}, the same
target-$5$ corridor becomes more rigid: the pole-$3$ singleton
vanishes, the pole-$4$ singleton is tied to the self-return singleton,
and the tail singleton is tied to neighboring target-$4$/target-$3$
transport channels. On the full visible
\texorpdfstring{$W^{(3)}$}{W3}/\texorpdfstring{$W^{(4)}$}{W4}/\texorpdfstring{$W^{(5)}$}{W5}
pairing locus of
Corollary~\ref{cor:winfty-stage5-target5-no-new-independent-data}, the
whole target-$5$ corridor carries no new independent coefficient, so
Proposition~\ref{prop:winfty-stage5-target4-pole5-w4-vanishing},
Proposition~\ref{prop:winfty-stage5-target3-pole5-w3-vanishing}, and
Proposition~\ref{prop:winfty-stage5-transport-cross-target-reduction}
collapse the remaining target-$4$/target-$3$ transport ladders to one
effective coefficient, and
Corollary~\ref{cor:winfty-stage5-effective-independent-frontier}
identifies the effective independent local reduction order there as the
single target-$4$ singleton \((3,5;4;0,4)\), with every other stage-$5$
higher-spin channel either zero or determined by it.
\ClaimStatusConditional{}
\end{proposition}

\begin{proof}
Parts \textup{(i)}--\textup{(iii)} follow from the packet decompositions already established:
Remark~\ref{rem:winfty-stage5-entry-transport} for the entry/transport split,
Remark~\ref{rem:winfty-stage5-target5-corridor} for the target-$5$ corridor,
and Remark~\ref{rem:winfty-stage5-transport-pole-profiles} for the transport ladder ordering.
The visible-pairing refinement combines
Propositions~\ref{prop:winfty-stage5-target5-pole3-pairing-vanishing}--\ref{prop:winfty-stage5-transport-cross-target-reduction}
with Corollary~\ref{cor:winfty-stage5-effective-independent-frontier}.
\end{proof}

\begin{proposition}[Principal stage-\texorpdfstring{$5$}{5}
target-\texorpdfstring{$5$}{5} corridor normal form]
\label{prop:winfty-stage5-principal-target5-no-new-independent-data}
\ClaimStatusConditional{}
The principal stage-$5$ target-$5$ corridor carries no new independent
coefficient precisely when the following identity package holds:
\[
\mathsf{C}^{\mathrm{DS}}_{3,4;5;0,2}(5)
=
-\frac{5}{4}\,\mathsf{C}^{\mathrm{DS}}_{3,5;4;0,4}(5)
=
\frac{5}{3}\,\mathsf{C}^{\mathrm{DS}}_{4,5;3;0,6}(5),
\]
\[
\mathsf{C}^{\mathrm{DS}}_{3,5;5;0,3}(5)=0,
\qquad
\mathsf{C}^{\mathrm{DS}}_{4,5;5;0,4}(5)=0.
\]
Equivalently, on the full visible
\texorpdfstring{$W^{(3)}$}{W3}/\texorpdfstring{$W^{(4)}$}{W4}/\texorpdfstring{$W^{(5)}$}{W5}
pairing locus the principal target-$5$ corridor carries no new
independent coefficient beyond the neighboring target-$4$/target-$3$
channels.
\end{proposition}

\begin{proposition}[Principal stage-\texorpdfstring{$5$}{5}
target-\texorpdfstring{$5$}{5} transport normal form]
\label{prop:winfty-stage5-principal-target5-transport-vanishing}
\ClaimStatusConditional{}
The target-$5$ transport part of the principal stage-$5$ normal form is
\[
\mathsf{C}^{\mathrm{DS}}_{3,5;5;0,3}(5)=0,
\qquad
\mathsf{C}^{\mathrm{DS}}_{4,5;5;0,4}(5)=0.
\]
Equivalently, on the full visible
\texorpdfstring{$W^{(3)}$}{W3}/\texorpdfstring{$W^{(4)}$}{W4}/\texorpdfstring{$W^{(5)}$}{W5}
pairing locus the principal target-$5$ corridor contracts to the tail
singleton.
\end{proposition}

\begin{proposition}[Principal stage-\texorpdfstring{$5$}{5} tail
singleton cross-target normal form]
\label{prop:winfty-stage5-principal-tail-cross-target-reduction}
\ClaimStatusConditional{}
The tail-cross-target part of the principal stage-$5$ normal form is
\[
\mathsf{C}^{\mathrm{DS}}_{3,4;5;0,2}(5)
=
-\frac{5}{4}\,\mathsf{C}^{\mathrm{DS}}_{3,5;4;0,4}(5)
=
\frac{5}{3}\,\mathsf{C}^{\mathrm{DS}}_{4,5;3;0,6}(5).
\]
Equivalently, on the full visible pairing locus the principal reduced
tail singleton adds no new independent coefficient beyond the
neighboring target-$4$/target-$3$ transport channels.
\end{proposition}

\begin{remark}[Target-$5$ corridor factorization]
\label{rem:winfty-stage5-principal-target5-factorization}%
Proposition~\ref{prop:winfty-stage5-principal-target5-no-new-independent-data}
holds if and only if both
Proposition~\ref{prop:winfty-stage5-principal-target5-transport-vanishing}
and
Proposition~\ref{prop:winfty-stage5-principal-tail-cross-target-reduction}
hold: the two transport vanishings together with the tail cross-target relation account for all five identities.
\end{remark}

\begin{proposition}[Principal residual stage-\texorpdfstring{$5$}{5}
front normal form]
\label{prop:winfty-stage5-principal-residual-front-one-coefficient}%
\label{prop:winfty-stage5-principal-residual-front-vanishings}%
\label{prop:winfty-stage5-principal-residual-cross-target-reduction}%
\label{prop:winfty-stage5-principal-residual-front-factorization}%
\ClaimStatusConditional{}
The residual-front part of the principal stage-$5$ normal form is
\[
\mathsf{C}^{\mathrm{DS}}_{5,5;4;0,6}(5)=0,
\qquad
\mathsf{C}^{\mathrm{DS}}_{4,5;4;0,5}(5)=0,
\qquad
\mathsf{C}^{\mathrm{DS}}_{3,5;3;0,5}(5)=0,
\]
\[
\mathsf{C}^{\mathrm{DS}}_{4,5;3;0,6}(5)
=
-\frac{3}{4}\,\mathsf{C}^{\mathrm{DS}}_{3,5;4;0,4}(5).
\]
Equivalently, outside the target-$5$ corridor, the principal stage-$5$
higher-spin packet carries one effective independent coefficient,
represented by the target-$4$ singleton \((3,5;4;0,4)\).
The three vanishings alone form the \emph{residual vanishing component}
and the cross-target relation
$\mathsf{C}^{\mathrm{DS}}_{4,5;3;0,6}(5)=-\tfrac{3}{4}\,\mathsf{C}^{\mathrm{DS}}_{3,5;4;0,4}(5)$
forms the \emph{residual cross-target reduction component};
the full statement holds if and only if both components hold.
\end{proposition}

\begin{proposition}[Principal stage-\texorpdfstring{$5$}{5}
one-coefficient normal-form package]
\label{prop:winfty-stage5-principal-one-coefficient-normal-form}
\ClaimStatusConditional{}
The principal stage-$5$ target packet satisfies the one-coefficient
normal form exactly when the following package of identities holds:
\[
\mathsf{C}^{\mathrm{DS}}_{3,4;5;0,2}(5)
=
-\frac{5}{4}\,\mathsf{C}^{\mathrm{DS}}_{3,5;4;0,4}(5),
\qquad
\mathsf{C}^{\mathrm{DS}}_{3,5;5;0,3}(5)=0,
\qquad
\mathsf{C}^{\mathrm{DS}}_{4,5;5;0,4}(5)=0,
\]
\[
\mathsf{C}^{\mathrm{DS}}_{5,5;4;0,6}(5)=0,
\qquad
\mathsf{C}^{\mathrm{DS}}_{4,5;4;0,5}(5)=0,
\qquad
\mathsf{C}^{\mathrm{DS}}_{3,5;3;0,5}(5)=0,
\qquad
\mathsf{C}^{\mathrm{DS}}_{4,5;3;0,6}(5)
=
-\frac{3}{4}\,\mathsf{C}^{\mathrm{DS}}_{3,5;4;0,4}(5).
\]
Equivalently, on the full visible
\texorpdfstring{$W^{(3)}$}{W3}/\texorpdfstring{$W^{(4)}$}{W4}/\texorpdfstring{$W^{(5)}$}{W5}
pairing locus the principal stage-$5$ higher-spin packet carries one
effective independent coefficient, represented by the target-$4$
singleton \((3,5;4;0,4)\). Equivalently again, this is the conjunction
of
Propositions~\ref{prop:winfty-stage5-principal-target5-no-new-independent-data}
and~\ref{prop:winfty-stage5-principal-residual-front-one-coefficient}.
This package is a principal-side input; it is not supplied by the
residue-side visible-pairing reductions alone.
\end{proposition}

\begin{proposition}[Principal stage-\texorpdfstring{$5$}{5} one-coefficient normal form
factors through the target-\texorpdfstring{$5$}{5} corridor and the residual front;
\ClaimStatusConditional]
\label{prop:winfty-stage5-principal-one-coefficient-factorization}
Assume the hypotheses of
Corollary~\ref{cor:winfty-stage5-effective-independent-frontier}. Then
the packaged principal stage-$5$ one-coefficient normal form holds if
and only if both of the following structural conditions hold:
\begin{enumerate}[label=\textup{(\roman*)}]
\item the target-$5$ corridor carries no new independent coefficient;
\item the residual front carries one effective independent
 coefficient.
\end{enumerate}
\end{proposition}

\begin{proof}
Combine Remark~\ref{rem:winfty-stage5-principal-target5-factorization} and Proposition~\ref{prop:winfty-stage5-principal-residual-front-factorization}; together they account for all seven identities in the normal form.
\end{proof}

\begin{proposition}[Stage-\texorpdfstring{$5$}{5} higher-spin comparison reduces to one
coefficient on the full visible pairing locus; \ClaimStatusConditional]
\label{prop:winfty-stage5-one-coefficient-reduction}
Assume the hypotheses of
Corollary~\ref{cor:winfty-stage5-effective-independent-frontier} and
the packaged principal stage-$5$ one-coefficient normal form.
Then the full stage-$5$ higher-spin comparison on
\(\mathcal{J}_5^{\mathrm{hs}}\) is equivalent to the single
target-$4$ singleton identity
\[
\mathsf{C}^{\mathrm{res}}_{3,5;4;0,4}(5)
=
\mathsf{C}^{\mathrm{DS}}_{3,5;4;0,4}(5).
\]
\end{proposition}

\begin{proof}
Corollary~\ref{cor:winfty-stage5-effective-independent-frontier} expresses every residue coefficient as a fixed multiple of \(\mathsf{C}^{\mathrm{res}}_{3,5;4;0,4}(5)\); matching this singleton matches all channels.
\end{proof}

\begin{proposition}[Stage-\texorpdfstring{$5$}{5} one-coefficient comparison on the full
visible pairing locus]
\label{prop:winfty-stage5-one-coefficient-comparison}
\ClaimStatusConditional{}
Assume the hypotheses of
Proposition~\ref{prop:winfty-stage5-one-coefficient-reduction}. Then
\[
\mathsf{C}^{\mathrm{res}}_{3,5;4;0,4}(5)
=
\mathsf{C}^{\mathrm{DS}}_{3,5;4;0,4}(5).
\]
Equivalently, on the full visible
\texorpdfstring{$W^{(3)}$}{W3}/\texorpdfstring{$W^{(4)}$}{W4}/\texorpdfstring{$W^{(5)}$}{W5}
pairing locus the stage-$5$ higher-spin problem is reduced to the
target-$4$ singleton \((3,5;4;0,4)\).
\end{proposition}

\begin{corollary}[Exact remaining stage-\texorpdfstring{$5$}{5} visible-pairing input
package; \ClaimStatusConditional]
\label{cor:winfty-stage5-exact-remaining-input}
Assume the hypotheses of
Corollary~\ref{cor:winfty-stage5-effective-independent-frontier}. Then
the full stage-$5$ higher-spin comparison on \(\mathcal{J}_5^{\mathrm{hs}}\)
holds on the full visible pairing locus if and only if the following
two inputs both hold:
\begin{enumerate}[label=\textup{(\roman*)}]
\item the packaged principal stage-$5$ one-coefficient normal form
 holds;
\item the target-$4$ singleton identity
 \[
 \mathsf{C}^{\mathrm{res}}_{3,5;4;0,4}(5)
 =
 \mathsf{C}^{\mathrm{DS}}_{3,5;4;0,4}(5)
 \]
 holds.
\end{enumerate}
Equivalently, the exact remaining local stage-$5$ visible-pairing
problem is the packaged principal one-coefficient normal form together
with the displayed singleton identity. The missing input is therefore
no longer residue-side transport reduction, but that packaged
principal normal form together with that single target-$4$ match.
Equivalently again, by
Proposition~\ref{prop:winfty-stage5-principal-one-coefficient-factorization},
this two-input package is the same as the pair of structural
conditions that the target-$5$ corridor carries no new independent
coefficient and the residual front carries one effective independent
coefficient, together with the singleton identity.
\end{corollary}

\begin{proof}
Necessity: the one-coefficient normal form of Corollary~\ref{cor:winfty-stage5-effective-independent-frontier} on the residue side forces the same on the principal side; the singleton identity is one matched channel. Sufficiency: \textup{(i)} supplies the hypothesis of Proposition~\ref{prop:winfty-stage5-one-coefficient-reduction}, and \textup{(ii)} is the single identity to which it reduces. The reformulation is Proposition~\ref{prop:winfty-stage5-principal-one-coefficient-factorization}.
\end{proof}

\begin{corollary}[Stage-\texorpdfstring{$5$}{5} higher-spin defect family collapses to one
representative defect on the full visible pairing locus]
\label{cor:winfty-stage5-one-defect-family}
Assume the hypotheses of
Corollary~\ref{cor:winfty-stage5-effective-independent-frontier},
and
the packaged principal stage-$5$ one-coefficient normal form.
Set
\[
D_5:=
\mathsf{C}^{\mathrm{res}}_{3,5;4;0,4}(5)
-
\mathsf{C}^{\mathrm{DS}}_{3,5;4;0,4}(5).
\]
Then every defect in the stage-$5$ higher-spin packet
\(\mathcal{J}_5^{\mathrm{hs}}\) is either zero or a fixed rational
multiple of \(D_5\). Explicitly,
\[
\begin{aligned}
\mathsf{C}^{\mathrm{res}}_{3,4;5;0,2}(5)
-
\mathsf{C}^{\mathrm{DS}}_{3,4;5;0,2}(5)
&=
-\frac{5}{4}\,D_5,\\
\mathsf{C}^{\mathrm{res}}_{3,5;5;0,3}(5)
-
\mathsf{C}^{\mathrm{DS}}_{3,5;5;0,3}(5)
&=0,\\
\mathsf{C}^{\mathrm{res}}_{4,5;5;0,4}(5)
-
\mathsf{C}^{\mathrm{DS}}_{4,5;5;0,4}(5)
&=0,
\end{aligned}
\]
\[
\begin{aligned}
\mathsf{C}^{\mathrm{res}}_{5,5;4;0,6}(5)
-
\mathsf{C}^{\mathrm{DS}}_{5,5;4;0,6}(5)
&=0,\\
\mathsf{C}^{\mathrm{res}}_{4,5;4;0,5}(5)
-
\mathsf{C}^{\mathrm{DS}}_{4,5;4;0,5}(5)
&=0,\\
\mathsf{C}^{\mathrm{res}}_{3,5;3;0,5}(5)
-
\mathsf{C}^{\mathrm{DS}}_{3,5;3;0,5}(5)
&=0,
\end{aligned}
\]
\[
\mathsf{C}^{\mathrm{res}}_{4,5;3;0,6}(5)
-
\mathsf{C}^{\mathrm{DS}}_{4,5;3;0,6}(5)
=
-\frac{3}{4}\,D_5.
\]
Equivalently, on that full visible pairing locus the whole stage-$5$
higher-spin defect family on \(\mathcal{J}_5^{\mathrm{hs}}\) vanishes
if and only if \(D_5=0\).
\ClaimStatusConditional{}
\end{corollary}

\begin{proof}
Subtract the residue-side one-coefficient normal form from the principal-side normal form.
\end{proof}

\begin{proposition}[Stage-\texorpdfstring{$5$}{5} entry-packet identities]
\label{prop:winfty-stage5-entry-identities}
\ClaimStatusConditional{}
Assume the hypotheses of
Remark~\ref{rem:winfty-stage5-entry-transport}. Then
\[
\mathsf{C}^{\mathrm{res}}_{s,t;u;m,n}(5)
=
\mathsf{C}^{\mathrm{DS}}_{s,t;u;m,n}(5)
\]
for every channel $(s,t;u;m,n)$ in the two-channel packet
$\mathcal{J}_5^{\mathrm{entry}}$. Equivalently, the singleton channel
identities
Propositions~\ref{prop:winfty-stage5-block-34}
and~\ref{prop:winfty-stage5-block-55}
both hold, in the reduction order dictated by
Remark~\ref{rem:winfty-stage5-entry-mixed-self}.
\end{proposition}

\begin{proposition}[Stage-\texorpdfstring{$5$}{5} mixed transport identities]
\label{prop:winfty-stage5-transport-identities}
\ClaimStatusConditional{}
Assume the hypotheses of
Remark~\ref{rem:winfty-stage5-entry-transport}. Then
\[
\mathsf{C}^{\mathrm{res}}_{s,t;u;m,n}(5)
=
\mathsf{C}^{\mathrm{DS}}_{s,t;u;m,n}(5)
\]
for every channel $(s,t;u;m,n)$ in the six-channel packet
$\mathcal{J}_5^{\mathrm{tr}}$. Equivalently, the three fixed-target
ladder identities
Propositions~\ref{prop:winfty-stage5-transport-target-3}
--\ref{prop:winfty-stage5-transport-target-5}
all hold. By
Remark~\ref{rem:winfty-stage5-transport-pole-profiles}, the
natural local reduction order is target~$5$, then target~$4$, then
target~$3$.
\end{proposition}

\begin{proposition}[Stage-\texorpdfstring{$5$}{5} transport target-\texorpdfstring{$3$}{3} ladder identities]
\label{prop:winfty-stage5-transport-target-3}
\ClaimStatusProvedElsewhere{}
Assume the hypotheses of
Remark~\ref{rem:winfty-stage5-entry-transport}. Then
\[
\mathsf{C}^{\mathrm{res}}_{s,t;u;m,n}(5)
=
\mathsf{C}^{\mathrm{DS}}_{s,t;u;m,n}(5)
\]
for every channel $(s,t;u;m,n)$ in the two-channel ladder
$\mathcal{J}_5^{\mathrm{tr},3}$ of
Remark~\ref{rem:winfty-stage5-transport-target-ladders}.
\end{proposition}

\begin{proposition}[Stage-\texorpdfstring{$5$}{5} transport target-\texorpdfstring{$4$}{4} ladder identities]
\label{prop:winfty-stage5-transport-target-4}
\ClaimStatusProvedElsewhere{}
Assume the hypotheses of
Remark~\ref{rem:winfty-stage5-entry-transport}. Then
\[
\mathsf{C}^{\mathrm{res}}_{s,t;u;m,n}(5)
=
\mathsf{C}^{\mathrm{DS}}_{s,t;u;m,n}(5)
\]
for every channel $(s,t;u;m,n)$ in the two-channel ladder
$\mathcal{J}_5^{\mathrm{tr},4}$ of
Remark~\ref{rem:winfty-stage5-transport-target-ladders}.
\end{proposition}

\begin{proposition}[Stage-\texorpdfstring{$5$}{5} transport target-\texorpdfstring{$5$}{5} ladder identities]
\label{prop:winfty-stage5-transport-target-5}
\ClaimStatusConditional{}
Assume the hypotheses of
Remark~\ref{rem:winfty-stage5-entry-transport}. Then
\[
\mathsf{C}^{\mathrm{res}}_{s,t;u;m,n}(5)
=
\mathsf{C}^{\mathrm{DS}}_{s,t;u;m,n}(5)
\]
for every channel $(s,t;u;m,n)$ in the two-channel ladder
$\mathcal{J}_5^{\mathrm{tr},5}$ of
Remark~\ref{rem:winfty-stage5-transport-target-ladders}. Together
with Proposition~\ref{prop:winfty-stage5-block-34}, this matches every
channel in the target-$5$ corridor of
Remark~\ref{rem:winfty-stage5-target5-corridor}; after the singleton
tail input of Remark~\ref{rem:winfty-stage5-reduced-tail-singleton}
is matched, this is exactly the residual continuation identified in
Remark~\ref{rem:winfty-stage5-target5-residual} and the mixed
transport mechanism identified in
Remark~\ref{rem:winfty-stage5-target5-transport-mechanism}.
Equivalently, the two singleton transport identities
Propositions~\ref{prop:winfty-stage5-transport-target5-35}
and~\ref{prop:winfty-stage5-transport-target5-45}
both hold, in the pole order dictated by
Remark~\ref{rem:winfty-stage5-target5-transport-singletons}.
\end{proposition}

\begin{proposition}[Stage-\texorpdfstring{$5$}{5} target-\texorpdfstring{$5$}{5} transport singleton from
\texorpdfstring{$W^{(3)}W^{(5)}$}{W3W5}]
\label{prop:winfty-stage5-transport-target5-35}
\ClaimStatusProvedElsewhere{}
Assume the hypotheses of
Remark~\ref{rem:winfty-stage5-entry-transport}. Then
\[
\mathsf{C}^{\mathrm{res}}_{3,5;5;0,3}(5)
=
\mathsf{C}^{\mathrm{DS}}_{3,5;5;0,3}(5).
\]
This is the first singleton in the residual target-$5$ continuation of
Remark~\ref{rem:winfty-stage5-target5-transport-singletons},
equivalently the pole-$3$ $W^{(5)}$-projection in
$W^{(3)}(z)W^{(5)}(w)$.
\end{proposition}

\begin{proposition}[Stage-\texorpdfstring{$5$}{5} target-\texorpdfstring{$5$}{5} transport singleton from
\texorpdfstring{$W^{(4)}W^{(5)}$}{W4W5}]
\label{prop:winfty-stage5-transport-target5-45}
\ClaimStatusProvedElsewhere{}
Assume the hypotheses of
Remark~\ref{rem:winfty-stage5-entry-transport}. Then
\[
\mathsf{C}^{\mathrm{res}}_{4,5;5;0,4}(5)
=
\mathsf{C}^{\mathrm{DS}}_{4,5;5;0,4}(5).
\]
This is the second singleton in the residual target-$5$ continuation of
Remark~\ref{rem:winfty-stage5-target5-transport-singletons},
equivalently the pole-$4$ $W^{(5)}$-projection in
$W^{(4)}(z)W^{(5)}(w)$.
\end{proposition}

\begin{proposition}[Stage-\texorpdfstring{$5$}{5}
\texorpdfstring{$(3,4)$}{(3,4)} higher-spin block identity]
\label{prop:winfty-stage5-block-34}
\ClaimStatusConditional{}
Assume the hypotheses of
Corollary~\ref{cor:winfty-stage5-higher-spin-packet}. Assume further
that the visible stage-$5$ quotient identifies its \(W^{(5)}\)-line
with the principal Drinfeld--Sokolov \(W^{(5)}\)-line, with the same
strong-generator sign, the normalized metric
\(\langle W^{(5)},W^{(5)}\rangle=c/5\), and the
Blumenhagen--Eholzer--Flohr--Honecker--Hornfeck--H\"ubel
\(W^{(3)}W^{(4)}W^{(5)}\) coupling. Then
\[
\mathsf{C}^{\mathrm{res}}_{3,4;5;0,2}(5)
=
\mathsf{C}^{\mathrm{DS}}_{3,4;5;0,2}(5).
\]
This is exactly the reduced tail comparison input isolated in
Remark~\ref{rem:winfty-stage5-reduced-tail-singleton}, the exact
missing comparison mechanism isolated in
Remark~\ref{rem:winfty-stage5-tail-mechanism}, and the
first singleton channel in the target-$5$ corridor of
Remark~\ref{rem:winfty-stage5-target5-corridor}.
\end{proposition}

\begin{proof}
On the normalized principal \(W_5\) quotient fix the strong-generator
sign of \(W^{(5)}\) and write \(g_{345}\) for the
\(W^{(3)}W^{(4)}W^{(5)}\) coupling. In the
Blumenhagen--Eholzer--Flohr--Honecker--Hornfeck--H\"ubel
normalization \cite{BEHFH96},
\[
g_{345}^{\,2}
=
\frac{1680\,c^2(5c+22)(2c-1)}
{(c+24)(7c+68)(3c+46)(10c+197)} .
\]
The visible two-point metric is
\(\eta_{55}=\langle W^{(5)},W^{(5)}\rangle=c/5\), and the invariant
three-point tensor is
\[
\langle W^{(3)}W^{(4)}W^{(5)}\rangle
=
\frac{c}{5}\,g_{345}.
\]
The residue coefficient is obtained by raising the \(W^{(5)}\)-index:
\[
\mathsf{C}^{\mathrm{res}}_{3,4;5;0,2}(5)
=
\eta^{55}\langle W^{(3)}W^{(4)}W^{(5)}\rangle
=
\frac{5}{c}\cdot\frac{c}{5}\,g_{345}
=g_{345}.
\]
The principal Drinfeld--Sokolov OPE coefficient
\(\mathsf{C}^{\mathrm{DS}}_{3,4;5;0,2}(5)\) is the same top-pole
coefficient \(g_{345}\) in the displayed normalization. Therefore the
two coefficients agree. Changing the sign of \(W^{(5)}\) changes both
sides by the same sign; the equality is a statement after the visible
\(W^{(5)}\)-line is identified with the principal \(W^{(5)}\)-line.
\end{proof}

\begin{computation}[Exact stage-\texorpdfstring{$5$}{5}
\texorpdfstring{$(3,4)$}{(3,4)} block coefficient check]
\label{comp:winfty-stage5-block34-coefficient}
The coefficient equality of
Proposition~\ref{prop:winfty-stage5-block-34} is checked by the exact
engine \texttt{compute/lib/winfty\_stage5\_block34.py}. The test
surface \texttt{compute/tests/test\_winfty\_stage5\_block34.py}
compares three independent paths:
\begin{enumerate}[label=\textup{(\roman*)}]
\item the metric contraction
\(\eta^{55}\langle W^{(3)}W^{(4)}W^{(5)}\rangle=g_{345}\);
\item the principal \(W_5\) bootstrap square
\[
g_{345}^{\,2}
=
\frac{1680\,c^2(5c+22)(2c-1)}
{(c+24)(7c+68)(3c+46)(10c+197)};
\]
\item the existing stage-$5$ visible normal-form ratios
\[
\mathsf{C}_{3,4;5;0,2}
=-\frac54\,\mathsf{C}_{3,5;4;0,4},
\qquad
\mathsf{C}_{4,5;3;0,6}
=-\frac34\,\mathsf{C}_{3,5;4;0,4}.
\]
\end{enumerate}
At \(c=10,24,26,123\) the exact rational square agrees with the two
pre-existing \(W_5\) engines; at \(c=\tfrac12\) the factor \(2c-1\)
kills \(g_{345}\) on both sides.
\end{computation}

\begin{proposition}[Stage-\texorpdfstring{$5$}{5}
\texorpdfstring{$(3,5)$}{(3,5)} higher-spin block identities]
\label{prop:winfty-stage5-block-35}
\ClaimStatusConditional{}
Assume the hypotheses of
Corollary~\ref{cor:winfty-stage5-higher-spin-packet}. Then
\[
\mathsf{C}^{\mathrm{res}}_{s,t;u;m,n}(5)
=
\mathsf{C}^{\mathrm{DS}}_{s,t;u;m,n}(5)
\]
for every channel $(s,t;u;m,n)$ in the three-channel block
$\mathcal{J}_5^{(3,5)}$ of
Remark~\ref{rem:winfty-stage5-higher-spin-subblocks}.
\end{proposition}

\begin{proposition}[Stage-\texorpdfstring{$5$}{5}
\texorpdfstring{$(4,5)$}{(4,5)} higher-spin block identities]
\label{prop:winfty-stage5-block-45}
\ClaimStatusConditional{}
Assume the hypotheses of
Corollary~\ref{cor:winfty-stage5-higher-spin-packet}. Then
\[
\mathsf{C}^{\mathrm{res}}_{s,t;u;m,n}(5)
=
\mathsf{C}^{\mathrm{DS}}_{s,t;u;m,n}(5)
\]
for every channel $(s,t;u;m,n)$ in the three-channel block
$\mathcal{J}_5^{(4,5)}$ of
Remark~\ref{rem:winfty-stage5-higher-spin-subblocks}.
\end{proposition}

\begin{proposition}[Stage-\texorpdfstring{$5$}{5}
\texorpdfstring{$(5,5)$}{(5,5)} self higher-spin block identity]
\label{prop:winfty-stage5-block-55}
\ClaimStatusConditional{}
Assume the hypotheses of
Corollary~\ref{cor:winfty-stage5-higher-spin-packet}. Then
\[
\mathsf{C}^{\mathrm{res}}_{5,5;4;0,6}(5)
=
\mathsf{C}^{\mathrm{DS}}_{5,5;4;0,6}(5).
\]
\end{proposition}

\begin{proposition}[First higher-spin bar-vs-DS packet beyond
\texorpdfstring{$\mathcal{I}_4$}{I4}]
\label{prop:winfty-stage5-higher-spin-identities}
\ClaimStatusConditional{}
Assume the hypotheses of
Corollary~\ref{cor:winfty-stage5-residue-eight-channel}. Then
\[
\mathsf{C}^{\mathrm{res}}_{s,t;u;m,n}(5)
=
\mathsf{C}^{\mathrm{DS}}_{s,t;u;m,n}(5)
\]
for every channel $(s,t;u;m,n)$ in the explicit packet
$\mathcal{J}_5^{\mathrm{hs}}$ of
Corollary~\ref{cor:winfty-stage5-higher-spin-packet}. This is the
next finite list of higher-spin bar-vs-Drinfeld--Sokolov identities
after the stage-$4$ four-channel packet. Equivalently, the four
subblock identities
Propositions~\ref{prop:winfty-stage5-block-34}
--\ref{prop:winfty-stage5-block-55}
all hold, or, equivalently again, the entry-packet and mixed-transport
identities
Propositions~\ref{prop:winfty-stage5-entry-identities}
and~\ref{prop:winfty-stage5-transport-identities}
both hold. On the full visible
\texorpdfstring{$W^{(3)}$}{W3}/\texorpdfstring{$W^{(4)}$}{W4}/\texorpdfstring{$W^{(5)}$}{W5}
pairing locus,
Proposition~\ref{prop:winfty-stage5-principal-one-coefficient-normal-form}
together with Proposition~\ref{prop:winfty-stage5-one-coefficient-reduction}
reduce this whole packet identity to the single scalar identity of
Proposition~\ref{prop:winfty-stage5-one-coefficient-comparison},
and hence through the packaged principal one-coefficient normal form.
\end{proposition}

\begin{corollary}[Visible stage-\texorpdfstring{$5$}{5} local network collapses
to one nontrivial singleton under principal normal form;
\ClaimStatusConditional]
\label{cor:winfty-stage5-visible-network-collapse}%
\label{cor:winfty-stage5-visible-conjecture-network-collapse}%
Assume the hypotheses of
Corollary~\ref{cor:winfty-stage5-effective-independent-frontier},
and
the packaged principal stage-$5$ one-coefficient normal form.
Then on the full visible
\texorpdfstring{$W^{(3)}$}{W3}/\texorpdfstring{$W^{(4)}$}{W4}/\texorpdfstring{$W^{(5)}$}{W5}
pairing locus:
\begin{enumerate}[label=\textup{(\roman*)}]
\item the packet-level stage-$5$ conditions
 \textup{(}\ref{prop:winfty-stage5-higher-spin-identities},
 \ref{prop:winfty-stage5-entry-identities},
 \ref{prop:winfty-stage5-transport-identities},
 \ref{prop:winfty-stage5-transport-target-3},
 \ref{prop:winfty-stage5-transport-target-4},
 \ref{prop:winfty-stage5-block-34}--\ref{prop:winfty-stage5-block-45}\textup{)}
 are each equivalent to the single target-$4$ singleton identity
 \[
 \mathsf{C}^{\mathrm{res}}_{3,5;4;0,4}(5)
 =
 \mathsf{C}^{\mathrm{DS}}_{3,5;4;0,4}(5);
 \]
\item the target-$5$ ladder conditions
 \textup{(}\ref{prop:winfty-stage5-transport-target-5},
 \ref{prop:winfty-stage5-transport-target5-35},
 \ref{prop:winfty-stage5-transport-target5-45},
 \ref{prop:winfty-stage5-block-55}\textup{)} become automatic.
\end{enumerate}
In particular, once the packaged principal one-coefficient normal form
is granted, the local stage-$5$ condition tree below
\(\mathcal{J}_5^{\mathrm{hs}}\) has exactly one nontrivial comparison
input, namely the target-$4$ singleton \((3,5;4;0,4)\).
\end{corollary}

\begin{proof}
Under the one-coefficient normal form on both sides, each block
condition is either automatic (self-block, target-$5$ ladder: all
channels vanish on both sides) or equivalent to the singleton identity
\(\mathsf{C}^{\mathrm{res}}_{3,5;4;0,4}=\mathsf{C}^{\mathrm{DS}}_{3,5;4;0,4}\)
via the fixed rational multiples \(-\tfrac54\) (mixed-entry), \(1\)
(representative), and \(-\tfrac34\) (target-$3$). The entry,
transport, and packet conditions are conjunctions of these, hence each
equivalent to the singleton identity.
\end{proof}

\begin{corollary}[Visible stage-\texorpdfstring{$5$}{5} local defect classes under principal
normal form]
\label{cor:winfty-stage5-visible-defect-classes}
Assume the hypotheses of
Corollary~\ref{cor:winfty-stage5-visible-network-collapse},
and set
\[
D_5:=
\mathsf{C}^{\mathrm{res}}_{3,5;4;0,4}(5)
-
\mathsf{C}^{\mathrm{DS}}_{3,5;4;0,4}(5).
\]
Then the downstream local stage-$5$ comparison surfaces on the full
visible
\texorpdfstring{$W^{(3)}$}{W3}/\texorpdfstring{$W^{(4)}$}{W4}/\texorpdfstring{$W^{(5)}$}{W5}
pairing locus split into the following exact defect classes:
\begin{enumerate}[label=\textup{(\roman*)}]
\item the mixed-entry block
 \textup{(}\ref{prop:winfty-stage5-block-34}\textup{)} carries defect
 \[
 \mathsf{C}^{\mathrm{res}}_{3,4;5;0,2}(5)
 -
 \mathsf{C}^{\mathrm{DS}}_{3,4;5;0,2}(5)
 =
 -\frac{5}{4}\,D_5;
 \]
\item the representative target-$4$ class consists of the singleton
 comparison \textup{(}\ref{prop:winfty-stage5-one-coefficient-comparison}\textup{)},
 the target-$4$ ladder \textup{(}\ref{prop:winfty-stage5-transport-target-4}\textup{)}, and the
 \texorpdfstring{$(3,5)$}{(3,5)} block
 \textup{(}\ref{prop:winfty-stage5-block-35}\textup{)}, with representative
 defect exactly \(D_5\);
\item the target-$3$ class consists of the target-$3$ ladder
 \textup{(}\ref{prop:winfty-stage5-transport-target-3}\textup{)} and the
 \texorpdfstring{$(4,5)$}{(4,5)} block
 \textup{(}\ref{prop:winfty-stage5-block-45}\textup{)}, with representative
 defect
 \[
 \mathsf{C}^{\mathrm{res}}_{4,5;3;0,6}(5)
 -
 \mathsf{C}^{\mathrm{DS}}_{4,5;3;0,6}(5)
 =
 -\frac{3}{4}\,D_5;
 \]
\item the target-$5$ ladder
 \textup{(}\ref{prop:winfty-stage5-transport-target-5},
 \ref{prop:winfty-stage5-transport-target5-35},
 \ref{prop:winfty-stage5-transport-target5-45},
 \ref{prop:winfty-stage5-block-55}\textup{)}
 are automatic, so their defects vanish.
\end{enumerate}
Equivalently, every nonautomatic local stage-$5$ visible-pairing
surface lies in one of the three exact defect classes
\(-\tfrac54 D_5\), \(D_5\), or \(-\tfrac34 D_5\).
\ClaimStatusConditional{}
\end{corollary}

\begin{proof}
The fixed ratios from Corollary~\ref{cor:winfty-stage5-visible-network-collapse} give defects \(-\tfrac54 D_5\), \(D_5\), and \(-\tfrac34 D_5\); automatic channels have zero defect.
\end{proof}

\begin{theorem}[All-stages \texorpdfstring{$\mathcal{W}_\infty$}{W-infinity}
MC4 rigidity reduction criterion;
\ClaimStatusConditional]%
\label{thm:winfty-all-stages-rigidity-closure}%
\index{MC4!all-stages closure}%
\index{W-infinity@$\mathcal{W}_\infty$!all-stages closure}%
Assume the hypotheses of
Proposition~\textup{\ref{prop:winfty-mc4-frontier-package}}, and assume
further that at each stage~$N$ the stage-$N$ quotient of the
$\mathcal{W}_\infty^{\mathrm{fact}}$ target carries visible generators
$W^{(2)},\ldots,W^{(N)}$ that are even Virasoro-primary fields with
the Virasoro field $T:=W^{(2)}$ acting by the standard Ward identity.
Assume, in addition, the principal-stage identification input:
for every $N\ge4$ the visible quotient is identified, after fixing
central charge and strong-generator signs, with the principal
Drinfeld--Sokolov stage $W_N$ and its normalized OPE constants.
Then for every $N \ge 4$, all stage-$N$ defects vanish:
\[
\mathsf{D}_{s,t;u;m,n}(N)
:=
\mathsf{C}^{\mathrm{res}}_{s,t;u;m,n}(N)
-
\mathsf{C}^{\mathrm{DS}}_{s,t;u;m,n}(N)
=
0
\qquad
\text{for all $(s,t;u;m,n) \in \mathcal{I}_N$.}
\]
In particular, under this added identification input, the stage-$4$
six-defect block, the stage-$5$
identity network
\textup{(}Propositions~\textup{\ref{prop:winfty-stage5-principal-target5-no-new-independent-data}}%
--\textup{\ref{prop:winfty-stage5-principal-one-coefficient-normal-form}},
\textup{\ref{prop:winfty-stage5-one-coefficient-comparison}}\textup{)},
and all higher-stage incremental packets $\mathcal{J}_{N+1}$ vanish.
\end{theorem}

\begin{proof}
The principal-stage identification is exactly the missing comparison
map between the residue calculus and the Drinfeld--Sokolov stage.
After central charge and strong-generator signs are fixed, it
identifies every visible OPE coefficient with the corresponding
principal coefficient on the finite packet~$\mathcal{I}_N$. Hence each
defect
\(\mathsf{C}^{\mathrm{res}}_{s,t;u;m,n}(N)
-\mathsf{C}^{\mathrm{DS}}_{s,t;u;m,n}(N)\)
vanishes.

For \(N=4\), the principal constants are the explicit rational
functions and Virasoro-target values of
Proposition~\ref{prop:w4-ds-ope-explicit}; the added identification is
therefore precisely the input required by
Propositions~\ref{prop:winfty-stage4-visible-diagonal-normalization},
\ref{prop:winfty-stage4-ward-inheritance}, and
\ref{prop:winfty-stage4-visible-borcherds-transport}. For \(N=5\),
the same identification is the principal one-coefficient normal-form
package of
Proposition~\ref{prop:winfty-stage5-principal-one-coefficient-normal-form}
together with the singleton comparison
Proposition~\ref{prop:winfty-stage5-one-coefficient-comparison}.
Without this finite-stage identification, the theorem supplies a
reduction criterion rather than a closure theorem.
\end{proof}

\begin{computation}[Full stage-\texorpdfstring{$4$}{4} reduction chain check]
\label{comp:winfty-stage4-reduction-chain}
\index{MC4!W-infinity stage-4 reduction chain}
The complete reduction chain from the residual packet
$\mathcal{J}_4$ to the completion packet of
Proposition~\ref{prop:winfty-mc4-frontier-package} has been
checked computationally
(\texttt{compute/lib/w4\_stage4\_coefficients.py}, $85$~tests):
\begin{itemize}
\item \emph{Primaryity}
 (Corollary~\ref{cor:winfty-ds-stage4-top-pole-packet}):
 $|\mathcal{J}_4| = 29 \to |\mathcal{J}_4^{\mathrm{top}}| = 7$,
 eliminating $22$ sub-leading entries; seven top-pole elements
 checked against the manuscript display;
\item \emph{Parity}
 (Corollary~\ref{cor:winfty-ds-stage4-parity-packet}):
 $|\mathcal{J}_4^{\mathrm{top}}| = 7
 \to |\mathcal{J}_4^{\mathrm{par}}| = 6$,
 removing exactly $(4,4,3,5)$ (odd-spin self-OPE);
\item \emph{OPE blocks}
 (Corollary~\ref{cor:winfty-ds-stage4-ope-blocks}):
 $6 = 1 + 3 + 2$ as $(3,3)$, $(3,4)$, $(4,4)$ blocks;
\item \emph{Mixed-self split}
 (Corollary~\ref{cor:winfty-ds-stage4-mixed-self-split}):
 three self-coupling scalars $+$ three mixed channels;
\item \emph{Virasoro elimination}
 (Corollary~\ref{cor:winfty-ds-stage4-five-plus-zero}):
 $\mathsf{C}^{\mathrm{DS}}_{4,4;2;0,6} = 2$ (universal $T$-coupling)
 and $\mathsf{C}^{\mathrm{DS}}_{3,4;2;0,5} = 0$ (mixed vanishing)
 split the stage-$4$ packet into $4$ higher-spin channels plus
 $2$ proved principal Virasoro-target values.
\end{itemize}
The full chain $29 \to 7 \to 6 = 4 + 2$ matches every intermediate
cardinality and every explicit element of the manuscript displays.
\end{computation}

\begin{proposition}[Explicit \texorpdfstring{$W_4$}{W4} OPE structure constants {\cite{Hornfeck93}};
\ClaimStatusConditional]
\label{prop:w4-ds-ope-explicit}
In the principal Drinfeld--Sokolov reduction
$W_4 = W(\mathfrak{sl}_4, f_{\mathrm{prin}})$ with generators
$T$ (spin~$2$), $W^{(3)}$ (spin~$3$), $W^{(4)}$ (spin~$4$)
normalized so that
$\langle W^{(s)}(z)\,W^{(s)}(0)\rangle = (c/s)\,z^{-2s}$,
the four higher-spin stage-$4$ principal targets of
Corollary~\textup{\ref{cor:winfty-ds-stage4-five-plus-zero}} are
given by the following rational functions of the central charge~$c$:
\begin{align}
c_{334}^{\,2}
&= \frac{42\,c^2(5c+22)}{(c+24)(7c+68)(3c+46)},
\label{eq:c334-explicit}
\\[4pt]
c_{444}^{\,2}
&= \frac{112\,c^2(2c-1)(3c+46)}{(c+24)(7c+68)(10c+197)(5c+3)},
\label{eq:c444-explicit}
\\[4pt]
\mathsf{C}_{3,4;3;0,4}^{\,2}
&= \tfrac{9}{16}\,c_{334}^{\,2},
\label{eq:c343-explicit}
\\[4pt]
\mathsf{C}_{3,4;4;0,3}^{\,2}
&= \tfrac{5}{7}\,c_{334}^{\,2}.
\label{eq:c344-explicit}
\end{align}
The two proved Virasoro-target identities in
Corollary~\textup{\ref{cor:winfty-ds-stage4-five-plus-zero}} are
confirmed:
\begin{equation}
\mathsf{C}^{\mathrm{DS}}_{4,4;2;0,6} = 2,
\qquad
\mathsf{C}^{\mathrm{DS}}_{3,4;2;0,5} = 0.
\label{eq:w4-virasoro-target-identities-values}
\end{equation}
\end{proposition}

\begin{proof}
The structure constants follow from the vertex-algebra Jacobi
identity (Borcherds identity) applied to the three-generator $W_4$
algebra. Formula~\eqref{eq:c334-explicit} is the bootstrap result
of Hornfeck~\cite{Hornfeck93}; formula~\eqref{eq:c444-explicit}
then follows from the $(W^{(4)}, W^{(4)}, W^{(4)})$ Jacobi identity.
The inter-coefficient relations~\eqref{eq:c343-explicit}
and~\eqref{eq:c344-explicit} are metric-adjoint and Borcherds-identity
consequences of~\eqref{eq:c334-explicit}, respectively; see
Bouwknegt--Schoutens~\cite{Bouwknegt-Schoutens}
and Blumenhagen \textit{et al.}~\cite{BEHFH96}.
The universal $T$-coupling
$\mathsf{C}_{4,4;2;0,6}=2$ follows from
Proposition~\ref{prop:winfty-ds-self-t-coefficient}, and the mixed
Virasoro vanishing
$\mathsf{C}_{3,4;2;0,5}=0$ from
Proposition~\ref{prop:winfty-ds-mixed-virasoro-ds-zero}.
\end{proof}

\begin{corollary}[Principal stage-\texorpdfstring{$4$}{4} higher-spin packet from two
primitive square classes; \ClaimStatusConditional]
\label{cor:w4-ds-stage4-square-class-reduction}
In the setting of
Proposition~\ref{prop:w4-ds-ope-explicit}, the higher-spin part of the
principal stage-$4$ target packet is determined at signless
square-class level by the pair
\[
c_{334}^{\,2},
\qquad
c_{444}^{\,2}.
\]
More precisely, the two mixed higher-spin channels are forced by
$c_{334}^{\,2}$:
\[
\mathsf{C}_{3,4;3;0,4}^{\,2}
= \tfrac{9}{16}\,c_{334}^{\,2},
\qquad
\mathsf{C}_{3,4;4;0,3}^{\,2}
= \tfrac{5}{7}\,c_{334}^{\,2}.
\]
Consequently, after removing the proved Virasoro-target values
\[
\mathsf{C}^{\mathrm{DS}}_{4,4;2;0,6}=2,
\qquad
\mathsf{C}^{\mathrm{DS}}_{3,4;2;0,5}=0,
\]
the principal stage-$4$ DS packet does not require four primitive
higher-spin targets: it consists of the two primitive square classes
$c_{334}^{\,2},c_{444}^{\,2}$ together with two mixed square-class
descendants forced by~$c_{334}^{\,2}$.
\end{corollary}

\begin{proof}
The explicit formulas
\eqref{eq:c343-explicit} and~\eqref{eq:c344-explicit} express the two
mixed higher-spin squares directly in terms of~$c_{334}^{\,2}$, while
\eqref{eq:w4-virasoro-target-identities-values} records the two
proved Virasoro-target values. Hence the displayed pair
$c_{334}^{\,2},c_{444}^{\,2}$ is the entire primitive higher-spin
square-class input on the principal side.
\end{proof}

\begin{remark}[Exact stage-4 identity packet]
\label{rem:winfty-stage4-exact-identity-packet}
Proposition~\ref{prop:winfty-mc4-frontier-package} and
Proposition~\ref{prop:w4-ds-ope-explicit} together show that the first
genuinely higher-spin MC4 packet is no longer extraction of principal
Drinfeld--Sokolov data. That side is already explicit. The live
stage-$4$ task is the six bar-side identities
\[
\mathsf{C}^{\mathrm{res}}_{s,t;u;m,n}(4)
=
\mathsf{C}^{\mathrm{DS}}_{s,t;u;m,n}(4)
\]
on the packet
\[
\Bigl\{
\begin{gathered}
(3,3;4;0,2),\,
(4,4;4;0,4),\\
(3,4;3;0,4),\,
(3,4;4;0,3),\,
(4,4;2;0,6),\\
(3,4;2;0,5)
\end{gathered}
\Bigr\}.
\]
For the first four channels, the DS targets are fixed by the explicit
formulas \eqref{eq:c334-explicit}--\eqref{eq:c344-explicit}, up to the
chosen strong-generator signs. The last two channels are fixed exactly
by \eqref{eq:w4-virasoro-target-identities-values}. Thus the first open higher-spin
packet is an exact six-entry identity problem on~$\mathcal{I}_4$, not a
further principal-side extraction problem.
\end{remark}

\begin{computation}[Full \texorpdfstring{$W_4$}{W4} OPE coefficient extraction]
\label{comp:w4-ds-ope-extraction}
\index{MC4!W4 OPE coefficient extraction}
The explicit structure constants of
Proposition~\ref{prop:w4-ds-ope-explicit} have been checked
computationally
(\texttt{compute/lib/w4\_ds\_ope\_extraction.py}, $121$~tests):
\begin{itemize}
\item all four squared structure constants evaluated at levels
 $k = 0, 1, 2, -2, -3$ (avoiding poles);
\item inter-coefficient ratios
 $\mathsf{C}_{3,4;3;0,4}^2/c_{334}^2 = 9/16$ and
 $\mathsf{C}_{3,4;4;0,3}^2/c_{334}^2 = 5/7$ checked symbolically;
\item Feigin--Frenkel duality $k \leftrightarrow -k-2h^\vee$
 checked: $c(k) + c(-k-8) = 246$;
\item unitarity non-negativity ($c_{334}^2 \geq 0$, $c_{444}^2 \geq 0$)
 confirmed at all $\mathcal{W}(N, N{+}1)$ minimal models
 for $4 \leq N \leq 19$;
\item both proved principal Virasoro-target values
 $\mathsf{C}^{\mathrm{DS}}_{4,4;2;0,6} = 2$ and
 $\mathsf{C}^{\mathrm{DS}}_{3,4;2;0,5} = 0$ confirmed.
\end{itemize}
\end{computation}

The completion-closure theory and the $\mathcal{W}_\infty$ seed-packet
analysis above establish the structural content of MC4: the strong
filtration axiom produces completed bar-cobar quasi-isomorphisms,
and the finite packet calculus isolates the exact stagewise residue
identities. The remainder of this chapter returns
to the foundational algebraic layer: the role of Maurer--Cartan
elements in the curved setting, the genus-by-genus obstruction
analysis, the bar complex as an algebra over the modular operad, and
the separation of reconstruction from duality that underlies the
entire monograph.

\subsection{Maurer--Cartan elements and deformation theory}
\label{sec:maurer-cartan-curved}

\begin{definition}[Maurer--Cartan element in curved context]
\label{def:mc-element-curved}
\label{ex:maurer-cartan}
Let $(A,\{\mu_n\}_{n\geq0})$ be a complete filtered curved
$A_\infty$ algebra, with $\mu_n$ of degree $2-n$ and with all sums
convergent in the filtration. A \emph{Maurer--Cartan element} is an
element $\alpha\in F^1A^1$ satisfying
\begin{equation}
\label{eq:mc-equation}
\sum_{n \geq 0}\mu_n(\alpha^{\otimes n})=0.
\end{equation}
This is the $A_\infty$ convention. For the ordered convolution dg Lie
algebra the same condition is written
$D\Theta+\tfrac12[\Theta,\Theta]=0$.
\end{definition}

\begin{theorem}[Twisting by MC elements {\cite{LV12}}; \ClaimStatusProvedElsewhere]
\label{thm:twisting-mc}
\label{ex:mc-deformations}
Given an MC element $\alpha$, the twisted curved $A_\infty$ structure is
\begin{equation}
\mu_n^\alpha(a_1,\ldots,a_n)
=
\sum_{i_0,\ldots,i_n\geq0}
\mu_{n+i_0+\cdots+i_n}
(\alpha^{\otimes i_0},a_1,\alpha^{\otimes i_1},
\ldots,a_n,\alpha^{\otimes i_n}),
\end{equation}
with the Koszul signs determined by the suspension convention of
Appendix~\ref{app:signs-shifts}.

The twisted structure $(A, \{\mu_n^\alpha\})$ is again a curved $A_\infty$ algebra, with new
curvature:
\begin{equation}
\mu_0^\alpha=\sum_{n\geq0}\mu_n(\alpha^{\otimes n}).
\end{equation}

If $\alpha$ is an MC element, then $\mu_0^\alpha = 0$, so the twisted structure is
\emph{uncurved}.
\end{theorem}

\begin{remark}[MC elements as uncurving data]
\label{rem:mc-uncurving}
If $(A, \{\mu_n\})$ is curved with $\mu_0 \neq 0$, an MC element
$\alpha \in A^1$ satisfying $\sum_{n \geq 0} \mu_n(\alpha^{\otimes n}) = 0$
produces the twisted structure $(A, \{\mu_n^\alpha\})$ with $\mu_0^\alpha = 0$:
the MC equation is the condition that twisting by $\alpha$ removes the curvature.
If no MC element exists, the curvature persists: the raw CDG
fiberwise differential over $\overline{\mathcal{M}}_g$ has
$d_{\mathrm{fib}}^2 = \kappa \cdot \omega_g \neq 0$. Ordinary
cohomology is then taken only after replacing it by the
period-corrected square-zero total differential; otherwise the
comparison lives in the coderived comodule category.
\end{remark}

\subsection{Obstruction hierarchy}
\label{sec:curved-summary}

The first three regimes of Convention~\ref{conv:regime-tags}
are distinguished by whether $d_0^{\,2} = 0$ holds without
correction and whether completion is needed.
For the standard positive-energy vertex-operator families whose anomaly
is a vacuum scalar, the curvature $\mu_0$ is central
($\mu_0 \propto \mathbf{1}$), so
$m_1^2(a) = [\mu_0, a] = 0$ and $m_1$ is a strict
differential on~$A$ in every regime.

\begin{center}
\renewcommand{\arraystretch}{1.3}
\begin{tabular}{llll}
\textbf{Regime} & \textbf{Characterization} & \textbf{Bar differential} & \textbf{Example} \\
\hline
I.\;Quadratic & $d_0^{\,2} = 0$
 & $d_{\mathrm{bar}}^2 = 0$
 & $\mathcal{H}_k$,\; $V_\Lambda$ \\[2pt]
II.\;Curved & $d_0^{\,2} \neq 0$;\; finitely generated
 & $D_{\mathrm{tot}}^2 = 0$;\; $d_{\mathrm{fib}}^2 = \kappa \cdot \omega_g$
 & $\mathrm{Vir}_c$,\; $\widehat{\mathfrak{g}}_k$ \\[2pt]
III.\;Filtered & non-quadratic;\; completion needed
 & $\widehat{d}_{\mathrm{bar}}$ on $\varprojlim_r \overline{B}(A)/F^r$
 & $\mathcal{W}_3$ \\
\end{tabular}
\end{center}

\noindent
Regime~I: $d_0^{\,2} = 0$, so
$d_{\mathrm{bar}}^2 = 0$ on the nose; the classical
bar-cobar adjunction
(Theorem~\ref{thm:bar-cobar-adjunction}) holds without
qualification. This includes both algebras with no simple
poles ($\mathcal{H}_k$: $d_{\mathrm{bracket}} = 0$) and
algebras whose bracket satisfies the Jacobi identity
($V_\Lambda$: $d_{\mathrm{bracket}}^2 = 0$).
Regime~II: $d_0^{\,2} \neq 0$, but the algebra is finitely
generated and the bar differential on the
\emph{period-corrected total} complex satisfies
$d_{\mathrm{bar}}^2 = 0$; the fiberwise differential over
$\overline{\mathcal{M}}_g$ acquires geometric curvature
$d_{\mathrm{fib}}^2 = \kappa(\mathcal{A}) \cdot \omega_g$;
the corrected differential $D^{(g)}$ of
Theorem~\ref{thm:quantum-diff-squares-zero} absorbs this.
Regime~III: the OPE involves composite fields, forcing a
filtration and the completed bar complex
$\widehat{\overline{B}}(A)$ with the strong filtration
axiom
$\mu_r(F^{i_1}, \ldots, F^{i_r}) \subset
F^{i_1 + \cdots + i_r}$
(Theorem~\ref{thm:completed-bar-cobar-strong}).

\section{\texorpdfstring{Curved $A_\infty$ structures: strict vs.\ homotopy nilpotence}{Curved A-infinity structures: strict vs homotopy nilpotence}}
\label{sec:strict-vs-homotopy}

For chiral algebras $\mathcal{A}$ with quantum corrections at genus $g$,
the curvature $\mu_0$ controls two distinct squares:
\begin{enumerate}[label=\textup{(\roman*)}]
\item \emph{Algebraic:}\; $m_1^2(a) = [\mu_0, a]$.
 When $\mu_0 \in Z(\mathcal{A})$, this vanishes for all~$a$, so $m_1$ is a strict differential on~$\mathcal{A}$. All standard chiral algebras (Heisenberg, Kac--Moody, Virasoro, $\mathcal{W}$-algebras) satisfy this: $\mu_0$ is proportional to the vacuum vector $\mathbf{1}$, which is central.
\item \emph{Geometric:}\; $\dfib^{\,2} = \kappa(\mathcal{A})\cdot\omega_g \neq 0$ whenever $\kappa(\mathcal{A}) \neq 0$
 (Convention~\ref{conv:higher-genus-differentials}).
 Even when $\mu_0$ is central, the fiberwise bar differential over~$\overline{\mathcal{M}}_g$ is \emph{not} an ordinary nilpotent differential; its square is controlled by the scalar invariant~$\kappa(\mathcal{A})$.
 Thus $\dfib$ defines the curved CDG structure, not an uncurved
 dg-coalgebra. Only the period-corrected differential $\Dg{g}$ gives
 the square-zero coalgebra differential used for ordinary cohomology.
\end{enumerate}
Under the hypotheses of Theorem~\ref{thm:quantum-diff-squares-zero},
the total corrected differential $\Dg{g}$ absorbs the geometric
curvature and satisfies $\Dg{g}^{\,2} = 0$.

\subsection{Mathematical foundations: three regimes}
\label{subsec:three-regimes}

The three regimes below correspond to the first three levels of the
regime hierarchy (Convention~\ref{conv:regime-tags}): quadratic
(Regime~I), curved-central (Regime~II), and filtered-complete
(Regime~III). The fourth level (programmatic) lies beyond the
scope of this curved bar-cobar analysis.

\subsubsection{\texorpdfstring{Regime I: strict differential ($d^2 = 0$ strict)}{Regime I: strict differential (d2 = 0 strict)}}
\label{def:strict-dg}

Regime~I is the quadratic locus of Convention~\ref{conv:regime-tags}: the bar differential satisfies $d_{\mathrm{bar}}^2 = 0$ on the nose, with no curvature or completion. The canonical examples (Heisenberg, free fields, lattice vertex algebras) are developed in \S\ref{sec:curved-koszul-i-adic} above and in Part~\ref{part:standard-landscape}.

\subsubsection{\texorpdfstring{Regime II: curved differential ($m_1^2 = [\mu_0, -]$, central curvature)}{Regime II: curved differential (m1 squared = commutator with mu0, central curvature)}}

\begin{theorem}[Central curvature and the square-zero total model;
\ClaimStatusProvedHere]\label{thm:central-implies-strict}
Let $(\mathcal{A},\{\mu_n\}_{n\geq0})$ be a complete filtered curved
chiral $A_\infty$ algebra whose curvature satisfies $\mu_1(\mu_0)=0$
and
\begin{equation}
\mu_0 \in Z(\mathcal{A}) := \{a \in \mathcal{A} \mid \mu_2(a \otimes b) = (-1)^{|a||b|}\mu_2(b \otimes a)
\text{ for all } b\}
\end{equation}
for the binary product. Then centrality kills the internal curvature
commutator:
\begin{equation}\label{eq:curved-m1-squared}
\mu_1^2(a)=\mu_2(\mu_0,a)-\mu_2(a,\mu_0)=0.
\end{equation}
An ordinary square-zero bar differential follows only after one of two
additional operations has been supplied:
\begin{enumerate}[label=\textup{(\roman*)}]
\item a convergent MC twist $\alpha$ with
$\mu_0^\alpha=\sum_{n\ge0}\mu_n(\alpha^{\otimes n})=0$; or
\item the period-corrected total differential $\Dg{g}$ of
Convention~\ref{conv:higher-genus-differentials}, whose boundary
terms satisfy the modular-operad relation of
Theorem~\ref{thm:quantum-diff-squares-zero}.
\end{enumerate}
In either case the chosen bar differential satisfies
\begin{equation}
d_{\text{bar}}^2 = 0 \quad \text{strictly}
\end{equation}
as an ordinary differential. The raw fiberwise CDG operator is not
included in this conclusion: it may still satisfy
$\dfib^{\,2}=\kappa(\cA)\omega_g$ when $\kappa(\cA)\neq0$.
\end{theorem}

\begin{proof}
The $n=1$ curved $A_\infty$ relation gives
\[
\mu_1^2(a)=\mu_2(\mu_0,a)-\mu_2(a,\mu_0)
=[\mu_0,a]_{\mu_2}.
\]
The closedness relation $\mu_1(\mu_0)=0$ is the $n=0$ relation.
Centrality of $\mu_0$ kills the commutator, proving
\eqref{eq:curved-m1-squared}.

The residue part of the genus-$0$ bar differential squares to zero by
the Arnold relation
\[
\eta_{ij}\wedge\eta_{jk}
+\eta_{jk}\wedge\eta_{ki}
+\eta_{ki}\wedge\eta_{ij}=0
\]
and anticommutes with the internal differential by the Borcherds
Leibniz identity for the chiral product. Thus the ordinary bar square
is controlled by the remaining curvature term.

In case~\textup{(i)}, twisting by $\alpha$ replaces the curvature by
\[
\mu_0^\alpha=\sum_{n\ge0}\mu_n(\alpha^{\otimes n})=0.
\]
The twisted $A_\infty$ relations are therefore uncurved, and the
associated bar coderivation squares to zero. In case~\textup{(ii)},
the total operator is the modular corrected differential $\Dg{g}$.
The identity $\Dg{g}^{\,2}=0$ is exactly
Theorem~\ref{thm:quantum-diff-squares-zero}: in convolution language it
is the MC equation
$D\Theta_\cA+\tfrac12[\Theta_\cA,\Theta_\cA]=0$ together with the
modular boundary relation. No sign computation for a bare insertion
operator $d_{\mathrm{correction}}$ proves square-zero by itself; the
square-zero assertion belongs to the MC twist or to the corrected total
model.
\end{proof}

\begin{remark}[Non-central curvature]\label{rem:non-central-curvature}
If the curvature $\mu_0$ is not central, then the raw curved
differential has nonzero curvature commutator and nilpotence holds
only up to homotopy: the bar complex requires
homotopy coherent structures, coderived filtrations whose
associated-graded differential is square-zero, and the full
$A_\infty$/$L_\infty$ framework at every level. Ordinary bar
cohomology and an unqualified bar spectral sequence are not defined
for the raw curved differential in this regime. The standard
positive-energy vertex-operator families with vacuum-scalar anomaly
have central curvature, so the internal commutator vanishes and the
period-corrected total differential is square-zero; the fiberwise CDG
operator remains curved when
$\kappa(\cA)\neq0$. See Remark~\ref{rem:voa-central-curvature} for
the precise mechanism.
\end{remark}

\begin{remark}[Central curvature in vertex operator algebras]\label{rem:voa-central-curvature}
For vertex operator algebras, the vacuum axiom guarantees that the vacuum vector
$\mathbf{1}\in V$ is central: $Y(\mathbf{1},z) = \mathrm{id}_V$. In the curved
$A_\infty$ framework, the curvature element $\mu_0$ arising from the genus-$g$
quantum correction is proportional to $\mathbf{1}$ (it equals the central charge or
level times $\mathbf{1}$; see Examples~\ref{ex:heisenberg-strict}--\ref{ex:w3-strict}).
Since $\mathbf{1}$ acts by the identity on every module, $\mu_0$ lies in $Z(\mathcal{A})$,
and Theorem~\ref{thm:central-implies-strict} yields
$d_{\mathrm{bar}}^2=0$ for the period-corrected total bar
differential. The raw fiberwise CDG differential still has square
$\kappa(\cA)\omega_g$ when the modular characteristic is nonzero.

For more general chiral algebras without a vacuum axiom (for instance, chiral
algebras on higher-dimensional varieties, or non-unital factorization algebras),
the curvature need not be central, and $d^2=0$ may hold only up to homotopy.
Similarly, chiral algebras arising from field theories with non-central anomalies
(such as gravitational anomalies that act non-trivially on the operator algebra)
fall outside the scope of Theorem~\ref{thm:central-implies-strict}.
\end{remark}

\subsubsection{\texorpdfstring{Regime III: general homotopy coherent ($d^2 \sim 0$ via homotopy)}{Regime III: general homotopy coherent (d2 0 via homotopy)}}

\begin{definition}[Homotopy coherent differential]\label{def:homotopy-coherent}
A \emph{homotopy coherent differential} on a graded space $V$ consists of:
\begin{itemize}
\item $d_1: V \to V[1]$ (the ``differential'')
\item $h: V \to V[-1]$ (a homotopy)
\item Satisfying: $d_1^2 = [d_1, h]$ (not zero, but homotopic to zero)
\item Plus higher coherence homotopies $h_2, h_3, \ldots$ ad infinitum
\end{itemize}
\end{definition}

This is the setting of Lurie's $(\infty, 1)$-categories~\cite{HTT},
derived algebraic geometry (To\"en--Vezzosi, Lurie), and non-curved
$A_\infty$ or $L_\infty$ structures. The standard positive-energy
vertex-operator families in this chapter have central curvature
(Remark~\ref{rem:voa-central-curvature}), so this level of generality
is used only for non-central anomaly surfaces.

\subsection{Application to chiral algebras}

\subsubsection{\texorpdfstring{Heisenberg algebra (level $k$)}{Heisenberg algebra (level k)}}

\begin{example}[Heisenberg: strict nilpotence]\label{ex:heisenberg-strict}
The Heisenberg algebra $\mathcal{H}_k$ has OPE $J(z)J(w) = k/(z-w)^2 + \text{reg.}$, curvature $\mu_0 = k \cdot \mathbf{1}$, and $\mu_0 \in Z(\mathcal{H}_k)$ since $\mathbf{1}$ is central.

The period-corrected total bar differential satisfies
$d_{\text{bar}}^2 = 0$ strictly.
At genus~$1$ the raw fiberwise operator has curvature
\[
\dfib^{\,2}=k\cdot\omega_1\cdot\mathrm{id}
\]
in the notation of Convention~\ref{conv:higher-genus-differentials}.
The corrected operator is
\[
\Dg{1}=\dfib+d_{\mathrm{per}}^{(1)}
\]
with the prime-form period term included. The identity
$(\Dg{1})^2=0$ is the genus-$1$ case of
Theorem~\ref{thm:quantum-diff-squares-zero}; it is not proved by
discarding a term such as $\omega_1^2$. Centrality of
$k\mathbf{1}$ kills the internal commutator, while the period
correction cancels the geometric curvature.
\end{example}

\subsubsection{\texorpdfstring{Affine Kac--Moody (level $k$)}{Affine Kac--Moody (level k)}}

\begin{example}[Kac--Moody: strict nilpotence]\label{ex:kac-moody-strict}
For $\widehat{\mathfrak{g}}_k$ with OPE $J^a(z)J^b(w) = k\kappa^{ab}/(z-w)^2 + f^{abc}J^c(w)/(z-w) + \text{reg.}$, the curvature is $\mu_0 = k \cdot K$ where $K$ generates the one-dimensional center.

Again the period-corrected total differential satisfies
$d_{\text{bar}}^2 = 0$ strictly.
At genus~$g$, the correction involves:
\[\mu_0^{(g)} = \kappa(\widehat{\fg}_k) \cdot \lambda_g \in H^2(\mathcal{M}_g, Z(\widehat{\mathfrak{g}}_k))\]
where $\lambda_g$ is the Hodge class and $\kappa(\widehat{\fg}_k) = \dim(\fg) \cdot (k+h^\vee)/(2h^\vee)$ is the modular characteristic (not the level~$k$; see Theorem~\ref{thm:modular-characteristic}).

Since $\mu_0^{(g)}$ is central, all period-corrected higher-genus
bar differentials square to zero strictly; the raw fiberwise
differential retains the scalar curvature
$\kappa(\widehat{\fg}_k)\lambda_g$.
\end{example}

\subsubsection{\texorpdfstring{Virasoro algebra (central charge $c$)}{Virasoro algebra (central charge c)}}

\begin{example}[Virasoro: curved but strict]\label{ex:virasoro-strict}
The Virasoro algebra $\text{Vir}_c$ has stress tensor $T$ with OPE
\[T(z)T(w) = \frac{c/2}{(z-w)^4} + \frac{2T(w)}{(z-w)^2} + \frac{\partial T(w)}{z-w} + \cdots\]
and curvature $\mu_0 \propto \mathbf{1}$ from the central charge (the proportionality constant depends on the $A_\infty$ model; the modular characteristic is $\kappa(\mathrm{Vir}_c) = c/2$). Since $\mathbf{1}$ is central, $\mu_0 \in Z(\mathrm{Vir}_c)$.

The Virasoro algebra has higher-order corrections
($m_3(T \otimes T \otimes T) \neq 0$ from the cubic Schwarzian
term), but the curvature $\mu_0 \propto \mathbf{1}$ remains central,
so the period-corrected total bar differential satisfies
$d_{\text{bar}}^2 = 0$ strictly. The central charge~$c$ is a
quantum correction that breaks classical conformal invariance in a
central way: it does not break associativity of the OPE algebra.
\end{example}

\subsubsection{\texorpdfstring{$W_3$ algebra}{W-3 algebra}}

\begin{example}[\texorpdfstring{$W_3$}{W3}: filtered, still strict]\label{ex:w3-strict}
The $W_3$ algebra has generators $L$ (conformal weight~$2$) and $W$
(conformal weight~$3$, primary), with non-linear OPE: $W(z)W(w)$
involves composite operators. The curvature element satisfies
$\mu_0 \propto \mathbf{1}$ (with higher-filtration central corrections
from the non-linear OPE), so all contributions lie in the center
$Z(\mathcal{W}_3)$ and the period-corrected total bar differential
satisfies $d_{\text{bar}}^2 = 0$ strictly. The algebra
is not quadratic and requires filtered structure
(Theorem~\ref{thm:filtered-to-curved}), but centrality of the
curvature is independent of the filtration regime. The completed
bar complex $\widehat{\bar{B}}(W_3) = \varprojlim_n \bar{B}(W_3)/I^n$
(Appendix~\ref{app:nilpotent-completion}) inherits strict nilpotence.
\end{example}

\subsection{Obstruction theory: genus-by-genus analysis}

\begin{theorem}[Strict nilpotence at genus zero; \ClaimStatusProvedHere]\label{thm:genus-zero-strict}
Let $\mathcal{A}$ be a chiral algebra with central curvature. Then at genus~$0$,
the bar differential satisfies $d_0^2 = 0$ strict.
\end{theorem}

\begin{proof}
At genus 0, the bar differential is
$d_0 = d_{\text{internal}} + d_{\text{residue}}$
with no quantum corrections ($\mu_0 = 0$ at genus 0).
By the Arnold relations (Theorem~\ref{thm:arnold-three}), $d_0^2 = 0$.
\end{proof}

\begin{theorem}[Strict nilpotence for the corrected genus tower; \ClaimStatusProvedHere]\label{thm:genus-induction-strict}
\label{ex:mc-periods}
\textup{[Regime: curved-central\linebreak
\textup{(}Convention~\textup{\ref{conv:regime-tags})}.]}

Let $\mathcal{A}$ be a chiral algebra whose genus-$g$ curvature
$\mu_0^{(g)}$ is central for every $g \geq 1$ under consideration.
Assume the period-corrected components $d_g$ are defined on the
completed genus tower and satisfy the modular-operad boundary
compatibilities below. Then the corrected full bar differential
$d_{\mathrm{full}} = \sum_{g \geq 0} d_g$ satisfies
$d_{\mathrm{full}}^2 = 0$ strict.
Here each $d_g$ denotes the genus-$g$ \emph{corrected component} of
the modular decomposition, not the raw fiberwise differential
$\dfib$ of Convention~\textup{\ref{conv:higher-genus-differentials}}.
\end{theorem}

\begin{proof}
The full bar differential decomposes as
$d_{\mathrm{full}} = d_0 + \sum_{g \geq 1} d_g$,
where $d_0 = d_{\mathrm{int}} + d_{\mathrm{res}}$ is the genus-$0$ bar
differential and each $d_g$ ($g \geq 1$) arises from the boundary strata of
$\overline{\mathcal{M}}_{g,n}$. We compute $d_{\mathrm{full}}^2$ by organizing
the terms by total genus.

\medskip
\emph{Step 1: The modular operad boundary relation.}
The collection $\{\overline{\mathcal{M}}_{g,n}\}_{g,n}$ forms a modular operad
whose composition maps are given by gluing along boundary divisors. The
boundary of $\overline{\mathcal{M}}_{g,n}$ is a normal crossings divisor with
strata indexed by dual graphs $\Gamma$. Codimension-$2$ strata arise from
refinements $\Gamma' \to \Gamma$ of dual graphs. The fundamental topological
identity $\partial^2 = 0$ on the stratification of $\overline{\mathcal{M}}_{g,n}$
states that each codimension-$2$ stratum appears with opposite signs from the two
codimension-$1$ strata that contain it. This is equivalent to the \emph{modular
operad axiom}: the composition in the modular operad is associative.

\medskip
\emph{Step 2: Genus-$0$ nilpotence ($d_0^2 = 0$).}
This is Theorem~\ref{thm:genus-zero-strict}: the Arnold relations
(Theorem~\ref{thm:arnold-three}) provide the codimension-$2$ cancellation on
$\overline{C}_n(X) = \overline{\mathcal{M}}_{0,n+1}$.

\medskip
\emph{Step 3: Cross-terms ($d_0 d_g + d_g d_0 = 0$).}
Each $d_g$ inserts a genus-$g$ correction $\mu_0^{(g)} \otimes \omega_g$,
where $\omega_g \in H^*(\overline{\mathcal{M}}_{g,n})$ and
$\mu_0^{(g)} \in Z(\mathcal{A})$ by hypothesis. The cross-term
$d_0 d_g + d_g d_0$ corresponds to codimension-$2$ boundary strata of
$\overline{\mathcal{M}}_{g,n+k}$ that involve one genus-$0$ degeneration
(a collision of points on the same component) and one genus-$g$ operation.
By the modular operad axiom (Step~1), these strata cancel pairwise:
\[
 d_0 d_g + d_g d_0 = \sum_{\substack{\text{codim-2 strata} \\ \text{(genus-0/genus-$g$)}}}
 (\pm 1)\, [\text{stratum}] = 0.
\]
Centrality of $\mu_0^{(g)}$ is essential here: it ensures that the algebraic
insertion $\mu_0^{(g)} \otimes (-)$ commutes with the internal differential
$d_{\mathrm{int}}$, so that the signs from the boundary orientation match the
signs from the Koszul sign rule.

\medskip
\emph{Step 4: Pure higher-genus terms ($\sum_{g_1 + g_2 = g} d_{g_1} d_{g_2} = 0$).}
The product $d_{g_1} d_{g_2}$ inserts two corrections
$\mu_0^{(g_1)}$ and $\mu_0^{(g_2)}$ into the bar complex. Since both lie in
$Z(\mathcal{A})$, the order of insertion is irrelevant: the insertions commute.
The sum over all decompositions $g_1 + g_2 = g$ corresponds to the codimension-$2$
boundary strata of $\overline{\mathcal{M}}_{g,n}$ that arise from two
\emph{separating} degenerations (splitting the surface into components of genera
$g_1$ and $g_2$). Non-separating degenerations at genus $g$ (which increase the
genus of a single component by~$1$) also contribute; these correspond to
self-gluing operations in the modular operad.

For each pair of codimension-$1$ divisors $D_1, D_2$ meeting in a codimension-$2$
stratum $D_1 \cap D_2$, the boundary-of-boundary identity assigns opposite signs
to the two orderings $D_1 \to D_1 \cap D_2$ and $D_2 \to D_1 \cap D_2$.
Centrality ensures the algebraic insertions respect these orientations (no
additional sign from non-commutativity), and so:
\[
 \sum_{g_1 + g_2 = g} d_{g_1} d_{g_2}
 = \sum_{\substack{\text{codim-2 strata of}\\
 \overline{\mathcal{M}}_{g,n}}} (\pm 1)\, [\text{stratum}] = 0.
\]

\medskip
\emph{Step 5: Conclusion.}
Collecting all terms:
\[
 d_{\mathrm{full}}^2 = d_0^2
 + \sum_{g \geq 1} (d_0 d_g + d_g d_0)
 + \sum_{g \geq 2}\sum_{g_1 + g_2 = g} d_{g_1} d_{g_2}
 = 0 + 0 + 0 = 0.
\]
The first term vanishes by Step~2, the second by Step~3, and the third by Step~4.
This completes the proof.
\end{proof}

\begin{remark}[Computational check]\label{rem:genus-induction-computational-check}
The proof above uses the modular operad structure of
$\{\overline{\mathcal{M}}_{g,n}\}$ in an essential way. The computation
also checks the conclusion through genus~$5$ for all standard examples
(Heisenberg, Kac--Moody, Virasoro, $W$-algebras), providing independent
confirmation.
\end{remark}

The proof of Theorem~\ref{thm:genus-induction-strict} relies at every step on
the modular operad axioms for $\{\overline{\mathcal{M}}_{g,n}\}$. The next
statement crystallizes this dependence into an explicit algebraic statement: the bar
complex \emph{is} an algebra over the Feynman transform of the commutative
modular operad.

\begin{theorem}[Bar complex as algebra over the modular operad;
 \ClaimStatusProvedHere]\label{thm:bar-modular-operad}%
\index{bar complex!algebra over modular operad}%
\index{modular operad!bar complex as algebra over}%
\index{Feynman transform!commutative modular operad}%
For a chiral algebra $\cA$, the collection
$\{\barB^{(g,n)}(\cA)\}_{2g-2+n>0}$ is an algebra over the Feynman transform
$\mathsf{F}\mathrm{Com}$ of the commutative modular operad in the sense of
Getzler--Kapranov~\textup{\cite{GetzlerKapranov98}}. Concretely:
\begin{enumerate}[label=\textup{(\roman*)}]
\item For each stable graph $\Gamma$ of type $(g,n)$, there is a composition
 map
 \[
 \circ_\Gamma \colon
 \bigotimes_{v \in V(\Gamma)} \barB^{(g_v,n_v)}(\cA)
 \longrightarrow \barB^{(g,n)}(\cA)
 \]
 given by contracting internal edges via the propagator $P_\cA$.
\item These composition maps satisfy the associativity and equivariance axioms
 of modular operad algebras: for any refinement $\Gamma' \to \Gamma$ of stable
 graphs, $\circ_\Gamma \circ (\otimes_v\, \circ_{\Gamma'_v})
 = \circ_{\Gamma' \circ \Gamma}$, and for $\sigma \in \mathrm{Aut}(\Gamma)$
 the maps are $\sigma$-equivariant.
\item The identity $\partial^2 = 0$ on $\barB^{(g,n)}(\cA)$ is a formal
 consequence of the boundary-of-boundary relation $\partial^2 = 0$ on the
 stratification of $\overline{\mathcal{M}}_{g,n}$.
\end{enumerate}
\end{theorem}

\begin{proof}
Part~(i) is the stable-graph bar differential construction of
Lemma~\ref{lem:stable-graph-d-squared}: each edge $e \in E(\Gamma)$ contributes
a contraction against the propagator $P_\cA$, and the composition
$\circ_\Gamma$ is the iterated contraction over all internal edges.

Part~(ii) follows from the associativity of edge contraction on stable graphs.
If $\Gamma' \to \Gamma$ is a refinement (collapsing a subset of edges of
$\Gamma'$ to produce $\Gamma$), then contracting in two stages (first the
edges internal to each $\Gamma'_v$, then the remaining edges of $\Gamma$) gives
the same result as contracting all edges of $\Gamma'$ at once. Equivariance
under automorphisms is immediate from the symmetric monoidal structure of the
target category.

Part~(iii) is the content of Theorem~\ref{thm:prism-higher-genus}: the bar
differential $d_{\mathrm{full}} = \sum_\Gamma \circ_\Gamma$ satisfies
$d_{\mathrm{full}}^2 = 0$ precisely because each codimension-$2$ boundary
stratum of $\overline{\mathcal{M}}_{g,n}$ appears with opposite signs from the
two codimension-$1$ strata containing it.
\end{proof}

\begin{remark}\label{rem:bar-modular-operad-consequences}%
\index{genus tower!as modular operad corollary}%
Theorem~\ref{thm:bar-modular-operad} identifies the modular operad
algebra structure with the Getzler--Kapranov compatibility of
edge contractions. All genus-tower structure
(the decomposition $d_{\mathrm{full}} = \sum_g d_g$, the
cross-genus cancellations in the proof of
Theorem~\ref{thm:genus-induction-strict}, and the recursive computation of bar
cohomology by genus) is a formal corollary of the bar complex being an
$\mathsf{F}\mathrm{Com}$-algebra.
The $\mathsf{F}\mathrm{Com}$-algebra structure is the strict
incarnation of the homotopy modular coalgebra on $\barB(\cA)$;
the Feynman transform provides the homotopy-invariant formulation
(Getzler--Kapranov~\cite[\S5]{GetzlerKapranov98}).
\end{remark}

\begin{theorem}[Bar-cobar as a factorization homotopy equivalence;
\ClaimStatusConditional]
\label{thm:quillen-equivalence-chiral}
\index{Quillen equivalence!bar-cobar|textbf}
\index{bar-cobar!Quillen equivalence}
\index{model structure!conilpotent coalgebras}
For a Koszul chiral algebra~$\cA$
\textup{(}Definition~\textup{\ref{def:chiral-koszul-pair})},
the bar-cobar adjunction
\[
\Omega_\kappa\colon
\mathrm{conil}\text{-}\mathrm{dg}\text{-}\cC^{\textup{!\textasciigrave}}_{\mathrm{ch}}\text{-}\mathrm{coalg}
\;\rightleftarrows\;
\mathrm{dg}\text{-}\mathrm{Ch\text{-}alg}
\;:\!B_\kappa,
\qquad
\Omega_\kappa\dashv B_\kappa
\]
is a homotopy equivalence on the finite-type
conilpotent-complete factorization model. The pointwise operadic input
is the Vallette/Loday--Vallette bar-cobar theorem
\textup{\cite[Theorem~2.1]{Val16}; \cite[Theorem~11.4.1]{LV12}}; the
chiral $\Dmod(X)$-enriched statement is obtained only after placing the
chiral operad in the Francis--Gaitsgory--Rozenblyum factorization
model structure
\textup{\cite[Chapter~IV.5, Theorem~3.1.2]{GR17}}. In that ambient:
\begin{enumerate}[label=\textup{(\roman*)}]
\item Weak equivalences are maps whose cobar construction is a
 quasi-isomorphism on the finite-type square-zero lane, respectively
 maps detected by coderived comodules on the curved lane.
\item Cofibrant coalgebras are generated by conilpotent cofree
 factorization coalgebras.
\item Fibrant algebra objects are represented by quasi-free
 factorization algebras in the chosen model.
\item $B_\kappa$ preserves fibrations and weak equivalences
 after transport to the factorization model.
\item $\Omega_\kappa$ preserves cofibrations and weak equivalences on
 the same finite-type conilpotent-complete surface.
\end{enumerate}
In particular, on this finite-type conilpotent-complete surface, the
homotopy category of dg chiral algebras is equivalent to the homotopy
category of conilpotent chiral coalgebras:
$\mathrm{Ho}(\mathrm{dg}\text{-}\mathrm{Ch\text{-}alg})
\simeq
\mathrm{Ho}(\mathrm{conil}\text{-}\mathrm{dg}\text{-}
\cC^{\textup{!\textasciigrave}}_{\mathrm{ch}}\text{-}\mathrm{coalg})$.
\end{theorem}

\begin{proof}[References]
Vallette's theorem supplies the conilpotent bar-cobar model structure
for cooperads in chain complexes over a field. The chiral tensor
product is pseudo-tensor before passing to factorization algebras, so
that theorem is not applied verbatim to $\Dmod(X)$-enriched cooperads.
The required chiral ambient is the factorization model of
Francis--Gaitsgory--Rozenblyum
\cite[Chapter~IV.5, Theorem~3.1.2]{GR17}, as recorded in
Theorem~$A^{\infty,2}$
(Theorem~\ref{thm:A-infinity-2}) and
Remark~\ref{rem:ainf2-heal-72}. Transporting the operadic
bar-cobar theorem through that ambient gives the displayed homotopy
equivalence. The algebra-side counit
$\psi_A\colon\Omega_\kappa B_\kappa(A)\to A$ is the bar-cobar
resolution map of Theorem~\ref{thm:bar-cobar-inversion-qi}; it is a
quasi-isomorphism on the Koszul locus. The coalgebra-side unit is
$\eta_C\colon C\to B_\kappa\Omega_\kappa(C)$.
\end{proof}

\begin{remark}[Relationship to Theorem~A]
\label{rem:quillen-verdier-connection}
\index{Quillen equivalence!Theorem A connection}
Theorem~\ref{thm:quillen-equivalence-chiral} is the
homotopy-categorical refinement of the geometric bar-cobar adjunction
and Verdier intertwining of
Theorem~A
(Theorem~\ref{thm:bar-cobar-isomorphism-main}): the Verdier
intertwining $\mathbb{D}_{\operatorname{Ran}} \barB(\cA_1) \simeq
(\cA_2)_\infty$ identifies the Verdier dual bar coalgebra with the
homotopy Koszul-dual algebra and becomes a Quillen-level statement about derived
functors, and the configuration-space integrals that define the bar
differential become the geometric realization of the model structure's
fibration data.
\end{remark}

\begin{remark}[Physical content of the Quillen equivalence]
\label{rem:quillen-physical}
\index{Quillen equivalence!physical interpretation}
The Quillen equivalence $\Omega_\kappa \dashv B_\kappa$ has a direct
physical reading. The bar functor~$B_\kappa$ passes from the on-shell
theory (the chiral algebra~$\cA$) to the off-shell theory (the bar
coalgebra~$\barB(\cA)$, which carries the full perturbative data
including gauge redundancies). The cobar functor~$\Omega_\kappa$
recovers the on-shell theory from the off-shell data. The Quillen
equivalence states that this passage is invertible: no information is
lost, precisely on the Koszul locus, where the bar spectral sequence
collapses at~$E_2$.

In renormalization language: Koszulness is the condition under which
perturbative quantization is invertible. A Koszul chiral algebra
admits a unique reconstruction from its perturbative data (bar
cohomology). Outside the Koszul locus, irreducible higher-loop
ambiguities may obstruct ordinary reconstruction. The four
shadow-depth classes G/L/C/M classify the complexity of this
reconstruction within the proved Koszul standard families; the
class~W boundary lies outside this finite-type Quillen statement.
\end{remark}

\begin{remark}[Hecke equivariance of Verdier duality]%
\label{rem:verdier-hecke-equivariance}%
\index{Verdier duality!Hecke equivariance}%
At genus~$1$, the Verdier duality functor~$\bD$ commutes
with Hecke correspondences~$T_p$ on~$\cM_{1,1}$
(Theorem~\ref{thm:hecke-verdier-commutation}):
the dual isogeny $\phi^\vee$ bijects degree-$p$ isogenies
from~$E$ with those from $E^\vee\simeq E$.
When $\cA\simeq\cA^!$, the Verdier involution commutes
with~$T_p$; combined with multiplicity one for newforms,
$Z_\cA$ decomposes into Hecke eigenforms, and the
constrained Epstein zeta factors into standard
$L$-functions
(Theorem~\ref{thm:self-dual-factorization}).
\end{remark}

\begin{corollary}[Rectification of $\mathrm{Ch}_\infty$-algebras;
\ClaimStatusProvedElsewhere]
\label{cor:rectification-ch-infty}
\index{rectification!Ch-infinity@$\mathrm{Ch}_\infty$-algebras}
Over chain complexes, every $\mathrm{Ch}_\infty$-algebra is
$\infty$-quasi-isomorphic to a strict algebra over the resolved
operad via the rectification functor
$\Omega_\kappa \tilde{B}_\iota$
\textup{\cite[Theorem~1.2]{Val16}}. For chiral algebras on a curve,
the same rectification is used only after transport to the
factorization ambient of Theorem~$A^{\infty,2}$. In particular, on the
finite-type conilpotent-complete factorization surface, the homotopy
category of strict chiral algebras and of $\mathrm{Ch}_\infty$-algebras
with $\infty$-morphisms are equivalent:
\[
\mathrm{Ho}(\mathrm{dg}\text{-}\mathrm{Ch\text{-}alg})
\;\simeq\;
\infty\text{-}\mathrm{Ch}_\infty\text{-}\mathrm{alg}\, /\,{\sim_h}.
\]
In the chain-operadic model,
\textup{\cite[Theorem~3.8]{Val16}} upgrades this to an equivalence of
$\infty$-categories. The chiral version uses the
Francis--Gaitsgory--Rozenblyum model structure before taking this
underlying $\infty$-category. This is the target for the DK-$5$
categorical lift
\textup{(}\S\textup{\ref{sec:concordance-three-pillars})}.
\end{corollary}

\begin{proof}[References]
\cite[Theorems~1.2, 3.7, and 3.8]{Val16} for the chain-operadic
rectification; \cite[Chapter~IV.5]{GR17} and
Theorem~\ref{thm:A-infinity-2} for the chiral factorization ambient.
\end{proof}

\subsection{Summary}

\begin{center}
\begin{tabular}{|>{\raggedright\arraybackslash}p{3.5cm}|>{\raggedright\arraybackslash}p{4cm}|>{\raggedright\arraybackslash}p{6cm}|}
\hline
\textbf{Regime} & \textbf{Condition} & \emph{Examples} \\
\hline
\hline
\textbf{Strict Nilpotence} & $\mu_0 \in Z(\mathcal{A})$ plus square-zero total correction &
\begin{itemize}[leftmargin=*]
\item Heisenberg $\mathcal{H}_k$
\item Kac--Moody $\widehat{\mathfrak{g}}_k$
\item Virasoro $\text{Vir}_c$
\item $W$-algebras $W_N$
\item Free fermions $\beta\gamma$
\end{itemize}
$d_{\text{bar}}^2 = 0$ strictly for the corrected total differential (Theorem~\ref{thm:genus-induction-strict}) \\
\hline
\textbf{Curved (Non-Central)} & $\mu_0 \notin Z(\mathcal{A})$ &
\begin{itemize}[leftmargin=*]
\item Hypothetical non-central extensions
\item Some deformed algebras
\end{itemize}
raw CDG square nonzero, need higher homotopies/coderived localization \\
\hline
\textbf{Homotopy Coherent} & No curvature, but $d^2 \sim 0$ only &
\begin{itemize}[leftmargin=*]
\item $(\infty,1)$-categorical structures
\item Derived geometry settings
\item Non-algebraic field theories
\end{itemize}
Requires full $A_\infty$ or $L_\infty$ framework \\
\hline
\end{tabular}
\end{center}

\begin{remark}[Consequences of strict nilpotence]\label{rem:why-strict-matters}
Strict nilpotence of the corrected total differential allows direct cohomology computation
(Theorem~\ref{thm:genus-induction-strict}) and ensures the bar-cobar
adjunction works without completion on the quadratic square-zero locus
(Theorem~\ref{thm:bar-cobar-adjunction}); homotopy nilpotence
requires spectral sequences that may not degenerate and completion
with attendant convergence issues
(Theorem~\ref{thm:completion-necessity}). Physically, strict
nilpotence means quantum corrections are controlled by central charges
(Theorem~\ref{thm:deformation-obstruction}); homotopy nilpotence
requires full renormalization group analysis. The standard
positive-energy vertex-operator families with the vacuum axiom, and
the unitary CFT chiral algebras whose anomaly is a scalar vacuum
class, have central curvature
(Remark~\ref{rem:voa-central-curvature};
Theorem~\ref{thm:genus-induction-strict}); after the
period-corrected total differential is chosen, they lie in the strict
regime.
\end{remark}

\subsection{Connection to literature}

\begin{theorem}[GLZ, Theorem 5.3; \ClaimStatusProvedElsewhere]\label{thm:glz-curved}
For a quadratic chiral algebra $\mathcal{A}$ with central curvature
$\mu_0 \in Z(\mathcal{A})$, complete for the chosen filtration, the
Koszul dual $\mathcal{A}^!$ exists as a curved cooperad. If the
curvature is killed by an MC twist, or if the period-corrected total
differential is used, the algebra-side counit
$\Omega(B(\mathcal{A})) \to \mathcal{A}$ is an isomorphism in the
ordinary derived category. Without that square-zero hypothesis the
same comparison is a Positselski coderived statement: its cone is
coacyclic in the second-kind category of curved comodules, not a raw
chain-level inverse~\textup{\cite{GLZ22,Positselski11}}.
\end{theorem}

Theorem~\ref{thm:central-implies-strict} provides the geometric
realization of the GLZ algebraic result.

\begin{theorem}[FG, Theorem 7.2.1; \ClaimStatusProvedElsewhere]
\label{thm:fg-factorization-bar-cobar}
For a factorization algebra $\mathcal{F}$ on a curve $X$,
$\textup{Fact}(X, \Omega(B(\mathcal{F}))) \simeq
\mathcal{F}$~\textup{\cite{FG12}}.
\end{theorem}

Combined with the explicit bar construction
(Definition~\ref{def:geometric-bar}), this identifies the geometric
bar differential with the standard one.

The Costello--Gwilliam MC equation $\delta I + \frac{1}{2}\{I,I\} = 0$
for the quantum effective action~$I$~\cite[Definition~3.2.1.1]{CG-vol2}
is equation~\eqref{eq:mc-equation} in the field theory context:
central charges in QFT correspond to central curvature in chiral
algebras, and both ensure that quantum corrections do not destroy
associativity.

\subsection{Computational corollaries}

\begin{corollary}[Bar cohomology computes Ext;
 \ClaimStatusProvedElsewhere]\label{cor:bar-computes-ext}
For a chiral algebra $\mathcal{A}$ with central curvature and a
square-zero total bar differential,
the bar construction computes self-Ext of the vacuum module:
\begin{equation}
H^*(\bar{B}(\mathcal{A}), d_{\text{bar}}) =
\operatorname{Ext}_{\mathcal{A}}^*(\omega_X, \omega_X).
\end{equation}
This is the chiral analogue of the classical bar-Ext
identification for algebras over operads
\textup{(}\cite{LV12}, Theorem~2.2.4; the chiral
adaptation is Theorem~\textup{\ref{thm:bar-cobar-isomorphism-main}}\textup{)}.
For a genuinely curved CDG bar object, the same formula is read after
passing to the coderived category of curved comodules.
\end{corollary}

\begin{corollary}[Koszul dual coalgebra {\cite{GK94}}; \ClaimStatusProvedElsewhere]\label{cor:koszul-dual-cooperad}
For quadratic $\mathcal{A}$ with central curvature and square-zero
total bar model, the bar
coalgebra $\barB(\cA)$ has cohomology
\begin{equation}
H^*(\bar{B}(\mathcal{A}))
\end{equation}
concentrated in bar degree~$1$ (the Koszul dual coalgebra), carrying a curved cooperad structure with comultiplication dual to~$m_2$ and curvature dual to~$\mu_0$. The Koszul dual \emph{algebra}~$\cA^!$ is characterized by Verdier intertwining: $\mathbb{D}_{\operatorname{Ran}}\barB_X(\cA) \simeq \cA^!_\infty$ (factorization \emph{algebra}, not coalgebra; Convention~\ref{conv:bar-coalgebra-identity}).
\end{corollary}

\begin{corollary}[Genus expansion convergence; \ClaimStatusProvedHere]\label{cor:genus-expansion-converges}
The genus expansion:
\begin{equation}
Z(\mathcal{A}) = \sum_{g=0}^\infty \hbar^{2g-2} Z_g(\mathcal{A})
\end{equation}
where $Z_g(\mathcal{A}) = H^*(\bar{B}^{(g)}(\mathcal{A}), \Dg{g})$, converges in the
$\hbar$-adic topology of the formal power series ring $\mathbb{C}[[\hbar]]$.
\end{corollary}

\begin{proof}
Each term $Z_g(\mathcal{A})$ is well-defined because
$\Dg{g}^{\,2} = 0$ strictly at genus $g$
(Theorem~\ref{thm:genus-induction-strict}), so the cohomology
$H^*(\bar{B}^{(g)}(\mathcal{A}))$ exists on the chosen square-zero
finite-type or completed surface.
The formal power series $\sum_{g \geq 0} \hbar^{2g-2} Z_g$ is an element of
$\hbar^{-2} \mathbb{C}[[\hbar^2]]$, which is a well-defined ring (Laurent series
with a pole of bounded order). Convergence in the $\hbar$-adic topology is
automatic: for every $N$, the partial sum $\sum_{g=0}^N \hbar^{2g-2} Z_g$
agrees with $Z(\mathcal{A})$ modulo $\hbar^{2N}$.

Here $\Dg{g}^{\,2} = 0$ holds \emph{strictly} at each genus for
the period-corrected total differential, so each $Z_g$ is an ordinary
cohomology object in that ambient, not the cohomology of the raw curved
fiberwise CDG operator. This is the strict nilpotence established in
Section~\ref{sec:strict-vs-homotopy}.
\end{proof}

\subsection{Perspectives on strict nilpotence}

\subsubsection{Physical perspective}

Associativity of the OPE algebra reflects locality: quantum
corrections enter as central charges that modify the normalization
but preserve associativity. The central charge~$c$ in the
Virasoro OPE $T(z)T(w) \sim (c/2)/(z-w)^4 + \cdots$ is zero
classically; its centrality ensures the Jacobi identity survives
quantization.

\subsubsection{Geometric perspective}

Kontsevich's formality theorem~\cite{Kon03} realizes the
$L_\infty$-quasi-isomorphism
$\text{Poly}_{\bullet}(M)[[t]] \xrightarrow{\sim}
\text{Dpoly}_{\bullet}(M)[[t]]$
via configuration space integrals $U_\Gamma = \int_{C_n(\mathbb{R}^d)} \omega_\Gamma$
satisfying $\sum_\Gamma U_\Gamma \cdot (\text{boundary terms}) = 0$
by Stokes' theorem. This is the genus-$0$ case of strict
nilpotence; at higher genus the same Stokes argument applies
because curvature terms are central and do not interfere with the
boundary calculations.

\subsubsection{Computational and functorial perspectives}

Strict nilpotence of the corrected total model holds degree by degree:
$d^2(a \otimes b) = 0$
by direct calculation, $d^2(a \otimes b \otimes c) = 0$ using
Arnold relations, and similarly through degree~$5$ using extended
Arnold relations; centrality of $\mu_0$ is used at each stage.
It also follows from functoriality:

\begin{theorem}[Functoriality of bar construction; \ClaimStatusConditional]\label{thm:bar-functorial-grothendieck}
The bar construction:
\[B: \mathrm{ChAlg}^{\mathrm{central,tot=0}} \to \mathrm{Coalg}\]
is a functor from chiral algebras with central curvature and
square-zero total differential to coalgebras, characterized by the
universal property:
\[\mathrm{Hom}_{\mathrm{ChAlg}}(\Omega(C),\mathcal{A}) \simeq
\mathrm{Hom}_{\mathrm{Coalg}}(C,B(\mathcal{A}))\]

For a central curved CDG model before passing to a square-zero total
differential, the same construction lands in curved CDG-coalgebras
and is interpreted coderivedly.
\end{theorem}

\begin{proof}
The adjunction follows from functoriality, the bar-cobar Hom identity,
and the square-zero total model.

\emph{Functoriality.}
By Theorem~\ref{thm:bar-functorial}, the bar construction $B$ is functorial: a morphism $f: \mathcal{A} \to \mathcal{B}$ of chiral algebras induces $B(f): B(\mathcal{A}) \to B(\mathcal{B})$ compatible with the coalgebra structures. The category $\mathrm{ChAlg}^{\mathrm{central,tot=0}}$ consists of chiral algebras whose curvature $\mu_0$ lies in the center $Z(\mathcal{A})$ and whose chosen total bar differential squares to zero.

\emph{Adjunction.}
The bar-cobar adjunction $\Omega \dashv B$ is established in
Theorem~\ref{thm:bar-cobar-adjunction}. For a coalgebra~$C$ and a
chiral algebra~$\mathcal{A}$, a coalgebra map
$\phi\colon C \to B(\mathcal{A})$ is equivalently a twisting morphism
from~$C$ to~$\mathcal{A}$, hence a chiral algebra map
$\tilde{\phi}\colon \Omega(C) \to \mathcal{A}$. Conversely,
$\tilde{\phi}$ restricts to the generators of~$\Omega(C)$ and extends
uniquely to the coalgebra map $C \to B(\mathcal{A})$ by the cofree
property of the bar coalgebra.

\emph{$d^2 = 0$ from the chosen square-zero total model.}
The bar differential $d_{\bar{B}}$ on $B(\mathcal{A})$ satisfies
$d_{\bar{B}}^2 = 0$ in this category because the square-zero total
differential is part of the hypothesis. Centrality of
$\mu_0 \in Z(\mathcal{A})$ kills the internal commutator
$[\mu_0,-]$, while Theorem~\ref{thm:central-implies-strict} supplies
the closedness and boundary-cancellation hypotheses for the total
correction. Before those hypotheses are imposed, $B(\mathcal{A})$ is
a curved CDG-coalgebra and the bar-cobar comparison is
coderived/comodule-detected.
\end{proof}

\begin{remark}[Dependencies]
This proof uses three ingredients:
\begin{enumerate}[label=(D\arabic*)]
\item Concrete bar-functor functoriality (Theorem~\ref{thm:bar-functorial}).
\item The bar-cobar adjunction framework (Theorem~\ref{thm:bar-cobar-adjunction}).
\item Central curvature plus square-zero total correction implies strict nilpotence (Theorem~\ref{thm:central-implies-strict}).
\end{enumerate}
\end{remark}

The adjunction exists independently of centrality. Centrality, together
with the square-zero total correction, is the hypothesis that permits
ordinary cohomology instead of the curved coderived category.

\subsection{Strict versus homotopy nilpotence}
\label{subsec:strict-vs-homotopy-nilpotence}

\begin{remark}[Strict vs homotopy nilpotence]
For chiral algebras $\mathcal{A}$ arising from vertex operator algebras
(satisfying the vacuum axiom $Y(\mathbf{1},z) = \mathrm{id}$) or from
conformal field theories whose anomalies are central, the corrected
total bar differential satisfies $d_{\text{bar}}^2 = 0$ strictly at
each genus in the corrected tower. The ingredients are: (1) the curvature $\mu_0$ is
proportional to the vacuum $\mathbf{1}$, hence central by the vacuum
axiom (Remark~\ref{rem:voa-central-curvature}); (2) on the affine
genus-$0$ screen, Arnold relations ensure $d_{\mathrm{residue}}^2 = 0$,
while on $\mathbb{P}^1$ and for $g\geq1$ this residue-square
computation also uses the projective and period relations of
Proposition~\ref{prop:arnold-higher-genus}; (3) the Leibniz rule gives
the anticommutation of $d_{\mathrm{internal}}$ and
$d_{\mathrm{residue}}$; (4) the modular operad boundary relations on
$\overline{\mathcal{M}}_{g,n}$ extend this through the corrected genus tower
(Theorem~\ref{thm:genus-induction-strict}). This covers the corrected
total models of the standard examples. For non-unital chiral algebras
or theories with non-central anomalies, the raw CDG square may vanish
only after higher homotopies and coderived localization.
\end{remark}

%% ================================================================
%% RECONSTRUCTION, DUALITY, AND THE CASIMIR TRANSGRESSION
%% ================================================================

\section{Reconstruction, duality, and the Casimir transgression}
\label{sec:reconstruction-duality-casimir}
\index{reconstruction!versus duality}
\index{Casimir transgression|textbf}
\index{Sugawara coupling!bar complex}

There are two natural operations on the bar complex of a chiral
algebra, and they are not the same. The first applies the cobar
functor directly: $\Omegach(\barBch(\cA)) \to \cA$. This is
\emph{reconstruction}: it recovers~$\cA$ from its collision data,
and Theorem~B says it is a quasi-isomorphism on the Koszul locus.
The second interposes Verdier duality on the Ran space:
$\mathbb{D}_{\Ran}\barBch(\cA)$. This is
\emph{duality}: it produces the Koszul dual partner~$\cA^!$, not
$\cA$ itself. The distinction between reconstruction and duality is
the bar-cobar analogue of the distinction between the identity and
the Fourier transform: both act on the same data, but they extract
complementary information.

A new phenomenon appears when two chiral algebras are coupled.
When two chiral algebras are tensored with a nontrivial mixed OPE
(the Sugawara coupling, the diagonal GKO coset), the bar differential
acquires a mixed-color component whose first obstruction is
\emph{universal}: it depends only on primarity, not on the
Lie type or level. After this universal piece is contracted away,
what survives is the \emph{Casimir transgression}: a single
quadratic relation that cuts out the coupled dual as a strict
complete intersection with no higher $A_\infty$~corrections.

\subsection{Reconstruction versus Verdier--cobar duality}
\label{subsec:reconstruction-vs-duality}

The five-object distinction
\[
\cA,\qquad
\barBch(\cA),\qquad
\cA^i = H^*(\barBch(\cA)),\qquad
\cA^! = (\cA^i)^\vee,\qquad
Z^{\mathrm{der}}_{\mathrm{ch}}(\cA)
\]
is the firewall for the whole bar-cobar theory.
Theorems~A and~B are about the reconstruction functor $R_X$: they
recover~$\cA$ from its collision data. Koszul duality (the
interplay of $\cA$ with its dual partner~$\cA^!$) is about the
\emph{distinct} functor~$K_X$. The derived chiral centre is a third
operation, Hochschild rather than Verdier or cobar. Every formula in
the shadow obstruction tower, every modular characteristic, and every
complementarity theorem depends on keeping these operations apart. The
following definition and theorem make the separation precise.

\begin{definition}[The two functors]
\label{def:reconstruction-vs-duality-functors}
\index{reconstruction functor $R_X$}
\index{Verdier duality functor $K_X$}
For an augmented chiral algebra~$\cA$, define
\[
R_X(\cA) := \Omegach_X \barBch_X(\cA),
\qquad
K_X(\cA) := \mathbb{D}_{\Ran}\barBch_X(\cA).
\]
We call $R_X$ the \emph{reconstruction functor} and $K_X$ the
\emph{Verdier duality functor}.
\end{definition}

\begin{theorem}[Reconstruction versus duality; \ClaimStatusConditional]
\label{thm:reconstruction-vs-duality}
Assume~$\cA$ lies on the finite chiral Koszul locus and
$(\cA, \cA^\sharp)$ is a finite chiral Koszul pair. Then
\[
R_X(\cA) \simeq \cA,
\qquad
K_X(\cA) \simeq \cA^\sharp.
\]
\end{theorem}

\begin{proof}
The first statement is the bar-cobar counit on the Koszul locus
(Theorem~\ref{thm:bar-cobar-inversion-qi}).
For the second, finite chiral Koszulity supplies the Verdier
identification
$\mathbb{D}_{\Ran}\barBch_X(\cA) \simeq (\cA^\sharp)_\infty$ (factorization \emph{algebra}, not coalgebra).
The strict comparison $(\cA^\sharp)_\infty \simeq \cA^\sharp$ on the
finite Koszul locus gives
\[
K_X(\cA)
= \mathbb{D}_{\Ran}\barBch_X(\cA)
\simeq (\cA^\sharp)_\infty
\simeq \cA^\sharp.
\qedhere
\]
\end{proof}

\begin{remark}[Five objects under Verdier duality and Hochschild cochains]
\label{rem:four-objects}
\index{Verdier duality!algebra versus coalgebra}
The bar complex $\barBch(\cA)$ is a factorization \emph{coalgebra}.
Its Verdier dual $\mathbb{D}_{\Ran}\barBch(\cA)$ is naturally a
factorization \emph{algebra}: the Verdier dual of a coalgebra
carries a multiplication, not a comultiplication. The intertwining
$\mathbb{D}_{\Ran}\barBch(\cA) \simeq \cA^!_\infty$ identifies
this \emph{algebra} with the homotopy Koszul dual algebra~$\cA^!_\infty$, whose underlying complex is~$\barBch(\cA^!)$,
via the canonical D-module pairing on holonomic objects
(Theorem~\ref{thm:verdier-bar-cobar}):
\begin{center}
\renewcommand{\arraystretch}{1.2}
\begin{tabular}{lll}
\textbf{Object} & \textbf{Type} & \textbf{Operation} \\
\hline
$\barBch(\cA)$ & coalgebra & coproduct from collision splitting \\
$\mathbb{D}_{\Ran}\barBch(\cA)$ & algebra
 & multiplication = Verdier dual of coproduct \\
$\barBch(\cA^!)$ & coalgebra & coproduct from $\cA^!$-collisions \\
$\cA^{\scriptstyle \text{\normalfont !\textasciigrave}} = H^*(\barBch(\cA))$ & coalgebra & bar cohomology \\
$\cA^! = (\cA^{\scriptstyle \text{\normalfont !\textasciigrave}})^\vee$ & algebra
 & Koszul dual algebra
 \\
$Z^{\mathrm{der}}_{\mathrm{ch}}(\cA)$ & algebra object
 & chiral Hochschild cochains $C^\bullet_{\mathrm{ch}}(\cA,\cA)$
\end{tabular}
\end{center}
The reconstruction functor $R_X = \Omegach_X \circ \barBch_X$
recovers~$\cA$ (an algebra). The duality functor $K_X =
\mathbb{D}_{\Ran} \circ \barBch_X$ produces
$\cA^!$ (an algebra): Verdier duality converts the coalgebra to an
algebra, and the finite Koszul comparison identifies that algebra with
the strict dual partner. The derived centre
$Z^{\mathrm{der}}_{\mathrm{ch}}(\cA)$ is
$C^\bullet_{\mathrm{ch}}(\cA,\cA)$, the Hochschild closed-sector algebra acting on
$\cA$; it is not a bar coalgebra, not bar cohomology, and not the
Verdier-Koszul dual. The
distinction: $\Omegach_X(\barBch_X(\cA)) \simeq \cA$, while
$\mathbb{D}_{\Ran}\barBch_X(\cA) \simeq \cA^!$.
\end{remark}

\begin{theorem}[Recognition theorem for finite chiral Koszul pairs;
\ClaimStatusConditional]
\label{thm:recognition-koszul-pairs}
\index{Koszul pair!recognition theorem}
Let $\cA$ and~$\cB$ be augmented chiral algebras satisfying
finite-type and convergence hypotheses. The following are
equivalent:
\begin{enumerate}[label=\textup{(\roman*)}]
\item $(\cA,\cB)$ is a finite chiral Koszul pair.
\item There are quasi-isomorphisms
 $K_X(\cA) \simeq \cB$ and $K_X(\cB) \simeq \cA$.
\end{enumerate}
In particular, $K_X$ is a contravariant involution on the finite
chiral Koszul locus: $K_X(K_X(\cA)) \simeq \cA$.
\end{theorem}

\begin{proof}
The implication (i)$\Rightarrow$(ii) is
Theorem~\ref{thm:reconstruction-vs-duality}. Conversely,
suppose $K_X(\cA) \simeq \cB$ and $K_X(\cB) \simeq \cA$.
Then the two quasi-isomorphisms are precisely
$\mathbb{D}_{\Ran}\barBch_X(\cA)\simeq\cB$ and
$\mathbb{D}_{\Ran}\barBch_X(\cB)\simeq\cA$. Dualizing once more gives
weak equivalences between the corresponding bar coalgebras, so the
Koszul twisting morphisms supplied by the bar construction identify
$(\cA,\cB)$ as a finite chiral Koszul pair.
\end{proof}

\subsection{The mixed-boundary spectral sequence}
\label{subsec:mixed-boundary-sseq}
\index{spectral sequence!mixed-boundary}
\index{coupled chiral tensor product}

Fix two augmented chiral algebras $\cA$ and $\cB$ and form the
chiral tensor product $\cA \otimes_{\mathrm{ch}} \cB$. Color
$\cA$-insertions red and $\cB$-insertions blue on
$\overline{\Conf}_{m,n}(X)$.

\begin{definition}[Mixed boundary filtration]
\label{def:mixed-boundary-filtration}
Let $D_{\mathrm{mix}} \subset \overline{\Conf}_{m,n}(X)$ be the
union of boundary strata where at least one red point and one blue
point collide. Define the mixed-boundary filtration by
\[
F^p_{\mathrm{mix}}
:= I(D_{\mathrm{mix}})^p \cdot
 \barBch_X(\cA \otimes_{\mathrm{ch}} \cB).
\]
The bar differential decomposes as
$d = d_\cA + d_\cB + d_{\mathrm{mix}}$, where $d_\cA$ and $d_\cB$
preserve the filtration and $d_{\mathrm{mix}}$ lowers it by one.
\end{definition}

\begin{theorem}[Mixed-boundary spectral sequence;
\ClaimStatusProvedHere]
\label{thm:mixed-boundary-sseq}
Assume the mixed-boundary filtration is exhaustive, complete, and
strongly convergent weight-by-weight. There is a convergent
spectral sequence
\[
E_1
\;\cong\;
H^\bullet_{\mathrm{bar}}(\cA)
 \;\widehat{\otimes}\;
H^\bullet_{\mathrm{bar}}(\cB)
\;\Longrightarrow\;
H^\bullet_{\mathrm{bar}}(\cA \otimes_{\mathrm{ch}} \cB),
\]
whose $d_1$-differential is induced by the mixed residue operator.
\end{theorem}

\begin{proof}
On the associated graded, mixed collisions are invisible: the colored
bar complex separates into the completed tensor product of the pure
red and pure blue bar complexes. The $d_1$-operator is the
first-order effect of $d_{\mathrm{mix}}$, namely the mixed residue.
\end{proof}

\begin{theorem}[PBW-regular tensor theorem; \ClaimStatusProvedHere]
\label{thm:pbw-regular-tensor}
\index{tensor product!regular mixed OPE}
If the mixed OPE between $\cA$-fields and $\cB$-fields is regular
\textup{(}no singular terms\textup{)} and both are chiral Koszul by
the PBW criterion, then $\cA \otimes_{\mathrm{ch}} \cB$ is chiral
Koszul and its bar cohomology coalgebra factorizes:
\[
C_{\cA \otimes_{\mathrm{ch}} \cB}
\;\cong\;
C_\cA \;\widehat{\otimes}\; C_\cB.
\]
If Verdier duality commutes with $\widehat{\otimes}$ in this regime,
the partner algebra factorizes as well:
$(\cA \otimes_{\mathrm{ch}} \cB)^\sharp
\simeq \cA^\sharp \otimes_{\mathrm{ch}} \cB^\sharp$.
\end{theorem}

\begin{proof}
Regular mixed OPE means $d_{\mathrm{mix}} = 0$. The spectral
sequence collapses at $E_1$, giving
$H^\bullet_{\mathrm{bar}}(\cA \otimes_{\mathrm{ch}} \cB)
\cong C_\cA \widehat{\otimes} C_\cB$. Koszul concentration
of both factors implies concentration of the product.
\end{proof}

\subsection{The Sugawara coupling and the universal primarity differential}
\label{subsec:sugawara-coupling}
\index{Sugawara construction!coupled bar complex}

The first genuinely coupled family is the semidirect Sugawara
extension. Let $\mathfrak{g}$ be simple with dual Coxeter
number~$h^\vee$, and let $k \neq -h^\vee$ be a noncritical level.
The Sugawara conformal vector
$L_\mathfrak{g}(z) = \tfrac{1}{2(k+h^\vee)}
\sum_i {:}u_i(z)u^i(z){:}$
produces the semidirect chiral algebra
\[
S_{\mathfrak{g},k}
:= \mathrm{Vir}_{c(k)} \ltimes \widehat{\mathfrak{g}}_k,
\qquad
c(k) = \frac{k\,\dim\mathfrak{g}}{k + h^\vee}.
\]
The mixed $\lambda$-bracket
$[L_\lambda J^a] = (\partial + \lambda)J^a$
says every current is primary of conformal weight~$1$.

\begin{theorem}[Universal first mixed Sugawara differential;
\ClaimStatusProvedHere]
\label{thm:universal-sugawara-d1}
On the one-Virasoro sector
$\eta \otimes C_\mathfrak{g}
\subset H^\bullet_{\mathrm{bar}}(\mathrm{Vir}_{c(k)})
\,\widehat{\otimes}\,
H^\bullet_{\mathrm{bar}}(\widehat{\mathfrak{g}}_k)$,
the first mixed differential is
\begin{equation}\label{eq:universal-d1-sugawara}
d_1^{\mathrm{mix}}(\eta \otimes x)
= 1 \otimes (\partial x + Nx),
\end{equation}
where $\partial$ is the translation coderivation and $N$ is the
total conformal-weight Euler operator on the affine dual
coalgebra~$C_\mathfrak{g}$.
\end{theorem}

\begin{proof}
The mixed differential records the singular part of the mixed OPE
between $L$ and the currents. The primality relation
$L_{(0)} J^a = \partial J^a$, $L_{(1)} J^a = J^a$,
$L_{(n)} J^a = 0$ for $n \geq 2$
contributes exactly two terms from each mixed collision: the
simple-pole residue ($\partial$) and the double-pole coefficient
(conformal weight). On arbitrary current composites the latter is
the Euler operator~$N$, because conformal weights add under normal
ordering and $\partial$ raises weight by one.
\end{proof}

\begin{corollary}[Universality; \ClaimStatusProvedHere]
\label{cor:sugawara-universality}
The operator $d_1^{\mathrm{mix}} = \partial + N$ depends on neither
the structure constants of~$\mathfrak{g}$ nor the value of~$k$,
only on the existence condition $k \neq -h^\vee$.
\end{corollary}

\begin{proof}
The preceding theorem identifies the mixed differential with the sum
of the translation coderivation and the conformal-weight Euler
operator. Its proof uses only the Virasoro primarity relations
$L_{(0)}J^a=\partial J^a$ and $L_{(1)}J^a=J^a$, together with the
fact that conformal weights add under normal ordering. No structure
constants of~$\mathfrak{g}$ enter, and the level appears only through
the existence of the Sugawara field, namely $k \neq -h^\vee$.
\end{proof}

\begin{proposition}[Contractibility on positive current weight;
\ClaimStatusProvedHere]
\label{prop:sugawara-contraction}
\index{contracting homotopy!Sugawara primarity}
On the reduced positive current-weight sector
$\eta \otimes \bar{C}_\mathfrak{g}$, the operator
$Q_{\mathrm{Sug}} := \partial + N$ is invertible with Neumann
series inverse
\[
Q_{\mathrm{Sug}}^{-1}
= \sum_{r \geq 0} (-N^{-1}\partial)^r \, N^{-1},
\]
which converges in the weight completion because $N^{-1}\partial$
is topologically nilpotent. Consequently the reduced mixed sector
is contractible.
\end{proposition}

\begin{proof}
On $\bar{C}_\mathfrak{g}$, the operator~$N$ is strictly positive
(every reduced current class has positive conformal weight), hence
invertible. The coderivation~$\partial$ raises weight by one, so
$N^{-1}\partial$ is topologically nilpotent, and the Neumann series
converges. A homotopy
$h(1 \otimes x) = \eta \otimes Q_{\mathrm{Sug}}^{-1} x$,
$h(\eta \otimes x) = 0$
satisfies $d_1^{\mathrm{mix}} h + h\, d_1^{\mathrm{mix}} = \mathrm{id}$.
\end{proof}

\begin{remark}[The surviving class]
\label{rem:surviving-conformal-class}
The contraction does \emph{not} kill
$\eta \otimes 1 \in \eta \otimes C_\mathfrak{g}$,
because $1 \notin \bar{C}_\mathfrak{g}$ and
$(\partial + N)(1) = 0$.
The vacuum-tensored conformal class survives the primarity stage
and sources the next transgression.
\end{remark}

\subsection{The Casimir transgression complex}
\label{subsec:casimir-transgression-complex}
\index{Casimir transgression complex|textbf}

The first mixed differential encodes only primarity. The first
genuinely \emph{new} coupled information is the Sugawara relation
$L = \nu_{\mathfrak{g},k}$, where
$\nu_{\mathfrak{g},k}
= \tfrac{1}{2(k+h^\vee)} \sum_i u_i(-1) u^i(-1)\lvert 0\rangle$.
Its bar shadow is the Casimir class
$\Omega_{\mathfrak{g},k}
:= [\nu_{\mathfrak{g},k}] \in C_\mathfrak{g}$.

\begin{construction}[Minimal Casimir-transgression complex]
\label{cons:minimal-casimir-transgression}
Define
\[
\mathfrak{C}_{\mathfrak{g},k}
:= \bigl(C_\mathfrak{g} \otimes \Lambda(\eta),\;
 d\eta = \Omega_{\mathfrak{g},k},\;
 d|_{C_\mathfrak{g}} = 0\bigr).
\]
This is \emph{minimal} in the sense that the entire
positive-current primarity sector has been contracted away; the
only remaining Virasoro datum is the conformal generator~$\eta$.
\end{construction}

\begin{theorem}[Sugawara Casimir transgression; \ClaimStatusProvedHere]
\label{thm:sugawara-casimir-transgression}
In the minimal coupled model, $d_2(\eta) = \Omega_{\mathfrak{g},k}$.
\end{theorem}

\begin{proof}
After the primarity reduction, $\eta$ survives because it is
tensored with the vacuum. The genuine Sugawara quotient imposes
$L_{-2}\lvert 0\rangle
= \tfrac{1}{2(k+h^\vee)} \sum_i u_i(-1)u^i(-1)\lvert 0\rangle$,
whose right-hand side is exactly~$\Omega_{\mathfrak{g},k}$.
Passing to the coupled retract, the first new differential is
the class of this relation, hence $d_2(\eta) = \Omega_{\mathfrak{g},k}$.
\end{proof}

\begin{theorem}[Homology of the minimal Casimir-transgression
complex; \ClaimStatusProvedHere]
\label{thm:casimir-transgression-homology}
Let $A_\mathfrak{g} := C_\mathfrak{g}^{\vee,\mathrm{cont}}$ be the
continuous dual, identified with the completed symmetric algebra
$\widehat{\operatorname{Sym}}
(\bigoplus_{n \geq 1} \mathfrak{g}^\vee_{(n)})$.
Let $q_{\mathfrak{g},k}
:= \tfrac{1}{2(k+h^\vee)} \sum_i x_i^{(1)} x^{i,(1)}$
be the dualized Casimir element. Then:
\begin{enumerate}[label=\textup{(\roman*)}]
\item $q_{\mathfrak{g},k}$ is a non-zero-divisor in the
 integral domain~$A_\mathfrak{g}$;
\item the one-constraint Koszul complex
 $\mathcal{K}(A_\mathfrak{g};\, q_{\mathfrak{g},k})
 = (A_\mathfrak{g}[\theta],\, d\theta = q_{\mathfrak{g},k})$
 has homology
 $H^\bullet(\mathcal{K}) \cong A_\mathfrak{g}/(q_{\mathfrak{g},k})$,
 concentrated in degree~$0$;
\item the minimal coupled dual is the continuous dual of the
 quadric:
 $H_\bullet(\mathfrak{C}_{\mathfrak{g},k})
 \cong (A_\mathfrak{g}/(q_{\mathfrak{g},k}))^{\vee,\mathrm{cont}}$.
\end{enumerate}
\end{theorem}

\begin{proof}
For (i), the completed symmetric algebra
$A_\mathfrak{g} = \widehat{\operatorname{Sym}}
(\bigoplus_{n \geq 1} \mathfrak{g}^\vee_{(n)})$
is an integral domain: if $f, g \neq 0$, then
$\mathrm{wt}_{\min}(fg) = \mathrm{wt}_{\min}(f) + \mathrm{wt}_{\min}(g)$
by the ordinary domain property of polynomial algebras on the
associated graded. Since $q_{\mathfrak{g},k} \neq 0$ for
$k \neq -h^\vee$, it is a non-zero-divisor.

For (ii), the Koszul complex has
$H^0 = \operatorname{coker}(A_\mathfrak{g}
\xrightarrow{\cdot q} A_\mathfrak{g})
= A_\mathfrak{g}/(q_{\mathfrak{g},k})$
and $H^1 = \ker(\cdot q) = \operatorname{Ann}(q) = 0$,
using (i).

Part (iii) follows by continuous dualization.
\end{proof}

\begin{theorem}[Quadric rigidity: no higher $A_\infty$~corrections;
\ClaimStatusProvedHere]
\label{thm:casimir-quadric-rigidity}
\index{rigidity!quadric}
\index{formality!Casimir-transgression complex}
The minimal Casimir-transgression complex is \emph{formal}: the
projection
$\pi \colon \mathcal{K}(A_\mathfrak{g};\,q_{\mathfrak{g},k})
\to A_\mathfrak{g}/(q_{\mathfrak{g},k})$
is a quasi-isomorphism of dg~algebras. By
Kadeishvili's theorem, the transferred minimal $A_\infty$-structure
on the homology is strict:
\[
m_1 = 0,
\qquad
m_2 = \text{ordinary multiplication},
\qquad
m_n = 0 \;\;(n \geq 3).
\]
For the minimal coupled Sugawara model, there are no higher
transgressions beyond the quadratic Casimir equation.
\end{theorem}

\begin{proof}
The projection $\pi$ sends $\theta \mapsto 0$ and is a dg~algebra
map. Since $A_\mathfrak{g}/(q)$ with
zero differential and ordinary multiplication is already a strict
dg~algebra model, the transferred $A_\infty$~structure is strict.
\end{proof}

\begin{remark}[What is and is not proved]
\label{rem:casimir-what-proved}
This solves the \emph{minimal} coupled problem exactly. It does
not prove that the entire coupled bar complex of the full Sugawara
quotient retracts to~$\mathfrak{C}_{\mathfrak{g},k}$. What it
proves is that once the coupled problem has been reduced to the
single surviving conformal generator and the single quadratic
Casimir relation, no further higher structure remains. The
appearance of genuine higher operations is tied not to coupling as
such, but to \emph{singular coupling}: the failure of regularity.
\end{remark}

\subsection{The complete-intersection transgression principle}
\label{subsec:ci-transgression}
\index{complete intersection!transgression}

The one-constraint theorem generalizes to arbitrary
codimension.

\begin{theorem}[Complete-intersection transgression principle;
\ClaimStatusProvedHere]
\label{thm:ci-transgression-principle}
Let $A$ be a weight-complete commutative domain and
$\mathbf{q} = (q_1, \ldots, q_r)$ a regular sequence. Then
\[
H^\bullet\!\bigl(
 A[\theta_1, \ldots, \theta_r],\;
 d\theta_i = q_i,\;
 d|_A = 0
\bigr)
\;\cong\;
A/(q_1, \ldots, q_r),
\]
the projection is a quasi-isomorphism of dg~algebras, and the
minimal $A_\infty$-model on the homology is strict.
\end{theorem}

\begin{proof}
Induction on~$r$. The case $r = 1$ is
Theorem~\ref{thm:casimir-transgression-homology}(ii). For the
inductive step, set $A' = A/(q_1, \ldots, q_{r-1})$.
Because~$\mathbf{q}$ is regular, $q_r$ is a non-zero-divisor
in~$A'$, and the $r = 1$ case applies.
\end{proof}

\subsection{Diagonal GKO transgression}
\label{subsec:gko-transgression}
\index{GKO coset!bar complex}
\index{diagonal Casimir!transgression}

The diagonal Goddard--Kent--Olive coupling exhibits the same
structural pattern. For a simple~$\mathfrak{p}$ with diagonal
embedding
$\widehat{\mathfrak{p}}_{m+m'} \hookrightarrow
V^m(\mathfrak{p}) \otimes V^{m'}(\mathfrak{p})$,
the coset Virasoro class $\eta_{\mathrm{GKO}}$ has coupled part
\[
\Omega_\Delta
:= -\frac{1}{m + m' + h^\vee_\mathfrak{p}}
\left[
 \sum_i u_i(-1)\lvert 0\rangle \otimes u^i(-1)\lvert 0\rangle
\right].
\]

\begin{theorem}[Diagonal GKO transgression; \ClaimStatusProvedHere]
\label{thm:gko-transgression}
In the minimal coupled GKO model,
$d_2(\eta_{\mathrm{GKO}}) = \Omega_\Delta$.
The dualized GKO constraint
$q_\Delta
= -\tfrac{1}{m+m'+h^\vee_\mathfrak{p}}
 \sum_i x_i^{\prime,(1)} x^{i,\prime\prime,(1)}$
is a non-zero-divisor in the completed tensor product
$A_\Delta = \widehat{\operatorname{Sym}}
(\bigoplus_{n \geq 1} \mathfrak{p}^{\vee\prime}_{(n)}
\oplus \bigoplus_{n \geq 1} \mathfrak{p}^{\vee\prime\prime}_{(n)})$.
Consequently, the minimal coupled GKO dual is the strict
quadric~$A_\Delta/(q_\Delta)$ with no higher $A_\infty$
corrections.
\end{theorem}

\begin{proof}
The primarity reduction is identical to the Sugawara case. The
surviving class $\eta_{\mathrm{GKO}}$ transgresses to the diagonal
Casimir cross term. The domain property of~$A_\Delta$ (an integral
domain as a completed polynomial algebra) gives regularity of
$q_\Delta$, and
Theorem~\ref{thm:casimir-transgression-homology} applies.
\end{proof}

\begin{remark}[Narrative thread: from separation to rigidity]
\label{rem:casimir-transgression-position}
The theorems of this section form a single argument in five steps,
not a collection of independent results.

\emph{Step~1} (Theorem~\ref{thm:reconstruction-vs-duality}):
$R_X \neq K_X$: reconstruction and duality are distinct
operations on the bar complex. This sits at the foundation of
the algebraic engine: Theorems~A and~B are about~$R_X$, while
the Koszul pair structure of
Chapter~\ref{chap:koszul-pairs} is about~$K_X$.

\emph{Step~2}
(Theorem~\ref{thm:mixed-boundary-sseq}):
the mixed-boundary spectral sequence extends the engine to coupled
tensor products, separating mono-color from mixed-color data.

\emph{Step~3}
(Theorem~\ref{thm:universal-sugawara-d1} and
Proposition~\ref{prop:sugawara-contraction}):
the universal primarity differential $Q_{\mathrm{Sug}} = \partial + N$
contracts away all reduced mixed data; it depends on neither the Lie
type nor the level.

\emph{Step~4}
(Theorem~\ref{thm:sugawara-casimir-transgression} and
Theorem~\ref{thm:casimir-transgression-homology}):
the surviving class transgresses to the Casimir quadric
$\Omega_{\mathfrak{g},k}$, cutting out the coupled dual as a strict
complete intersection.

\emph{Step~5}
(Theorem~\ref{thm:casimir-quadric-rigidity} and
Theorem~\ref{thm:ci-transgression-principle}):
the quadric is formal (no higher $A_\infty$ corrections), and the
principle generalizes to arbitrary codimension via regular sequences.

The conceptual upshot: genuine higher $A_\infty$ corrections in
coupled chiral bar complexes arise from \emph{singular}
coupling (the failure of the transgression classes to form a
regular sequence), not from coupling per~se. The entire sequence
is a Maurer--Cartan element in the mixed-color
convolution algebra, and the Casimir
transgression is the projection of that MC element to the
quadratic stratum.
\end{remark}

\section{Curved Bar-Cobar at CY-$2$ / CY-$3$}
\label{sec:bccurved-supplement}

\begin{remark}[CY-$2$ curved bar-cobar at K3]
\label{rem:bccurved-K3-cy2}
At $\cC = D^b\mathrm{Coh}(K3)$, curved bar-cobar produces the Mukai-Heisenberg
$\mathcal{H}_{\mathrm{Muk}}$ at $c = 24$ on $\Lambda_{\mathrm{Muk}} = \mathrm{II}_{4,20}$ as the
$E_2$-chiral output of $\Phi_2$, with shift $[2]$ and the projective $\widetilde{M}_{24}$
anomaly cocycle of order $6$ in $H^3(M_{24},U(1))\supseteq\mathbb Z/12$ (generalised-Mathieu-moonshine
't~Hooft anomaly, Gaberdiel--Persson--Ronellenfitsch--Volpato 2013; not a Schur multiplier, since
$M_{24}$'s is trivial per ATLAS). The curvature term comes from the Mukai pairing's signature
$(4, 20)$: $c_+ = 4$ is the positive-signature contribution to the Koszul conductor
$K^{\kappa_{\mathrm{ch}}} = 2c_+ = 8$.
\end{remark}

\begin{remark}[Curved bar-cobar at $K3 \times E$ and the Hall--Drinfeld double]
\label{rem:bccurved-K3xE-hall-drinfeld}
At $K3 \times E$ ($d = 3$) the curved bar-cobar output is a
recognition target for the K3 chiral Hall--Drinfeld double
$\mathcal{D}_\hbar(\mathcal{Y}^{\mathrm{Hall}}_\hbar(\mathrm{CoHA}_{K3 \times E}))$
(Vol~III) with Schiffmann--Vasserot shuffle presentation attached to the Manin
pair $(\mathfrak{g}_{\Delta_5}, \mathfrak{n}_+^{\mathrm{imag}} \oplus \mathfrak{h}^{\mathrm{imag,rk\,23}})$;
signature-$(2,1)$ Cartan forbids Manin-triple Lagrangian decomposition.
The finite Hall/CoHA source, Mukai/Siegel pairing,
PBW/no-extra-relations check, radical quotient, parity fixture,
completion, inverse-limit, and Heegner divisor comparison are gates for
recognition. Vol~I does not construct an intrinsic
$H^2(\mathfrak{g}_{\Delta_5})$ classification; the Igusa cusp form is
the Heegner target against which the completed source is compared.
\end{remark}

\section{K3 Hall Curvature and the Igusa Cusp Form}
\label{sec:bcac-deep}

\begin{remark}[Curved bar-cobar adjunction for $\mathbf H_{\Delta_5}$]
\label{rem:bcac-deep-1}
Assume the K3 Hall candidate has passed the source-recognition gates
and write the resulting completed object as $\mathbf H_{\Delta_5}$. The
curved bar-cobar adjunction on this object is controlled by the
curvature element
\[
w^{\mathrm{chir}} \;=\; \hbar^2\cdot K^{\kappa_{\mathrm{ch}}}\cdot \mathbf{1}_{\mathbf H_{\Delta_5}} \;=\; -\mathbf{1}_{\mathbf H_{\Delta_5}}\quad\text{at $\hbar^2 = -1/8$},
\]
so that the curved bar $B^w(\mathbf H_{\Delta_5})$ admits a strict
quasi-isomorphism
$\Omega^w(B^w(\mathbf H_{\Delta_5}))\simeq\mathbf H_{\Delta_5}$
as curved coalgebras-algebras only on the curvature-cancelled
Lusztig surface. Off that surface the curvature is non-trivial and
the adjunction is genuinely curved: the comparison is
pro-object/weight-completed and coderived, detected by comodules
rather than by a raw chain inverse.
\end{remark}

\begin{remark}[Curved bar in the Hall source category]
\label{rem:bcac-deep-2}
A candidate curved bar $B^w(\mathbf H_{\Delta_5})$ is recognized from
finite Hall/CoHA windows of $\mathrm{CoHA}_{K3\times E}$ only after the
pairing, PBW/no-extra-relations, radical, parity, completion,
inverse-limit, and Heegner gates. On each finite window its differential
has the form
\[
d_{B^w} \;=\; d_{\mathrm{Hall}} + d_{\mathrm{Sieg}} + \hbar\, d_{\mathrm{assoc}} + \hbar^2\, d_{\Phi_{10}/\eta^{24}} + w^{\mathrm{chir}}\cdot\mathrm{Id},
\]
where the last piece is the curvature contribution. The closedness
$(d_{B^w})^2 = 0$ in the square-zero total model \emph{requires}
$\hbar^2\cdot K^{\kappa_{\mathrm{ch}}} = -1$,
reproducing the universal identity as the curvature-cancellation
condition. Away from this equality $(B^w,d_{B^w},w^{\mathrm{chir}})$
is a curved CDG object. Primary-lit: Positselski 2011 (curved bar),
Schiffmann-Vasserot 2012, Drinfeld 1990.
\end{remark}

\begin{remark}[Mukai pairing as curvature deformation parameter]
\label{rem:bcac-deep-3}
The Mukai pairing signature $(c_+,c_-) = (4,20)$ on
$\Lambda_{\mathrm{Muk}} = II_{4,20}$ enters the curved bar-cobar
adjunction as the deformation parameter governing the
$\mathrm{Sp}_4(\mathbb Z)$-action: the positive-cone contribution
$c_+ = 4$ controls the Mukai-doubling $K^{\kappa_{\mathrm{ch}}} = 2c_+ = 8$,
while the negative-cone contribution $c_- = 20$ controls the timelike
Cartan direction of the Fake Monster BKM algebra at the Borel cusp.
The curvature element $w^{\mathrm{chir}} = -1$ at Lusztig reflects
the unique stable point in this two-parameter deformation. Primary-lit:
Mukai 1987, Borcherds 1992, Gritsenko-Nikulin 1996--98.
\end{remark}

\section{The Curved Bar-Cobar Pair on
the Humbert stratification}
\label{sec:bcac-supplement}

The curved CDG-coalgebra surface
of Theorem~\ref{thm:chiral-positselski-weight-completed} carries the
K3 chiral bialgebra $\mathbf{H}_{\Delta_5}$ with curvature
$w^{\mathrm{chir}}$. On the open Koszul locus of the Humbert
stratification of $\overline{\mathcal{A}_2}$ this curvature cancels
exactly, so the curved adjunction is a strict chain-level inversion.
On a Humbert wall strict inversion is obstructed; the surviving
statement is weight-completed and coderived, with the
$A_\infty$-deformation providing the minimal comodule-detected
homotopy at $H_3$.

\begin{remark}[Curved bar-cobar strict inversion on
$U = \overline{\mathcal{A}_2}\setminus(H_1\cup H_3\cup H_4)$]
\label{rem:bcac-1}
\index{curved bar-cobar adjunction!strict inversion on Koszul locus}
\index{Koszul locus!curvature cancellation}
On the Koszul locus $U$ of
Theorem~\ref{thm:climax-W18-global-inversion}, the curvature
element of
Remark~\ref{rem:bcac-deep-1},
\[
 w^{\mathrm{chir}}
 \;=\; \hbar^2 \cdot K^{\kappa_{\mathrm{ch}}}
 \cdot \mathbf{1}_{\mathbf{H}_{\Delta_5}}
 \;=\; -\mathbf{1}_{\mathbf{H}_{\Delta_5}}
 \quad \text{at}\quad \hbar^2 = -1/8,\ K^{\kappa_{\mathrm{ch}}} = 8,
\]
cancels exactly against the identity piece of the bar
differential, making the curved CDG-coalgebra
$(B^w(\mathbf{H}_{\Delta_5}), d_{B^w}, w^{\mathrm{chir}})$
strictly closed: $(d_{B^w})^2 = 0$ without the residual
$w\cdot\mathrm{Id}$ piece. The curved cobar
$\Omega^w(B^w(\mathbf{H}_{\Delta_5}))|_U \simeq
\mathbf{H}_{\Delta_5}|_U$ is therefore a chain-level
quasi-isomorphism. This is the
chain-level upgrade of
Remark~\ref{rem:bcac-deep-1} from its K3-chiral ambient
statement.
\end{remark}

\begin{remark}[Curvature perturbation on $H_1$: $\mathbb{Z}/8$-band]
\label{rem:bcac-2}
\index{curved bar-cobar!$H_1$ curvature perturbation}
\index{Bruinier Chern class!$\mathbb{Z}/8$-monodromy}
On a formal neighbourhood of $H_1$, the curvature
$w^{\mathrm{chir}}$ picks up an additional Bruinier-Chern
contribution from the order-$8$ monodromy
\cite[Prop.~5.1]{Bruinier2002LNM}: writing
$\zeta = e^{2\pi i/8}$ for the primitive $8$-th root of unity, the
curvature on the $\mathbb{Z}/8$-eigenspace $B^{(k)}$ at $H_1$ is
$w^{(k)} = \zeta^k\cdot w^{\mathrm{chir}}$, so the curved bar
differential acquires the $\mathbb{Z}/8$-twist
$d_{B^w}|_{H_1} = \bigoplus_k (d + \zeta^k\mathrm{Id}\cdot
w^{\mathrm{chir}})$. Each eigenspace piece is bounded; the
direct-sum bar coalgebra fails Mittag--Leffler precisely on the
phase mismatch between eigenspaces, while the weight-completed bar
coalgebra restores it (Lemma~\ref{lem:weight-filtration-basics}).
Strict chain-level curved bar-cobar is lost at $H_1$; the
weight-completed coderived analogue survives.
\end{remark}

\begin{remark}[Curvature perturbation on $H_3$: three-loop
Costello $c_3 = -8$ and $A_\infty$ homotopy]
\label{rem:bcac-3}
\index{curved bar-cobar!$H_3$ three-loop obstruction}
\index{Costello correction!$c_3 = -8$}
\index{$A_\infty$ deformation!Positselski perturbation lemma}
On a formal neighbourhood of $H_3$, the $\hbar^3$-order of the
curved bar-cobar pair acquires the three-loop obstruction
\[
 \phi^{(3)}|_{H_3}
 \;=\; -8 \cdot c_{\Phi_{10}/\eta^{24}}|_{H_3},
\]
the Costello--Gaiotto/Gritsenko--Nikulin normalized value
\cite{Costello2013,CostelloGaiotto2018,
GritsenkoNikulin1998}). The curved bar differential acquires the
order-$\hbar^3$ correction
\[
 d_{B^w}|_{H_3}
 \;=\; d_{\mathrm{Hall}} + d_{\mathrm{Sieg}}
 + \hbar\, d_{\mathrm{assoc}}
 + \hbar^2\, d_{\Phi_{10}/\eta^{24}}
 + \hbar^3\, (-8)\, d_{\Phi_{10}/\eta^{24}}|_{H_3}
 + w^{\mathrm{chir}}\cdot\mathrm{Id},
\]
which is non-zero as a $3$-cocycle on the weight-completed
Chevalley--Eilenberg complex. The strict chain-level curved
bar-cobar $(B^w, \Omega^w)$ fails at $H_3$; the
$A_\infty$-deformation installs the explicit homotopy
$h^{(3)}|_{H_3}$ with
$[d_{B^w}, h^{(3)}|_{H_3}] = \phi^{(3)}|_{H_3}$,
constructed by the Positselski perturbation lemma
\cite[\S\S3--4]{Positselski11}. The resulting $A_\infty$-curved
bar-cobar pair has a coacyclic cone in the completed coderived
category and requires the $h^{(3)}$ homotopy as extra data; it is not
a raw chain-level inverse on the wall.
\end{remark}

\begin{remark}[Curvature perturbation on $H_4$: Kuga--Satake
$\mathbb{Z}/2$-gerbe]
\label{rem:bcac-4}
\index{curved bar-cobar!$H_4$ Kuga--Satake twist}
\index{etale gerbe!$\mathbb{Z}/2$-band on $H_4$}
On a formal neighbourhood of $H_4$, the Gauss--Manin monodromy
carries a $\mathbb{Z}/2$-twist
\cite[Thm.~9.3]{Bruinier2002LNM} from the real-multiplication
involution $\sqrt{2}\mapsto -\sqrt{2}$. The curved bar coalgebra
$B^w(\mathbf{H}_{\Delta_5})|_{H_4}$ splits as a sum of $\pm
1$-isotypic components
$B^w(\mathbf{H}_{\Delta_5})|_{H_4} = B^{w, +}\oplus B^{w, -}$
exchanged by the involution. Morrison's global Torelli theorem
\cite{Morrison84} lifts this $\mathbb{Z}/2$-twist through the
Kuga--Satake construction \cite{KugaSatake67,Deligne72} to an
\'etale $\mathbb{Z}/2$-gerbe $\mathcal{G}_{H_4}\to H_4$ whose
banding cocycle in $H^2_{\text{\'et}}(H_4, \mathbb{Z}/2)$
trivialises the twist on \'etale neighbourhoods but not globally.
Strict chain-level curved bar-cobar descends through
$\mathcal{G}_{H_4}$ only at the pro-object coderived level; the
weight-completed curved bar-cobar
(Theorem~\ref{thm:chiral-positselski-weight-completed}) handles
the gerbe banding automatically.
\end{remark}

\begin{theorem}[Curved bar-cobar trichotomy on
$\overline{\mathcal{A}_2}$]
\label{thm:bcac-trichotomy}
\index{curved bar-cobar adjunction!trichotomy}
\index{bar-cobar inversion!global chain-level}
Assume the K3 Hall/CoHA source has passed the finite-source, pairing,
PBW/no-extra-relations, radical, parity, completion, inverse-limit, and
Heegner comparison gates, and write the recognized completed object as
$\mathbf{H}_{\Delta_5}$. The curved bar-cobar adjunction
$(\Omega^w\dashv B^w)$ on it satisfies the three
ambient-qualified equivalences:
\begin{enumerate}
\item[\textup{(a)}] $\Omega^w(B^w(\mathbf{H}_{\Delta_5}))|_U
\simeq \mathbf{H}_{\Delta_5}|_U$ as a \emph{strict chain-level
quasi-isomorphism} on the Koszul locus $U$, with curvature
$w^{\mathrm{chir}} = -1$ cancelling against the identity piece
of the bar differential
(Remark~\ref{rem:bcac-1}).
\item[\textup{(b)}] $\Omega^w\circ B^w \simeq \mathrm{id}$
\emph{weight-completed in}
$D^{\mathrm{co}}_{\mathrm{ch}}(\overline{\mathcal{A}_2})$ across
the full base, with weight-completed witness
$\widehat{\bar B}^{\mathrm{ch}}(\mathbf{H}_{\Delta_5})$
(Theorem~\ref{thm:chiral-positselski-weight-completed}).
\item[\textup{(c)}] On a formal neighbourhood of $H_3$, strict
chain-level inversion fails with obstruction
$\phi^{(3)}|_{H_3} = -8\cdot c_{\Phi_{10}/\eta^{24}}|_{H_3}$; the
$A_\infty$-deformation installs the homotopy $h^{(3)}|_{H_3}$ of
Remark~\ref{rem:bcac-3} to produce a coacyclic-cone repair in the
completed coderived $A_\infty$ ambient
(cross-ref Theorem~\ref{thm:climax-W18-global-inversion}~(c)).
\end{enumerate}
\end{theorem}

\begin{proof}
Theorem~\ref{thm:climax-W18-global-inversion} supplies the same
trichotomy in the chiral bar-complex ambient; the present theorem
transports it to the curved CDG-coalgebra ambient. For statement~(a),
the curvature cancellation
$w^{\mathrm{chir}} + \mathrm{Id} = 0$ at $\hbar^2 = -1/8$,
$K^{\kappa_{\mathrm{ch}}} = 8$
(Remark~\ref{rem:bcac-deep-1}) is uniform on $U$ because the
Mittag--Leffler condition holds uniformly on $U$
(Lemma~\ref{lem:cclimax-H1-failure} shows it fails only at
$H_1$, Lemma~\ref{lem:cclimax-H4-failure} only at $H_4$,
and the three-loop coboundary of
Lemma~\ref{lem:cclimax-H3-failure} is non-zero only at $H_3$;
away from their union the strict quasi-isomorphism is
well-posed). Statement~(b) is the weight-completed Theorem~B of
Theorem~\ref{thm:chiral-positselski-weight-completed} applied to
$\mathbf{H}_{\Delta_5}$; see
Theorem~\ref{thm:climax-W18-global-inversion}~(b). Statement~(c)
is the $A_\infty$ Positselski perturbation at $H_3$
(Remark~\ref{rem:bcac-3}).

Primary inputs: Positselski $2011$
(CDG-coalgebra perturbation) \cite{Positselski11};
Gritsenko--Nikulin $1998$ (denominator restriction)
\cite{GritsenkoNikulin1998};
Bruinier $2002$ LNM~$1780$ (Heegner--Chern reciprocity)
\cite{Bruinier2002LNM}; Morrison $1984$ (global Torelli for K3)
\cite{Morrison84}.
\end{proof}

\begin{remark}[Relation to the Vol.~III
$\kappa_{\mathrm{BKM}} = c_N(0)/2$ identity]
\label{rem:bcac-5}
\index{universal trace identity!chain-level footing}
The trichotomy
(Theorem~\ref{thm:bcac-trichotomy}) provides the
ambient-qualified footing for the universal trace identity
$\kappa_{\mathrm{BKM}}(X) = \tfrac{1}{2} K(A_X)$ of
Equation~\eqref{eq:trace-identity-platonic}: the identity is
defined on the K3-fibered Class-A locus where the chiral bar
complex $B^w(A_X)$ satisfies the trichotomy, and the
$\phi^{(3)}|_{H_3}$ obstruction tracks the three-loop correction
to the 1-loop BKM determinant. At the strict chain level on the
Koszul locus $U$, the identity reduces to a finite-weight check;
weight-completed, it extends across Humbert walls; at $H_3$ the
$A_\infty$-deformation carries the $c_3 = -8$ correction
faithfully. The chapter's foundation (establishing the
curvature cancellation at the Lusztig point) and the chain-level
(establishing the Humbert stratification) therefore together
provide the strict-open and coderived-wall scaffold on which the Vol.~III
identification $\kappa_{\mathrm{BKM}} = c_N(0)/2$ of
Chapter~\ref{V3-ch:k3-chiral-bialgebra-platonic} rests.
\end{remark}

\section{Extension/obstruction formula for curved bar-cobar across each
admissible Humbert divisor $H_n$ ($n\equiv 0,3\pmod 4$)}
\label{sec:bcac-extension-obstruction-admissible}

The trichotomy theorem
(Theorem~\ref{thm:bcac-trichotomy}) records curved bar-cobar
in three ambient-qualified lanes on the leading three walls
$\{H_1,H_3,H_4\}$. The same extension/obstruction mechanism exists
across every admissible Humbert divisor
$H_n\subset\overline{\mathcal{A}_2}$ with $n\equiv 0,3\pmod 4$, and
is governed by the explicit formula below on the formal neighbourhood
of $H_n$.

\subsection{The curved nearby-cycle functor $\psi_n^w$}
\label{subsec:bcac-curved-nearby-cycle}

\begin{definition}[Curved nearby cycle of the curved bar coalgebra]
\label{def:bcac-curved-nearby-cycle}
\index{nearby cycle!curved bar coalgebra!along admissible $H_n$}
\index{CDG-coalgebra!curved nearby cycle}
Let $(B^w(A),d_{B^w},w^{\mathrm{chir}})$ denote the curved bar
CDG-coalgebra on $A=\mathbf{H}_{\Delta_5}$. For each admissible
$n\equiv 0,3\pmod 4$, the \emph{curved nearby-cycle functor}
along $H_n$ is the triple
\[
 \psi_n^w(A)
 \;:=\;
 \Bigl(
 \mathrm{gr}^V_\bullet\,B^w(A)\big|_{H_n},\;
 \mathrm{gr}^V_\bullet\, d_{B^w}\big|_{H_n},\;
 \mathrm{gr}^V_\bullet\, w^{\mathrm{chir}}\big|_{H_n}
 \Bigr)
\]
obtained by graded-piece restriction along the
Kashiwara--Malgrange $V$-filtration
\cite{Kashiwara1983,Malgrange1983}. The curvature element
$\mathrm{gr}^V_0 w^{\mathrm{chir}}|_{H_n}$ takes the value
\[
 \mathrm{gr}^V_0 w^{\mathrm{chir}}\big|_{H_n}
 \;=\;
 -1 + c_{-2,1}(n,r_n)\cdot c_{\Phi_{10}/\eta^{24}}\big|_{H_n}
 \cdot [H_n]^\vee,
\]
where the correction term is the residue of the Bruinier lift
\cite{Bruinier2002LNM,GritsenkoNikulin1998}. The cases $n=1,3,4$
recover the Costello-corrected values
$-2,-8,+2$ respectively (cross-ref
Remark~\ref{rem:cclimax-W17-costello-link} in
Chapter~\ref{chap:climax-platonic}).
\end{definition}

\begin{lemma}[Curved Maurer--Cartan on the nearby cycle]
\label{lem:bcac-curved-MC-on-nearby-cycle}
\index{Maurer--Cartan!curved nearby cycle}
The triple $\psi_n^w(A)$ of
Definition~\ref{def:bcac-curved-nearby-cycle} satisfies the curved
Maurer--Cartan equation
\[
 (\mathrm{gr}^V_\bullet d_{B^w}|_{H_n})^2
 \;+\;
 [\mathrm{gr}^V_\bullet w^{\mathrm{chir}}|_{H_n},\,-]
 \;=\;
 0
\]
on $\mathrm{gr}^V_\bullet B^w(A)|_{H_n}$; equivalently, $\psi_n^w(A)$
is a CDG-coalgebra in the sense of
\cite[\S4]{Positselski11}.
\end{lemma}

\begin{proof}
The curved bar coalgebra $(B^w(A),d_{B^w},w^{\mathrm{chir}})$
satisfies $d_{B^w}^2+[w^{\mathrm{chir}},-]=0$ on
$\widehat{\mathcal{O}}_{H_n}$
(Remark~\ref{rem:bcac-deep-1}). The $V$-filtration is
preserved by both $d_{B^w}$ and the action of $w^{\mathrm{chir}}$
(by
Lemma~\ref{lem:bar-differential-preserves-monodromy-filtration}
of Chapter~\ref{ch:theorem-B-scope-platonic}; the curvature acts
by identity-times-scalar and trivially preserves any filtration),
hence the curved MC equation descends to graded pieces.
\end{proof}

\subsection{Explicit extension/obstruction formula}
\label{subsec:bcac-explicit-extension-obstruction}

\begin{theorem}[Extension/obstruction formula for curved bar-cobar
across $H_n$]
\label{thm:bcac-extension-obstruction-across-Hn}
\index{extension!curved bar-cobar!across $H_n$}
\index{obstruction!curved bar-cobar!explicit formula}
\index{Positselski perturbation!curved chain level}
Let $n\equiv 0,3\pmod 4$ be admissible and let
$(B^w(A),d_{B^w},w^{\mathrm{chir}})|_{\widehat{U}_{H_n}}$ denote
the curved bar coalgebra on the punctured formal neighbourhood.
The obstruction to extending the strict chain-level curved bar-cobar
quasi-isomorphism $\Omega^w(B^w(A))\xrightarrow{\sim}A$ from
$\widehat{U}_{H_n}$ to $\widehat{\mathcal{O}}_{H_n}$ is the class
\[
\boxed{\quad
 [\mathrm{obs}^w_{H_n}]
 \;=\;
 c_{-2,1}(n,r_n)\cdot
 \bigl[c_{\Phi_{10}/\eta^{24}}|_{H_n}\bigr]
 \;\in\;
 H^1\!\bigl(
 \mathrm{Hom}_{\mathrm{CDG}}(B^w(A),\Omega^w B^w(A))|_{\widehat{\mathcal{O}}_{H_n}}
 \bigr),
\quad}
\]
where $c_{-2,1}(n,r_n)$ is the $n$-th Fourier coefficient of the
Bruinier lift $\phi_{-2,1}$
(Theorem~\ref{thm:tbsp-lim1-is-fourier-coefficient} in
Chapter~\ref{ch:theorem-B-scope-platonic}). The extension is
strictly at chain level if and only if $[\mathrm{obs}^w_{H_n}]=0$;
otherwise the extension is $A_\infty$-corrected with explicit
homotopy $\phi^{(n)}$ satisfying
\[
 [d_{B^w},\,\phi^{(n)}]\big|_{H_n}
 \;=\;
 c_{-2,1}(n,r_n)\cdot
 c_{\Phi_{10}/\eta^{24}}\big|_{H_n}
 \quad
 \text{in}\quad
 \mathrm{Hom}_{\mathrm{CDG}}(B^w(A),\Omega^w B^w(A)).
\]
In the second case the output is a coderived/comodule-detected
extension with coacyclic cone, not a strict chain-level inverse.
\ClaimStatusProvedHere{}
\end{theorem}

\begin{proof}
Step 1 (reduction to Tor-class). Apply
\cite[Prop.~3.5]{Positselski11}: the strict chain-level extension
obstruction for a CDG-coalgebra quasi-isomorphism across a closed
substack is the $\mathrm{Tor}_1$-class in
$\mathrm{Tor}_1^{\widehat{\mathcal{O}}_{H_n}}(\Omega^w B^w A,A)$.
This identifies
$[\mathrm{obs}^w_{H_n}]$ with the class of
Theorem~\ref{thm:tbsp-Tor-obstruction-per-wall}.

Step 2 (identification with Fourier coefficient). By
Theorem~\ref{thm:tbsp-lim1-is-fourier-coefficient}, this class
equals $c_{-2,1}(n,r_n)\cdot[H_n]^\vee$. Under the de Rham
$\to$ CDG comparison map of
\cite[\S7]{Positselski11}, $[H_n]^\vee$ is the Chern class of
$\Phi_{10}/\eta^{24}$ restricted to $H_n$, i.e.
$c_{\Phi_{10}/\eta^{24}}|_{H_n}$
(Harvey--Moore threshold formula
$-\log|\Phi_{10}|^2=\sum_{\text{inst}}\cdots$
\cite{HarveyMoore1995}).

Step 3 (Positselski perturbation lemma). If
$c_{-2,1}(n,r_n)\neq 0$, the Positselski perturbation lemma
\cite[\S\S3--4]{Positselski11} constructs the explicit homotopy
$\phi^{(n)}$: start from a strict lift of the CDG-coalgebra map on
$\widehat{U}_{H_n}$; attempt to extend across $H_n$; the
boundary of the extension is precisely
$c_{-2,1}(n,r_n)\cdot c_{\Phi_{10}/\eta^{24}}|_{H_n}$; the
perturbation lemma produces $\phi^{(n)}$ with this boundary.

Step 4 (three leading walls). For $n=1,3,4$, specialise:
$c_{-2,1}(1,1)=-2$, $c_{-2,1}(3,1)=-8$, $c_{-2,1}(4,0)=2$
(Remark~\ref{rem:tbsp-leading-three-walls-first-coefficients} in
Chapter~\ref{ch:theorem-B-scope-platonic}), recovering the
three obstruction classes of
Lemmas~\ref{lem:cclimax-H1-failure}--\ref{lem:cclimax-H4-failure}
in Chapter~\ref{chap:climax-platonic}.
\end{proof}

\begin{remark}[Chain-level comparison with the three-loop
$\phi^{(3)}=-8$ formula]
\label{rem:bcac-three-loop-comparison}
\index{three-loop obstruction!chain-level factorisation}
The three-loop coboundary
$\phi^{(3)}|_{H_3}=-8\cdot c_{\Phi_{10}/\eta^{24}}|_{H_3}$ of
Theorem~\ref{thm:bcac-trichotomy}(c) factors as
\[
 \phi^{(3)}|_{H_3}
 \;=\;
 c_{-2,1}(3,1)\cdot c_{\Phi_{10}/\eta^{24}}|_{H_3}
 \;=\;
 (-8)\cdot c_{\Phi_{10}/\eta^{24}}|_{H_3},
\]
where the coefficient $-8$ is the $(3,1)$-Fourier coefficient of
$\phi_{-2,1}$. Thus the Costello correction
\cite{Costello2013,CostelloGaiotto2018} has the chain-level
Fourier-coefficient reading:
$\phi^{(n)}|_{H_n}$ is uniformly
$c_{-2,1}(n,r_n)\cdot c_{\Phi_{10}/\eta^{24}}|_{H_n}$ at every
admissible wall, with $c_{-2,1}(3,1)=-8$ as the leading three-loop
instance.
\end{remark}

\subsection{Čech descent across the admissible Humbert
stratification}
\label{subsec:bcac-cech-descent}

\begin{lemma}[Triple-intersection cocycle splits on
$H_n\cap H_m$]
\label{lem:bcac-triple-intersection-cocycle-splits}
\index{wall-of-walls!cocycle splitting!curved bar-cobar}
For $n,m\equiv 0,3\pmod 4$ admissible and distinct, on the
wall-of-walls intersection $H_n\cap H_m$ (when non-empty) the
commutator class
$[\gamma_n,\gamma_m]\in H^2_{\text{\'et}}(H_n\cap H_m,\mathbb{Z}/N_{nm})$
(where $N_{nm}=\mathrm{lcm}(\ell_n,\ell_m)$) splits canonically
via the Schmid--Cattani--Kaplan commuting-$\mathfrak{sl}_2$-triples
construction \cite{CattaniKaplanSchmid86}. In particular, the
\v Cech cocycle condition on the cover
$\{\widehat{\mathcal{O}}_{H_n}\}_{n\ \text{admissible}}\cup
\{U^{\mathrm{adm}}\}$ is satisfied.
\end{lemma}

\begin{proof}
On $H_n\cap H_m$, the two monodromies $T_n,T_m$ commute
(both are pullbacks from the abelianised $V$-filtration), so the
pair $(T_n,T_m)$ lifts to a commuting pair of quasi-unipotent
operators with Jordan blocks of total size $\ell_n\cdot\ell_m$.
The Schmid--Cattani--Kaplan theorem
\cite[\S4.2]{CattaniKaplanSchmid86} constructs a canonical pair of
commuting $\mathfrak{sl}_2$-triples
$(e_n,f_n,h_n),(e_m,f_m,h_m)$ such that the total weight filtration
splits as a bi-filtration. The commutator cocycle $[\gamma_n,\gamma_m]$
is the obstruction to extending this splitting to the nilpotent
operator $N_n+N_m$; on commuting $\mathfrak{sl}_2$-triples the
splitting extends automatically, so the cocycle is cohomologous to
zero.
\end{proof}

\begin{theorem}[Global curved bar-cobar inversion via \v Cech
descent]
\label{thm:bcac-global-cech-descent}
\index{curved bar-cobar!global inversion!\v Cech descent}
\index{global inversion theorem!curved lane}
Let
$\mathcal{U}_{\mathcal{A}_2}
=\{U^{\mathrm{adm}}\}\cup\{\widehat{\mathcal{O}}_{H_n}\}_{n\ \text{admissible}}$
be the open cover of $\overline{\mathcal{A}_2}$ by the Koszul
locus and the formal neighbourhoods of admissible Humbert
divisors. The curved bar-cobar adjunction
$(\Omega^w\dashv B^w)$ on $A=\mathbf{H}_{\Delta_5}$ assembles via
\v Cech descent on $\mathcal{U}_{\mathcal{A}_2}$ to a global
curved bar-cobar equivalence
\[
 \Omega^w(B^w(A))
 \;\simeq\;
 A
 \quad\text{in}\quad
 D^{\mathrm{co}}_{\mathrm{ch}}(\overline{\mathcal{A}_2}),
\]
assembled as follows:
\begin{itemize}
\item on $U^{\mathrm{adm}}$: strict chain-level quasi-isomorphism
with curvature cancellation $w^{\mathrm{chir}}+\mathrm{Id}=0$
(Theorem~\ref{thm:bcac-trichotomy}(a));
\item on each $\widehat{\mathcal{O}}_{H_n}$: weight-completed
coderived with monodromy-refined filtration, $A_\infty$-corrected
by $\phi^{(n)}=c_{-2,1}(n,r_n)\cdot
c_{\Phi_{10}/\eta^{24}}|_{H_n}$
(Theorem~\ref{thm:bcac-extension-obstruction-across-Hn});
\item on double intersections
$\widehat{\mathcal{O}}_{H_n}\cap U^{\mathrm{adm}}$: restriction to
the generic fibre identifies both lanes
(Mittag--Leffler restriction);
\item on wall-of-walls
$\widehat{\mathcal{O}}_{H_n}\cap\widehat{\mathcal{O}}_{H_m}$: the
Deligne--Malcev bi-unipotent \v Cech cocycle splits by
Lemma~\ref{lem:bcac-triple-intersection-cocycle-splits}.
\end{itemize}
\ClaimStatusProvedHere{}
\end{theorem}

\begin{proof}
Lane assembly: the open lane is a strict chain-level
quasi-isomorphism, while each wall lane is a completed coderived
equivalence with coacyclic cone, by Theorem~\ref{thm:bcac-trichotomy}
and Theorem~\ref{thm:bcac-extension-obstruction-across-Hn}. \v Cech
cocycle condition on double
intersections: on $\widehat{\mathcal{O}}_{H_n}\cap U^{\mathrm{adm}}$
the restriction of the weight-completed lane to the generic fibre
is Mittag--Leffler strict (by
Lemma~\ref{lem:per-isotypic-ML-Hn} of
Chapter~\ref{ch:theorem-B-scope-platonic}) and equals the strict
chain-level lane on $U^{\mathrm{adm}}$. On
$\widehat{\mathcal{O}}_{H_n}\cap\widehat{\mathcal{O}}_{H_m}$,
Lemma~\ref{lem:bcac-triple-intersection-cocycle-splits} delivers
the cocycle splitting. Triple-intersection cocycle condition: on
$\widehat{\mathcal{O}}_{H_n}\cap\widehat{\mathcal{O}}_{H_m}
\cap\widehat{\mathcal{O}}_{H_\ell}$, the Schmid--Cattani--Kaplan
triples are commuting (three pairwise-commuting
$\mathfrak{sl}_2$-triples), producing a triple-commuting splitting;
the \v Cech 2-cocycle condition reduces to the Jacobi identity on
the three commutator brackets, automatic on commuting triples.
Primary: Positselski $2011$ \cite{Positselski11};
Cattani--Kaplan--Schmid $1986$ \cite{CattaniKaplanSchmid86};
Bruinier $2002$ \cite{Bruinier2002LNM};
Saito $1988$ \cite{Saito1988}.
\end{proof}

\begin{remark}[Independent checks for
Theorem~\ref{thm:bcac-global-cech-descent}]
\label{rem:bcac-cech-descent-independent-checks}
Four checks: (1) direct computation at each leading
wall $H_1,H_3,H_4$ via
Lemmas~\ref{lem:cclimax-H1-failure}--\ref{lem:cclimax-H4-failure}
in Chapter~\ref{chap:climax-platonic}; (2) identification
with the global inversion theorem
(Theorem~\ref{thm:tbsp-global-inversion-all-admissible} in
Chapter~\ref{ch:theorem-B-scope-platonic}) lane (iv), which is
this theorem stated in chiral bar ambient (here stated in curved
CDG-coalgebra ambient); (3) Chern-class identification with
$\kappa_{\mathrm{BKM}}(X)=\tfrac{1}{2}K(A_X)$ on the K3-fibered
Class-A locus (Vol.~III trace identity); (4) limiting-case
recovery of Koszul duality on $U^{\mathrm{adm}}$ by restriction.
\end{remark}
